% 本文件由 LaTeX 排版 Agent 根据有效 translation.json 自动生成。
% author=Clement of Rome; work_id=0808; title=克莱孟讲道集

\chapter{克莱孟讲道集}

% structure: p0001 source=h1 target=omitted reason=page-title-already-rendered-by-parent

\Emph{引言书信}\newline{}\Emph{讲道一}\newline{}\Emph{讲道二}\newline{}\Emph{讲道三}\newline{}\Emph{讲道四}\newline{}\Emph{讲道五}\newline{}\Emph{讲道六}\newline{}\Emph{讲道七}\newline{}\Emph{讲道八}\newline{}\Emph{讲道九}\newline{}\Emph{讲道十}\newline{}\Emph{讲道十一}\newline{}\Emph{讲道十二}\newline{}\Emph{讲道十三}\newline{}\Emph{讲道十四}\newline{}\Emph{讲道十五}\newline{}\Emph{讲道十六}\newline{}\Emph{讲道十七}\newline{}\Emph{讲道十八}\newline{}\Emph{讲道十九}\newline{}\Emph{讲道二十}\par

\SourceRecord{Translated by Thomas Smith.；Ante-Nicene Fathers；Vol. 8.；Edited by Alexander Roberts, James Donaldson, and A. Cleveland Coxe.；Buffalo, NY: Christian Literature Publishing Co.,；1886.；<http://www.newadvent.org/fathers/0808.htm>.}

% structure: p0001 source=h1 target=section reason=page-title

\section{导入信函}

% structure: p0002 source=h2 target=subsection reason=relative-heading-rank; base=h2

\subsection{伯多禄致雅各伯书}

\Person{伯多禄}致\Person{雅各伯}，至圣教会的主教，在万有之父之下，藉耶稣基督，恒祝平安。\par

% structure: p0004 source=h3 target=subsubsection reason=relative-heading-rank; base=h2

\subsubsection{第一章 保留的教义}

我的弟兄，我知道你热切渴望那有益于我们众人的事，我恳求并劝请你，不要将我寄给你的宣讲录交给任何外邦人，也不要在考验之前交给我们本族的人；但若有人经过考验而被证明是配得的，就照梅瑟将\Emph{他的书}交付给继承他座位的那七十人的方式，传给他。因此，这种谨慎的果实直到如今依然可见。因为他的同胞处处遵守同一君主制和政体规则，绝不可能以其他方式思想，也不能被多义圣经引离正道。因为按照传授给他们的规则，他们努力纠正圣经中的不一致之处，若有人因不知传统而困惑于先知们的各种言论。因此，他们禁止任何人教导，除非他先学会了如何使用圣经。这样，他们中间有一位天主、一条法律、一个希望。\par

% structure: p0006 source=h3 target=subsubsection reason=relative-heading-rank; base=h2

\subsubsection{第二章 对伯多禄教义的曲解}

因此，为使同样的事也发生在我们当中，如同发生在那七十人身上一样，请将我的宣讲录以同样的入教奥秘交给我们的弟兄们，使他们能教导那些愿意参与教导的人；因为若不这样，我们真理的话语就会被分裂成许多意见。我知道这事，并非作为先知，而是已经看见这邪恶的开端。因为有些外邦人拒绝了我的法律宣讲，依附于我那敌人的某些无法无天和轻浮的宣讲。这些人甚至在我还活着的时候，就试图通过种种不同的解释来篡改我的话语，以废除法律；好像我自己也有这种心思，却没有自由宣告它——绝无此事！因为这样做就是反对天主藉梅瑟所讲的法律，而我们的主曾为它的永恒延续作证；因为祂这样说：\BibleQuote{天地要过去，但法律的一撇一画也绝不能过去。}祂说这话，是为使一切得以实现。但这些人，我不知道他们如何自称知道我的心思，竟敢于解释他们从我这里听到的话语，比我这说话的人更聪明，告诉他们的慕道者这就是我的意思——其实我从未想过。然而，如果我还在世，他们就敢这样歪曲我，那么后来的人将更加胆大妄为！\par

% structure: p0008 source=h3 target=subsubsection reason=relative-heading-rank; base=h2

\subsubsection{第三章 入教仪式}

因此，为避免发生这样的事，我祈祷并恳求你，不要将我寄给你的宣讲录交给任何人，无论本国人或外族人，都要先经过考验；但若有人经过考验而被证明是配得的，就按照梅瑟的入教仪式，他将\Emph{他的书}交付给继承他座位的七十人的方式，传给他；这样他们可以持守信仰，到处传授真理的准则，按我们的传统解释一切；免得他们因无知而自陷，因按自己心思的臆测而陷入错误，并将他人带入同样的毁灭之坑。凡我认为好的事，我已经公正地向你指出；你认为好的事，我的主，也请体面地执行。祝安。\par

% structure: p0010 source=h3 target=subsubsection reason=relative-heading-rank; base=h2

\subsubsection{第四章 关于领受此书者的誓言}

1. 因此，\Person{雅各伯}读了这信后，召集了长老们；他向他们宣读后说：\Quote{我们的\Person{伯多禄}严正而合宜地嘱咐我们关于确立真理的事，我们不应将寄给我们的宣讲录随便交给任何人，只应交给善良、虔诚、愿意教导、受过割礼且忠信的人。而且不应一次全部交付给他；若他在第一次的事上显为不明智，就不应将其余的交托给他。因此，他必须经受不少于六年的考验。然后，按照梅瑟的入教仪式，\Emph{交付书籍的人}应带他到一条河流或一个泉源，就是活水之处，义人的重生在那里发生，并要他——不是宣誓（因为这是不合法的）——而是站在水边起誓，正如我们自己重生时，为了不犯罪而被要求做的那样。}\par

2. 他应说：\InnerQuote{我呼请天地、水，以及其中包含的一切，此外还有那渗透一切、没有它我便不能呼吸的空气作证：我将永远服从那给我宣讲录的人；而他所给我的那些书，我决不以任何方式传给任何人，无论是通过抄写、交付抄写、交给抄写员，或由我自己或由他人，或通过任何其他入教仪式、诡计、方法，或因疏忽保管、放在\Emph{任何人}面前、允许他\Emph{观看}，或以任何方式方法传给任何人；除非我确信某人配得，如同我自己曾被判断一样，甚至更配，并且经过不少于六年的考验；但只传给一位虔诚、良善、被选来教导的人，正如我领受的一样，我将交付给他，并照我主教的意愿行事。}\par

3. \Quote{但另一方面，即使他是我的儿子或兄弟，或朋友，或以任何方式与我有亲属关系，若他不配，我也不会赐予他这份恩惠，因为那是不合适的；我不会因阴谋而畏惧，也不会因礼物而软化。但如果我一旦认为交给我宣讲的书卷不真实，我便不会如此传给他们，而要交还。当我外出时，我会随身携带它们，无论我拥有其中哪些。如果我不打算随身携带，我不会让它们留在我家中，而会存放在我的主教那里，他持有同样的信德，且出自\Emph{像我一样}的人。如果我不幸生病，并面临死亡，且我无子女，我也会同样行事。但如果我死后有一个不配或尚不能胜任的儿子，我也同样行事。因为我会将它们存放在我的主教那里，以便我的儿子长大后，若配得上这份信托，他可以按照这项约定的条款，作为他父亲的遗赠交给他。}\par

4. \Quote{并且我将如此行，我再次呼请天、地、水——万物所包裹在其中——为证，此外还有遍及一切的大气，没有它我便不能呼吸，我愿永远顺服那交给我这些宣讲书卷的人，并在一切事上遵守我所承诺的，或甚至更多。因此，我若遵守此盟约，便与圣者同分；但我若违背所立之约，愿宇宙与我为敌，遍及一切的以太与统驭万有的天主也与我为敌，没有谁能高于祂，没有谁比祂更伟大。即使我将来认识另一位神，我现在也指着祂发誓——无论祂存在与否——我必不另有所行。此外，我若说谎，则生时被咒，死时被诅，且受永罚。}\par

\Quote{此后，让受托付的人与他一同分享饼和盐。}\par

% structure: p0016 source=h3 target=subsubsection reason=relative-heading-rank; base=h2

\subsubsection{第五章 接受的誓词}

\Person{雅各伯}如此说完，长老们恐惧战兢。因此\Person{雅各伯}见他们极其害怕，便说：\Quote{弟兄们，同仆们，请听我说。若我们将书卷不分给所有人，恐怕被一些胆大妄为的人所败坏，或因解释而曲解，正如你们听说有人已做的那样，那么就连真正寻求真理的人也将永远在错误中游荡。因此，这些书卷还是由我们保管更好，并且我们只以先前所提的一切谨慎交给那些愿意虔诚生活并拯救他人的人。但若有人在接受此誓词后另有所行，他很有理由应受永罚。因为那导致他人毁灭的人，怎能自己不灭亡？} 长老们因此对\Person{雅各伯}的话感到欣慰，喊道：\Quote{愿那预知万有、仁慈揀选您作我们主教的主受赞颂！} 说完这话，我们都站起来，向圣父和万有的天主祈祷，愿荣耀归于祂，直到永远。阿们。\par

% structure: p0018 source=h2 target=subsection reason=relative-heading-rank; base=h2

\subsection{\Person{伯多禄}致\Person{雅各伯}书}

\Person{克莱孟}致\Person{雅各伯}主，主教中之主教，因天主的眷顾治理\Place{耶路撒冷}——希伯来人的圣教会，以及各地卓越建立的教会，并诸位长老、执事及其余弟兄：愿平安常与你们同在。\par

% structure: p0020 source=h3 target=subsubsection reason=relative-heading-rank; base=h2

\subsubsection{第一章 \Person{伯多禄}殉道}

我主，请告知您：\Person{西满}——他因真信德及最稳固的教义基础而被分别出来，作为教会的基石，并为此被\Person{耶稣}亲自以祂真理的口命名为\Person{伯多禄}，我们主的初果，宗徒之首；圣父首先向他启示了圣子；\Person{基督}有充分理由祝福了他；蒙召的、被选的、与\Emph{基督}一同坐席及同行的卓越而可嘉的门徒——他既是最适宜的人，受命光照世界较暗的部分，即西方，并得以完成此事——我何以赘言，不愿提及悲伤之事，但无论多么不情愿，也必须告诉您——他因对人类的至深爱心，甚至来到\Place{罗马}，公开而明确地作证，反对那与他为敌的恶者；他宣讲有一位善良的君王统御全世界，同时以其天主启示的教义拯救人，而他本人，因暴力，以现世生命换取了永生。\par

% structure: p0022 source=h3 target=subsubsection reason=relative-heading-rank; base=h2

\subsubsection{第二章 \Person{克莱孟}的祝圣}

但那时，当他即将离世，弟兄们聚集起来，他突然抓住我的手，站起身来，在教会面前说：\Quote{弟兄们，同仆们，请听我说。既然我被主和导师\Person{耶稣基督}——祂的宗徒——所教导，我的死期临近，我将按手在这位\Person{克莱孟}身上，立他作你们的主教；我将我的宣讲之座托付给他——他自始至终与我同行，因此听了我所有的讲道——总而言之，他分享了我所有的试炼，已在信德中显为坚固；我见他比众人更虔诚、仁爱、纯洁、博学、贞洁、善良、正直、宽厚，并慷慨地背负着一些望教者的忘恩负义。因此，我授予他束缚与释放的权力，使凡他在世上所定夺的，在天上也必被决定。因为他应束缚当束缚的，释放当释放的，深知教会的规矩。因此你们要听从他，因为你们知道，那使真理的主持者难过的人，便是得罪\Person{基督}，冒犯圣父。因此他不得存活；所以为主持者的人应持有医师的位置，而不应珍视无理性野兽的愤怒。}\par

% structure: p0024 source=h3 target=subsubsection reason=relative-heading-rank; base=h2

\subsubsection{第三章 不愿就任主教}

当他这样说时，我向他下跪，恳求他，辞让这荣誉和席位的权柄。但他回答说：\Quote{关于此事，不要问我；因为我认为这样是好的，尤其在你辞让之时。因为这席位不需要一个傲慢、贪求占据它的人，而需要一个行为虔诚、精通\Emph{天主}的言语的人。但是，如果你能给我指出一个\Emph{比你更好}的人，一个陪我旅行更多、听过我更多讲论、更好学习了教会规章的人，我就不会强迫你违背你的意愿去行善。但你不能向我指出你的更好者，因为你是我所拯救的群众中精选的初果。然而，再想一想：如果你因害怕罪的危险而不接受教会的管理，那么你肯定在犯罪更多，当你能够帮助那些仿佛在海上处于危险中的虔诚人却不愿如此做时，你只顾及自己的利益，而不顾众人的共同益处。但是，你完全明白你应当承担这危险，而我为众人的帮助不断要求你。因此，你越早同意，你就能越早解除我的焦虑。}\par

% structure: p0026 source=h3 target=subsubsection reason=relative-heading-rank; base=h2

\subsubsection{第四章　赏报的酬报}

\Quote{但是，我自己，哦，克莱孟，也知道从那无知的群众之中所指定给你的忧苦、焦虑、危险和责难；你将能高贵地承受这些，仰望天主所赐予你忍耐的大赏报。但也请公平地和我一同思考：基督什么时候需要你的帮助？是现在，当那邪恶者向祂的新娘宣战的时候；还是将来，当祂胜利地统治、不再需要进一步帮助的时候？难道对任何一个稍微有理解力的人，这不明显是现在吗？因此，要怀着完全的好意，现在在紧迫的时刻急忙站在这位良善之王的战线上，祂的特性是在胜利后赐予巨大的赏报。因此，要欣然接受监督的职务；而且更是及时，因为你从我学习了教会的管理，为了那投靠我们的弟兄们的安全。}\par

% structure: p0028 source=h3 target=subsubsection reason=relative-heading-rank; base=h2

\subsubsection{第五章　一项训诫}

\Quote{然而，我愿在众人面前提醒你，为了众人的缘故，有关管理的事宜。你应当生活无咎，以最大的热诚摆脱一切生活的挂虑，既不做担保人，也不做律师，更不卷入任何其他世俗的事务。因为基督不愿指派你作世俗事务的法官、仲裁者，或现世生活的调停人，免得你因被局限在人的现世挂虑中，而无法由真理之道将人中的善者与恶者分别出来。但要让门徒们彼此履行这些职务，而不是把你从能拯救的讲论中拖开。因为正如同你承担世俗的挂虑，并疏忽你所受命去做的事是邪恶的，同样，每个平信徒若不彼此扶持，即使在世俗的需要上，也是罪。如果所有人都不懂得采取措施使你免于那些你应当免于的挂虑，就让他们从执事那里学习；这样你就能常常拥有教会的挂虑，既为妥善管理，也为持守真理之言。}\par

% structure: p0030 source=h3 target=subsubsection reason=relative-heading-rank; base=h2

\subsubsection{第六章　主教的职责}

\Quote{现在，如果你被世俗挂虑所占据，你就会欺骗你自己和你的听众。因为由于被事务占据而不能指出有益之事，你应当受罚，因为你没有教导有益之事；而他们因没有学习，应当因无知而丧亡。所以，你应当没有挂虑地主持他们，以便适时发出能够拯救他们的言语；让他们听从你，知道真理的使者在地上所捆绑的，在天上也被捆绑；他所释放的，也被释放。但你要捆绑应当捆绑的，释放应当释放的。这些以及类似的事，就是与你作为主席相关的事。}\par

% structure: p0032 source=h3 target=subsubsection reason=relative-heading-rank; base=h2

\subsubsection{第七章　长老的职责}

\Quote{关于长老，要接受这些\Emph{指导}。首先，让他们及早为年轻人成婚，防患于青春情欲的纠缠。但也不可忽略那些已年长者的婚姻，因为情欲甚至在有些老人身上也很强烈。因此，为使淫乱不发生在你们中间，不给你们带来大瘟疫，要谨慎并搜察，免得任何时候奸淫之火在你们中间被秘密点燃。因为奸淫是非常可怕的事，它在刑罚中位居第二，第一是归于那些在谬误中的人，即使他们保持贞洁。因此，你们作为教会的长老，要使基督的净配（我所说的净配是指教会身体）保持贞洁；因为如果她被她尊贵的净配发现是贞洁的，她将获得极大的尊荣；而你们，作为婚礼的宾客，将获得极大的称赞。但如果她被发现有罪，她本身将被驱逐；而你们将受罚，如果她的罪是因你们的疏忽而犯的。}\par

% structure: p0034 source=h3 target=subsubsection reason=relative-heading-rank; base=h2

\subsubsection{第八章　\Quote{待众人行善？}}

\Quote{因此，首要之事，要谨慎持守贞洁；因为在天主眼中，邪淫已被视为可憎之物。但邪淫有多种形式，就如\Person{克莱孟}亲自向你们解释的。第一是奸淫，即男人不应只享有自己的妻子，或女人不应只享有自己的丈夫。若有人保持贞洁，他也能行仁爱，并因此获得永久的怜悯。因为正如奸淫是大恶，仁爱便是至善。因此，要以庄重和怜悯的眼目爱所有弟兄，为孤儿尽父母之职，为寡妇尽丈夫之职，以各样仁慈供给她们所需，为适龄者安排婚姻，为无职业者通过工作提供必要支持；给工匠工作，给无能者施舍。}\par

% structure: p0036 source=h3 target=subsubsection reason=relative-heading-rank; base=h2

\subsubsection{第九章. \Quote{愿弟兄之爱常存}}

\Quote{但我知道，如果你们将爱铭刻在心，就必行这些事；而唯一的适宜方法就是共同分享食物。因此，要留意常常彼此款待，照你们的能力，以免缺失。因为这是行善的原因，而行善则是救恩的原因。所以，你们所有人要将你们的供应品共同献给天主内众弟兄，知道你们付出暂时之物，将获得永恒之物。更要喂养饥饿者，给口渴者喝，给赤身者衣；探望病人；向那些在监狱里的人显出自己，尽力帮助他们，并欢欢喜喜地接待陌生人到你们家中。然而，不再详述，仁爱必教导你们行一切善事，就如憎恨人向那些不愿得救的人暗示恶行。}\par

% structure: p0038 source=h3 target=subsubsection reason=relative-heading-rank; base=h2

\subsubsection{第十章. \Quote{凡是诚实的}}

\Quote{让有争议的弟兄们不要由世俗权威审判；而务必由教会长老和好他们，甘心服从他们。此外，要逃避贪婪，因为它能藉暂时利益的借口剥夺你们永恒的福乐。谨慎保持你们的秤、量器、砝码和买卖之物公平。在受托之事上要忠诚。再者，如果你们心中对从天主而来的审判有不可磨灭的记忆，就必持守这些及类似的事直到终了。因为谁若确信生命终结时必有公义天主的审判——他如今只是长久忍耐和良善，为使善人将来永远享受不可言喻的福乐，而罪人作为恶者将得受永不可言喻的惩罚——他还会犯罪吗？确实，这些事就是如此，若不是真理的先知曾预言并起誓必成就，怀疑倒是合理的。}\par

% structure: p0040 source=h3 target=subsubsection reason=relative-heading-rank; base=h2

\subsubsection{第十一章. 要消除疑惑}

\Quote{因此，作为真先知的门徒，要除去带来恶行的三心二意，热切从事善行。但如果你们中有人对我所说将要发生的事有所怀疑，若他关心自己的灵魂，就应毫不羞愧地承认，并由主教使之满意。但如果他已正确相信，就让他的言行满有把握，如同逃离那定罪的大火，进入天主永恒的善国。}\par

% structure: p0042 source=h3 target=subsubsection reason=relative-heading-rank; base=h2

\subsubsection{第十二章. 执事的职责}

\Quote{此外，教会的执事要明智地巡行，如同主教的眼睛，仔细查问教会每位成员的行为，\Emph{查明}谁将要犯罪，以便由主教用劝诫加以制止，或许他可免于犯罪。他们要约束不守规矩的人，使他们不致停止聚会听道，以便能以真理之言对抗因世俗遭遇和邪恶交往从各方落在心上的忧虑；因为这些忧虑若长久不处理，就成了燃料。他们也要学习谁正受身体疾病的苦，并将他们带到不知情的会众面前，好使他们去探望，按主教的判断供给所需。即使他们这样做是出于无知，也没有什么不对。执事们要关注这些及类似的事。}\par

% structure: p0044 source=h3 target=subsubsection reason=relative-heading-rank; base=h2

\subsubsection{第十三章. 教理导师的职责}

\Quote{让教理导师先受教，然后教导；因为这是关乎人灵魂的工作。因为教授圣言的人必须适应学习者不同的判断。因此，教理导师必须博学、无可指责、经验丰富、经过考验，正如你们知道\Person{克莱孟}就是这样的人，他将在我的教导之后作你们的导师。因为现在要我详述细节是太过分了。然而，如果你们同心合意，就能抵达安息之港，那里有伟大君王的平安之城。}\par

% structure: p0046 source=h3 target=subsubsection reason=relative-heading-rank; base=h2

\subsubsection{第十四章. 教会的船}

\Quote{因为教会的全部事业就像一艘大船，载着来自许多地方、渴望居住于善国王城的人，在暴风中航行。因此，让天主作你们的船长；让舵手比作基督，大副比作主教，水手比作执事，见习船员比作教理导师，众弟兄的群众比作乘客，世界比作大海；逆风比作诱惑、迫害和危险；各样的苦难比作波浪；陆风和飓风比作欺骗者和假先知的言论；海角和礁石比作高位上发出可怕威胁的审判官；两海交汇处和荒野之地比作不合理的人和那些怀疑真理应许的人。让伪善者被视为海盗。此外，把强大的漩涡、塔尔塔罗斯的卡律布狄斯、凶杀般的沉船和致命的毁灭视为罪恶。因此，为要顺风航行，安全抵达所希望之城的港口，要祈祷以求垂听。但祈祷由善行而变得可闻。}\par

% structure: p0048 source=h3 target=subsubsection reason=relative-heading-rank; base=h2

\subsubsection{第十五章. 航行中的事件}

\Quote{因此，让乘客保持安静，坐在自己的位置上，免得因混乱而造成颠簸或倾斜。让船上工作人员注意航程。让执事们不要忽略他们所受托的任何事；让长老们像水手一样，细心地安排每个人所需要的事。让主教像领航员一样，警醒地思量舵手所说的话。让基督，即救主，被爱戴如舵手，并唯独相信祂所说的事；让众人向天主祈求一帆风顺的航行。让那些航行的人预料到各种磨难，如同航行在浩瀚而汹涌的大海上——这世界：有时沮丧、受迫害、分散、饥饿、口渴、赤身露体、被围困；有时又联合、聚集、安息；但也晕船、眩晕、呕吐，即承认罪恶，如同产生疾病的胆汁——我指的是源于苦毒的罪恶，以及由无序的情欲积累的邪恶，通过承认它们，如同通过呕吐，你们得以摆脱疾病，借由审慎达到健康的安全。}\par

% structure: p0050 source=h3 target=subsubsection reason=relative-heading-rank; base=h2

\subsubsection{第十六章　主教的辛劳与赏报}

\Quote{但你们所有人都要知道，主教比你们所有人更辛劳；因为你们每人只受自己的苦，而他却受自己的苦和每个人的苦。因此，哦，克肋孟，你要根据自己的能力，作为每个人的助手主持大局，关心所有人的忧虑。因此我知道，在你承担管理之责时，我不是施予恩惠，而是领受恩惠。但要鼓起勇气，慷慨地承受它，因为知道天主会在你进入安息之港时赏报你——最大的福分，一份不能被夺走的赏报，与你为众人的安全承担更多辛劳相称。因此，即使许多兄弟因你的崇高正义而恨你，他们的恨意绝不能伤害你，但义人的天主的爱将大大地有益于你。因此，要努力摆脱源于不义的赞美，并因正义的管理而获得来自基督的有益赞美。}\par

% structure: p0052 source=h3 target=subsubsection reason=relative-heading-rank; base=h2

\subsubsection{第十七章　人民的职责}

说了这番话以及更多之后，他再次看着群众，说：\Quote{而你们，我亲爱的兄弟和同仆，也要在一切事上顺从真理的引导者，因为知道，那使他忧伤的人，没有接受基督——他的座位已受托于基督；而没有接受基督的人，将被视为轻视了父；因此他将被逐出善的王国。为此，要努力参加所有的集会，免得你们作为逃兵，因使你们的首领灰心而招致罪的指控。因此，你们所有人首先要思想与他有关的事，因为知道那恶者，因你们每个人而更加敌对，单独与他争战。所以你们要努力以爱对待他，彼此仁慈，顺从他，以便他得安慰，你们得救。}\par

% structure: p0054 source=h3 target=subsubsection reason=relative-heading-rank; base=h2

\subsubsection{第十八章　\Quote{如同外邦人和税吏}}

\Quote{但有些事你们自己也应当考虑，因为他因阴谋而不能公开说话。例如：若他对某人有敌意，不要等他说话；不要与那人结盟，而要明智地遵从主教的意愿，与他所敌对的人为敌，不与他所不交谈的人交谈，以便每个人都想得你们为友，而与他修和并得救，听从他的教导。但若有人仍与他所敌对的人为友，与他所不交谈的人交谈，那么他自己也是那些要败坏教会的人之一。因为他虽在身体上与你们同在，但在判断上不与你们同在，他是反对你们的；他比外来的公开敌人更坏，因为表面上的友谊却分散了内部的人。}\par

% structure: p0056 source=h3 target=subsubsection reason=relative-heading-rank; base=h2

\subsubsection{第十九章　克肋孟的任命}

说了这些话之后，他在众人面前按手在我身上，并强迫我坐在他自己的座位上。当我坐下后，他立刻对我说：\Quote{我恳求你，在在场所有兄弟面前，每当我离开此生（我必会离开）时，你要将你从童年起的推理经过，以及从一开始直到现在你与我同行的旅程，以及我在每个城市宣讲的讲道，和\Emph{看见}我的作为，简要地写信给主的兄弟雅各伯。最后你还要如我之前所说的，将我的死讯告诉他。因为当他得知我虔诚地经历了我所应受的苦难时，那件事不会使他过度忧伤。当他得知在我之后，教师的座位不是交给一个无学问的人，或一个对生命之言无知，或不知道教会规矩的人时，他将得到最大的安慰。因为欺骗者的言论毁灭了许多听者的灵魂。}\par

% structure: p0058 source=h3 target=subsubsection reason=relative-heading-rank; base=h2

\subsubsection{第二十章　克肋孟的顺服}

因此，我的主雅各伯，我应许并按所吩咐的去行，已经将他在每个城市中大部分讲道按章节写成书，这些书已经写给您，并由他亲自送去作为凭证；我也就这样寄给您，题写为\Quote{\WorkTitle{克肋孟——伯多禄通俗讲道摘要}}。然而，我即将按所吩咐的开始陈述它们。\par

\SourceRecord{Translated by Thomas Smith.；Ante-Nicene Fathers；Vol. 8.；Edited by Alexander Roberts, James Donaldson, and A. Cleveland Coxe.；Buffalo, NY: Christian Literature Publishing Co.,；1886.；<http://www.newadvent.org/fathers/080800.htm>.}

% structure: p0001 source=h1 target=section reason=page-title

\section{讲道一}

% structure: p0002 source=h2 target=subsection reason=relative-heading-rank; base=h2

\subsection{第一章 童年的疑问}

我\Person{克莱孟}，身为罗马公民，自少年时代起便能操守贞洁，我的心灵自幼便将内心的情欲引向沮丧与痛苦。因为我有一个习惯——不知从何而来——经常思考关于死亡的问题：当我死时，我将不再存在，也没有任何人会记得我，而无限的时间将所有人的一切带入遗忘吗？那时我将没有存在，也不认识那些存在者；既不知晓，也不被知晓，既未曾存在，也将不存在吗？世界是受造的吗？在它受造之前有什么吗？如果它一直存在，它也将继续存在；但如果它是受造的，它也将被解散。在它解散之后，还会再有什么吗？除非是寂静与遗忘？或者将有某物存在，是目前无法想象的。\par

% structure: p0004 source=h2 target=subsection reason=relative-heading-rank; base=h2

\subsection{第二章 从恶中生出善}

我不断地思索这些以及类似的问题——不知从何而起——感到极度悲痛，以致面色苍白，日渐消瘦；而最可怕的是，若有时我想抛开这种无益的思考，痛苦反而更加剧烈；我因此悲伤，却不知我心中有一位美丽的住客，即我的思想，它将为我带来有福永生的因由，这是我后来经验知道的，并感谢天主，万有的主。因为正是这起初折磨我的思想，迫使我寻求并发现真理；然后我怜悯那些我起初因无知而冒然称为有福的人。\par

% structure: p0006 source=h2 target=subsection reason=relative-heading-rank; base=h2

\subsection{第三章 困惑}

于是，从童年起，我深陷于这样的思考中，为了求得确定的知识，我常去哲学家的学园。但我看到的无非是学说的树立与推翻，争辩、求胜、三段论技巧和假设的伎俩；有时一种\Emph{观点}占上风——例如，灵魂是不朽的；有时则说灵魂是会死的。因此，若有时不朽的学说胜出，我就欢喜；若会死的学说胜出，我就悲伤。我更加沮丧，因为我无法满意地确立任何一个学说。然而，我发现讨论中的观点是依据其辩护者而取为真或假，并非按其实质显现。因此，我发现接受与否并不取决于讨论主题的真实性质，而是根据辩护者的能力来证明观点的是非，我因此更加困惑。我从灵魂深处叹息，因为我既无法确立任何东西，也无法摆脱对这些事的思考，尽管如前所述我想摆脱。\par

% structure: p0008 source=h2 target=subsection reason=relative-heading-rank; base=h2

\subsection{第四章 更大的困惑}

我再次生活在疑惑中，自忖道：\Quote{既然事情清楚，我何苦徒劳呢？如果我死后不再存在，那么在我还存在时，不应现在就自寻烦恼。因此我将把我的悲伤留到那\Emph{日子}，那时我不再存在，也就不会感到悲伤。但如果我将存在，现在无故自扰又有何益呢？}紧接著另一个推理攻击我：我说，\Quote{如果我没有虔敬地生活，到那时我岂不是要遭受比现在更糟的痛苦？按照某些哲学家的学说，我不是要被交付给\Person{皮里弗勒格同}和\Person{塔尔塔罗斯}，像\Person{西绪福斯}、\Person{提堤俄斯}、\Person{伊克西翁}或\Person{坦塔罗斯}那样，在\Person{哈得斯}永远受罚吗？}但我又回答说：\Quote{但这些事并不存在。}我又说：\Quote{但若有呢？}因此我说，\Quote{既然事不确定，更稳妥的办法是让我虔敬地生活。}但我为了正义，如何能克制肉体的享乐，而我所期待的不过是一个不确定的希望呢？我既不完全确信什么是取悦天主的正义之事，也不知道灵魂是不朽还是会死。我既找不到任何确凿的学说，也无法停止这样的争论。\par

% structure: p0010 source=h2 target=subsection reason=relative-heading-rank; base=h2

\subsection{第五章 一个决定}

那么，我该做什么呢，除非如此？我要去埃及，与神庙的圣师和先知交好；我要寻找并找到一位术士，用重金贿赂他，让他招魂，这称为招魂术，仿佛我要询问它一些事务。而询问的目的在于了解灵魂是否不朽。但灵魂回答它是不朽的，并不能通过它说话或我听见而给我知识，只能通过它被看见；这样，用我自己的眼睛看见它，我就能从它显现的事实本身获得自足而恰当的确证，知道它存在；从此耳闻的不确定言辞再也不能推翻眼睛所定夺之事。然而，我把这个计划告诉了一位身为哲学家的同伴；他劝我不要冒险，并提出了许多理由。\Quote{因为如果，}他说，\Quote{灵魂不听从术士，你将带着罪恶的良心生活，因为你已经违反了禁止做这些事的法律。但如果灵魂听从了他，那么除了你带着罪恶的良心生活之外，我认为你也不会因这次尝试而增进虔敬之事。因为人们说，天主对于打扰已脱离肉身的灵魂的人是愤怒的。}我听到这话后，确实更不愿做这事了，但我没有放弃原来的计划；我因无法实施而苦恼。\par

% structure: p0012 source=h2 target=subsection reason=relative-heading-rank; base=h2

\subsection{第六章 从\Place{犹太}传来的消息}

为了不以长篇大论向你们讨论这些事，当我忙于这样的思考和行为时，有一个传言，始于春天，在提庇留凯撒统治时期，渐渐四处传播，并如同天主真正的喜讯一样传遍世界，无法将天主的旨意沉寂于无声。因此它到处变得越来越大、越来越响，声称在犹太地有一个人，从春季开始，向犹太人宣讲不可见的天主的国，并说他们中凡愿意改变生活方式的人，都将享受这国。为了使人们相信他所说的话充满神性，他仅凭命令就行了许多奇妙的神迹和征兆，好像从天主那里获得了权能。因为他使聋子听见，瞎子看见，瘸子行走，扶起弯腰的人，驱除各种疾病，赶走所有魔鬼；甚至远望他的麻风病人，也被他治愈离去；死者被带到那里，就复活了；没有他做不到的事。随着时间推移，因为更多人到来，那事——我现在不说传言，而是——真理变得更大更强；因为如今终于在各处有集会，商议并探询那出现者是谁，以及他的目的是什么。\par

% structure: p0014 source=h2 target=subsection reason=relative-heading-rank; base=h2

\subsection{第七章 罗马的福音}

然后同一年，在秋季，有一个人站在公共场所喊道：\Quote{罗马人啊，请听。天主之子已在犹太地降临，向所有愿意的人宣告永生，只要他们按照派遣他的圣父的旨意生活。因此，你们要把生活方式从坏的改为好的，从暂时的事物改为永恒的事物；因为你们要知道，只有一位天主，他在天上，你们不义地居住在他的世界，在他公义的眼前。但如果你们改变，并按照他的旨意生活，那么，重生进入另一个世界，成为永恒，你们将享受他那不可言说的美物。但如果你们不信，你们的灵魂在身体解体后，将被抛入火狱，在那里永远受罚，并为那无益的行为而悔改。因为每个人的悔改期限是今世。}我因此听到这些话时，感到悲伤，因为在如此众多的人群中，听到这样的宣告，却没有一个人说：我要去犹太地，好知道这个告诉我们这些话的人是否讲真话，即天主之子为了美好而永生的希望，揭示派遣他的圣父的旨意，已经来到犹太地。因为他们所说的他宣讲的事不是小事：他断言有些人的灵魂\Emph{本身}是不朽的，将享受永恒的美物；而另一些人的灵魂被抛入不灭的火中，将永远受罚。\par

% structure: p0016 source=h2 target=subsection reason=relative-heading-rank; base=h2

\subsection{第八章 离开罗马}

当我这样谈论别人时，也责备自己说：我为什么责备别人，而我自己却犯了同样粗心大意的罪？但我将急忙前往犹太地，先安排好我的事务。当我这样下定决心后，却出现了长时间的拖延，因为我的世俗事务难以处理。于是，我进一步思考生命的本性，它如何用希望迷惑人，设下圈套陷害那些急于行动的人，是的，以及我被希望抛来抛去被夺去了多少时间，而我们人在忙碌中死去，我便放下所有的事务，赶往波图斯；到达港口后，上了一艘船，被逆风带到了亚历山大里亚，而非犹太地；因天气恶劣滞留那里，与哲学家们交往，告诉他们关于那在罗马出现之人的传闻和言论。他们回答说，他们确实对在罗马出现的人一无所知；但关于那生在犹太地、据传说是天主之子的人，他们从许多来自那里的人那里听说过，并了解到他用一句话所做的一切奇妙之事。\par

% structure: p0018 source=h2 target=subsection reason=relative-heading-rank; base=h2

\subsection{第九章 巴尔纳伯的宣讲}

当我说我希望遇到一个见过他的人时，他们立即带我去见一个人，说：\Quote{这里有一个人，不仅认识他，而且也是他的同乡，一个希伯来人，名叫巴尔纳伯，他自称是他的门徒之一；他住在这附近，并乐意向愿意的人宣告他许诺的条件。}于是我同他们去了；到了以后，我站在他旁边的人群中听他说话；我察觉到他不是用辩证的技巧讲真理，而是简朴地、无准备地陈述他亲眼所见、亲耳所听那显现的天主之子所行所说的事。甚至从周围的人群中，他也引出许多见证人，见证他所讲述的奇迹和言论。\par

% structure: p0020 source=h2 target=subsection reason=relative-heading-rank; base=h2

\subsection{第十章 哲学家的刁难}

\Quote{然而，当群众对他如此朴实地讲述的事情持好感时，哲学家们受其世俗学识的驱使，开始嘲笑他、戏弄他，责备他以过分的狂妄，运用诡辩的庞大武库。但他不理睬他们的闲言碎语，也不参与他们精微的质询，而是毫不尴尬地继续说着他的话。于是他们中有一人问道：为什么蚊子虽然如此微小，有六条腿，还有翅膀；而大象，最大的走兽，却没有翅膀，只有四条腿？但是他，在问题提出之后，重新开始他被中断的讲论，仿佛已经回答了问题，又继续他原来的讲论，只是在每次中断后加上这个开场白：我们只受命告诉你们那派遣我们来者的言语和奇妙作为；我们不以逻辑论证，反而从你们中间提出许多见证人，他们就在旁边，我记得他们的面容，如同活生生的图像。选择相信或不相信这些充分的见证，留给你们。但我不会停止向你们宣告对你们有益的事；因为保持沉默对我来说是损失，而不相信对你们是毁灭。其实我本可以回答你们无聊的问题，如果你们是出于热爱真理而发问。但关于蚊子和大象结构不同的原因，不适合告诉那些不认识万有的天主的人。}\par

% structure: p0022 source=h2 target=subsection reason=relative-heading-rank; base=h2

\subsection{第十一章 \Person{克莱孟}的热忱}

\Quote{天主早已预见到你们不配，所以安排得当，使他的旨意不能被你们接受，这从在场有判断力的人眼中看来是明显的。因为如今已派出传报他旨意的使者，他们不炫耀文法技巧，而以朴素且富有技巧的言语阐述他的旨意，使凡听见的人都能明白所说的话，并且不带任何嫉妒的心态，仿佛不愿将之提供给所有人；你们来到这里，不但不明白什么对你们有益，反而嘲笑真理，自害己身，这真理与蛮族为伍，定你们的罪，而你们因自己的邪恶和它言语的朴实，拒绝接纳它，以免被证实只是喜爱言语，而非喜爱真理和智慧的人。你们这些不能言善辩的人，还要学说话到几时？因为你们的许多言论还不如一句话有价值。那么，如果如他所说的，将有一场审判，你们那希腊群众同心合意，将会说什么呢？\InnerQuote{啊，天主，你为什么没有向我们宣告你的旨意？}如果你们被认为配得回答，岂不将被告知这个？\InnerQuote{我早在创世之前就已认识将要存在的所有性格，按照各人的功过预先对待了每一个人，并没有公开；但我愿意给我所投奔的人充分保证事实如此，并解释为何从一开始，在最初的时代，我没有容许我的旨意被公开宣扬；如今在世界终末，我派了使者宣扬我的旨意，他们却被那些不愿受益、并故意拒绝我友谊的人侮辱和嘲弄。}哦，大错！宣讲者甚至面临生命的危险，而这是由那些被召获得救恩的人所行的。}\par

% structure: p0024 source=h2 target=subsection reason=relative-heading-rank; base=h2

\subsection{第十二章 \Person{克莱孟}对众人的斥责}

\Quote{如果从一开始，不配的人就被召获得救恩，那么对我使者的这种错误对待就会从开始就针对所有人。因为这些人现在所做的错误事，是为了证明我正义的预知：我不曾从一开始就选择将那值得尊敬的言语无益地暴露于公开的轻蔑，这是好的；我决定抑制它，因为它是可敬的，不是对从开始就配得的人——因为我也将它传授给了他们——而是对如同你们所看到的那样不配的人——那些恨我、也不愿爱自己的人。现在，停止嘲笑这个人，要恭敬地聆听他的宣告，或者让听众中任何愿意的人回答。不要像恶狗一样狂吠，以混乱的喧闹震聋那些将要得救之人的耳朵，你们这些不义和憎恨天主的人，颠倒救恩的方法以至不信。你们藐视那位被派来向你们谈论天主神性的人，你们如何能够获得宽恕？你们这样对待一个人，即使他说的不是真理，你们也应当因他对你们的善意而接待他。}\par

% structure: p0026 source=h2 target=subsection reason=relative-heading-rank; base=h2

\subsection{第十三章 \Person{克莱孟}受\Person{巴尔纳伯}教导}

在我讲这些话以及其他类似的话时，人群中激起了很大的骚动；有些人可怜\Person{巴尔纳伯}，同情我；但另一些人毫无理智，对我切齿痛恨。但傍晚已经来临，我拉着\Person{巴尔纳伯}的手，强行把他带回我的住处，尽管他不愿意，我硬让他留在那里，以免有人对他下手。过了几天，他简要地教导了我真道，是几天内所能教导的，之后他说他要赶往犹太地守节，并且希望将来与本国的人来往。\par

% structure: p0028 source=h2 target=subsection reason=relative-heading-rank; base=h2

\subsection{第十四章 \Person{巴尔纳伯}的离开}

但我显然看出他神色不安。因为当我对他说：\Quote{只要你把从那位显现者所听来的话陈述给我，我必以言辞润色，宣讲天主的计划；你若如此，数日内我必与你同航，因我极其渴望前往犹太地，或许终生与你同住。}——他听后答道：\Quote{你若愿查考我们的事，并学习于你有益的事，就立刻与我同航。若不然，我现在就把我的住处和你想见之人的住所指示给你，以便你愿来时能找到我们。因我明日要启程回家。}我见他无法被说服，便随他直到港口；从他那里得知他答应指示给我的居所方向后，对他说：\Quote{若非明日我要收回一笔欠款，我必立即与你同航。但我很快会赶上你。}说完这话，将他托付给船上的掌舵者后，我忧伤地返回，念及他是一位极好且亲爱的朋友。\par

% structure: p0030 source=h2 target=subsection reason=relative-heading-rank; base=h2

\subsection{第十五章 引见伯多禄}

但过了数日，我未能收回全部欠款，为求迅速，我舍弃了余额，视之为阻碍，便也启航往犹太地去，十五日后抵达凯撒勒雅·斯特拉托尼斯。上岸寻找住处时，得知有一位名叫伯多禄的，是那位在犹太地显现并施行奇事异能的先生最受尊敬的门徒，将于次日与基特城的撒玛黎雅人西满进行言语辩论。我闻此讯，便请求指示他的住处；一得知，便站在门前。屋内的人看见我，议论我是谁，从何而来。看，巴尔纳伯出来了；他一看见我，就拥抱我，极其欢喜，流泪不止。他拉着我的手，领我到伯多禄那里，对我说：\Quote{这就是伯多禄，我告诉你他在天主的智慧中为最大者，我常向他提起你。所以请随意进去，因我已如实告诉他你的优点；同时，也向他透露了你的心意，所以他自己也很想见你。因此，当我亲手将你献给他时，我给了他一份厚礼。}他边说边介绍我，说：\Quote{伯多禄，这就是克莱孟。}\par

% structure: p0032 source=h2 target=subsection reason=relative-heading-rank; base=h2

\subsection{第十六章 伯多禄的问候}

那位有福的人一听到我的名字，便跳起来亲吻我；让我坐下后，立即说：\Quote{你接待了真理的宣讲者巴尔纳伯，荣耀了永生天主，此举高尚，你毅然无愧，不惧怕粗野群众的怨恨。你将有福。因你如此恭敬地接待了真理的使者，真理本身也将使你——一个异乡人——成为她自己城邑的居民。如此，你将大大喜乐，因为你如今施予了少许恩惠，即善言的美意。你将承受永恒且不能失去的福分。你不必费心向我详述你的生活方式；因为诚实的巴尔纳伯已将关于你的一切都告诉了我，几乎每日都称赞你。简单告诉你，如对真正的朋友，眼下的事：与我们同行，除非有什么妨碍你，从一城到另一城，直到罗马本地，领受我将宣讲的真理之言。你若愿说什么，请说吧。}\par

% structure: p0034 source=h2 target=subsection reason=relative-heading-rank; base=h2

\subsection{第十七章 所提问题}

于是我从头陈述我的旨趣，以及我如何耗费自己于难题上，以及我初始向你揭露的一切，以免重复同样的事。然后我说：\Quote{我准备好与你同行；因我不知为何，欣然愿意如此。但我愿先确知真理，好知道灵魂是有死的还是永生的；若它永恒，是否要为自己在此所行之事受审判。此外，是否有任何正义且中悦天主的事；世界是否受造，为何目的受造；是否将解体；若将解体，是变得更好，还是归于无有。}不一一细述，我说我愿意学习这些事及由此而生的事。他答道：\Quote{克莱孟，我将很快传授你关于存在之物的知识；即便现在，请听。}\par

% structure: p0036 source=h2 target=subsection reason=relative-heading-rank; base=h2

\subsection{第十八章 无知的原因}

\Quote{天主的旨意曾以多种方式被遮蔽。首先，有邪恶的教导、邪恶的交往、可怕的社交、不当的言谈、错误的偏见。由此产生错谬，然后是胆大妄为、不信、淫乱、贪婪、虚荣；以及其他成千上万的邪恶，犹如浓烟充满房屋一般充满了世界，遮蔽了居住在世上的人们的眼目，不容他们举目从天主所赐的描绘中认识创造主，并知道何事中悦祂。因此爱好真理者理应从内心呼求援助，以爱真理的理性呼唤，使那住在充满烟雾的房屋中的人可以走近并开门，让外界的阳光进入屋内，驱散内部的烟火。}\par

% structure: p0038 source=h2 target=subsection reason=relative-heading-rank; base=h2

\subsection{第十九章 真先知}

\Quote{现在，我将那位作为助手的人称为真先知；唯独祂能够光照人的灵魂，使我们能用双眼看见永生救恩的道路。否则这是不可能的，正如你自己也知道，因为你刚才说过，每一种教义都被建立又被推翻，同一教义被认为是真或假，全凭倡导者的能力；因此教义不是按其本来面目显现，而是从那些倡导者那里取得真或假的外貌。为此，宗教的全部事务需要一位真先知，好让他按事物的本来面目告诉我们，以及我们必须如何相信一切。因此，首先需要用每个先知性的标记来检验先知，一旦确定他是真的，此后就要在一切事上相信他，而不是对他的各别言论作判断，而要接受它们为确定的，这些言论虽然是通过表面的信德被接受，却是通过确实的判断。因为通过我们最初的证明和严格的全面审查，一切事物都是以正确的理性接受的。因此，在一切事之前，必须寻求真先知，因为没有祂，人不可能得到任何确定性。}\par

% structure: p0040 source=h2 target=subsection reason=relative-heading-rank; base=h2

\subsection{第二十章 伯多禄对克肋孟的满意}

同时，他藉着向我解释祂是谁、如何找到祂，并将祂真正可寻得的事实摆在我面前，使我满意，表明真理藉着先知的论述比眼所见之物更清晰地传入耳中；以致我惊异，奇怪为何没有人看见那些人人寻求的东西，尽管它们就在眼前。然而，我遵他的命令写了这篇关于先知的论述后，他使人将那卷书从凯撒勒雅·斯特拉托尼斯寄给你，说他受你委托每年寄送他的论述和事迹。因此，在他首先只论及真理的先知的第一天，他就在各方面坚定了我；然后他这样说：\Quote{从此以后，你要留意我与另一边的人之间的讨论；即使我处于不利地位，我也不怕你会怀疑已经传授给你的真理，因为我深知似乎是我被击败，并非先知传授给我们的教义。不过，我希望在与有悟性的人的探讨中不会处于不利地位——我指的是那些热爱真理的人，他们能辨别什么是似是而非、人为雕琢、悦耳动听的论述，什么是未经修饰、纯朴、只信赖藉着它们\Emph{所传递的}真理的论述。}\par

% structure: p0042 source=h2 target=subsection reason=relative-heading-rank; base=h2

\subsection{第二十一章 不可动摇的信念}

当他说完这些话，我回答说：\Quote{现在，我感谢天主；因为正如我希望被说服，祂已恩赐了我。然而，就我而言，你可以完全放心，我绝不会怀疑；甚至，如果你自己有一天想要使我离开先知的教义，你也不可能做到，因为我深知我所领受的。并且不要以为我承诺我永不会怀疑是一件大事；因为我自己，或任何听过你关于先知的论述的人，都不可能怀疑真教义，因为我们首先听到了并了解了先知宣告的真理。因此，要坚信天主所愿的教义；因为一切邪恶的计谋都被战胜了。因为对于预言，无论是辞令的艺术、诡辩的把戏、三段论还是任何其他手段，都不能占任何优势；也就是说，如果听过真先知的人确实渴求真理，并且不只是为了真理的假象而注意其他事物。所以，我的主伯多禄，不要不安，好像你已将最大的益处献给了一个无知的人；因为你将之献给了一个察觉这一恩惠的人，一个不能被他所托付的真理引诱离开的人。因为我知道，这是那些人们希望迅速领受、而不愿缓慢达成的事物之一。因此，我知道我不应该因为\Emph{我获得它的}快速而轻视所托付给我的、无与伦比的、唯一安全的事物。}\par

% structure: p0044 source=h2 target=subsection reason=relative-heading-rank; base=h2

\subsection{第二十二章 感恩}

当我这样说完了，伯多禄说：\Quote{我感谢天主，既为你的得救，也为我的满意。因为我确实很高兴知道你看懂了预言有多么伟大。既然，如你所说，如果我自己有一天——愿天主禁止——想要让你改信另一种教义，我将无法说服你，那么从明天起你就开始参加我与对手的讨论吧。明天我将与西满术士有一场讨论。}说完这话，他自己私下用了餐，然后命令我也用餐；祝福了食物，吃饱后谢了恩，然后向我解释了这件事，继续说道：\Quote{愿天主使你凡事都能像我，并且领洗后与我同席共餐。}说完这话，他吩咐我去休息；因为此时我的身体确实需要睡眠。\par

\SourceRecord{Translated by Thomas Smith.；Ante-Nicene Fathers；Vol. 8.；Edited by Alexander Roberts, James Donaldson, and A. Cleveland Coxe.；Buffalo, NY: Christian Literature Publishing Co.,；1886.；<http://www.newadvent.org/fathers/080801.htm>.}

% structure: p0001 source=h1 target=section reason=page-title

\section{第二篇讲道}

% structure: p0002 source=h2 target=subsection reason=relative-heading-rank; base=h2

\subsection{第一章 伯多禄的随从}

\Emph{因此}，次日，我\Person{克莱孟}，黎明前从睡眠中醒来，得知\Person{伯多禄}已经起身，正在与他的随从谈论关于天主的敬拜（他们共有十六人，我认为最好列出他们的名字，正如我后来所了解到的，好叫你们也知道他们是谁。其中第一个是\Person{匝凯}，他曾是税吏，和他的兄弟\Person{索福尼亚}；\Person{若瑟}和他的义兄弟\Person{米该亚}；还有孪生的\Person{多默}和\Person{厄里厄则尔}；以及司铎\Person{埃乃阿}和\Person{拉匝禄}；此外还有\Person{厄里叟}，和\Person{匝弗尔}之子\Person{本雅明}；以及建筑者\Person{鲁彼禄}和\Person{匝加利亚}；还有雅米尼亚人\Person{阿纳尼雅}和\Person{哈盖}；以及朋友\Person{尼则达}和\Person{阿桂拉}），——于是我进去向他问安，并应他的请求坐下。\par

% structure: p0004 source=h2 target=subsection reason=relative-heading-rank; base=h2

\subsection{第二章 健康的心灵寓于健康的身体}

他中断了正在进行的讨论，向我保证，并道歉为何没有唤醒我让我听他讲论，理由是航行的不适。他希望这种不适消散，所以容我睡觉。\Quote{因为，}他说，\Quote{每当灵魂被某种身体上的需要所分心时，它就不能很好地领会呈现给它的教导。因此，我不愿意与那些因灾祸而极度悲伤、或过度愤怒、或陷入爱情狂热、或遭受身体疲惫、或为生活忧虑所苦、或被其他任何痛苦折磨的人交谈；他们的灵魂，如我所说，因沮丧而同情受苦的身体，也因此占据了自己的理智。}\par

% structure: p0006 source=h2 target=subsection reason=relative-heading-rank; base=h2

\subsection{第三章 有备无患}

\Quote{并且不要有人说：\InnerQuote{那么，向处于困境的人提供安慰和劝告不恰当吗？}对此我回答说：如果确实有人能够，就让他提供；但如果不，就让他等待时机。因为我知道凡事都有其适当的季节。因此，用加强灵魂的话语提前应对邪恶是恰当的；这样，一旦任何邪恶临到他们，心智预先用正确的论点武装起来，就能承受所遭遇的：因为那时心智在斗争的关键时刻要知道向那位以良好忠告帮助他的人求助。}\par

% structure: p0008 source=h2 target=subsection reason=relative-heading-rank; base=h2

\subsection{第四章 一项请求}

\Quote{不过，我听说，哦\Person{克莱孟}，在\Place{亚历山大里亚}，\Person{巴尔纳伯}曾完美地向你阐释了关于预言的道理。是不是这样？}我回答说：\Quote{是的，而且非常好。}然后\Person{伯多禄}：\Quote{因此，现在没有必要用你已经知道的教导来占用时间，这时间可以用于其他你不知道的教导。}然后我说：\Quote{你说得对，哦\Person{伯多禄}。但请恩准我这个决心一直跟随你的人，持续地向一个欢喜的听众阐释这位先知的道理。因为，正如我从\Person{巴尔纳伯}所学到的，离开了他，就不可能认识真理。}\par

% structure: p0010 source=h2 target=subsection reason=relative-heading-rank; base=h2

\subsection{第五章 认识真先知的卓越性}

\Quote{就你而言，纠正的过程已经结束了，因为你知道那无误预言的伟大，没有它，没有人能接受那最有益的事。因为在存在或可能存在的事物中，有许多和各种祝福，其中最蒙福的一切——无论是永生，还是永远的健康，或完美的理解，或光明，或喜乐，或不朽，或任何其他存在的或在事物本性中能至善的事物——若不首先认识事物的本来面目，就无法拥有；而这种知识除非首先认识真理的先知，否则无法获得。}\par

% structure: p0012 source=h2 target=subsection reason=relative-heading-rank; base=h2

\subsection{第六章 真先知}

\Quote{现在，真理的先知是那位始终知道一切事物的——已过的事是它们曾经的样，现在的事是它们现在的样，未来的事是它们将来的样；无罪的、仁慈的、唯独受托宣告真理的。读一读，你们就会发现那些自以为靠自己找到了真理的人，是\Emph{受了欺骗}。因为宣告真理是先知的特性，正如带来白昼是太阳的特性。因此，所有渴望认识真理，却没有幸运从他那里学到的人，都没有找到，而是寻求着死去。因为一个从自己的无知中寻求真理的人，怎能找到真理？即使他找到了，他也不认识，把它当作不是而错过。他也不能从另一个类似地许诺从他无知中给予知识的人那里获得真理；除非是道德知识和诸如此类的事，这些可以通过理性认识，理性给每个人这样的知识：他不应该亏待他人，因为他不希望\Emph{自己}受亏待。}\par

% structure: p0014 source=h2 target=subsection reason=relative-heading-rank; base=h2

\subsection{第七章 独自寻求真理无益}

\Quote{因此，所有曾经寻求真理、相信自己能够找到的人，都陷入了陷阱。这就是希腊的哲学家和较开明的外邦人所遭受的。因为，他们专注于可见的事物，对不可见的事物凭猜想作出判断，认为任何时刻呈现给他们的\Emph{那样}的事就是真理。就像那些知道真理的人一样，他们仍在寻求真理，拒绝一些呈现给他们的假设，抓住另一些，好像他们知道，其实他们不知道，什么是真什么是假。他们关于真理教条化，即使是那些寻求真理的人，他们不知道寻求真理的人不能从自己的流浪中学到真理。因为，如我所说，即使真理站在他旁边，他也认不出，因为他不认识她。}\par

% structure: p0016 source=h2 target=subsection reason=relative-heading-rank; base=h2

\subsection{第八章 真理的检验}

\Quote{而那说服每一个想从自身学习之人的，绝非真实者，而是令人愉悦者。因此，既然此人以此为悦，彼人以彼为悦，便有人将此当作真理，有人将彼当作真理。但真理是先知所认可的，而非各人所喜爱的。因为若那令人愉悦者即为真实，则那唯一者将成为众多，这是不可能的。故此，希腊的语文家——而非哲学家——凭猜测探究事物，凭靠臆想创立了众多且各异的教条，以为假设间恰当的连贯便是真理，却不知当他们为自己设下虚妄的开端时，其结论便与开端一致了。}\par

% structure: p0018 source=h2 target=subsection reason=relative-heading-rank; base=h2

\subsection{第九章 \Quote{世上的软弱之物。}}

\Quote{因此，人应放下其余一切，唯独将自己交托于真理的先知。我们全都能判断他是否为先知，纵然我们完全不学无术、是诡辩的新手、不谙几何、未习音乐。因为天主关怀众人，使关于祂自身的发现对众人更为容易，为使外邦人并非无力，希腊人也并非不能寻见祂。故此，关于祂的发现是容易的；情形如下：—}\par

% structure: p0020 source=h2 target=subsection reason=relative-heading-rank; base=h2

\subsection{第十章 先知的考验}

\Quote{若他是一位先知，且能知晓世界如何造成、其中的事物、以及直至终结将要发生的事，若他曾预言我们任何事，而我们确知它已完全应验，我们便容易相信那祂\Emph{说}将要发生的事，从已发生的事推断；我们相信祂，我所说，不仅因为祂知道，更因为祂预知。那么，即使理解力极其有限的人，岂不也明白：关于天主所喜悦的事，我们应当超越众人地相信那位超越众人所知者，即使祂未曾学习？因此，若有人不愿将认识真理的能力赋予这样一位——我指因祂内圣神的神性而拥有预知的那位——反而将认识的能力赋予其他任何人，难道他不是缺乏理解吗？因为他将那连先知也不愿赋予的认识能力，赋予了一位非先知者？}\par

% structure: p0022 source=h2 target=subsection reason=relative-heading-rank; base=h2

\subsection{第十一章 无知、知识、预知}

\Quote{因此，首要之事，我们当以完全的判断，借着先知的应许来考验先知；一旦确知祂是先知，我们便当无疑地遵循祂教导的其他话语；并凭着对盼望之事的信心，我们当依据最初的判断行事，深知那位告诉我们这些事者，其本性不会说谎。因此，若祂后来所说的任何事在我们看来似乎说得不好，须知并非祂说错了，而是我们理解不当。因为无知不能判断知识，同样，知识也不能真正判断预知；但预知为无知者提供知识。}\par

% structure: p0024 source=h2 target=subsection reason=relative-heading-rank; base=h2

\subsection{第十二章 真先知的教导}

\Quote{因此，亲爱的\Person{克莱孟}，若你愿知关于天主的事，你须从祂独一者学习，因为唯独祂知晓真理。若其他人知道任何事，也是从祂或祂的门徒领受的。这便是祂的教导和真实的宣告：唯有一位天主，世界是祂的化工；祂全然公义，必定在某一时刻按各人的行为报应各人。}\par

% structure: p0026 source=h2 target=subsection reason=relative-heading-rank; base=h2

\subsection{第十三章 未来的赏罚}

\Quote{因为凡说天主按其本性是公义的，必然也相信人的灵魂是不死的：否则祂的公义何在？有些敬虔生活的人曾遭恶待，有时被暴虐地处死，而另一些全然不虔、纵情奢侈的人却如常人一般死去。既然完美矛盾的天主是善的，也是公义的，除非灵魂在脱离肉体后是不死的，否则祂不能被称为公义；这样，恶人在死后，因在此世已享福乐，便在地狱中为自己的罪受罚；而善人在此世已因自己的罪受罚，则在那时，如同在义人的怀抱中，成为善事的承继者。既然天主是公义的，我们便完全清楚：有审判，且灵魂是不死的。}\par

% structure: p0028 source=h2 target=subsection reason=relative-heading-rank; base=h2

\subsection{第十四章 正义与不义}

\Quote{但若有人，依照这位撒玛黎雅人\Person{西满}的意见，不肯承认天主是公义的，那么人还能将正义或其可能性归于谁呢？因为若万有的根源没有正义，就必须认为在仿佛是果实的人性中也不可能找到正义。若在人身上能找到正义，那在天主身上岂不更多？但若正义无处可寻，无论在天主或人身上，那么不义也不能存在。然而确有正义这回事，因为不义之名取自正义的存在；所谓不义，乃是当正义与其比较而被发现与其相反时所得的名称。}\par

% structure: p0030 source=h2 target=subsection reason=relative-heading-rank; base=h2

\subsection{第十五章 成对}

\Quote{因此，天主教导人关于万有实体的真理，祂本是独一的，将一切本原区分为成对和相对者，祂从起初就是独一唯一的天主，创造了天和地、昼和夜、光和火、太阳和月亮、生命和死亡。但在这些之中，祂只使人能自我控制，有成为义或不义的适应性。祂也为人类设定了各种组合的形态，首先将微小的事物放在他们面前，然后将伟大的事物放在后面，例如世界和永恒。但现在的世界是暂时的；将来的世界是永恒的。首先是无知，然后是知识。祂也如此安排了预言的领导者。因为，既然现在的世界是女性的，如同母亲生育她儿女的灵魂，但将来的世界是男性的，如同父亲\Emph{从母亲那里}接收他的儿女，因此在这世界里就来了一连串的先知，作为将来世界的儿子，并且有对人的知识。如果虔敬的人明白了这个奥秘，他们就不会误入歧途；但即使是现在，他们也应该知道那现在迷惑所有人的西满，是错误和欺骗的同谋。现在，预言规则的教义如下。}\par

% structure: p0032 source=h2 target=subsection reason=relative-heading-rank; base=h2

\subsection{第十六章  人的道路与天主的道路相反}

\Quote{如同在起初，独一的天主，像右手和左手，先造了天，后造了地，祂也按次序设立了所有的组合；但对于人，祂不再这样做，而是变化了所有的组合。因为，既然从祂那里，更伟大的事物先来，较劣等的事物后来，我们在人身上却发现相反的情形——先是更坏的，后是更优的。因此，从按天主的肖像受造的亚当，首先生出不义的加音，然后生出义人亚伯尔。再次，从你们当中称为德乌卡利翁的那一位，两种灵体被派出，即不洁的和纯洁的，先是黑色的乌鸦，然后是白色的鸽子。从亚巴郎，我们民族的族长，也生出两个首生者——先是依市玛耳，后是依撒格，蒙天主降福。同样，从依撒格自己，又生出两个——厄撒乌亵渎者，和雅各伯虔诚者。所以，在出生上为首，如同世界上的首生者，是大司祭\Emph{亚郎}，然后是立法者\Emph{梅瑟}。}\par

% structure: p0034 source=h2 target=subsection reason=relative-heading-rank; base=h2

\subsection{第十七章  先是更坏的，后是更好的}

\Quote{同样，关于厄里亚的组合，本应来到，却自愿延到另一时间，因为决定以后方便时享受它。因此，在妇人所生的人中，他先来；在人间子中，他后来。按此顺序，有可能看出西满属于哪个系列——他在我之前来到外邦人中，而我在他之后来到，如同光明临到黑暗，知识临到无知，医治临到疾病。这样，正如真先知告诉我们的，必须先有一个假先知从某个欺骗者而来；然后，同样，在圣所被移去之后，真福音必须被秘密传播，以纠正将要出现的异端。此后，在末期，假基督必须先来，然后我们的耶稣必被显示确实是基督；此后，永恒的光明升起，一切黑暗之物必消失。}\par

% structure: p0036 source=h2 target=subsection reason=relative-heading-rank; base=h2

\subsection{第十八章  关于西满术士的错误}

\Quote{既然，如我所说，有些人不知道组合的规则，因此他们不知道谁是我的先行者西满。因为，如果他被知道，就不会被相信；但现在，他不被知道，反而被不当相信；尽管他的行为是憎恨者的行为，他却被爱；尽管他是敌人，却被当作朋友接待；尽管他是死亡，却被当作救主渴求；尽管他是火，却被当作光明看待；尽管他是欺骗者，却被当作真理的言者相信。}\par

然后我克莱孟听见这话，说：\Quote{我请问，谁是这样的一位欺骗者？我愿意知道。} 伯多禄接着说：\Quote{若你愿意学习，你能够从那些人也使我准确知道关于他一切事情的人那里得知。}\par

% structure: p0039 source=h2 target=subsection reason=relative-heading-rank; base=h2

\subsection{第十九章  尤斯塔，一名皈依者}

\Quote{在我们中间有一位尤斯塔，叙利腓尼基人，民族是客纳罕人，她的女儿被重病压迫。她来到我们的主面前，喊叫并恳求祂医治她的女儿。但祂，也被我们请求，说：\InnerQuote{医治外邦人是不合法的，他们因使用各种肉食和习俗而像狗，而天国的筵席已给了以色列子民。} 但她听见这话，并恳求像狗一样分得从这桌上掉下的碎屑，她以改变自己，像天国子民一样生活，获得了她女儿按她所求的医治。因为她是外邦人，若停留在同样的生活道路上，祂不会医治她，因医治她作为外邦人是不合法的。}\par

% structure: p0041 source=h2 target=subsection reason=relative-heading-rank; base=h2

\subsection{第二十章  因信仰而被离婚}

\Quote{因此，她采用了按法律的生活，与被治好的女儿一起被丈夫赶出家门，因为他的情感反对我们。但她忠于自己的承诺，且家境富裕，自己守寡，却把女儿嫁给了一个依恋于真信仰的穷人。并且，为了女儿而禁绝结婚，她买了两个男孩并教育他们，把他们当作儿子。他们从小与西满术士一起受教育，学会了关于他的一切事。因为他们的友谊如此之深，以致于他们在他愿意与他们联合的一切事上与他联合。}\par

% structure: p0043 source=h2 target=subsection reason=relative-heading-rank; base=h2

\subsection{第二十一章  尤斯塔的养子，西满的同伙}

\Quote{这些人遇到在此寄居的匝凯，从他那里接受了真理之言，并对自己先前的革新悔改，立即揭发西满在一切事上与他同谋。我一到此寄居，他们就与他们的养母一同来到我这里，由他（\Emph{匝凯}）引荐给我，从此他们一直与我同在，享受真理的教导。}伯多禄说完这话，就派人去叫他们，并吩咐他们应准确地向\Person{我}报告关于西满的一切。于是他们呼求天主为证，决不在任何事上作假，便开始了陈述。\par

% structure: p0045 source=h2 target=subsection reason=relative-heading-rank; base=h2

\subsection{第二十二章 西满的教义}

阿奎拉首先开始这样讲话：\Quote{最亲爱的兄弟，请听，好让你准确认识这个人的一切：他是谁，是什么人，从哪里来；他所行的是什么事，以及他如何和为何行这些事。这个西满是安多尼和辣黑耳的儿子，是撒玛黎雅人，来自吉塔村，该村离城有六舍尼远。他在亚历山大城受了严格的训练，精通魔术，野心勃勃，希望被尊为某种最高的权能，甚至比创造世界的天主还大。有时他暗示自己是基督，自称为\OriginalTerm{站立者}{Standing One}。他用这个称呼暗示他要永远站立，没有任何腐败的原因使他的身体倒下。他既不说创造世界的是至高者，也不相信死人会复活。他拒绝耶路撒冷，以革黎斤山代替它。他不用我们的基督，而自称基督。他凭自己的臆测解释法律；他确实说将有审判，但他并不期待它。因为如果他相信自己将被天主审判，他就不会敢对天主本身不敬。因此，有些人不了解他，以宗教为外衣，败坏真理之事，而那些忠于相信他所说某种意义上的希望和审判的人，便遭到毁灭。}\par

% structure: p0047 source=h2 target=subsection reason=relative-heading-rank; base=h2

\subsection{第二十三章 西满是洗礼者若翰的门徒}

\Quote{但他开始涉及宗教教义，情形如下：有一位若翰，是\OriginalTerm{每日施洗者}{day-baptist}，按照组合的方式，他是我们主耶稣基督的前驱；正如主有十二位宗徒，代表太阳的十二个月份，同样，这位若翰也有三十名首领，满足月亮的月份计数；其中有一位名叫海伦娜的妇女，使这一点也不乏救恩计划的意义。因为女人是半个男人，补足了\OriginalTerm{三十之数}{triacontad}的不完全；正如月亮，它的运行并不完成月份的完整周期。但这三十人中，若翰最看重、最尊崇的是西满；他之所以在若翰死后没有成为首领，原因如下：——}\par

% structure: p0049 source=h2 target=subsection reason=relative-heading-rank; base=h2

\subsection{第二十四章 竞选策略}

\Quote{他因练习魔术而前往埃及，若翰被杀后，多西托斯渴望领导权，便谎称西满已死，并继任其位。但西满不久后返回，并坚称此位属于自己，当他遇到多西托斯时，他没有要求此位，因为他知道一个人若获得了超出预期之权，便无法被移除。因此，他假装友好，暂时屈居多西托斯之下，位居第二。但几天后，他加入三十位同门之中，开始诋毁多西托斯，说他未能正确传授训诲。他说他这样说，并非因他不愿正确传授，而是出于无知。有一次，多西托斯察觉到西满这种巧妙的指控正在削弱众人对他的看法，以至于他们不认为他是\OriginalTerm{站立者}{Standing One}，他便怒气冲冲地来到惯常的聚会地点，找到西满，用棍子打他。但棍子似乎穿过了西满的身体，如同穿过烟雾一般。多西托斯因此惊慌，对他说：\InnerQuote{如果你是站立者，我也要朝拜你。}然后西满说他是。多西托斯知道他自己不是站立者，便俯伏朝拜；他与二十九位首领联合，将西满提升到他原本的尊位；这样，没过几天，当西满站立时，多西托斯自己却倒下跌死了。}\par

% structure: p0051 source=h2 target=subsection reason=relative-heading-rank; base=h2

\subsection{第二十五章 西满的欺骗}

\Quote{但西满正与海伦娜四处游荡，甚至直到现在，如你所见，仍在煽动百姓。他说他将这个海伦娜从最高天带到了这个世界；她是王后，作为\OriginalTerm{全生者}{all-bearing being}，也是智慧；为了她，希腊人和蛮族争战，眼前只是真理的一个形象；因为那真正的真理当时与首位神同在。此外，他通过巧妙地解释某些由希腊神话编造出来的类似事物，欺骗了许多人；特别是他行了许多显著的奇事，以致如果我们不知道他是用魔术行这些事，我们自己也会被骗。但当初我们是他的同工，只要他行这些事而不损害宗教事业；如今他疯狂地开始试图欺骗那些虔诚的人，我们便离开了他。}\par

% structure: p0053 source=h2 target=subsection reason=relative-heading-rank; base=h2

\subsection{第二十六章 他的邪恶}

\Quote{因为他甚至开始犯谋杀罪，正如他自己向我们透露的那样，像朋友对朋友一样，说通过可怕的咒语将孩子的灵魂与身体分离，作为他展示任何他喜欢之事的助手，并画了那男孩的肖像，将它放在他睡觉的内室里，说他曾用神圣的技艺用空气造出了那男孩，并画了他的肖像，又将他送回了空气。他还解释说他如此行了这事。他说人的第一个灵魂，变成了热的性质，吸引并吸入了周围的空气，如同葫芦的形状；然后当它在灵的形态内时，将它变成了水；他说他将其中的空气变成了血的性质，这血因灵的稠密而无法倾泻，并使血液凝固成肉；然后肉如此凝固后，他展示了一个不是由土而是由空气造的人。这样，他说服自己能够造出一种新人，他说他逆转了变化，又将他恢复为空气。当他告诉别人这事时，他被相信了；但我们这些在场参加他仪式的人却虔诚地不信。因此我们谴责了他的不敬，并离开了他。}\par

% structure: p0055 source=h2 target=subsection reason=relative-heading-rank; base=h2

\subsection{第二十七章　他的承诺}

当阿奎拉说了这话后，他的兄弟尼塞塔斯说：\Quote{有必要，我们的兄弟克莱孟，提一下阿奎拉遗漏的事。因为首先，天主作证，我们没有参与任何不敬的工作，而是看着他施行；只要他做无害的事并展示它们，我们也高兴。但是，当为了欺骗虔诚的人，他说他借神性行出了那些由魔术做成的事时，我们就不再容忍他，尽管他给了我们许多承诺，特别是我们的雕像应被认为配得上在庙宇中有一席之地，我们应被认为是神，被群众崇拜，被国王尊荣，被认为配得上公共荣誉，并拥有无限的财富。}\par

% structure: p0057 source=h2 target=subsection reason=relative-heading-rank; base=h2

\subsection{第二十八章　徒劳的劝告}

\Quote{这些事，以及比这些更大的事，他许诺给我们，条件只是我们与他交往，并对他的行径的邪恶保持沉默，好使他的欺骗计划得以成功。但我们仍不同意，反而劝告他停止这样的疯狂，对他说：\InnerQuote{我们，西满，记念我们自幼以来的友谊，出于对你的爱，给你好的劝告。停止这种企图。你不能成为神。要敬畏那真是天主的。要知道你是人，你生命的时间是短暂的；即使你获得巨大的财富，甚至成为国王，你短暂的生命中也只能享受很少的东西，而邪恶得来的东西很快消失，并为冒险者招致永罚。因此我们劝告你敬畏天主，每个人的灵魂都必须为他在这里所行的行为受审判。}}\par

% structure: p0059 source=h2 target=subsection reason=relative-heading-rank; base=h2

\subsection{第二十九章　灵魂的不朽}

\Quote{他听了这话就笑了；当我们问他为什么因我们给他好劝告而笑时，他回答说：\InnerQuote{我笑你们愚蠢的假设，因为你们相信人的灵魂是不朽的。}然后我说：\InnerQuote{西满，我们不奇怪你试图欺骗我们，但我们惊讶于你甚至欺骗自己。告诉我，西满，即使没有别人完全确信灵魂是不朽的，至少你和我们应该确信：你曾将一个灵魂从人体中分离出来，并与它交谈，向它发命令；我们则在场，听了你的命令，并清楚地见证了所命令的事成就了。}然后西满说：\InnerQuote{我知道你们的意思；但你们对于你们所辩论的事一无所知。}尼塞塔斯说：\InnerQuote{如果你知道，就说出来；但如果你不知道，不要以为我们可以被你的话欺骗，你说你知道而我们不知道。因为我们并非如此幼稚，以致你能在我们心中播下一种狡猾的怀疑，使我们以为你知道一些不可言说的事，从而你能抓住并控制我们，通过欲望的束缚来约束我们。}}\par

% structure: p0061 source=h2 target=subsection reason=relative-heading-rank; base=h2

\subsection{第三十章　一个论证}

\Quote{然后西满说：\InnerQuote{我知道你们知道我分离了一个灵魂从一个身体中；但我知道你们不知道，服侍我的并不是死者的灵魂，因为它不存在，而是一个冒充灵魂的恶魔在工作。}然后尼塞塔斯说：\InnerQuote{我们一生听过许多难以置信的事，但我们不期望听到比这话更愚蠢的话。因为如果一个恶魔冒充死者的灵魂，那么灵魂有什么用，以至于它要与身体分离？我们不是亲自在场，并听见你从身体中召唤灵魂吗？怎么当召唤一个时，另一个未受召唤的却服从，好像受了惊吓？而且你自己，当我们有时问为什么交流有时停止时，你不是说灵魂在世上完成了它要在身体中度过的时限后，就去了阴府吗？你还补充说，自杀者的灵魂不容易被允许前来，因为它们回到阴府后就被看守着。}}\par

% structure: p0063 source=h2 target=subsection reason=relative-heading-rank; base=h2

\subsection{第三十一章　一个困境}

尼塞塔如此说完，阿基拉自己接着说道：\Quote{我只愿向您请教这一点，西满：所招来的是灵魂还是魔鬼？它惧怕什么，以致不敢藐视咒语？于是西满说：\InnerQuote{它知道若不服从，必受惩罚。}阿基拉说：\InnerQuote{因此，若灵魂被招来，便也有审判。若灵魂是不朽的，便肯定有审判。如您所言，那些因邪恶勾当被招来的，若不服从而受罚，您怎么还敢强迫它们？因为被强迫者因不服从而受罚，您岂不惧怕？您迄今未因自己的行径受罚，并不奇怪，因为审判尚未到来；那时您将为强迫他人所作的恶行承受惩罚，而被强迫者为尊重导致恶行的誓言而行事，当被宽恕。}他听后大怒，威胁要处死我们，如果我们不闭口不提他的行径。}\par

% structure: p0065 source=h2 target=subsection reason=relative-heading-rank; base=h2

\subsection{第三十二章 西满的奇事}

阿基拉如此说完，我克肋孟问道：\Quote{那么，他行的是什么奇事？}他们告诉我，他使雕像行走，在火上翻滚而不被烧，有时飞行；他以石头变饼；他变成蛇；他把自己变成山羊；他有两副面孔；他把自己变成金子；他打开锁着的门；他熔化铁器；在宴席上，他制造各种形状的影像。在他家中，他使菜肴仿佛自行出现侍奉他，不见捧持者。我听了这些话感到惊奇，但许多人作证说，他们曾在场并亲眼看见这些事。\par

% structure: p0067 source=h2 target=subsection reason=relative-heading-rank; base=h2

\subsection{第三十三章 成对论}

这些事如此说完，杰出的伯多禄也亲自发言：\Quote{兄弟们，你们必须认识结合律的真理，凡不偏离它的人便不会误入歧途。因为，如我们所说，我们看到万物皆成对且相反；夜在先，然后有昼；先有无知，然后有知识；先有疾病，然后有医治；同样，错误的事先进入我们的生活，然后真理降临，犹如医生降临疾病。因此，当上天宠爱的民族即将从埃及人的压迫中被救赎时，首先藉着亚郎手中变为蛇的棍杖产生疾病，然后藉着梅瑟的祈祷引入医治。如今，当外邦人即将从偶像迷信中被救赎时，统治他们的邪恶已预先派遣它的同盟，如另一条蛇，就是你们所见的这个西满，他行奇事以惊人欺人，而非行医治的征兆以使人皈依获救。因此，你们也当从所行的奇迹判断行奇迹者：行奇迹者的品性和奇迹的性质。若他行无益的奇迹，便是邪恶的代理人；若他行有益的事，便是良善的领导者。}\par

% structure: p0069 source=h2 target=subsection reason=relative-heading-rank; base=h2

\subsection{第三十四章 无用与仁慈的奇迹}

\Quote{所以，那些是无用的征兆——你们所说的西满所行的。但我说的那些：使雕像行走，在炭火上翻滚，变成龙，变成山羊，在空中飞翔，以及所有此类事，都不是为了医治人，其本性是欺骗许多人。但慈悲的真理所行的奇迹是仁慈的，如你们所听到的主所行的，以及我在祂之后藉祈祷所行的；你们中许多人曾目睹这些事：有人摆脱各种疾病，有人脱离魔鬼，有人手复原，有人脚复原，有人重见光明，有人恢复听觉，以及其他一切怀着仁慈之心所能行的事。}\par

% structure: p0071 source=h2 target=subsection reason=relative-heading-rank; base=h2

\subsection{第三十五章 讨论推迟}

伯多禄如此说完，天将破晓时，匝凯进来致候我们，对伯多禄说：\Quote{西满将讨论推迟到明天，因为今天是他每十一天一次的安息日。}伯多禄回答他：\Quote{告诉西满：随你何时愿意；要知道，靠着天主的眷顾，只要你愿意，我们随时准备好见你。}匝凯听了这话，便出去回复。\par

% structure: p0073 source=h2 target=subsection reason=relative-heading-rank; base=h2

\subsection{第三十六章 一切皆好}

但他（伯多禄）见我心烦意乱，便问原因；得知是因讨论推迟，他说：\Quote{明白了世界是在天主的善眷顾下运行的人，亲爱的克肋孟，就不会为任何发生的事烦恼；他考虑到统治者的眷顾使事物顺利进行。因此，他知道天主是公义的，且以好良心生活；他懂得用正确的理性从灵魂中抖落任何临到的烦恼，因为一旦完全，它必归于某种未知的善。所以，现在不要让术士西满推迟讨论使你忧愁；这或许出于天主的眷顾，为你的益处。因此，我要坦率地向你说话，如同你是我特别的朋友。}\par

% structure: p0075 source=h2 target=subsection reason=relative-heading-rank; base=h2

\subsection{第三十七章 敌营中的间谍}

\Quote{我们当中有些人假装作为同伴跟随西满，仿佛被他那极其无神的错误所说服，为要探知他的意图并向我们透露，使我们能够有利地对抗这个可怕的人。如今我已从他们那里得知他将在讨论中使用什么论据。知道了这事，我一方面感谢天主，另一方面因讨论推迟而为你庆幸；因为你将在讨论之前被我教导，关于他将用以毁灭无知者的那些论据，你将能够安然聆听，没有跌倒的危险。}\par

% structure: p0077 source=h2 target=subsection reason=relative-heading-rank; base=h2

\subsection{第三十八章 法律的败坏}

\Quote{因为圣经中为此增添了许多反对天主的不实之言。先知梅瑟奉天主的命令，将法律连同注释传给某些被拣选的人，大约七十位，以便他们也能教导那些愿意学习的百姓；不久后，成文的法律中又添加了一些反对创造天地及其中万物的天主之法律的虚假内容；那恶者竟敢为某种正义目的而行此事。此事发生得合情合理，为使那些胆敢听从反对天主之言论的人被定罪，并使那些因爱天主而不仅不信反对祂的言论，甚至不愿听到它们的人——即使这些言论可能属实——得以显明，他们判断：在信仰上冒险，远比因亵渎的言语而怀着恶毒的良心活着更为稳妥。}\par

% structure: p0079 source=h2 target=subsection reason=relative-heading-rank; base=h2

\subsection{第三十九章 战术}

\Quote{因此，据我所知，西门打算公开露面，谈论那些为了试探而增添到圣经中的反对天主的章节，为要尽可能引诱许多可怜的人远离天主的爱。因为我们不愿公开说这些章节被增添到圣经中，否则我们会使无知的群众困惑，从而成就这个邪恶西门的目的。因为他们尚未具备分辨能力，就会逃避我们，视我们为不敬虔者；或者，仿佛不仅亵渎的章节是虚假的，他们甚至会远离圣言。因此，我们不得不假装赞同这些虚假的章节，并反过来就它们向他提问，以使他陷入困境，并在私下里对那些经受信仰考验、心志良善的人解释这些反对天主的章节；而做到这一点只有一种方法，且是简短的方法。那就是：}\par

% structure: p0081 source=h2 target=subsection reason=relative-heading-rank; base=h2

\subsection{第四十章 预备训导}

\Quote{凡是对天主所说或所写的，都是虚假的。但我说这话是真实的，不仅为名声，更为真理；等我的讲论稍进一步，我就会使你信服。因此，我最亲爱的克莱孟，你不必因西门在本次讨论之前插入了一天的间隔而难过。因为今天，在讨论之前，你将接受有关增添到圣经中的章节的教导；然后，在讨论中，就关于唯一良善的天主、也是世界的创造者这一主题，你应当不分心。但在讨论中，你甚至会惊讶：不虔敬的人竟罔顾圣经中为天主所说的大量言论，却只盯着那些反对祂的言论，并欣然提出这些；于是听众因无知而相信了反对天主的事，成为祂国度的弃民。因此，借着这一延后，你学习了圣经的奥秘，获得了不犯罪得罪天主的\Emph{手段}，将无比喜乐。}\par

% structure: p0083 source=h2 target=subsection reason=relative-heading-rank; base=h2

\subsection{第四十一章 为求教诲，而非争辩}

那时我克莱孟听了这话，说：\Quote{我实在喜乐，并感谢天主，祂在一切事上都做得美好。然而，祂知道除了万物都是为天主之外，我无法有其他想法。所以不要以为我问问题是因为怀疑有关天主的话或将要论及的话，而是为要学习，以便自己也能教导那些诚心愿意学习的人。因此请告诉我，增添到圣经中的虚假内容是什么，它们是如何成为虚假的。}然后伯多禄回答说：\Quote{即使你没有问我，我也会按次序继续，并如我所应许的，向你解释这些事。那么，你就学习圣经在诸多方面如何曲解祂，好使你在遇到时能够识别。}\par

% structure: p0085 source=h2 target=subsection reason=relative-heading-rank; base=h2

\subsection{第四十二章 天主的正确观念为圣德所必需}

\Quote{但我将要告诉你的，作为例证已经足够。但我亲爱的克莱孟，我不认为任何对天主拥有哪怕一点点爱意和真诚的人，能够接受甚至听到那些反对祂的言论。因为一个人若认为有许多神，而不是只有一位，他怎能拥有\OriginalTerm{独尊的灵魂}{monarchic soul}，并成为圣洁呢？但即使只有一位，谁若在祂身上发现许多缺陷，还会热心追求圣洁呢？因为他会希望万物的元始因自身本性的缺陷，不会追究他人的罪行。}\par

% structure: p0087 source=h2 target=subsection reason=relative-heading-rank; base=h2

\subsection{第四十三章 论天主属性的先验论证}

\Quote{因此，我们决不可相信万物的主，那创造天地及其中万物的，与别人分享祂的统治，或者说祂说谎。因为若祂说谎，那么谁讲真理呢？或者祂试验如同无知；那么谁预知呢？若祂反复思量并改变主意，谁在理解上完美、在计划上恒久呢？若祂嫉妒，谁超越竞争？若祂使人心刚硬，谁使人智慧？若祂使人眼瞎耳聋，谁赐人视觉听觉？若祂偷窃，谁施行公义？若祂嘲笑，谁真诚？若祂软弱，谁全能？若祂不义，谁正义？若祂创造邪恶之物，谁将创造良善之物？若祂行恶，谁将行善？}\par

% structure: p0089 source=h2 target=subsection reason=relative-heading-rank; base=h2

\subsection{第四十四章 同上（续）}

\Quote{然而，若祂渴望这肥沃的山丘，那么万物属于谁呢？若祂是虚假的，那么谁才是真实的呢？若祂居住在帐幕中，谁又是无界限的呢？若祂喜爱油脂、祭献、供物和奠祭，那么谁是无所需、圣洁、纯净而完美的呢？若祂喜悦烛火和灯台，那么是谁将光体安置在天空呢？若祂住在阴影、黑暗、暴风和烟雾中，谁是那光照宇宙的光呢？若祂伴随着号角、呐喊、标枪和箭矢而来，谁是为众所期盼的安宁呢？若祂喜爱战争，谁又渴望和平呢？若祂制造邪恶，谁创造美好？若祂无情，谁又是爱人者？若祂不忠于自己的承诺，谁还可信？若祂喜爱恶人、奸夫和杀人者，谁会是公正的审判者？若祂改变心意，谁是坚定不变的？若祂拣选恶人，谁又站在善人一边？}\par

% structure: p0091 source=h2 target=subsection reason=relative-heading-rank; base=h2

\subsection{第四十五章 当如何思想天主}

\Quote{因此，我的儿子克莱孟，要当心不可对天主存有别的想法，只能认为祂是唯一的天主、主和圣父，良善而正义，创造者，宽容，仁慈，维系者，恩主，安排对人类的爱，劝勉纯洁，不朽且使人不朽，无可比拟，住在善人的灵魂中，不能被容纳却又被容纳，祂将大世界固定于空间中如同中心，祂铺展诸天，坚固大地，储存众水，布置星辰于天空，使泉源涌流于大地，产生断层，举起山岳，为海洋设定界限，命令风和气流，祂以策谋的神灵安然保全那被无尽海洋包围的身体。}\par

% structure: p0093 source=h2 target=subsection reason=relative-heading-rank; base=h2

\subsection{第四十六章 审判将要来临}

\Quote{这就是我们的审判者，我们应当仰望祂，并规范自己的灵魂，凡事想祂所喜悦，说祂的好话，深信祂以忍耐显明众人的顽固，惟独祂是良善的。祂在万有终结之时，将作为公正的审判者，坐在每一个妄行不当之事的人面前。}\par

% structure: p0095 source=h2 target=subsection reason=relative-heading-rank; base=h2

\subsection{第四十七章 一个切题的问题}

我克莱孟听了这话，便说：\Quote{这真是虔诚；这真是敬虔。}我又说：\Quote{因此，我想知道，为何圣经写了这类事情？因为我记得您说过，这是为了定那些胆敢相信任何反对天主之话的人的罪。但既然您允许，我们斗胆在您的命令下询问：最亲爱的伯多禄，若有人对我们说：\InnerQuote{圣经是真实的，尽管对你而言，那些反对天主的话似乎是虚假的}，我们该如何回答他？}\par

% structure: p0097 source=h2 target=subsection reason=relative-heading-rank; base=h2

\subsection{第四十八章 一个特别的案例}

于是伯多禄回答说：\Quote{你的询问很好；因为这对你的得救有益。所以请听：既然圣经中有许多反对天主的话，而时间紧迫，因天色已晚，你可以随意提出任何一件事，我将解释它，表明它是虚假的，不仅因为它反对天主，而且因为它确实是虚假的。}于是我回答说：\Quote{我希望知道，当圣经说天主无知时，您如何能表明祂知道？}\par

% structure: p0099 source=h2 target=subsection reason=relative-heading-rank; base=h2

\subsection{第四十九章 归谬法}

于是伯多禄回答说：\Quote{你向我们提出了一个容易回答的问题。不过，请听，天主如何一无所不知甚至预知。但首先回答我问你的问题。写圣经并讲述世界如何被造、说天主不预知的那个人，他是人吗？}我说：\Quote{他是人。}伯多禄回答说：\Quote{那么，作为一个人，他如何可能确切知道世界是如何被造的，并且知道天主不预知呢？}\par

% structure: p0101 source=h2 target=subsection reason=relative-heading-rank; base=h2

\subsection{第五十章 一个令人满意的回答}

那时我已隐约明白，笑着说他是一位先知。伯多禄说：\Quote{那么，既然他作为人，因从天主领受了预知而一无所不知，那么那位将预知恩赐赐予人的天主，祂自己岂能无知呢？}我说：\Quote{您说得对。}伯多禄说：\Quote{再跟我进一步。既然我们承认天主预知一切，那么那些说祂无知的圣经必然是虚假的，而那些说祂知道的圣经才是真实的。}我说：\Quote{必然如此。}\par

% structure: p0103 source=h2 target=subsection reason=relative-heading-rank; base=h2

\subsection{第五十一章 权衡在天平上}

于是伯多禄说：\Quote{因此，若有些圣经是真实的，有些是虚假的，我们的主说\InnerQuote{你们要做精明的兑换商}，这话说得很有道理，因为圣经中有一些真话，也有一些伪作。对于因虚假圣经而犯错的人，祂恰当地指出了他们错误的原因，说：\InnerQuote{你们因此错了，因为不知道圣经中真实的事；为此你们也不认识天主的能力。}}我说：\Quote{\Emph{您说得}非常出色。}\par

% structure: p0105 source=h2 target=subsection reason=relative-heading-rank; base=h2

\subsection{第五十二章 否认圣人的过犯}

于是伯多禄回答说：\Quote{确实，我有充分的理由既不相信任何反对天主的事，也不相信任何反对律法中所记载的义人的事，认为那些都是不敬的想象。因为我深信，亚当并非罪人，他是天主亲手塑造的；诺厄并非醉酒，他在全世界中被视为义人；亚巴郎并不同时拥有三个妻子，他因节制而被认为配得众多后裔；雅各伯也没有与四人——其中两人是姐妹——交往，他是十二支派之父，并预示了我们主的来临；梅瑟也不是杀人犯，也没有从偶像崇拜的司祭那里学习审判——他向全世界颁布了天主的法律，并因公正的审判而被见证为忠信的管家。}\par

% structure: p0107 source=h2 target=subsection reason=relative-heading-rank; base=h2

\subsection{第五十三章 会议的结束}

\Quote{但这些事和类似的事，我将在适当的时候给你们解释。至于其余的事，既然你们看到，夜幕已经降临，今天所说的就足够了。但无论何时，无论何事，你们都可以大胆地问我们，我们会立即乐意解释。} 这样说完后，他站起身来。然后，进食之后，我们便就寝，因为黑夜已经降临。\par

\SourceRecord{Translated by Thomas Smith.；Ante-Nicene Fathers；Vol. 8.；Edited by Alexander Roberts, James Donaldson, and A. Cleveland Coxe.；Buffalo, NY: Christian Literature Publishing Co.,；1886.；<http://www.newadvent.org/fathers/080802.htm>.}

% structure: p0001 source=h1 target=section reason=page-title

\section{讲道词 三}

% structure: p0002 source=h2 target=subsection reason=relative-heading-rank; base=h2

\subsection{第一章 辩论之晨}

因此，两天过去，第三天破晓时，我克肋孟与其余同伴约在第二遍鸡叫时起身，准备与西蒙辩论，发现灯仍亮着，伯多禄跪着祈祷。他结束祈祷后转身，见我们已准备好听讲，便说：\par

% structure: p0004 source=h2 target=subsection reason=relative-heading-rank; base=h2

\subsection{第二章 西蒙的计谋}

\Quote{我希望你们知道，那些按照我们的安排与西蒙交往、以便了解他的意图并将其报告给我们、好使我们能对抗他各种邪恶诡计的人，已派人来告诉我，西蒙今天正如他所安排的，准备好到众人面前，并从圣经证明：创造天地及其中万物的那一位并非至高的天主，而是另有未知的、至高的天主，他以不可言喻的方式作万神之神；并且他派遣了两位神，一位是创造世界者，另一位是颁布律法者。他说这些事，为要破坏那些敬拜创造天地的独一真天主之人的正确信仰。}\par

% structure: p0006 source=h2 target=subsection reason=relative-heading-rank; base=h2

\subsection{第三章 他的目的}

\Quote{我听到这事，怎能不灰心呢！因此我希望你们，我的弟兄们，与我交往的人，知道我内心无比忧伤，看到恶者警醒诱惑人，而人对自己的得救毫不在意。因为对于那些即将被说服、认为地上的偶像并非神明的外邦人，他设法引入许多其他神明的观念，好使他们若停止多神狂热，就被欺骗而说出反对天主唯一统御的话，甚至\Emph{比现在}更糟，以致他们仍不珍视与那独一统治相关的真理，永远不能获得怜悯。为了这个企图，西蒙来与我们作战，用圣经中虚假的章节武装自己。更可怕的是，他竟敢从那些他\Emph{事实上}并不相信的先知口中，如此教条式地反对真天主。}\par

% structure: p0008 source=h2 target=subsection reason=relative-heading-rank; base=h2

\subsection{第四章 为外邦人设下的罗网}

\Quote{对我们这些从祖先手中继承了对创造万有之天主的崇拜以及能欺骗人的书卷奥秘的人，他无法得逞；但对于那些从小怀有多神思想、且不认识圣经虚妄的外邦人，他却大得逞。不仅他，任何人为外邦人讲述任何虚妄的、如梦般华丽的故事反对天主，都会被相信，因为他们自童年起心智就习惯接受反对天主的话。他们之中将只有少数人，如同众人中的少数，因诚实而不愿听任何反对创造万有之天主的恶言。只有这些外邦人才能得救。因此，你们中谁也不要完全抱怨西蒙或任何人，因为一切事都公正地发生，连圣经的虚妄也被善加用来作试验。}\par

% structure: p0010 source=h2 target=subsection reason=relative-heading-rank; base=h2

\subsection{第五章 谬误的用处}

此时我克肋孟听了这话，说：\Quote{我主，您怎么说连圣经的虚妄也被善加用来考验人呢？}他回答说：\Quote{圣经的虚妄被允许写成，是为了某种公义的理由，因邪恶的要求。当我说\InnerQuote{善加}时，意思是：在天主的计画中，恶者不爱天主不亚于善者，只在这一点上被善者超越：恶者因对深奥之事的爱，不宽恕因无知而不敬的人，渴望消灭不敬之人；但善者愿意给他们提供救治。因为善者希望所有人都藉悔改得医治，但只拯救认识祂的人。祂不医治不认识祂的人：并非祂不愿意，而是因为将预备给王国儿女的美善赐给那些因缺乏理智而像无理性动物的人，是不合法的。}\par

% structure: p0012 source=h2 target=subsection reason=relative-heading-rank; base=h2

\subsection{第六章 炼狱与地狱}

\Quote{这就是那独一天主的本性，祂创造世界，创造我们，赐给我们一切，只要任何人仍在虔敬的范围内，不亵渎祂的圣神，祂出于对那人的爱，就因爱而将灵魂带向自己。即使灵魂有罪，祂的本性是在它因所作的事受到适当惩罚后拯救它。但若有人否认祂，或以任何其他方式对祂不敬，而后悔改，他必因所犯的罪受罚，但必得救，因为他回头并活了。也许极度的虔敬和祈祷甚至能免于惩罚，无知被承认为悔改后罪过得赦的理由。但不悔改的人必被火罚消灭，尽管在其他一切方面他们至圣。但是，如我所说，到了指定时间，五分之一被永久之火惩罚者将被消灭。因为对独一天主不敬的人不能永远忍受。}\par

% structure: p0014 source=h2 target=subsection reason=relative-heading-rank; base=h2

\subsection{第七章 何为不敬？}

\Quote{对祂的不敬，是指在宗教上至死说另有别的神，无论更高或更低，或以任何方式说除了那真正存在的一位之外另有他神。因为真正存在的那位，正是人的身体所承载的形象；为了祂，诸天和众星虽本质上更优越，却服从于本质较低的人，因为统御者的形象。天主如此祝福人超过一切，好使人因着与恩惠的多少相称地爱那施恩者，藉此爱为来世得救。}\par

% structure: p0016 source=h2 target=subsection reason=relative-heading-rank; base=h2

\subsection{第八章 魔鬼的诡计}

\Quote{因此，人对天主的爱足以得救恩。恶者知道这一点；当我们急忙将使人不朽的对祂的爱播撒在那些出自外邦人、准备相信唯一的天主之人的灵魂中时，这恶者拥有足以毁灭无知者的兵器，急忙播撒许多神祇或至少一位更大者的假设，为使人们构想并信服那不智慧的事，如同犯奸淫之罪而死亡，并被逐出祂的国。}\par

% structure: p0018 source=h2 target=subsection reason=relative-heading-rank; base=h2

\subsection{第九章 圣经的不确定性}

\Quote{因此，凡愿意听任何反对天主独一统治的话的人，都配得被弃绝；但若有人胆敢听反对天主的话，如同信赖圣经，让他首先与我一同思考：若有人随己意形成适合他自己的教义，然后仔细查考圣经，他就能从圣经中提出许多证词支持他所形成的教义。那么，当每个人所愿的都能在圣经中找到时，怎么能信赖它们来反对天主呢？}\par

% structure: p0020 source=h2 target=subsection reason=relative-heading-rank; base=h2

\subsection{第十章 \Person{西满}的意图}

\Quote{因此，明天将要与我们公开辩论的\Person{西满}，大胆反对天主的独一统治，希望从这些圣经中提出许多论述，表明有许多神祇，并且有一位不是创造这世界的那位，而是高于祂；同时，他要提出许多来自圣经的证明。但我们也能轻易从中指出许多经文，表明唯有创造世界的那位是天主，除祂以外没有别的。但若有人愿意说别的，他也能按自己的喜好从中提出证明。因为圣经讲说各种各样的事，使那些不感恩地探求的人找不到真理，而只找到他所愿找的，真理为感恩的人保留——感恩就是保持我们对那创造我们存在之主的爱。}\par

% structure: p0022 source=h2 target=subsection reason=relative-heading-rank; base=h2

\subsection{第十一章 预测与预言之区别}

\Quote{因此，首先必须知道，除非来自真理的先知，否则无处可寻真理。但祂是真正的先知，常常知道一切，甚至所有人的思想，祂没有罪，因为确信于天主的审判。因此我们不应仅仅考虑祂的预知，而应考虑祂的预知是否能脱离其他原因而成立。因为医生预测某些事，有病人的脉象作为提交给他们的质料；有些人通过飞禽预测，有些人通过祭献，另一些人通过各种不同的质料提交给他们；然而这些都不是先知。}\par

% structure: p0024 source=h2 target=subsection reason=relative-heading-rank; base=h2

\subsection{第十二章 相同的主题}

\Quote{但若有人说这些预测所\Emph{显明}的预知与真正植入的预知相似，他就大受欺骗。因为他只宣告当前的事，并且若他讲真话，也仅此而已。然而，这些事对我也有用，因为它们证明预知是存在的。但唯一真先知的预知不仅知道当前的事，而且无限延伸预言直到未来的世界，且不需要任何解释，不隐晦、模棱两可地预言，使所说的话需要另一位先知来解释；而是清楚、简单地，如同我们的主和先知，藉着内在而常流的圣神，常常知道一切事。}\par

% structure: p0026 source=h2 target=subsection reason=relative-heading-rank; base=h2

\subsection{第十三章 先知知识的不变性}

\Quote{因此，祂满怀信心地陈述关于未来的事——我指苦难、地方、边界。因为祂是无瑕的先知，用祂灵魂的无垠之眼观看一切，知道隐藏的事。但若我们像许多人那样认为，即使是真先知，并非总是，而是有时当祂拥有圣神并藉着祂预知，但当祂没有时便无知——若我们这样设想，我们便欺骗了自己并误导了他人。因为这样的事属于那些被混乱之神疯狂感动的人——属于那些在祭台旁酩酊大醉、饱食脂油的人。}\par

% structure: p0028 source=h2 target=subsection reason=relative-heading-rank; base=h2

\subsection{第十四章 先知之神的恒常}

\Quote{因为若允许任何自称有预言的人在被他发现是假的情况下，仍能被相信他那时没有预知的圣神，那么要指控他是假先知便很困难；因为在他所说的许多事中，有几件应验了，于是他便被认为拥有圣神，尽管他将起初的事说在最后，将最后的事说在起初；将过去的事说成未来，将未来的事说成已过去；并且毫无次序；或是从别人那里借来并改动的事，有些是减少的、未成形的、愚昧的、模棱两可的、不雅的、隐晦的，宣扬一切无良心。}\par

% structure: p0030 source=h2 target=subsection reason=relative-heading-rank; base=h2

\subsection{第十五章 基督的预言}

\Quote{但我们的主并非如此预言；而是如我所说，凭着内在而常流的圣神是先知，并时时知道一切，祂满怀信心地陈述，如我先前所说，清楚地指出苦难、地方、指定的时间、方式、边界。因此，预言关于圣殿时，祂说：\InnerQuote{你们看见这些建筑吗？我实在告诉你们：这里将没有一块石头留在另一块上不被拆毁；这世代不会过去，直到毁灭开始。因为他们要来到，要坐在这里，要围攻它，并在这里杀戮你们的儿女。}同样，祂用明白的话说出立即要发生的事，我们现在可以用眼睛看见，为使应验在那些听到这话的人中间。因为真理的先知发出证明之言，是为祂听众的信德。}\par

% structure: p0032 source=h2 target=subsection reason=relative-heading-rank; base=h2

\subsection{第十六章 结合之教义}

\Quote{然而，谬误的宣讲者众多，他们有一个首领，就是邪恶的首领；正如真理的先知只有一位，他也是虔诚的首领，将在自己的时期，使所有纯洁的人作他的先知。但人受骗的主要原因，在于他们事先不理解联合的道理，我将每天私下向你们简要阐述此事；因为详述便太长了。但愿你们成为我所说话语的喜爱真理的审判者。}\par

% structure: p0034 source=h2 target=subsection reason=relative-heading-rank; base=h2

\subsection{第十七章　亚当是否拥有圣神}

\Quote{但如今我将开始论述。天主创造了万物，若有人不许祂亲手所造的人拥有祂伟大而神圣的先知圣神，而竟将之归于一个出自不纯血统的人，这岂不是大错特错？我不认为他能得到宽恕，尽管他因谬误的圣经而思想反对万有的父的可怖之事。因为谁侮辱永恒之王的肖像及其所属，这罪便算作加于那按其肖像被造者身上。但有人说，当他犯罪时，圣神离开了他。若如此，圣神便与他一同犯罪；说这话的人怎能逃过危险？或许他犯罪后领受了圣神？那便是赐予不义者；公义何在？但它赐予义人及不义之人。这是最不义的事。因此，一切谬误，即使有千种推理支持，也要在很久之后遭到驳斥。}\par

% structure: p0036 source=h2 target=subsection reason=relative-heading-rank; base=h2

\subsection{第十八章　亚当并非无知}

\Quote{不要受骗。我们的父亲一无所不知；事实上，即使公开流传的法律，为了不配之人的缘故，以无知之罪指控他，却将渴求知识的人打发到他那里，说：\InnerQuote{问问你的父亲，他必告诉你；你的长老，他们必指示你。}这位父亲，这些长老应当被询问。但你们未曾询问，王国属于谁的时代，预言的宝座属于谁，尽管祂亲自指出自己，说：\InnerQuote{经师和法利塞人坐在梅瑟的座位上；凡他们对你们所说的，你们都要遵行。}听从他们，祂说，如同受托拥有天国钥匙者，即知识，唯有知识能开启生命之门，唯有通过它才能进入永生。但确实，祂说，他们持有钥匙，却不许愿意进入的人进入。}\par

% structure: p0038 source=h2 target=subsection reason=relative-heading-rank; base=h2

\subsection{第十九章　基督的王权}

\Quote{因此，我说，祂自己如同父亲为儿女般从座位上站起，宣示那些从起初秘密交付给配得之人的事，将怜悯延伸至外邦人，怜恤众人的灵魂，却疏忽了自己的亲属。因为祂被视为配得来世的君王，便与那藉预选篡夺了现今王国的人作战。使祂极度忧伤的是，正是那些祂为他们如同为儿子争战的人，因无知而攻击祂。然而祂仍爱那些恨祂的人，为不信者哭泣，祝福那些毁谤祂的人，为那些与祂为敌的人祈祷。祂不仅作为父亲如此行，还教导祂的门徒也当如此，对众人以兄弟相待。我们的父如此行，我们的先知如此行。这合情合理：祂应作祂儿女的君王；藉父亲对儿女的爱，以及儿女对父亲内在的尊敬，可产生永久的和平。因为当善人治理时，被治理者因治理者而有真正的喜乐。}\par

% structure: p0040 source=h2 target=subsection reason=relative-heading-rank; base=h2

\subsection{第二十章　唯一的先知基督在不同时代显现}

\Quote{但请留意我开头关于真理的言论。若有人不许天主亲手所造的人拥有基督的圣神，却容许另一个出自不洁血统的人拥有它，他岂不犯了大不敬？若他不许他人拥有它，而说唯有那位从世界开端便改变形貌和名字、在世界多次显现，直至来到自己时代、为天主的工程受膏于怜悯，然后永远安息者，才拥有它，那便是最虔诚之举。祂的尊荣在于统管万有，在天空、大地和水域。此外，祂自己造了人，人有气息，那是灵魂不可言喻的外衣，为使人能不朽。}\par

% structure: p0042 source=h2 target=subsection reason=relative-heading-rank; base=h2

\subsection{第二十一章　否认吃禁果}

\Quote{祂自己是唯一的真先知，适当地给每个动物按其本性的功过命名，因为祂创造了它们。若祂给任一动物起名，那便是被造之物的名称，由造物者所赐。那么，祂如何还需要吃树上的果子，才能知道善恶，而祂本被命令不得吃它？但愚昧之人竟相信这事，以为一个无理性的野兽比创造万有的天主更有能力。}\par

% structure: p0044 source=h2 target=subsection reason=relative-heading-rank; base=h2

\subsection{第二十二章　男女}

\Quote{但一个伴侣与他一同被造，是女性的本性，与他不甚相同，如同性质与实体、月亮与太阳、火与光之间的差别。她，作为女性，掌管现今世界如与她相似者，被委任为第一位女先知，在一切妇人所生者中宣告预言。而另一位，作为人子，是男性，预告更美好的来世之事，如同男性。}\par

% structure: p0046 source=h2 target=subsection reason=relative-heading-rank; base=h2

\subsection{第二十三章　两种预言}

\Quote{那么让我们明白，有两种预言：一种是雄性的；且应规定，第一种，即雄性的，在降临的次序中被排在另一种之后；但第二种，即雌性的，被指定在成对降临中先来。因此，这第二种，既生于女子之中，作为此现世的女主管，希望被认为具有男子气概。因此，她偷窃雄性的种子，与她自己肉身的种子一同播种，便结出果实——即言语——完全作为她自己的。她应许要赐予现世的地上财富作为嫁妆，想要以快的换取慢的，以大的换取小的。}\par

% structure: p0048 source=h2 target=subsection reason=relative-heading-rank; base=h2

\subsection{第二十四章 女先知——误导者}

\Quote{然而，她不仅胆敢说和听有许多神，还相信自己就是其中之一，并希望成为她本性所不是的，抛弃她所有的，并作为一个女子在行经时献祭而沾血；然后她玷污触摸她的人。但当她怀胎生下暂时的君王时，她挑起战争，流许多血；那些想要从她那里学习真理的人，她通过告诉他们一切相反的、呈现许多各种服务，使他们总是寻求却找不到任何东西，直到死亡。因为从一开始，死亡的原因就临到瞎眼的人身上；因为她预言欺骗、模棱两可和歪曲，欺骗那些相信她的人。}\par

% structure: p0050 source=h2 target=subsection reason=relative-heading-rank; base=h2

\subsection{第二十五章 加音的名与本性}

\Quote{因此，这模棱两可的名字是她给她的长子起的，称他为\Emph{加音}，这个名字可以用两种方式解释；因为它既可解释为\Emph{占有}，也可解释为\Emph{嫉妒}，表示将来他要么嫉妒一个女人，要么嫉妒财产，要么嫉妒父母对她的爱。但如果这些都不是，那么他将被称为\Emph{占有}。因为她先占有了他，这对他也是有利的。因为他是个谋杀者、说谎者，并且不甘心在统治方面与他的罪和平共处。此外，那些从他而出的人依次是第一批通奸者。还有瑟琴、竖琴和制造战争器具的人。因此，他的后裔的预言充满通奸者和瑟琴，暗中通过享乐刺激战争。}\par

% structure: p0052 source=h2 target=subsection reason=relative-heading-rank; base=h2

\subsection{第二十六章 亚伯尔的名与本性}

\Quote{但他在人子中，那些灵魂天生就拥有预言、明确地作为雄性、指示来世希望的人，称自己的儿子为亚伯尔，这名字毫无歧义地译为\Emph{悲伤}。因为他指定他的儿子们为他们被欺骗的弟兄们悲伤。他应许来世的安慰时并不欺骗他们。当他说我们必须向唯一的天主祈祷时，他自己不说有众神，也不相信另一个说它们的人。他保守他所有的善，并不断增加。他憎恨祭献、流血和奠祭；他喜爱贞洁、纯洁、圣洁。他熄灭祭台的火，平息战争，教导虔诚的宣讲者智慧，洁净罪恶，批准婚姻，赞成节制，引导所有人走向贞洁，使人慷慨，规定正义，封印他们当中完美的，发布和平之言，明确预言，肯定地说，经常提及永罚之火，不断宣告天国，指示天上的财富，应许不褪色的光荣，显示因行为而得赦免。}\par

% structure: p0054 source=h2 target=subsection reason=relative-heading-rank; base=h2

\subsection{第二十七章 先知与女先知}

\Quote{何需多言？雄性全然是真理，雌性全然是虚假。但由雄性和雌性所生的，有时说真理，有时说虚假。因为雌性用她自己的血，如同红火，包裹雄性的白色种子，以外部骨头的支撑来维持自己的软弱，并因暂时的血肉之花而喜悦，以短暂的快乐损害判断的力量，引领大部分进入奸淫，从而剥夺他们在将来那卓越的新郎中的份。因为每一个人都是新娘，每当被真先知的全部真理之言所播种，他的理智就得到光照。}\par

% structure: p0056 source=h2 target=subsection reason=relative-heading-rank; base=h2

\subsection{第二十八章 属灵的奸淫}

\Quote{因此，宜听从唯一真理的先知，知道由别人所播种的话，带着奸淫的罪责，好像被新郎从祂的王国中逐出。但对于那些知道奥秘的人，死亡也是由属灵的奸淫产生的。因为每当灵魂被他人播种时，它就被圣神离弃，如同犯奸淫或通奸；因此，活的躯体，生命之圣神被抽离，就化为尘土，而灵魂在身体解体后，在审判时受到罪应有的惩罚；就如在人中，犯奸淫被捉的妇人，先被赶出房屋，然后被判受罚。}\par

% structure: p0058 source=h2 target=subsection reason=relative-heading-rank; base=h2

\subsection{第二十九章 信号发出}

当伯多禄正要向我们详细解释这奥秘之言时，匝凯来了，说：\Quote{伯多禄，现在确实是您出去参加讨论的时候了；因为一大群人挤在庭院里等着您；他们中间站着西满，像一个战争首领，由他的持矛士兵陪同。}伯多禄听了这话，命令我去祈祷，因为尚未领受带来救恩的圣洗，然后对那些已经成全的人说：\Quote{我们起来祈祷吧，愿天主以祂不失败的美善，帮助我为祂所造的人的救恩而努力。}说了这话，祈祷后，他走到庭院露天部分，那是一个大空间；有许多人为了看他而聚集，他的卓越使他们更加急切地赶来听。\par

% structure: p0060 source=h2 target=subsection reason=relative-heading-rank; base=h2

\subsection{第三十章 宗徒的问候}

因此，他站着看见所有群众都深默地注视着他，而\Person{西满术士}站在中间，就开始这样说道：\Quote{愿平安归于你们所有人，你们已准备好将你们的右手交给天主的真理，这真理是他伟大而无可比拟的恩赐，在今世，那位派遣我们来的、作为那至上益处之无误先知的那位，嘱托我们在训导的话语之前，以问候的方式向你们宣告，为的是如果你们中间有平安之子，平安就能藉我们的教导抓住他；但如果你们中有人不愿接受，那么我们就抖掉脚上的尘土作为见证——这尘土是我们辛劳所背负、带来给你们为使你们得救的——而往别人的居所和城市去。}\par

% structure: p0062 source=h2 target=subsection reason=relative-heading-rank; base=h2

\subsection{第三十一章 信德在于天主}

\Quote{我们确实告诉你们，在审判的日子，住在\Place{索多玛}和\Place{哈摩辣}之地，比住在不信之地更容易容忍。首先，因为你们自己没有保存合理的事；其次，因为你们听到了关于我们的事，却没有到我们这里来；第三，因为我们来到了你们这里，你们却不相信我们。因此，我们关心你们，主动祈祷愿我们的平安降临你们身上。所以，如果你们愿意拥有它，你们必须立即承诺不作恶，并且慷慨地忍受委屈；人的本性无法承受这些，除非它首先接受了那至上益处的知识，即认识那统御万有的主的公义本性，他保护并报复受委屈的人，并永远善待虔诚者。}\par

% structure: p0064 source=h2 target=subsection reason=relative-heading-rank; base=h2

\subsection{第三十二章 邀请}

\Quote{你们因此作为天主感恩的仆人，自己觉察到什么是合理的，就承担起那令他喜悦的生活方式，这样，你们爱他，也被他所爱，就可以永远享受美善。因为只有他最有可能赐予这美善，他使无变有，创造了诸天，奠定了大地，给海定了界限，储藏了\Place{阴府}中的事物，并用空气充满了所有地方。}\par

% structure: p0066 source=h2 target=subsection reason=relative-heading-rank; base=h2

\subsection{第三十三章 创造的化工}

\Quote{唯有他将那最初单一的简单实体转化为四种对立的元素。这样组合它们，他从中造出无数复合物，使它们转变为相反的本性，并混合起来，从而从对立物的结合中产生生活的愉悦。同样，唯有他，藉其旨意的\OriginalTerm{要有}{Fiat}，创造了天使和神灵的种族，使天上充满了居民；他也用星辰装饰了可见的穹苍，并为它们指定了路径，安排了它们的轨道。他使大地坚实以出产果实；他给海定了界限，在旱地上划出居住之地；他储藏\Place{阴府}中的事物，将其指定为灵魂的所在；他用空气充满了所有地方，使一切活物能安全呼吸以存活。}\par

% structure: p0068 source=h2 target=subsection reason=relative-heading-rank; base=h2

\subsection{第三十四章 创造的广大}

\Quote{啊，智慧天主的伟大之手，它运行万有！因为他造了无数种类的飞鸟，各式各样，彼此在各个方面都不同：我说的是它们的颜色、喙、爪、外貌、感官、声音等等。还有多少不同种类的植物，以无限多样的颜色、质地和香气为特征！还有多少陆上和水中的动物，无法尽述其形状、形态、栖息地、颜色、食物、感官、本性、数目！然后还有众多的高山、各样的石头、可怕的洞穴、泉源、河流、沼泽、海洋、港口、岛屿、森林，以及所有有人居住的世界和无人居住的地方！}\par

% structure: p0070 source=h2 target=subsection reason=relative-heading-rank; base=h2

\subsection{第三十五章 \Quote{这些是他途径的一部分。}}

\Quote{还有多少未知之事，超出了人的智慧！就是在我们所能理解的范围内，有谁知晓其界限？我说的是，天如何运转，星辰如何沿轨道运行，它们有什么形状，其存在的本质是什么，以及它们以太般的路径是什么。风的气流从何而来并被带到各处，且有不同的能量；泉源为何不停地涌出，河流奔流不息地注入大海，而\Emph{泉源}既不会枯竭，\Emph{大海}也不会被充满！无边的\Place{塔耳塔洛斯}那不可测的深度有多远！环绕万有的天穹被什么支撑！云如何从空气生出，又化为空气！雷电、雪、冰雹、雾、冰、风暴、阵雨、悬云的性情是什么！还有他如何创造植物和动物！这些事，他都精确地、不断地在它们无数种类中完善！}\par

% structure: p0072 source=h2 target=subsection reason=relative-heading-rank; base=h2

\subsection{第三十六章 统治受造物}

\Quote{因此，若有人用理性精确地察看万有，他就会发现天主造它们是为了人。因为降雨是为了果实，使人类能享用它们，动物也能得到喂养，从而对人有帮助。太阳照耀，使空气转为四季，每个季节都能给予人特殊的服务。泉源涌出，使人类有饮品。此外，在可能的范围内，谁是受造物的主宰？难道不是人吗？人接受了智慧去耕种大地、航行于海：使鱼、鸟和野兽成为他的猎物；探究星辰的轨道，开采大地，航海；建造城市，划定王国，制定法律，施行正义，认识那不可见的天主，知晓天使的名字，驱赶魔鬼，尽力用药物治疗疾病，找到对付毒蛇的咒语，理解相斥之物？}\par

% structure: p0074 source=h2 target=subsection reason=relative-heading-rank; base=h2

\subsection{第三十七章 \Quote{认识他，就是永生。}}

但是，人啊，如果你心存感恩，明白天主在一切事上都是你的恩主，你甚至可以成为不朽的，因着你的感恩，为你所造之物得以存续。如今你能成为不朽坏的，如果你承认你从前不认识的那一位，如果你爱你从前离弃的那一位，如果你唯独向那能惩罚或拯救你身体和灵魂的一位祈祷。因此，在一切之前，你要思考：没有人分享祂的统治，没有人与祂共享名号——即被称为神。因为唯独祂既是被称为神，也是神。也不可允许认为有另一位，或以那名称呼任何其他者。若有人胆敢如此，灵魂将受永罚。\par

% structure: p0076 source=h2 target=subsection reason=relative-heading-rank; base=h2

\subsection{第三十八章 西满的挑战}

伯多禄说了这些话后，西满在人群外面大声喊道：\Quote{你为什么说谎，欺骗站在你周围的无知群众，劝说他们不可认为有神并称他们为神，而犹太人中间通行的书却说有许多神？现在我愿意当着所有人的面，根据这些书与你讨论必须认为有神；首先表明你认为神的那一位，祂并非至高全能的\Emph{存有}，因为祂没有预知、不完全、有欠缺、不善良、并受许多无数痛苦情欲的支配。因此，当我照我所言从圣经中证明了这一点，就必然有另一位没有记载的神，祂有预知、完全、无欠缺、善良、远离一切痛苦情欲。但你们称为造物主的那一位却受相反之恶的支配。}\par

% structure: p0078 source=h2 target=subsection reason=relative-heading-rank; base=h2

\subsection{第三十九章 归咎于天主的缺陷}

\Quote{因此，亚当起初按祂的肖像被造时，是盲目的受造，据说不知善恶，被发现是悖逆者，被逐出乐园，受死亡的惩罚。同样，造他的那一位，因为祂不能看见各处，论及索多玛的毁灭时说：\BibleQuote{来，我们下去，看看他们是否照那达到我跟前的呼声行事；如果不是，让我知道。}这样，祂显明自己无知。论到亚当，祂说：\BibleQuote{让我们将他赶出去，免得他伸手摸生命树，吃了就永远活着。}祂说\Emph{免得}表明无知；将他赶出去，免得他吃了永远活着，也表明嫉妒。至于经上记载\BibleQuote{天主后悔造了人}，这既表明后悔又表明无知。因为这种反思是一种观点，借此一人因无知而想探究他所意愿之事的结局，或是因事情未如预期而悔恨者的作为。至于经上记载\BibleQuote{主闻到了甘饴的馨香}，这是有需要者的作为；祂喜爱肥脂是那不善良者的作为。而祂的试探，如经上所载\BibleQuote{天主试探了亚巴郎}，是恶者的作为，且不知试验的结果。}\par

% structure: p0080 source=h2 target=subsection reason=relative-heading-rank; base=h2

\subsection{第四十章 伯多禄的答复}

同样，西满从圣经中取了许多段落，似乎表明天主受各种软弱的支配。对此伯多禄说：\Quote{那邪恶且全然恶毒的人，喜欢在自己犯罪的事上控告自己吗？请回答我这个问题。}西满说：\Quote{他不喜欢。}伯多禄说：\Quote{那么，天主怎能是邪恶恶毒的呢？因为那些关于祂的邪恶之事已经普遍记载，却是出于祂自己的意愿而加上的！}西满说：\Quote{可能是对祂的指控是由另一个权势所写，并非出于祂的选择。}伯多禄说：\Quote{那么，我们首先来研究这个。如果祂确实如你先前所承认的那样，自愿控告自己，那么祂就不是邪恶的；但如果是出于另一个权势，就必须全力探究和调查，是谁将那位独一善者置于一切邪恶之下。}\par

% structure: p0082 source=h2 target=subsection reason=relative-heading-rank; base=h2

\subsection{第四十一章 \Quote{\OriginalTerm{问题状态}{Status Quæstionis}}}

西满说：\Quote{你显然回避听取圣经中对你的天主的指控。}伯多禄说：\Quote{在我看来，你本人正是在这样做；因为回避调查顺序的人，并不希望进行真正的探究。因此，我有条不紊地进行，并希望先考虑作者，显然是渴望走正路。}西满说：\Quote{首先承认，如果针对造物主的记载是真实的，祂并非超乎万有之上，因为根据圣经，祂受一切邪恶的支配；然后我们再来探究作者。}伯多禄说：\Quote{为了不显得我因不愿展开调查而反对你的无序，我回答你。我说：如果针对天主的记载是真实的，它们并不表明天主是邪恶的。}西满说：\Quote{你怎能这样主张？}\par

% structure: p0084 source=h2 target=subsection reason=relative-heading-rank; base=h2

\subsection{第四十二章 亚当是盲目的吗？}

那时，\Person{伯多禄}说：\Quote{因为写着与那些说祂坏话的言论相反的话；因此两者都无法确证。}那时，\Person{西满}说：\Quote{那么，对于那些说祂是恶的或说祂是善的圣经，如何确定真理呢？}那时，\Person{伯多禄}说：\Quote{凡圣经中与祂所造的受造界和谐一致的言论都是真的，凡与受造界相悖的都是假的。}那时，\Person{西满}说：\Quote{你怎能证明圣经自相矛盾？}而\Person{伯多禄}说：\Quote{你说\Person{亚当}被造时是瞎眼的，其实并非如此；因为祂不会向一个盲人指出知善恶树，并命令他不可尝它。}那时，\Person{西满}说：\Quote{祂指的是他的心智瞎了。}那时，\Person{伯多禄}说：\Quote{一个在尝树之前就与造他的那一位和谐一致，给所有动物起了适当名字的人，怎会在心智方面瞎眼呢？}那时，\Person{西满}说：\Quote{如果\Person{亚当}有预知，他为何不预知蛇会欺骗他的妻子？}那时，\Person{伯多禄}说：\Quote{如果\Person{亚当}没有预知，他如何给按出生时的将来行为给人类之子起名，称第一个为\Person{加音}（解释为\InnerQuote{嫉妒}），他因嫉妒杀了他的兄弟\Person{亚伯尔}（解释为\InnerQuote{悲伤}），因为他的父母为他，第一个被杀的人悲伤？}\par

% structure: p0086 source=h2 target=subsection reason=relative-heading-rank; base=h2

\subsection{第四十三章 天主的预知}

\Quote{但如果\Person{亚当}，作为天主的作品，有预知，那么创造他的天主更加有预知。而那写着天主思索，如同因无知而运用推理的话是假的；以及主试探\Person{亚巴郎}，为要知道他是否能忍受；以及所写的：\InnerQuote{我们下去，看看他们是否照那来到我这里的哀声行事；若不然，我也好知道。}并且，为免我的言论过长，凡归因于祂无知或任何邪恶的言论，都被其他肯定相反的言论推翻，证明是假的。但正因为祂确实预知，祂对\Person{亚巴郎}说：\InnerQuote{你确要知道，你的后裔必在不是他们自己的地上寄居；他们必奴役他们，恶待他们，苦害他们四百年。但他们所要服事的那国，我要审判；以后他们必带着许多财物从那里出来。但你要平平安安地到你列祖那里，得享长寿；到第四代，他们必回到这里，因为\Person{阿摩黎人}的罪孽尚未满盈。}}\par

% structure: p0088 source=h2 target=subsection reason=relative-heading-rank; base=h2

\subsection{第四十四章 天主的法令}

\Quote{但怎么？\Person{梅瑟}不是预先暗示百姓的罪，并预言他们分散在列国中吗？但如果祂把预知给了\Person{梅瑟}，祂自己怎能没有预知呢？但祂确实有预知。并且如果祂有预知，正如我们也已证明的，那么说祂思索、祂后悔、祂下去察看等等，都是过分的说法。凡事先被预知将要发生的事，都借由智慧的安排，无需后悔而成就。}\par

% structure: p0090 source=h2 target=subsection reason=relative-heading-rank; base=h2

\subsection{第四十五章 祭献}

\Quote{但祂不喜悦祭献，这由以下之事表明：那些贪恋肉食的人，一尝到就立刻被杀，并被葬在坟墓里，以致它被称为贪欲之墓。那么，起初就不喜欢宰杀动物、不愿它们被杀的那一位，并非因渴望祭献而设立祭献；祂从起初也没有要求祭献。因为祭献不能没有宰杀动物，初熟之物也不能献上。但祂创造了纯净的诸天，创造了太阳为众人带来光明，并为无数星辰指定了不变的运行秩序，祂又怎能居住在黑暗、烟雾和暴风中（因为这也是写着）？因此，\Person{西满}啊，天主的亲手笔迹——我指的是诸天——表明创造之主的旨意是纯洁而稳固的。}\par

% structure: p0092 source=h2 target=subsection reason=relative-heading-rank; base=h2

\subsection{第四十六章 对天主的毁谤}

\Quote{因此，那些控告造天之主的言论，既被旁边的相反言论所废止，又被受造界所驳斥。因为它们不是由先知的手所写。所以它们也显明与创造万有的天主之手相反。}那时，\Person{西满}说：\Quote{你怎能证明这一点？}\par

% structure: p0094 source=h2 target=subsection reason=relative-heading-rank; base=h2

\subsection{第四十七章 梅瑟的预知}

那时，\Person{伯多禄}说：\Quote{天主的法律是由\Person{梅瑟}赐下的，没有书写，传给七十位智者，以便代代相传，使治理得以延续。但在\Person{梅瑟}被接去之后，这法律是由某人书写的，而不是由\Person{梅瑟}。因为法律本身写着：\InnerQuote{\Person{梅瑟}死了；他们把他葬在\Place{培敖尔}的房屋附近，直到今日无人知道他的坟墓。}但\Person{梅瑟}怎能写\Person{梅瑟}死了呢？而在\Person{梅瑟}之后约五百年，这法律被发现放在所建的圣殿里；又过了约五百年，它被掳走，在\Person{拿步高}时代被烧毁。如此，在\Person{梅瑟}之后书写，且多次失传，这甚至表明\Person{梅瑟}的预知，因为他预见到它会消失，所以没有写它；而那些写它的人，因未预见到其消失而被显明是无知的，他们不是先知。}\par

% structure: p0096 source=h2 target=subsection reason=relative-heading-rank; base=h2

\subsection{第四十八章 真理的试金石}

那时，\Person{西满}说：\Quote{既然，如你所说，我们必须通过将有关天主的事与受造界比较来理解，那么如何可能识别法律中来自\Person{梅瑟}传统且真实、并与这些谎言混杂的其他事呢？}那时，\Person{伯多禄}说：\Quote{按照天主的眷顾，在成文法律中记载了一节经文，无可争议，为要清楚显示所写之事哪些是真、哪些是假。}那时，\Person{西满}说：\Quote{是哪一节？指给我们看。}\par

% structure: p0098 source=h2 target=subsection reason=relative-heading-rank; base=h2

\subsection{第四十九章 真先知}

\Person{伯多禄}便说：\Quote{我立刻告诉你们。在法律的第一卷书末尾写着：\BibleQuote{权杖必不离开犹大，领袖必不离他的两腿之间，直到那应得者来到；祂是万民的期待。}因此，若有人能理解那在犹大的权杖和领袖消失之后而来的、且为万民所期待的那一位，他就能藉这节经文认出祂确实已经来临；并相信祂的教导，从而知道圣经中哪些是真确的，哪些是虚假的。}然后\Person{西满}说：\Quote{我明白你所说的是指你的耶稣，即圣经所预言的那一位。因此，就当这是真的。那么，请告诉我们，祂如何教导你分辨圣经。}\par

% structure: p0100 source=h2 target=subsection reason=relative-heading-rank; base=h2

\subsection{第五十章 祂关于圣经的教导}

然后\Person{伯多禄}说：\Quote{关于真伪混杂的问题，我记得有一次祂责备撒杜塞人说：\BibleQuote{所以你们错了，因为你们不知道圣经的真事，由此你们也不认识天主的能力。}但是，如果祂指责他们不知道圣经的真事，那就显然圣经中有虚假的事。此外，正因为祂说：\InnerQuote{你们要作精明的兑换银钱者}，这是因为有真确的和虚假的话语。而当祂说：\InnerQuote{你们为什么不明白圣经中合理的事呢？}祂使那自愿正确判断的人的理解力更强。}\par

% structure: p0102 source=h2 target=subsection reason=relative-heading-rank; base=h2

\subsection{第五十一章 祂关于法律的教导}

\Quote{祂派遣到当时圣经的经师和教师那里，如同到那些知道当时法律真事的人那里，这是众所周知的。还有祂说：\BibleQuote{我来不是要废除法律，}然而祂似乎是在废除它，这表明祂所废除的并不属于法律。祂说：\BibleQuote{天地要过去，但法律的一笔一划也不能过去，}这表明在天地之前就要过去的事并不真正属于法律。}\par

% structure: p0104 source=h2 target=subsection reason=relative-heading-rank; base=h2

\subsection{第五十二章 基督的其他言论}

\Quote{既然天地依然存在，祭献、国度、以及由妇人所生者中的预言，以及诸如此类的事，都因不是天主的法令而过去了；因此祂说：\BibleQuote{凡我天父没有栽种的植物，都要被拔出来。}所以祂，作为真先知，说：\BibleQuote{我是生命之门；谁从我进入，就进入生命。}因为没有其他教导能够拯救。因此祂也呼喊道：\BibleQuote{凡劳苦的，都到我这里来，}即那些寻求真理而找不到的人；又说：\BibleQuote{我的羊听我的声音；}并在别处说：\BibleQuote{寻求，就会找到。}因为真理不在表面。}\par

% structure: p0106 source=h2 target=subsection reason=relative-heading-rank; base=h2

\subsection{第五十三章 基督的其他言论}

\Quote{而且有一个从天上来的见证声音说：\BibleQuote{这是我的爱子，我所喜悦的，你们要听祂。}除此之外，祂愿意更充分地指正那些先知是错的——就是他们声称曾从他们学习的那先知——祂宣称他们渴望真理而死，却没有学到真理，说：\BibleQuote{许多先知和君王渴望看见你们所看见的，听见你们所听见的；我实在告诉你们，他们既没有看见，也没有听见。}另外祂说：\BibleQuote{我就是梅瑟所预言的那一位，他说：主我们的天主要从你们弟兄中给你们兴起一位像我一样的先知；你们要听祂一切的话；凡不听那先知的，必要死亡。}}\par

% structure: p0108 source=h2 target=subsection reason=relative-heading-rank; base=h2

\subsection{第五十四章 其他言论}

\Quote{因此，没有祂的教导，就不可能获得救恩的真理，即使人永远在寻找不存在那东西的地方寻找。但救恩的真理过去就在，现在也在我们的耶稣的话语中。因此，祂知道法律的真实之事，就对撒杜塞人说——他们问梅瑟为何允许娶七个妻子——：\BibleQuote{梅瑟因你们心硬，才给你们这条诫命；但从起初并不是这样：因为起初创造人的天主，造了他们一男一女。}}\par

% structure: p0110 source=h2 target=subsection reason=relative-heading-rank; base=h2

\subsection{第五十五章 基督的教导}

\Quote{但对那些像圣经所教导那样认为天主起誓的人，祂说：\BibleQuote{你们的话，是就说是，非就说非；其他多说的，便是出于邪恶者。}对那些说亚巴郎、依撒格和雅各伯已死的人，祂说：\BibleQuote{天主不是死人的天主，而是活人的天主。}对那些认为天主试探人（如圣经所说）的人，祂说：\BibleQuote{那试探人的是那恶者，}祂自己也曾被试探。对那些认为天主不预知的人，祂说：\BibleQuote{因为你们的天父在你们祈求祂以前，已经知道你们需要这一切。}对那些相信（如圣经所说）天主不看顾一切的人，祂说：\BibleQuote{要在暗中祈祷，你的父在暗中看见，必要报答你。}}\par

% structure: p0112 source=h2 target=subsection reason=relative-heading-rank; base=h2

\subsection{第五十六章 基督的教导}

\Quote{对那些认为天主不是良善的（如圣经所说）的人，祂说：\BibleQuote{你们中间，哪个人儿子向他求饼，反而给他石头；或者求鱼，反而给他蛇？你们纵然不好，尚且知道把好东西给你们的儿女，何况你们的天父，岂不更将好东西赐给那些求祂、并遵行祂旨意的人！}但对那些肯定祂在圣殿里的人，祂说：\BibleQuote{不可指着天起誓，因为天是天主的宝座；也不可指着地起誓，因为地是祂的脚凳。}对那些认为天主喜悦祭献的人，祂说：\BibleQuote{天主喜欢仁爱，而不是祭献}——即认识祂自己，而不是全燔祭。}\par

% structure: p0114 source=h2 target=subsection reason=relative-heading-rank; base=h2

\subsection{第五十七章 基督的教导}

\Quote{但对那些被说服认为祂是邪恶的（如圣经所说）的人，祂说：\BibleQuote{不要称我为善，因为只有\Emph{一位}是善的。}又说：\BibleQuote{你们应当是良善和慈悲的，如同你们的天父，祂使太阳升起，光照恶人，也光照善人；降雨给义人，也给不义的人。}但对那些被误导而想象有多个神的人（如圣经所说），祂说：\BibleQuote{以色列，你要听：上主你的天主是唯一的主。}}\par

% structure: p0116 source=h2 target=subsection reason=relative-heading-rank; base=h2

\subsection{第五十八章 \Person{西满}的逃亡}

因此\Person{西满}因见\Person{伯多禄}正迫使他自己按\Person{耶稣}所教导的方式运用圣经，便不愿讨论进入有关天主的道理，尽管\Person{伯多禄}已将讨论转为问答形式，恰如\Person{西满}自己所要求的。然而，讨论持续了三天。当第四天破晓时，他暗中出发，一直到了\Place{腓尼基}的\Place{提洛}。过了不多几日，一些先遣者来对\Person{伯多禄}说：\Quote{\Person{西满}在\Place{提洛}行了许多大奇迹，搅乱了那里许多百姓；并且用许多诽谤使你被人憎恨。}\par

% structure: p0118 source=h2 target=subsection reason=relative-heading-rank; base=h2

\subsection{第五十九章 \Person{伯多禄}决心追随}

\Person{伯多禄}听了这话，次日夜间召集了听道群众；他们一到齐，他便说：\Quote{我正在前往那些说有许多神祇的万国，为教训并宣讲有一位天主，祂创造了天、地及其中万物，好使他们爱祂而得救之时，邪恶已抢先于我，并凭结合律将\Person{西满}差遣在我前，为的是这些人若停止说有许多神祇，并否认地上被称为神祇者，就会以为天上也有许多神祇；这样，因不感受独一统治的卓越，他们便要在永罚中丧亡。最可怕的是，由于真道具有无比的威力，他以诽谤先发制我，并劝他们甚至一开始就不要接待我；免得那诽谤者被证实自己实际上是一个魔鬼，而真道却被接受和相信。因此我必须迅速赶上他，免得这诬告因时间拖延而完全掌控所有人。}\par

% structure: p0120 source=h2 target=subsection reason=relative-heading-rank; base=h2

\subsection{第六十章 应指派继任者}

\Quote{因此，既然必须另选一人代替我填补我的位置，让我们大家同心合意祈求天主，使祂显明我们当中谁是最好的，使他坐在基督的座位上，虔诚地治理祂的教会。那么，谁当被选立呢？因为藉天主的旨意，那人是蒙福的：\BibleQuote{他的主人要委派他管理他的同伴，按时分粮给他们，不在心里想：我的主人必会迟延；也不动手打他的同伴，且与娼妓和酒徒一起吃喝。那仆人的主人要在意想不到的日子、不知道的时辰来到，把他斩成两半，使他和伪善者一同承受不信者的那份。}}\par

% structure: p0122 source=h2 target=subsection reason=relative-heading-rank; base=h2

\subsection{第六十一章 独一统治}

\Quote{但在座的人中，若有人能教导人的愚昧却退缩，只顾自己的安逸，就该听这句判词：\BibleQuote{可恶懒惰的仆人！你应把我的钱交给钱庄，我回来时就能连本带利收回。把这无用的仆人丢到外面的黑暗里去。}并且有充分理由；因为，祂说：\BibleQuote{人啊，考验我的话正如钱庄考验银钱一样。}所以信徒群众应当服从某一个人，以达和谐。因为趋向一人统治的形式（即独一统治）能使臣民因良好秩序而享平安；但若人人因好权而不愿服从一人，则必因分裂而全然败亡。}\par

% structure: p0124 source=h2 target=subsection reason=relative-heading-rank; base=h2

\subsection{第六十二章 服从带来和平}

\Quote{此外，还要让你们眼前发生的事说服你们：现今世上由于有许多国王，战事不断发生。因为每个人都把别人的统治作为战争的借口。但如果有一位普世的主宰，他没有理由发动战争，就会享有永久的和平。因此，简而言之，对于那些被认为配得永生的人，天主将在将来的世界里任命一位普世的君王，好藉独一统治而有不断的和平。所以，众人都应当跟从某一位领袖，尊崇他如同天主的肖像；而领袖应当熟悉进入圣城的道路。}\par

% structure: p0126 source=h2 target=subsection reason=relative-heading-rank; base=h2

\subsection{第六十三章 \Person{匝凯}被任命}

\Quote{但在座的人中，我除了\Person{匝凯}还能选谁呢？主也曾进到他家里并歇息，认为他配得救恩。}说了这话，\Person{伯多禄}按手在站在旁边的\Person{匝凯}身上，迫使他坐在自己的椅子上。但\Person{匝凯}俯伏在他脚前，恳求准许他辞去统治职务；同时许诺并说：\Quote{凡是统治者当做的事，我都愿意做；只求你不要给我这个名号；因为我害怕担任统治者的名号，它充满了苦涩的嫉妒和危险。}\par

% structure: p0128 source=h2 target=subsection reason=relative-heading-rank; base=h2

\subsection{第六十四章 主教职}

然后\Person{伯多禄}说：\Quote{你若为此害怕，不要被称为\Emph{统治者}，而要称为\Emph{被指派者}，因为主曾允许你这样称呼，当祂说：\BibleQuote{那仆人如果被他的主人\Emph{指派}管理他的同伴，是有福的。}但如果你希望别人全然不知你有管理权，那么在我看来，你似乎不知道：公认的首领权威对群众的尊重有极大的影响。因为每个人都服从那已接受权威的人，良心是极大的约束。难道你不清楚：你不应像外邦人的统治者那样统治，而应像仆人一样服事他们，像父亲对待受压迫者，像医生一样探望他们，像牧人一样守护他们——总之，要全力照顾他们的得救？你难道以为我不知道我迫使你承担的是什么劳苦，要你被众人评判，而任何人都不可能取悦众人？但那行善的人最可能取悦天主。因此，我恳求你，为天主、为基督、为弟兄们的得救、为他们的秩序以及你自己的益处，诚心承担。}\par

% structure: p0130 source=h2 target=subsection reason=relative-heading-rank; base=h2

\subsection{第六十五章}

% structure: p0131 source=h2 target=subsection reason=relative-heading-rank; base=h2

\subsection{\OriginalTerm{我不愿当主教}{Nolo Episcopari}}

\Quote{还要考虑另一件事：管理基督的教会所付出的劳苦和危险越大，赏报也越大。然而，对于能够管理却拒绝的人，惩罚也同样更大。因此，我知道你是我随从中最优秀的，我愿你在判断上善用主所托付给你的卓越能力，以便你能得到\Emph{做得好，忠信又善良的仆人}的称赞，而不致像那隐藏一塔冷通的人那样被指责并受罚。但如果你不愿被立为教会的好管家，就当另举一位比你更有学识、更忠信的人来代替。然而你不能这样做，因为你曾与主为伴，目睹祂的奇妙作为，并学会了教会的管理。}\par

% structure: p0133 source=h2 target=subsection reason=relative-heading-rank; base=h2

\subsection{第六十六章 不顺从的危险}

\Quote{你的工作是下令什么是适当的；弟兄们的工作则是服从，不可违命。因此，服从的人将得救，但违命的人将受主的惩罚，因为主持者受托代表基督。所以，对主持者的尊敬或轻视会传递给基督，再由基督传给天主。我说这话，是让这些弟兄们明白，他们因不顺从你所面临的危险；因为凡不服从你命令的，就是不服从基督；而不服从基督的，就是冒犯天主。}\par

% structure: p0135 source=h2 target=subsection reason=relative-heading-rank; base=h2

\subsection{第六十七章 教会职务人员的责任}

\Quote{因此，教会如同建在山上的城，必须有天主所嘉许的秩序和良好治理。特别是，主教作为首领，他所讲的话应当被听从；长老们应留意所吩咐的事得以实行。执事们巡行各处，照顾弟兄们的身体和灵魂，并向主教报告。其余的弟兄们都应耐心忍受委屈；但如果他们希望所受的委屈得到裁判，就应在长老面前和好；长老们应将这和好向主教报告。}\par

% structure: p0137 source=h2 target=subsection reason=relative-heading-rank; base=h2

\subsection{第六十八章 \Quote{婚姻常受尊重}}

\Quote{他们不仅应当劝勉年轻人结婚，也应当劝勉年长者结婚，以免燃烧的情欲因淫乱或奸淫给教会带来瘟疫。因为，在一切罪中，奸淫的邪恶最被天主憎恨；它不仅毁灭犯罪者本人，也毁灭与他一同吃喝交往的人。它就像狗的狂犬病一样，具有传染自己狂乱的本性。因此，为了贞洁，不仅长老们，连众人都应赶紧结婚。因为犯奸淫者的罪必然临到众人。所以，劝勉弟兄们贞洁，这是首要的爱德。因为这是灵魂的治疗。身体的滋养是休息。}\par

% structure: p0139 source=h2 target=subsection reason=relative-heading-rank; base=h2

\subsection{第六十九章 \Quote{不可停止聚会}}

\Quote{但如果你们爱弟兄们，就不要从他们那里索取什么，反倒要与他们分享你们所有的。要喂养饥饿的人，给口渴的人喝，给赤身露体的人穿，探望病人，尽你们所能帮助在监里的人，欢欢喜喜地接待陌生人到自己的住所，不要怨恨任何人。至于你们应当怎样虔诚，你们自己的心灵会借着正直的判断教导你们。但首要的是——若我真的需要告诉你们——要常常聚会，即使是每时每刻，尤其是在规定的聚会上。因为你们若这样做，就是在安全的围墙之内。因为混乱无序是灭亡的开端。因此，谁也不可因对弟兄的嫉妒而离开聚会。因为你们中若有人离开聚会，他就会被视为那拆散基督教会的人，并与奸淫者一同被赶出去。因为如同奸淫者受到他内心的邪灵影响，他以某种借口将自己分离，给恶者留下攻击自己的余地——如同一只待窃的羊，像在羊栈外被找到的人。}\par

% structure: p0141 source=h2 target=subsection reason=relative-heading-rank; base=h2

\subsection{第七十章 \Quote{听从主教}}

\Quote{然而，要听从你们的主教，不要厌倦给予他一切尊荣；要知道，将尊荣给他，就是归给基督，并由基督归给天主；而给予尊荣的人，必得百倍的回报。因此，要尊敬基督的宝座。因为你们甚至奉命尊敬梅瑟的座位，尽管坐在其上的被视为罪人。现在，我对你们说够了；我认为不必再对他说应当如何无可指责地生活，因为他是一位合格的门徒，属于那位也教导了我的那一位。}\par

% structure: p0143 source=h2 target=subsection reason=relative-heading-rank; base=h2

\subsection{第七十一章 基督徒的各种本分}

\Quote{但是，弟兄们，有一些事你们不必等着听，而应当自己思考什么是合情合理的。只有\Person{匝凯}完全献身，为你们劳苦，需要维持生计，不能处理自己的事务，他如何获得必要的供养呢？你们预先考虑他的生活，岂不合理吗？不要等他向你们请求，因为那是乞丐所为。他宁愿饿死也不愿这样做。难道你们不考虑\Quote{工人当得起他的工资}吗？谁也不能说：那么，白白赐予的圣言就被出售了吗？绝不可以！因为如果有人有生活来源却收取报酬，他就是出售圣言；但若没有生活来源的人为了生活而收取供养——正如主在晚餐时和在朋友间也收取，祂虽一无所有，却唯独祂是万有的主人——他便没有犯罪。因此，要适当地尊敬长老、要理讲授者、有用的执事、善度生活的寡妇、作为教会儿女的孤儿。凡是遇到紧急需要供应之处，大家应一同捐献。彼此以善待，为了你们自己的救恩，不要畏缩于忍受任何事物。}\par

% structure: p0145 source=h2 target=subsection reason=relative-heading-rank; base=h2

\subsection{第七十二章 授职}

他说了这些话后，将手按在\Person{匝凯}身上，说：\Quote{万有的统治者与主，圣父与天主，请以羊群守护牧人。祢是起因，祢是能力。我们是被助者，祢是助佑者、医师、救主、壁垒、生命、希望、避难所、喜乐、期待、安息。总之，祢是我们的一切。为得永久的救恩，祢合作、保存、保护。祢无所不能。因为祢是统治者的统治者，万主之主，诸王的管理者。赐予监督能力，使他能释放当释放的，束缚当束缚的。使他明智。请因祢的名，保护祢的基督的教会如同美丽的新娘。因为光荣永远属于祢。愿光荣归于圣父、圣子及圣神，直到永远。阿们。}\par

% structure: p0147 source=h2 target=subsection reason=relative-heading-rank; base=h2

\subsection{第七十三章　圣洗}

他说了这些话后，接着说：\Quote{你们中谁愿意受洗，从明天开始禁食，并每日被覆手，询问你们所愿意的事。因为我还要同你们一起停留十天。}三天后，他开始施洗，并召唤我、\Person{阿桂拉}和\Person{尼凯达斯}，对我们说：\Quote{我打算七天后前往\Place{提洛}，我希望你们今天就离开，并秘密地留宿在\Person{犹斯塔}的女儿、\Place{客纳罕}人\Person{贝勒尼基}那里，从她那里学习，并准确写信告诉我\Person{西满}在做什么。因为此事对我至关重要，好让我相应地做好准备。因此，你们即刻平安地出发吧。}我们便离开正在施洗的他，按照他的命令，先于他前往\Place{腓尼基}的\Place{提洛}。\par

\SourceRecord{Translated by Thomas Smith.；Ante-Nicene Fathers；Vol. 8.；Edited by Alexander Roberts, James Donaldson, and A. Cleveland Coxe.；Buffalo, NY: Christian Literature Publishing Co.,；1886.；<http://www.newadvent.org/fathers/080803.htm>.}

% structure: p0001 source=h1 target=section reason=page-title

\section{第四篇讲道}

% structure: p0002 source=h2 target=subsection reason=relative-heading-rank; base=h2

\subsection{第一章. \Person{贝尔尼刻}的款待}

\Emph{因此}我\Person{克肋孟}，辞别\Place{凯撒勒雅斯特拉托尼斯}，与\Person{尼塞塔斯}和\Person{阿桂拉}一起，进入了\Place{腓尼基}的\Place{提洛}。并遵照遣派我们的\Person{彼得}的指示，我们寄宿在\Person{贝尔尼刻}那里，她是\Person{犹斯塔}——那客纳罕妇人的女儿。她极其喜乐地接待了我们，并待我以极大的敬重，待\Person{阿桂拉}和\Person{尼塞塔斯}以深情，并如朋友般自由畅谈，因着喜乐，她礼貌地款待我们，并殷勤地催促我们歇息身体。因此，我觉察到她正试图让我们稍作耽搁，我便说：\Quote{你确实做得很好，忙于履行爱的本分；但敬畏我们的天主必须领先于此。因为我们正为许多灵魂而进行一场战斗，我们害怕将自己的安逸置于他们的救恩之前。}\par

% structure: p0004 source=h2 target=subsection reason=relative-heading-rank; base=h2

\subsection{第二章. \Person{西满}的作为}

\Quote{因为我们听说那术士\Person{西满}，在\Place{凯撒勒雅}与我们的主\Person{彼得}的辩论中被击败后，立即赶到了这里，正在做许多坏事。因为他违背真理，向所有敌人诽谤\Person{彼得}，并偷窃众人的灵魂。他本人是个术士，却称彼得为术士；他本人是个骗子，却宣称彼得是骗子。尽管在辩论中他全然被击败并逃走了，他却说自己得胜了；并且他不断指责他们不应听\Person{彼得}——好像他真是担心他们会被一个可怕的术士迷惑似的。}\par

% structure: p0006 source=h2 target=subsection reason=relative-heading-rank; base=h2

\subsection{第三章. 使命的目标}

\Quote{因此，我们的主\Person{彼得}得知这些事后，派遣我们来调查所报告给他的事情；如果确有其事，我们可以写信告诉他，使他能前来当面驳斥那人对他的指控。既然许多灵魂面临着危险，为此我们必须暂时忽略身体的安歇；我们想从你们这些住在这里的人口中确实得知，我们所听到的是否属实。请务必详细告诉我们。}\par

% structure: p0008 source=h2 target=subsection reason=relative-heading-rank; base=h2

\subsection{第四章. 西满的所作所为}

但\Person{贝尔尼刻}被问时，她说：\Quote{这些事确实如你所听到的；我还要告诉你一些关于这同一\Person{西满}的其他事，也许你还不知道。因为他每天都使整个城市震惊，在集市中间制造幻影和幽灵；当他行走时，雕像移动，许多阴影走在他前面，他说这些是死者的灵魂。许多试图证明他是骗子的人，他很快使他们与自己和好；之后，以宴请为幌子，杀了一头牛，给他们吃，他用各种疾病感染他们，并将他们交给恶魔支配。总之，他伤害了许多人，并被当作神，因此既被畏惧又被尊敬。}\par

% structure: p0010 source=h2 target=subsection reason=relative-heading-rank; base=h2

\subsection{第五章. 谨慎是勇德的上策}

\Quote{所以，我认为没有人能扑灭这样点燃的火。因为无人怀疑他的承诺；但每个人都断言事情确实如此。因此，为了避免你们暴露于危险中，我劝你们在\Person{彼得}来之前不要试图对抗他；惟有他能抵抗如此的能力，因为他是我们的主耶稣基督最受尊敬的门徒。我如此惧怕这人，以至如果他未曾在我主\Person{彼得}的辩论中被击败过，我会劝你们甚至说服\Person{彼得}本人不要去对抗\Person{西满。}}\par

% structure: p0012 source=h2 target=subsection reason=relative-heading-rank; base=h2

\subsection{第六章. 西满的离去}

于是我说道：\Quote{如果我们的主\Person{彼得}不知道他自己一个人就能胜过这能力，他就不会先派遣我们，命令我们秘密地探查关于\Person{西满}的事情，并写信给他。}然后，傍晚来临，我们吃了晚餐，就去睡觉。但到了早晨，\Person{贝尔尼刻}的一位朋友来，说\Person{西满}已乘船去了\Place{西顿}，并留下了\Person{阿皮翁·普雷斯托尼刻斯}——一位\Place{亚历山大}人，职业是文法学者，我认识他，他是我父亲的朋友；还有一位占星家，\Person{底奥斯颇里的阿努比翁}，以及\Person{阿忒诺多鲁斯}，一位雅典人，信奉\Person{厄丕枯禄}的学说。我们得知了这些关于\Person{西满}的事情后，在早晨写了信并派送给了\Person{彼得}，然后出去散步了。\par

% structure: p0014 source=h2 target=subsection reason=relative-heading-rank; base=h2

\subsection{第七章. 阿皮翁的问候}

而\Person{阿皮翁}遇见了我们，不仅带着刚才提到的两个同伴，还有大约三十人。他一看到我，就向我问安并亲吻我，说道：\Quote{这就是\Person{克肋孟}，我曾多次向你们谈起他的高贵出身和通识教育；因为他与\Person{提庇留凯撒}的家族有亲缘关系，并具备所有希腊学问，却被一个名叫\Person{彼得}的野蛮人引诱，按照犹太人的方式说话和行事。因此，我恳求你们与我一同努力使他归正。现在我当着你们的面问他。让他告诉我，既然他认为自己全心投入了虔诚，他抛弃本国的习俗，堕落至野蛮人的习俗中，难道不是最不虔敬的行为吗？}\par

% structure: p0016 source=h2 target=subsection reason=relative-heading-rank; base=h2

\subsection{第八章. 一次挑战}

我回答说：\Quote{我确实接受你对我友好的情感，但我对你的无知持异议。因为你的情感是友好的，因为你希望继续持守那些你认为是好的\Emph{习俗}。但你不准确的知识，却以友谊为借口，试图为我设下圈套。}阿比盎于是说：\Quote{在你看来，一个人遵守祖辈的习俗，并依照希腊人的方式判断，这就是无知吗？}然后我回答说：\Quote{一个渴望虔诚的人，不应完全遵守祖辈的习俗；而要遵守那些虔诚的习俗，抛弃那些不虔诚的。因为一个不虔诚父亲之子，若他愿意虔诚，就不应渴望追随他父亲的宗教。}阿比盎于是回答说：\Quote{那么如何？你说你的父亲是一个行为邪恶的人吗？}我说：\Quote{他不是行为邪恶，而是意见邪恶。}阿比盎说：\Quote{我想知道他的邪恶理解是什么。}我说：\Quote{因为他相信希腊人的虚假邪恶神话。}阿比盎问：\Quote{希腊人的这些虚假邪恶神话是什么？}我说：\Quote{就是关于诸神的错误意见，如果你愿意容忍我，你将与那些渴望学习的人一起听到。}\par

% structure: p0018 source=h2 target=subsection reason=relative-heading-rank; base=h2

\subsection{第九章　哲学家的不配目的}

\Quote{因此，在我们开始谈话之前，让我们先退到某个更安静的地方，在那里我将与你交谈。我之所以希望私下谈话，是因为多数人，甚至所有哲学家，都不是诚实地按照事物本相来判断。因为我们知道许多人，甚至那些以哲学自傲的人，是虚荣的，或者是为了利益而穿上哲学家的长袍，而不是为了美德本身；如果他们找不到他们从事哲学所追求的东西，就会转向嘲弄。因此，为了这样的人，让我们选择一个适合私下会谈的地方。}\par

% structure: p0020 source=h2 target=subsection reason=relative-heading-rank; base=h2

\subsection{第十章　凉爽的退隐处}

他们当中有一位富人，拥有一座常青植物的花园，说：\Quote{因为天气很热，让我们稍微退到城外的我的花园里去吧。}于是他们出去，坐在一个地方，那里有清澈凉凉的溪流，和各种树木的绿荫。我愉快地坐在那里，其他人围着我；他们静默着，不是用言语而是用热切的眼神向我表明，他们需要我对我的主张的证明。因此我继续这样说：——\par

% structure: p0022 source=h2 target=subsection reason=relative-heading-rank; base=h2

\subsection{第十一章　真理与习俗}

\Quote{希腊人啊，真理与习俗之间有一个很大的区别。因为真理是在被诚实寻求时发现的；但习俗，无论所接受的习俗的性质如何，是真是假，它不经判断而自我强化；接受它的人既不以之为真而喜悦，也不以之为假而忧伤。因为这样的人不是凭判断，而是凭偏见相信，将自己的希望寄托在之前生活之人的意见上，纯属偶然。而且，脱下祖先的衣服并不容易，即使它被证明是完全愚蠢和可笑的。}\par

% structure: p0024 source=h2 target=subsection reason=relative-heading-rank; base=h2

\subsection{第十二章　\OriginalTerm{命运}{genesis}}

\Quote{因此我说，希腊人的全部学问是一个邪恶恶魔的最可怕捏造。因为他们引入了许多他们自己的神，且是邪恶的，受制于各种情欲；这样，愿意做类似事情的人就不会感到羞耻，这属于人，以神话中诸神的邪恶不安生活为榜样。因为他不感到羞耻，这样的人就不提供任何悔改的希望。还有其他人引入了命运，即所谓的\OriginalTerm{命运}{genesis}，违背它没有人能遭受或做任何事。因此，这也与前者相似。因为任何人若认为除了\OriginalTerm{命运}{genesis}之外没有人有事可做或可受，就容易陷入罪恶；一旦犯罪，他就不因自己的不虔诚而悔改，以他受\OriginalTerm{命运}{genesis}驱使而做这些事作为自己的辩护。既然他不能纠正\OriginalTerm{命运}{genesis}，他就没有理由为自己所犯的罪感到羞耻。}\par

% structure: p0026 source=h2 target=subsection reason=relative-heading-rank; base=h2

\subsection{第十三章　命定}

\Quote{还有其他人引入一种无预见的命运，仿佛万物自行运转，没有一位主人的监督。但这样思考，如我们所说，是最严重的意见。因为，仿佛没有人在监督、预判并按各人的应得分配，他们就因无所畏惧而轻易行事，尽其所能。因此，持这样意见的人不容易，或者根本不会，过有德的生活；因为他们没有预见到可能使他们转变的危险。但你们所谓的野蛮犹太人的教义是最虔诚的，它引入一位作为这整个世界的圣父和创造者，本性善良而正直：善良，因祂宽恕悔改者的罪；正直，因祂按各人悔改后行为的配得而报应各人。}\par

% structure: p0028 source=h2 target=subsection reason=relative-heading-rank; base=h2

\subsection{第十四章　\Quote{合乎虔诚的教义}}

\Quote{这种教义，即使也是神话式的，只要是虔诚的，对今世生活也不是没有益处。因为每个人，在期待被全视的天主审判时，受到更大的激励趋向美德。但若这教义也是真实的，它会使那生活有德的人免于永罚，并赐予他从天主而来的永恒说不尽的福乐。}\par

% structure: p0030 source=h2 target=subsection reason=relative-heading-rank; base=h2

\subsection{第十五章　诸神的邪恶}

\Quote{但我回到希腊人的首要教义，即以故事形式陈述有许多神，且受制于各种情欲。为了不在清楚的事上花费太多时间，提及每一个所谓神的邪恶行为，我无法讲述他们所有的恋情：宙斯、波塞冬、普路托、阿波罗、狄俄尼索斯、赫拉克勒斯，以及他们每一个。这些你们自己并非不知道，也学会了他们的生活方式，受过希腊学问的教导，以便作为诸神的竞争者，你们可以做类似的事。}\par

% structure: p0032 source=h2 target=subsection reason=relative-heading-rank; base=h2

\subsection{第十六章　朱庇特的邪恶}

\Quote{但我将始于最尊贵的宙斯，他的父亲克洛诺斯，如你们所说，吞吃了自己的孩子，并用金刚石镰刀割掉其父乌拉诺斯的肢体，为那些热心于诸神奥秘的人树立了孝敬父母和爱护子女的榜样。而朱庇特本人捆绑了自己的父亲，将他囚禁在塔尔塔罗斯；他也惩罚其他诸神。为了那些想做不可言说之事的人，他生了墨提斯，并吞吃了她。但墨提斯是种子；因为吞吃一个孩子是不可能的。作为那些自渎之人的借口，他劫走了伽倪墨得斯。作为奸淫者的帮凶，他常被发现是个奸淫者。对于那些想与姐妹乱伦的人，他以与姐妹赫拉、得墨忒耳以及天上的阿佛洛狄忒（有人称之为多多那）的交合为榜样。对于那些想与女儿乱伦的人，从他的故事中有一个邪恶的例子，就是他与其女珀耳塞福涅乱伦。但他有无数的不虔敬行为，以至由于他极度的邪恶，那不虔敬之人竟相信他是神的故事。}\par

% structure: p0034 source=h2 target=subsection reason=relative-heading-rank; base=h2

\subsection{第十七章 \Quote{制造它们的人必定像他们}}

\Quote{你们会认为无知的人对这些幻想略微愤怒是合理的。但我们对那些有学问的人——其中一些自称文法家和智者，断言这些行为配得上诸神——该说些什么呢？因为他们自己放纵，就抓住这个神话借口，作为诸神的模仿者，自由地行丑事。}\par

% structure: p0036 source=h2 target=subsection reason=relative-heading-rank; base=h2

\subsection{第十八章 第二天性}

\Quote{因此，住在乡村的人比他们犯罪少得多，因为他们没有受那些教导，而敢于做这些事的人却受了教导，从邪恶的教训中学会了不虔敬。因为那些自幼年便通过这样的寓言学习文字的人，当他们的灵魂尚柔软时，就把那些被称为神的人的不虔敬行为移植到自己的心中；因此，当他们长大时，就结出果实，如同播在灵魂里的恶种。最糟的是，根深蒂固的污秽不易铲除，因为当他们成年时，这些污秽在他们看来是苦涩的。因为每个人都乐于保持童年养成的习惯；既然习惯的力量几乎不亚于本性，他们就难以改向起初未播在灵魂里的善事。}\par

% structure: p0038 source=h2 target=subsection reason=relative-heading-rank; base=h2

\subsection{第十九章 \Quote{无知即是福}}

\Quote{因此，年轻人不应满足于那些腐化的课程，中年人也应小心避免听闻希腊人的神话。因为关于他们的诸神的教训远比无知更糟，正如我们从乡村居民的例子所表明的，他们因未受希腊人教导而犯罪较少。实在，他们的这些寓言、表演和书籍应当避免，如果可能，甚至他们的城市也应回避。因为那些充满邪恶学问的人，甚至用气息感染与他们交往的人，如同用疯狂传染他们自己的情欲。最糟的是，他们中最有学问的人，就越多地偏离本性的判断。}\par

% structure: p0040 source=h2 target=subsection reason=relative-heading-rank; base=h2

\subsection{第二十章 哲学家的虚假理论}

\Quote{他们中甚至有些自称哲学家的人断言这样的罪是无所谓的，并说对那些愤怒于这类行为的人是无知的。因为他们说，这些事在本性上并非罪，而是由古代智者制定的法律所禁止的，那些智者知道人心因不稳定，在这些事上大受激动，彼此争战；为此，智者立法禁止这些事为罪。但这是一个可笑的假设。因为，如果这些事不是罪，何以会引起骚乱、谋杀和各种混乱？难道奸淫不带来生活上的缺陷和更多邪恶吗？}\par

% structure: p0042 source=h2 target=subsection reason=relative-heading-rank; base=h2

\subsection{第二十一章 奸淫的恶果}

\Quote{但有人说，如果一个人不知道他的妻子是奸妇，为什么他不愤怒、不狂怒、不发狂？为什么他不发动战争？可见这些事在本性上并非邪恶，而是人的不合理意见使它们可怕。但我认为，即使这些可怕的事不发生，妇女与奸夫交往，通常会抛弃丈夫，或若继续同居，就会谋害他，或把丈夫劳苦所得给予奸夫；并且当丈夫不在时，如果因奸夫而怀孕，就会因羞愧的定罪而试图毁灭腹中的胎儿，成为杀婴者；甚至可能在毁灭胎儿时自己也一同灭亡。但如果丈夫在家时她因奸夫怀孕生子，孩子长大后不认识自己的父亲，以为不是父亲的人为父；这样，不是父亲的人临终时将财产留给别人的孩子。奸淫还会自然产生多少其他邪恶！隐秘的罪恶我们不知道。因为狂犬病狗传染所有接触者以不可见的狂犬病，同样，奸淫的隐秘之恶，即使不被知晓，也阻断后裔。}\par

% structure: p0044 source=h2 target=subsection reason=relative-heading-rank; base=h2

\subsection{第二十二章 更卓越的道}

\Quote{但现在让我们略过这个。但我们都知道，一般人对奸淫过度愤怒，因此爆发战争、房屋倾覆、城邑被掳，以及无数其他邪恶。为此，我投奔犹太人的圣洁天主和法律，确信法律由天主的公义审判所制定，灵魂必按其所行的功过在某个时候接受报应。}\par

% structure: p0046 source=h2 target=subsection reason=relative-heading-rank; base=h2

\subsection{第二十三章 \Quote{我往何处去躲避祢的面？}}

我这样说了以后，阿皮翁打断了我的话。\Quote{什么！}他说，\Quote{难道希腊的法律不也禁止恶行、惩罚奸夫吗？}于是我说道：\Quote{那么，那些违反法律的希腊众神就该受惩罚。但我怎能自制，如果我认为众神自己首先实践了所有恶行，包括奸淫，却没有受罚？他们反倒更该受罚，因为他们并非欲望的奴隶。但如果他们也受欲望支配，他们又怎么是神呢？}然后阿皮翁说：\Quote{我们不要以众神，而以法官为鉴；看着他们，我们就会害怕犯罪。}我接着说：\Quote{这并不恰当，阿皮翁。因为着眼于人的人，因指望逃脱察觉而敢于犯罪；但把全视的天主摆在灵魂前的人，知道不能逃过祂的注意，就连暗中也必不犯罪。}\par

% structure: p0048 source=h2 target=subsection reason=relative-heading-rank; base=h2

\subsection{第二十四章  寓意}

阿皮翁听了这话，说道：\Quote{自从我听说你与犹太人交往以来，我就知道你的判断已疏离了。因为有人说得很好：\InnerQuote{邪恶的交往败坏良好的品行。}}我于是说：\Quote{因此良好的交往纠正邪恶的品行。}阿皮翁说道：\Quote{今天我完全满足于知道了你的立场，所以我也允许你先说。但明天，在此处，如果你同意，我将在这些朋友聚会时，当着他们的面，证明我们的众神既非奸淫者，也非杀人者，并非败坏孩童者，亦非与姐妹或女儿乱伦者。但古人只希望好学者得知奥秘，就用你所说的那些寓言将其掩盖。因为他们以生理学的方式谈论煮物，称其为宙斯，称时间为克罗诺斯，称水之常流性质为瑞亚。然而，正如我所承诺的，明天我要逐一向你们解释，在你们清晨聚集时，揭示事物的真相。}我对此回答说：\Quote{明天就照你所承诺的行吧。但现在，且听一些与你将要说的相反的话。}\par

% structure: p0050 source=h2 target=subsection reason=relative-heading-rank; base=h2

\subsection{第二十五章  次日的约定}

\Quote{如果众神的行为本是善的，却被邪恶的寓言掩盖，那么编织这掩盖的人就显出极大的邪恶，因为他用邪恶的叙事遮盖了高尚的事，使人无人效仿。但如果他们确实行了不虔敬的事，那么反该用善良的叙事遮盖它们，免得人们视他们为优越者，便同样开始犯罪。}我说了这些，在场的人显然开始对我所说的话抱有好感：因为他们再三恳切地要求我次日再来，然后就离去了。\par

\SourceRecord{Translated by Thomas Smith.；Ante-Nicene Fathers；Vol. 8.；Edited by Alexander Roberts, James Donaldson, and A. Cleveland Coxe.；Buffalo, NY: Christian Literature Publishing Co.,；1886.；<http://www.newadvent.org/fathers/080804.htm>.}

% structure: p0001 source=h1 target=section reason=page-title

\section{第五篇讲道}

% structure: p0002 source=h2 target=subsection reason=relative-heading-rank; base=h2

\subsection{第一章 \Person{阿皮翁}未现身}

因此，次日，在\Place{提洛}，按约定，我来到那安静之处，在那里见到了其余的人，还有一些其他人。我向他们致意。但未见\Person{阿皮翁}，便询问他缺席的原因；有人说他从昨晚起便身体不适。于是，当我说我们应立即出发去看望他时，几乎众人都恳求我先对他们讲论，然后我们再去见他。因此，既然众人看法一致，我便开始说：——\par

% structure: p0004 source=h2 target=subsection reason=relative-heading-rank; base=h2

\subsection{第二章 \Person{克莱孟}先前对\Person{阿皮翁}的了解}

\Quote{昨日我离开此地时，朋友们，我承认，因对即将与\Person{阿皮翁}进行的讨论深感忧虑，我彻夜未眠。无法入睡之际，我想起我在\Place{罗马}对他耍的一个把戏。情况如下：我从少年时期起，我\Person{克莱孟}就是个爱好真理、追寻有益灵魂之事的人，终日忙于构建和驳斥理论；但因未能找到任何完美之事，我心绪烦乱而卧病在床。当我卧病时，\Person{阿皮翁}来到\Place{罗马}，他是我父亲的朋友，便寄宿在我处；他听闻我卧病，因略通医术，便前来探问病因。但我深知此人极其憎恨犹太人，且曾著书反对他们，并与这位\Person{西蒙}结交，并非出于求学之心，而是因知他是撒玛黎雅人，且是犹太人的憎恨者，已出来反对犹太人，所以才与他结盟，以便从他那里学到反对犹太人的言论；——}\par

% structure: p0006 source=h2 target=subsection reason=relative-heading-rank; base=h2

\subsection{第三章 \Person{克莱孟}的计谋}

\Quote{我之前就知道\Person{阿皮翁}的这些事，他一问我生病的缘由，我便假装回答，说我正像年轻人那样心中痛苦烦乱。他听了说道：\InnerQuote{我儿，尽管向父亲一样畅所欲言：你的灵魂有何疾苦？}当我再次假装呻吟，因羞于谈论爱情，便以沉默和低垂的目光暗示了我想传达的意思。他见我如此，以为我是恋慕一女子，便说：\InnerQuote{人生中没有不可救助之事。因为我自己年轻时，曾恋慕一位极有才德的女子，不仅以为无法得到她，甚至不敢期望向她表白。然而，我遇到了一位精通魔法的埃及人，与之结交，向他吐露了恋情，他不仅助我成就了愿望，而且慷慨地尊重我，毫不迟疑地教给我一道咒语，藉此我得到了她；一得到她，通过他的秘密指导，因老师的慷慨而心悦诚服，我的恋情便治愈了。}}\par

% structure: p0008 source=h2 target=subsection reason=relative-heading-rank; base=h2

\subsection{第四章 \Person{阿皮翁}的承诺}

\Quote{他说：\InnerQuote{那么，如果你也像常人一样遭受此类事，就放心地向我坦诚吧；因为七日内我就能使你完全得到她。}我听了，看着自己心中的目标，便说：\InnerQuote{请原谅我，我并非完全相信魔法的存在；因为我已经试过许多夸下海口的人，都欺骗了我。然而，你的保证打动了我，使我心生希望。但当我细想起来，我担心魔鬼有时不听从术士的命令。}}\par

% structure: p0010 source=h2 target=subsection reason=relative-heading-rank; base=h2

\subsection{第五章 魔法理论}

\Quote{这时\Person{阿皮翁}说：\InnerQuote{你要承认，我比你更了解这些事情。然而，免得你以为从我这里听说的与你所讲的全无关系，我要告诉你，魔鬼为何必须服从术士的命令。因为正如士兵不能违抗将军，将军本身不能违抗君王——若有人对抗上司，便难逃惩罚——同样，魔鬼不能不服事作为他们将军的天使；当被因他们之名召唤时，魔鬼颤抖顺服，深知若违命将受重罚。但天使自身，当术士以他们统治者的名号召唤时，也服从，以免因抗命而被毁灭。因为除非所有有生命有理性的存在者都预见到统治者的惩罚，否则相互背叛，混乱便会发生。}}\par

% structure: p0012 source=h2 target=subsection reason=relative-heading-rank; base=h2

\subsection{第六章 疑虑}

\Quote{于是我说：\InnerQuote{那么，诗人与哲学家所说的那些事是否正确：在阴府中，恶人的灵魂因他们的罪恶而被审判受罚，如\Person{伊克西翁}、\Person{坦塔罗斯}、\Person{提堤俄斯}、\Person{西绪福斯}、\Person{达那俄斯}的女儿们，以及此地所有不敬神的人？再者，如果这些事不真，魔法又如何可能存在呢？}他告诉我阴府中确实如此后，我问他：\InnerQuote{我们为何不惧怕魔法，既然深信在阴府中奸淫要受惩罚？因为我不认为强迫不愿的女子行奸淫是正义之事；但若有人愿意去说服她，我不仅感谢他，而且乐意接受。}}\par

% structure: p0014 source=h2 target=subsection reason=relative-heading-rank; base=h2

\subsection{第七章 有区别的区别}

\Quote{这时\Person{阿皮翁}说：\InnerQuote{难道你不认为，无论是用魔法得到她，还是用言语欺骗她，都是一回事吗？}我便说：\InnerQuote{不完全相同；这两者相去甚远。因为用魔法之力胁迫不愿的女子，如同图谋侵害贞洁女子，将受最严厉的惩罚；但用言语说服她，把选择权交由她自己的意愿，并不强迫她。我认为，说服\Emph{一个女人}的人所受的惩罚，比强迫她的人要轻。因此，你若能说服她，我得到她后必将感谢你；否则，我宁死也不愿违她之意勉强她。}}\par

% structure: p0016 source=h2 target=subsection reason=relative-heading-rank; base=h2

\subsection{第八章 谄媚还是魔法}

\Quote{於是阿皮雍真的困惑了，說：\Quote{我該對你說什麼呢？因為你時而像一個被愛情攪擾的人，祈求得著她；轉眼間又像並不愛她似的，把恐懼看得比慾望更重。你認為若能說服她，你便無可指責，如同沒有罪；但若藉魔術得著她，你將招致懲罰。難道你不知道，每項行為的結局——即事情已經發生——才是被審判的，而達到目的的手段則不被考慮？如果你犯姦淫，是靠魔術成功的，你將被判為行惡；若靠說服，難道在姦淫上你就能免於罪嗎？}於是我說：\Quote{因著我的愛，我必須從可用的手段中選擇一個來獲得我所愛的人；我會盡可能選擇哄騙她而非使用魔術。但用諂媚說服她也不容易，因為那女人很有哲學頭腦。}}\par

% structure: p0018 source=h2 target=subsection reason=relative-heading-rank; base=h2

\subsection{第九章　一封情書}

\Quote{於是阿皮雍說：\Quote{我更希望能夠說服她，正如你所願，只要我們能與她交談。}我說：\Quote{那不可能。}阿皮雍便問是否可以寄一封信給她。我說：\Quote{那倒可以。}阿皮雍說：\Quote{就在今晚，我要寫一篇讚揚姦淫的文章，你從我這兒拿去送給她；我希望她會被說服而同意。}阿皮雍於是寫了那篇文章，交給了我；就在今晚我想起它，幸好我隨身帶著，連同我帶在身邊的其他文件。}說著，我把那篇文章展示給在場的人，並應他們的要求讀給他們聽；讀完後，我說：\Quote{諸位，這就是希臘人的教導，提供無憂無慮犯罪的寬廣許可。文章如下：——}\par

% structure: p0020 source=h2 target=subsection reason=relative-heading-rank; base=h2

\subsection{第十章　情人致所愛者}

\Quote{匿名，因愚人之法律。奉萬物首生者\OriginalTerm{厄洛斯}{Eros}之命，致敬。我知道你潛心哲學，為德行之故而崇尚高貴生活。但在所有神明中有誰比眾神更高貴？在人類中有誰比哲學家更高貴？因為唯獨他們知道哪些行為在本性上是善或惡，哪些並非如此卻因法律規定而被視為善惡。如今有些人以為所謂姦淫的行為是邪惡的，儘管它在各方面都是良善的。因為它是厄洛斯為增進生命而設立的。厄洛斯是眾神中最年長的。沒有厄洛斯，元素、神明、人類、無理性動物或任何其他事物都無法混合或生成。因為我們都是厄洛斯的工具。他藉著我們，是萬物受生的製造者，是居住在我們靈魂中的心智。因此，並非我們自己願意時，而是當他命令我們時，我們才渴望行他的旨意。但如果我們照他的旨意而渴望，卻為了所謂的貞潔而試圖克制慾望，我們所做的豈不是最大的不敬，竟反對眾神和人類中最古老的神？}\par

% structure: p0022 source=h2 target=subsection reason=relative-heading-rank; base=h2

\subsection{第十一章　\Quote{放縱情慾，行各樣污穢。}}

\Quote{但願所有門戶都向他敞開，所有有害且專橫的法律都被廢除，這些法律是由狂熱的人訂立的，他們受愚昧的支配，不願理解合理之事，此外還懷疑所謂的姦夫，難怪被宙斯本人透過米諾斯和拉達曼提斯以專橫的法律所嘲弄。因為住在我們靈魂中的厄洛斯是無法約束的；戀人的激情並非自願。因此，制定這些法律的宙斯本人也曾親近無數婦女；並且，據一些智者說，他有時與人類交合，作為生育子女的恩人。但對於那些他知道隱瞞身份會有益的人，他就改變形狀，以免使他們悲傷，或似乎違反他自己制定的法律。因此，你們這些辯論哲學的人，為了美好的生活，理應效法那些被公認為更高貴者，他們曾行過萬次性交。}\par

% structure: p0024 source=h2 target=subsection reason=relative-heading-rank; base=h2

\subsection{第十二章　朱庇特的愛情}

\Quote{為免浪費時間徒舉更多例子，我將先從宙斯本人——眾神與人類之父——的一些擁抱說起。因為數目眾多，無法列舉全部。請聽這位偉大朱庇特的愛情事蹟，他為避開愚昧之人的狂熱，屢次改變形體以隱瞞。首先，他想向智者表明姦淫並非罪過，當他即將結婚時——按眾人的看法，他明知自己是姦夫，在他的第一次婚姻中，但實際上並非如此——藉著一個表面上的罪，他成就了一場無罪的婚姻。因為他娶了自己的妹妹赫拉，假扮成杜鵑的翅膀；她生了赫柏和伊利提亞。因為他未曾與任何人交合而生墨提斯，赫拉亦然，生了伏爾甘。}\par

% structure: p0026 source=h2 target=subsection reason=relative-heading-rank; base=h2

\subsection{第十三章　朱庇特的愛情續}

\Quote{他把从克罗诺斯和塔拉塞所生的妹妹奸污，这是在他肢解克罗诺斯之后，由此生了厄洛斯和库普里斯，他们又称她为多多涅。然后他扮成萨堤尔与尼克透斯的女儿安提俄佩交合，生了安菲翁和泽托斯。他又化作丈夫安菲特律翁的形状拥抱其妻阿尔克墨涅，生了赫拉克勒斯。他又变为鹰接近阿斯克勒庇俄斯的女儿埃癸娜，生了埃阿科斯。他化作熊与福科斯的女儿阿玛尔忒亚同寝；又化作金雨落在阿克里西俄斯的女儿达那厄身上，由此生出珀耳修斯。他对吕卡翁的女儿卡利斯托变得如狮子般狂野，生了第二个阿尔库斯。他又用牛与福尼克斯的女儿欧罗巴交合，生出米诺斯、拉达曼堤斯和萨耳珀冬；又化作蚂蚁与阿刻罗俄斯的女儿欧律墨杜萨交合，生出密耳弥冬。他又化作兀鹰与赫耳塞俄斯的一个宁芙交合，生出西西里的古代智者们。他到罗得岛与地生的朱诺交合，生出帕耳该俄斯、克罗尼俄斯、库提斯。他又化作其丈夫福尼克斯的形状奸污了奥西阿，为他生出了安喀诺俄斯。他又化作天鹅或鹅与忒斯提俄斯的女儿涅墨西斯（也被认为是勒达）交合，生出了海伦；又化作星宿生出了卡斯托耳和波吕丢刻斯。与拉弥亚时他变为戴胜鸟。}\par

% structure: p0028 source=h2 target=subsection reason=relative-heading-rank; base=h2

\subsection{第十四章 朱庇特不加掩饰的恋情}

\Quote{他化作牧羊人使谟涅摩绪涅成为缪斯之母。他自燃其身，与卡德摩斯的女儿塞墨勒结婚，生出狄俄倪索斯。他化作龙奸污了自己的女儿珀耳塞福涅，她被认为是其兄弟普路托的妻子。他还与许多其他女人交合，并未改变形态；因为她们的丈夫并不视此为罪而对他怀有恶意，反而深知他与自己妻子交合是为他们慷慨生子，赐予他们赫耳墨斯们、阿波罗们、狄俄倪索斯们、恩底弥翁们以及我们提到过的其他人，因这父亲身份而极其俊美。}\par

% structure: p0030 source=h2 target=subsection reason=relative-heading-rank; base=h2

\subsection{第十五章 反自然的淫欲}

\Quote{且不在无休止的阐述中浪费时间，你将发现朱庇特与众神有无数结合。但愚昧之人称这些神的行为为通奸；甚至那些神也不以男色滥用为可耻之事，反而视之为合宜。例如，朱庇特自己恋慕伽倪墨得斯；波塞冬恋慕珀罗普斯；阿波罗恋慕喀倪剌斯、扎铿托斯、许阿铿托斯、福耳巴斯、许拉斯、阿德墨托斯、库帕里索斯、阿密克拉斯、特洛伊罗斯、廷那的布兰科斯、波特尼亚的帕罗斯、俄耳甫斯；狄俄倪索斯恋慕拉俄尼斯、安珀洛斯、许墨奈俄斯、赫耳玛佛洛狄托斯、阿喀琉斯；阿斯克勒庇俄斯恋慕希波吕托斯，赫淮斯托斯恋慕珀琉斯；潘恋慕达佛尼斯；赫耳墨斯恋慕珀耳修斯、克律萨斯、忒修斯、俄德律索斯；赫拉克勒斯恋慕阿布得罗斯、德律俄普斯、伊俄卡斯托斯、菲罗克忒忒斯、许拉斯、波吕斐摩斯、海蒙、科诺斯、欧律斯透斯。}\par

% structure: p0032 source=h2 target=subsection reason=relative-heading-rank; base=h2

\subsection{第十六章 对不贞的赞美}

\Quote{亲爱的，我已向你部分陈述了所有较著名神灵的恋情，好让你知道对这件事的狂热仅限于愚昧之人。因此他们是凡人，悲苦地度过一生，因为他们出于热忱，将众神视为卓越之事宣称为邪恶。因此，将来你有福了，效法众神而非世人。因为世人因自己所感受的，见你持守他们以为的贞洁，就赞美你，却不帮助你。但众神见你相似于他们，必既赞美又帮助你。}\par

% structure: p0034 source=h2 target=subsection reason=relative-heading-rank; base=h2

\subsection{第十七章 星座}

\Quote{因为且为我数算他们赏报了多少情妇，有些安置于星辰中，有些则祝福她们的子女和同伴。如此，宙斯使卡利斯托成为星座，称为小熊座，也有人称之为狗尾。波塞冬也因安菲特里忒将海豚置于天上；他又因米诺斯之女欧律阿勒的缘故，为其子俄里翁在天上赐予一席之地。狄俄倪索斯将阿里阿德涅的花冠制成星座，宙斯则赐予协助他抢掠伽倪墨得斯的鹰以及伽倪墨得斯本人以\InnerQuote{倒水者}的尊荣。他又为欧罗巴的缘故荣耀公牛；又将卡斯托耳、波吕丢刻斯和海伦赐予勒达，使他们成为星辰。又为达那厄的缘故赐予珀耳修斯；为卡利斯托的缘故赐予阿尔库斯。还有那处女，即\OriginalTerm{狄刻}{\GreekText{Δίκη}}，为忒弥斯的缘故；以及赫拉克勒斯为阿尔克墨涅的缘故。但我不再详述；因为一一细说众神因众多情妇而祝福了多少人，未免冗长。他们在与凡人的交合中，愚昧之人斥之为恶行，不知享乐乃是人类的大益。}\par

% structure: p0036 source=h2 target=subsection reason=relative-heading-rank; base=h2

\subsection{第十八章 哲学家为通奸辩护}

\Quote{然而如何？著名的哲学家不是也颂扬享乐，随心所欲与女人交合吗？其中第一位便是那位希腊的导师，福玻斯本人曾说：\InnerQuote{万民之中，苏格拉底最为智慧。}他不是教导说，在治理良好的城邦中妇女应共有吗？他不是将美貌的阿尔喀比亚德藏在他的哲学家长袍下吗？苏格拉底的门徒安提斯泰尼写道，不可放弃所谓通奸的必要性。甚至他的弟子第欧根尼，不也是公开与拉伊斯交合，并为此收受酬劳，将她扛在肩上招摇过市吗？伊壁鸠鲁不是颂扬享乐吗？亚里斯提卜不是涂抹香膏，完全委身于阿佛洛狄忒吗？芝诺不是暗示无差别，说神性遍及万物，为让聪明人知道，人与任何人交合如同与自身交合，因而禁止所谓通奸，或与母亲、女儿、姐妹、孩子交合，是多余的吗？克律西波斯在他的情书里提及阿尔戈斯的一座雕像，描绘赫拉与宙斯处于猥亵姿态。}\par

% structure: p0038 source=h2 target=subsection reason=relative-heading-rank; base=h2

\subsection{第十九章 情书结尾}

\Quote{我知道，对未通晓真理的人而言，这些事似乎可怕且极其卑劣；但对众神和希腊的哲学家，以及对那些受狄俄尼索斯和得墨忒耳秘仪启示的人而言，却并非如此。但首先，为了不浪费时间谈论所有神祇和所有哲学家的生活，请以两位为首要标记——众神中最伟大的宙斯，以及哲人中的苏格拉底。我在信中所提的其他事，请理解并留意，以免让你的爱人悲伤；因为，你若行事与神祇和英雄相悖，将被视为邪恶，并受相应的惩罚。但若你向每一位爱人献身，那么，作为神祇的效仿者，你将自他们那里获得恩惠。至于其余，最亲爱的，请记住我向你揭示的奥秘，并以书信告知你的选择。祝你安好。}\par

% structure: p0040 source=h2 target=subsection reason=relative-heading-rank; base=h2

\subsection{第二十章 对其的利用}

\Quote{因此，我从阿皮翁那里收到这封便笺，仿佛我真的要把它寄给一位心爱的人，便假装她已回信；次日，当阿皮翁来时，我把回信交给他，仿佛出自她手，内容如下：}\par

% structure: p0042 source=h2 target=subsection reason=relative-heading-rank; base=h2

\subsection{第二十一章 对阿皮翁来信的答复}

\Quote{我诧异，你称赞我智慧时，却像对待愚人般写信给我。因为，你为说服我顺从你的激情，便引用神祇神话中的例子，说厄洛斯是众神中最年长的，如你所言，超越众神和人类，竟不惧怕亵渎，为的是败坏我的灵魂并侮辱我的身体。厄洛斯并非众神的首领——我指的是那个与情欲相关的厄洛斯。因为他若自愿贪恋，那他自己便是自己的痛苦与惩罚；而自愿受苦者不可能是神。但他若非自愿地贪恋交合，并渗透我们的灵魂，如同经过我们身体的肢节，侵入我们的心智，那么驱使他去爱的力量便比他更大。再者，驱使他者若自身被另一种欲望所驱，那么更大于他的驱使者便出现。如此，我们便陷入无穷的爱之链条，这是不可能的。因此，既无驱使者，亦无受驱者；这不过是爱者本身的情欲冲动，因希望而增长，因绝望而消减。}\par

% structure: p0044 source=h2 target=subsection reason=relative-heading-rank; base=h2

\subsection{第二十二章 虚假的神话}

\Quote{那些不愿克制卑劣情欲的人，便污蔑神祇，意在表明神祇先行他们所作之事，从而使自己免于责备。因为那些被称为神祇的人，若为生育子女而行奸淫，而非出于放荡，那为何他们也玷污男子？但据说他们以星宿赠予情妇。难道在此之前没有星宿，直到因放纵，天空才被奸淫者以星宿装点？而那些化为星宿者的子女，又如何在地狱中受惩罚——阿特拉斯被负荷，坦塔罗斯受干渴折磨，西西弗斯推石，提提俄斯脏腑被刺穿，伊克西翁不停绕轮旋转？为何这些神圣的情人将所污辱的女子化为星宿，却未赐予后者同样的恩惠？}\par

% structure: p0046 source=h2 target=subsection reason=relative-heading-rank; base=h2

\subsection{第二十三章 并非神祇}

\Quote{那么，他们并非神祇，而是暴君的象征。因为在高加索山脉中，并非在天上，而是在地上，显露出一个坟墓，属于克洛诺斯，一个野蛮且吞食子女的人。此外，那个传说中同样吞食自己女儿墨提斯的放荡宙斯的坟墓，可在克里特岛见到；普路托和波塞冬的坟墓在阿刻戎湖；赫利俄斯的在阿斯特拉，塞勒涅的在卡莱，赫耳墨斯的在赫尔莫波利斯，阿瑞斯的在色雷斯，阿佛洛狄忒的在塞浦路斯，狄俄尼索斯的在忒拜，其余的在别处。无论如何，我所提及的这些人的坟墓都可见到；因为他们曾是凡人，且在这些事上，是邪恶之人和术士。否则，他们便不会成为暴君——我指传说中著名的宙斯和狄俄尼索斯——而是通过改变形貌，他们得以战胜任何所愿之人，为了他们设计的目的。}\par

% structure: p0048 source=h2 target=subsection reason=relative-heading-rank; base=h2

\subsection{第二十四章 若原则为善，便应贯彻}

\Quote{但若我们必须效仿他们的生活，让我们不仅模仿他们的奸淫，也模仿他们的宴会。因为克洛诺斯吞食了自己的子女，宙斯也同样吞食了自己的女儿。我还能说什么？珀罗普斯曾作为晚餐献给众神。因此，让我们在举行不圣洁的婚姻之前，先举办一场如同神祇的晚餐；这样，晚餐才配得上婚姻。但你绝不会同意此事；我也不会同意奸淫。此外，你以厄洛斯的愤怒恐吓我，视之为强大之神。厄洛斯并非神祇，按我的理解，他是为了生命延续，根据运行万物的上智安排，从活物的体质中产生的一种欲望，为使整个种族不致消亡，而因着快乐，从必死者之实体中生出另一人，藉由合法婚姻产生，使他知晓在年老时赡养自己的父亲。而奸淫所生者无法做到这一点，因为他们对生养他们的人没有天然的情感。}\par

% structure: p0050 source=h2 target=subsection reason=relative-heading-rank; base=h2

\subsection{第二十五章 结婚比欲火更佳}

\Quote{因此，既然情欲之渴望是为了延续和合法繁衍，正如我所说，父母应为子女的贞洁着想，藉由读圣洁之书以教导，预先以卓越的言论培养他们，来引导这欲望；因为习惯是第二天性。此外，应时常提醒他们法律所定的惩罚，以恐惧为缰绳，使他们不致奔逐于邪恶的快乐。同时，在欲望萌发之前，应以婚姻满足青春期的自然冲动，首先说服他们不要注视其他女子的美貌。}\par

% structure: p0052 source=h2 target=subsection reason=relative-heading-rank; base=h2

\subsection{第二十六章 答复的结束}

\Quote{因为我们的心，每当欣喜地印上爱人的形像，总像是在镜中看着那模样，回忆便折磨它；若不能遂愿，它就设法去得到；若真得到了，它反更增，如同火焰添了木柴，尤其在激情兴起之前，没有畏惧印在爱者的灵魂上。因为正如水灭火，畏惧也熄灭了不合理的欲望。因此，我从一个犹太人那里学会了理解和做天主喜悦的事，就不会被你们谎言的寓言诱入奸淫。但愿天主帮助你们，照你们的愿望和努力保持贞洁，并给那因爱火燃烧的灵魂提供良药。}\par

% structure: p0054 source=h2 target=subsection reason=relative-heading-rank; base=h2

\subsection{第二十七章 仇恨的理由}

\Quote{阿皮翁听了这假装的回答后，说：\InnerQuote{我恨犹太人难道没有理由吗？现在这里有个犹太人遇见了她，使她皈依了他的宗教，并说服她守贞洁，从此她再也不可能与其他男子交合了；因为这些人把天主当作一切行为的普世监察者，在贞洁上极其坚持，仿佛无法向祂隐藏似的。}}\par

% structure: p0056 source=h2 target=subsection reason=relative-heading-rank; base=h2

\subsection{第二十八章 骗局被坦白}

\Quote{我听了这话，对阿皮翁说：\InnerQuote{现在我要向你坦白实情。我并未恋慕那女人，也未恋慕任何人，因为我的灵魂完全专注于其他欲望，即对真道的探究。直到如今，尽管我考察了许多哲学家的学说，却未倾向其中任何一个，除了犹太人的教义——他们有个商人曾在罗马寄居，卖麻布衣服，一次幸运的相遇将独一天主的教义简单地摆在我面前。}}\par

% structure: p0058 source=h2 target=subsection reason=relative-heading-rank; base=h2

\subsection{第二十九章 阿皮翁的怨恨}

\Quote{然后，阿皮翁从我这里听了实情，带着他对犹太人毫无道理的仇恨，既不知道也不想知道他们的信仰是什么，毫无理性地发怒，立即默默离开了罗马。自那以后，这是我第一次与他见面，因此我自然料到他如今会发怒。不过，我要当着你们的面问他，对于那些被称为神祇的人，他有什么可说的；他们的生平被传说充满各种情欲，时常向民众传颂，以效仿他们；而除了我所说的人的激情外，他们的坟墓也在不同地方被展示。}\par

% structure: p0060 source=h2 target=subsection reason=relative-heading-rank; base=h2

\subsection{第三十章 承诺的讨论}

其他人听我说了这些事，并想知道后来如何，就陪我去拜访阿皮翁。我们找到他时，他刚洗过澡，坐在一张摆满食物的桌前。因此我们很少询问关于神祇的事。但他，我猜是明白了我们的愿望，答应第二天他会就神祇发表一些见解，并且指定了同一个地方给我们交谈。他一答应，我们就感谢他，各自回家了。\par

\SourceRecord{Translated by Thomas Smith.；Ante-Nicene Fathers；Vol. 8.；Edited by Alexander Roberts, James Donaldson, and A. Cleveland Coxe.；Buffalo, NY: Christian Literature Publishing Co.,；1886.；<http://www.newadvent.org/fathers/080805.htm>.}

% structure: p0001 source=h1 target=section reason=page-title

\section{第六篇讲道}

% structure: p0002 source=h2 target=subsection reason=relative-heading-rank; base=h2

\subsection{第一章 \Person{克肋孟}遇见\Person{阿皮翁}}

\Emph{到了}第三天，我与朋友们来到\Place{提洛}约定的地方，看见\Person{阿皮翁}坐在\Person{阿努比翁}和\Person{雅典诺多洛斯}之间，还有许多其他有学识的人与他一同等候我们。我毫不惊慌，向他们问好，然后在\Person{阿皮翁}对面坐下。过了一会儿，他开始说道：\par

\Quote{我想从这一点开始，尽快直接切入问题。在我儿\Person{克肋孟}加入我们之前，我这里的友人\Person{阿努比翁}和\Person{雅典诺多洛斯}——他们昨天曾聆听你的讲论——向我报告了你所说的话，即当我在\Place{罗马}拜访你时，你假装友爱，而我却对诸神提出了许多虚假的指控：我指控他们犯有鸡奸、淫乱和各种乱伦行为。但是，我儿，你应当知道，我写那些关于诸神的事并不是认真的，而是出于对你的爱隐藏了真理。现在，如果你乐意，你可以从我这里听到那真理。}\par

% structure: p0005 source=h2 target=subsection reason=relative-heading-rank; base=h2

\subsection{第二章 神话不可按字面理解}

\Quote{古代最智慧的人，那些通过辛勤劳作通晓一切真理的人，将知识之路向那些不配且不喜受教于神圣事物的人隐藏起来。因为实际上并非真的从\Person{乌剌诺斯}及其母亲\Person{盖亚}生出了十二个孩子，如神话所计的那样：六个儿子——\Person{俄刻阿诺斯}、\Person{科俄斯}、\Person{克瑞俄斯}、\Person{许珀里翁}、\Person{伊阿珀托斯}、\Person{克洛诺斯}；六个女儿——\Person{忒亚}、\Person{忒弥斯}、\Person{谟涅摩叙涅}、\Person{得墨忒耳}、\Person{忒堤斯}和\Person{瑞亚}。也不是\Person{克洛诺斯}用金刚石镰刀阉割了他的父亲\Person{乌剌诺斯}，如你所说，并将那部分扔进海里；也不是\Person{阿佛洛狄忒}从流出的血滴中诞生；也不是\Person{克洛诺斯}与\Person{瑞亚}结合，并吞食了他首生的儿子\Person{普路同}，因为\Person{普罗米修斯}的某种说法使他担心从他所生的孩子会比他更强大，并剥夺他的王位；也不是他以同样的方式吞食了他的第二个孩子\Person{波塞冬}；也不是当\Person{宙斯}随后出生时，他的母亲\Person{瑞亚}将他藏了起来，当\Person{克洛诺斯}要\Person{宙斯}以便吞食他时，她给了他一块石头代替；也不是这块石头被吞下后，挤压了先前被吞食的那些人，迫使他们出来，所以先被吞食的\Person{普路同}先出来，然后是\Person{波塞冬}，然后是\Person{宙斯}；也不是\Person{宙斯}——如故事所说——被他母亲的智慧所救，升到天上，剥夺了他父亲的王位；也不是他惩罚了他父亲的兄弟；也不是他降临尘世贪恋凡间女子；也不是他与他的姐妹、女儿、嫂嫂相交，犯下可耻的鸡奸；也不是他吞食了他的女儿\Person{墨提斯}，以便从她那里使\Person{雅典娜}从自己的脑中诞生（并从他的大腿生出\Person{狄俄倪索斯}，据说他被\Person{提坦}撕碎）；也不是他在\Person{珀琉斯}和\Person{忒提斯}的婚礼上设宴；也不是他将\Person{厄里斯}（不和女神）排除在婚礼之外；也不是\Person{厄里斯}因受辱而在此方制造了宴饮者之间争吵和不和的机会；也不是她从\Place{赫斯珀里得斯}的花园里取了一个金苹果，在上面写上\InnerQuote{给最美的}。然后他们编造说，\Person{赫拉}、\Person{雅典娜}和\Person{阿佛洛狄忒}找到了苹果，并为此争吵，来到\Person{宙斯}面前；\Person{宙斯}没有为她们裁决，而是由\Person{赫耳墨斯}将她们送到牧人\Person{帕里斯}那里，让他评判她们的美貌。但并没有那样的女神评判；也不是\Person{帕里斯}将苹果给了\Person{阿佛洛狄忒}；也不是\Person{阿佛洛狄忒}因受此尊荣，就以赐予\Person{海伦}为妻来尊荣他。因为女神所赐的尊荣绝不可能成为一场普世战争的借口——而那战争导致受尊荣者——他本人与\Person{阿佛洛狄忒}的种族有近亲关系——的毁灭。但是，我儿，正如我所说，这样的故事有一种独特的哲学意义，可以通过寓言的方式展现出来，连你听了也会惊叹。} 我说：\Quote{我恳求你不要拖延折磨我。} 他说：\Quote{不要害怕；我不会浪费时间，立即开始。}\par

% structure: p0007 source=h2 target=subsection reason=relative-heading-rank; base=h2

\subsection{第三章 \Person{阿皮翁}开始解释神话}

\Quote{曾有一个时期，除了混沌和秩序混乱元素的混杂——它们只是简单地堆积在一起——之外，别无他物。本性见证此事，伟大的人物也认为情况如此。在这些伟人中，我将向你介绍那在智慧上超越所有人的一位，即\Person{荷马}，他在谈及原初混乱时说道：\InnerQuote{但愿你全部变成水和土；}暗示万物皆源于此，并且当水和土的元素分离时，万物回归其最初状态，即混沌。而\Person{赫西俄德}在\WorkTitle{神谱}中说：\InnerQuote{确实，混沌是最初生成的存在。}现在，通过\InnerQuote{生成}，他显然意指混沌是生成的，即有开端，而非无始永存。而\Person{奥尔甫斯}将混沌比作一个蛋，其中蕴含着原初元素的混乱混合。这个\Person{奥尔甫斯}称为蛋的混沌，为\Person{赫西俄德}所假定，是有开端的，产生于无限物质，并以如下方式起源。}\par

% structure: p0009 source=h2 target=subsection reason=relative-heading-rank; base=h2

\subsection{第四章 混沌的起源}

\Quote{这种物质分为四种，且具有生命，是一个整个无限的深渊，可以说，在永恒之流中，没有秩序地漂流，时时刻刻形成无数但无效的组合（因此由于缺乏秩序又再次解散）；它确实成熟了，但无法被约束以产生生物。偶然有一次，这个无限的海洋，以其本性被自然运动所驱动，有序地从相同流向相同（折回自身），如同漩涡，混合物质，使得从每种物质中，宇宙中间（如同在模具的漏斗中）流下了恰好对产生生物最有用和最合适的物质。这被万有承载的漩涡带下，吸引了周围的灵，并且如此受孕，非常肥沃，形成了一个分离的物质。正如通常在水中形成气泡一样，周围的一切都有助于这个球状圆球的受孕。然后，在它自身受孕之后，被围绕它的天主之灵托起，也许产生了最伟大的东西；可以说，一件作品，具有从整个无限深渊受孕而来的生命，形状像蛋，速度像鸟。}\par

% structure: p0011 source=h2 target=subsection reason=relative-heading-rank; base=h2

\subsection{第五章 克洛诺斯与瑞亚的解释}

\Quote{现在你必须把克洛诺斯理解为\OriginalTerm{时间}{chronos}，把瑞亚理解为水样物质的\OriginalTerm{流动}{rheon}。因为整个物质体在产生如蛋般的球状、包罗万象的\OriginalTerm{天}{ouranos}之前，漂流了一段时间，起初充满了生产性的骨髓，因此能够从自身产生各种元素和颜色，而从一种物质和一种颜色产生各种形态。正如孔雀蛋似乎只有一种颜色，但潜在地具有将要出现的动物所有颜色一样，这个从无限物质受孕的活蛋，由基底的、永恒流动的物质推动，产生许多不同的形式。因为在圆周内，一个既是男性又是女性的活物，由内住的天主之灵的技术形成。俄耳甫斯称它为\OriginalTerm{法涅斯}{Phanes}，因为它\OriginalTerm{显现}{phaneis}时，宇宙从它发出光辉，带着那最荣耀的元素——火——在潮湿中完善的光泽。这并非不可信，因为在萤火虫中，大自然让我们看到一种潮湿的光。}\par

% structure: p0013 source=h2 target=subsection reason=relative-heading-rank; base=h2

\subsection{第六章 法涅斯与普路托}

\Quote{这个蛋，即最初的物质，稍稍变热，被里面的活物打破，然后形成并出现了某种东西；正如俄耳甫斯所说的，当\InnerQuote{宽大的蛋被打破}时，等等。因此，藉那\OriginalTerm{显现}{phaneis}者并出现者的强大力量，这个球体获得了连贯性并保持了秩序，而它本身，可以说，坐在天顶上，在那里以难以言喻的奥秘永远放射光芒。但留在球体内部的生产性物质，分离了万物的实质。因为首先，它的下部，就像渣滓一样，由于自身的重量而下沉；他们称之为\OriginalTerm{普路托}{Pluto}，因为其厚重、重量和大量底层的物质（\OriginalTerm{多}{polu}），称它为\OriginalTerm{哈得斯}{Hades}和死者之王。}\par

% structure: p0015 source=h2 target=subsection reason=relative-heading-rank; base=h2

\subsection{第七章 波塞冬、宙斯与墨提斯}

\Quote{那么，当他们说这种原初物质，虽然最肮脏粗糙，被克洛诺斯（即时间）吞噬时，这应从自然角度理解为它向下沉。而在第一次沉淀之后汇聚起来并漂浮在第一种物质表面的水，他们称之为\OriginalTerm{波塞冬}{Poseidon}。然后剩下的，最纯洁最高贵的一切，因为它是透明的火，他们称之为\OriginalTerm{宙斯}{Zeus}，因其\OriginalTerm{炽热}{zeousa}的本性。既然火上升，它没有被时间或克洛诺斯吞噬而下降；但是，正如我所说，火的物质，由于具有生命并自然上升，直接飞入空气，空气因其纯净而非常智能。因此，藉着他自身的热，宙斯——即炽热的物质——提起了残留在底层潮湿中的东西，即那非常强大和神圣的灵，他们称之为\OriginalTerm{墨提斯}{Metis}。}\par

% structure: p0017 source=h2 target=subsection reason=relative-heading-rank; base=h2

\subsection{第八章 帕拉斯与赫拉}

\Quote{当这到达以太的顶端时，被它（可以说潮湿与热混合）吞噬；并引起它不停的跳动，产生了智能，他们称之为\OriginalTerm{帕拉斯}{Pallas}，来自这种\OriginalTerm{跳动}{pallesthai}。这就是艺术智慧，以太工匠藉此建造了整个世界。而从遍及万物的宙斯，即这非常热的以太，\OriginalTerm{空气}{aer}延伸到我们大地；他们称之为\OriginalTerm{赫拉}{Hera}。因此，因为它出现在以太（最纯净的物质）之下（正如女人在纯洁方面逊色），当两者比较优劣时，她正确地被视为宙斯的姐妹，就起源于同一物质而言，但作为他的配偶，因为这个像妻子一样低等。}\par

% structure: p0019 source=h2 target=subsection reason=relative-heading-rank; base=h2

\subsection{第九章 阿尔忒弥斯}

\Quote{我们理解\OriginalTerm{赫拉}{Hera}是大气的一个愉悦调和，因此她非常多产；但\OriginalTerm{雅典娜}{Athena}，即他们所谓的\OriginalTerm{帕拉斯}{Pallas}，被认为是处女，因为由于高热她无法产生任何东西。类似地，\OriginalTerm{阿尔忒弥斯}{Artemis}也被解释：他们把她当作空气的最深层，因此称她为处女，因为她由于极端寒冷无法承载任何东西。那由上层和下层蒸汽产生的困扰而醉酒的混合物，他们称之为\OriginalTerm{狄俄倪索斯}{Dionysus}，因为扰乱理智。地下的水，其本性是一，但流经所有地脉，分为许多部分，他们称之为\OriginalTerm{奥西里斯}{Osiris}，因为被切割成碎片。他们把\OriginalTerm{阿多尼斯}{Adonis}理解为有利的季节，\OriginalTerm{阿佛洛狄忒}{Aphrodite}为交配和生殖，\OriginalTerm{得墨忒耳}{Demeter}为大地，\OriginalTerm{珀耳塞福涅}{Proserpine}为种子；有些把\OriginalTerm{狄俄倪索斯}{Dionysus}理解为葡萄树。}\par

% structure: p0021 source=h2 target=subsection reason=relative-heading-rank; base=h2

\subsection{第十章 所有此类故事皆为寓意}

\Quote{并且我必须请求你们认为所有这些故事都包含某种寓意。请将\Person{阿波罗}视为巡行不息的太阳（\OriginalTerm{巡行者}{peri-polôn}），他是\Person{宙斯}的一个儿子，也被称为\Person{密特拉}，因为他完成一年的周期。而这位无所不在的\Person{宙斯}的所谓变化必须被看作是季节的众多变迁，而他那数不清的妻子们你们必须理解为岁月或世代。因为从以太之中发出、经过空气而运行的能力，依次与所有的岁月和世代结合，不断地使它们变化，从而产生或毁坏庄稼。成熟的果实被称为他的孩子们，某些季节的荒芜则被归因于不法的结合。}\par

% structure: p0023 source=h2 target=subsection reason=relative-heading-rank; base=h2

\subsection{第十一章 克莱孟以前就听过这些}

当\Person{阿皮翁}这样进行寓意解释时，我陷入了沉思，似乎没有跟上他的话。于是他打断了话头，对我说：\Quote{如果你不跟上我所说的，我还有什么必要说呢？} 我回答说：\Quote{请不要以为我不懂你说的话。我完全理解；而且正因如此，这不是我第一次听到这些。为了让你知道我对这些事情并不陌生，我将概述你所说的，并按顺序补充你遗漏的那些故事的寓意解释，这些解释是我从别人那里听到的。} \Person{阿皮翁}说：\Quote{请这样做。}\par

% structure: p0025 source=h2 target=subsection reason=relative-heading-rank; base=h2

\subsection{第十二章 阿皮翁解释的概要}

我回答说：\Quote{我目前不特别谈论那个活蛋，它是由无限物质经过幸运的组合而孕育的，当它破碎时，雌雄同体的\Person{法内斯}跳了出来，正如有些人所说。关于这一切我只说一点，直到这个破碎的球体获得凝聚性，其中仍留有一些类似骨髓的物质；我将简要叙述凭借这种物质在其中发生的一切及后续。因为从\Person{克罗诺斯}和\Person{瑞亚}出生，如你所说——即通过时间和物质——首先是\Person{普路托}，他代表沉淀下来的沉积物；然后是\Person{波塞冬}，中间的液体物质，漂浮在下方较重的物体之上；第三个孩子——即\Person{宙斯}——是以太，是至高者。他没有被吞食；而是作为一种火性的能力，自然上升，一跃而至最高的以太。}\par

% structure: p0027 source=h2 target=subsection reason=relative-heading-rank; base=h2

\subsection{第十三章 克罗诺斯与阿芙洛狄忒}

\Quote{而\Person{克罗诺斯}的锁链是天地结合，正如我听说别人寓意解释的那样；他的阉割是元素的分裂与分离；因为万物都按其各自的本性被切断和分开，以便每一种类能自行排列。时间不再生出任何东西；但由时间生出的事物，按照自然法则，产生它们的后继者。而从海中出现的\Person{阿芙洛狄忒}是从水分中涌现出的丰饶实体，温暖的精神与之混合，引起那种性欲，并完善了世界之美。}\par

% structure: p0029 source=h2 target=subsection reason=relative-heading-rank; base=h2

\subsection{第十四章 佩琉斯与忒提斯、普罗米修斯、阿喀琉斯与波吕克塞娜}

\Quote{而那场婚宴，即\Person{宙斯}为水仙女\Person{忒提斯}与英俊的\Person{佩琉斯}举行盛宴的场合，包含这样的寓意——让你知道，\Person{阿皮翁}，你不是我唯一听到这类事情的人。宴会就是世界，十二位是命运的天柱，称为黄道带。\Person{普罗米修斯}是远见（\OriginalTerm{远见}{prometheia}），万物由此而生；\Person{佩琉斯}是黏土（\OriginalTerm{黏土}{pelos}），即从大地聚集并与水混合以产生人的东西；从这两者（即水和土）的混合中，第一个后代不是被生育出来，而是被塑造完整，称为\Person{阿喀琉斯}，因为他从未将嘴唇（\OriginalTerm{嘴唇}{cheile}）凑到乳房上。正值年华正茂，他被一箭射死，当时他渴望拥有\Person{波吕克塞娜}，即某种与真理不同且疏远（\OriginalTerm{疏远}{xene}）的东西，死亡从他的脚伤悄然侵入。}\par

% structure: p0031 source=h2 target=subsection reason=relative-heading-rank; base=h2

\subsection{第十五章 帕里斯的审判}

\Quote{然后\Person{赫拉}、\Person{雅典娜}、\Person{阿芙洛狄忒}、\Person{厄里斯}、苹果、\Person{赫尔墨斯}、审判和牧羊人都隐藏着如下意义：\Person{赫拉}代表尊严；\Person{雅典娜}，男子气概；\Person{阿芙洛狄忒}，快乐；\Person{赫尔墨斯}，语言，它解释（\OriginalTerm{解释性的}{hermeneutikos}）思想；牧羊人\Person{帕里斯}，不理性的兽性激情。如果在盛年之时，理性——灵魂的牧羊人——变得兽性，不顾自身利益，拒绝与男子气概和节制同流，只选择快乐，将奖品仅给予欲望，讨价还价要从欲望那里得到能使它快乐的东西，那么这样错误判断的人将选择快乐以致自毁并毁掉朋友。\Person{厄里斯}是嫉妒的恶意；赫斯珀里得斯的金苹果也许是财富，连像\Person{赫拉}那样节制的人有时也被诱惑，像\Person{雅典娜}那样有男子气概的人也变得嫉妒，从而做出不合身份的事，而像\Person{阿芙洛狄忒}那样的灵魂之美在文雅的伪装下被摧毁。简而言之，在所有人中，财富引发邪恶的不和。}\par

% structure: p0033 source=h2 target=subsection reason=relative-heading-rank; base=h2

\subsection{第十六章 赫拉克勒斯}

\Quote{而\Person{赫拉克勒斯}，他杀死了带领并守护财富的蛇，是真正的哲学理性，它摆脱了一切邪恶，走遍世界，探访人的灵魂，惩罚它所遇到的一切——即像凶猛狮子、胆小牡鹿、野蛮野猪或多头水蛇的人；\Person{赫拉克勒斯}其他所有传说般的劳动也都隐晦地指向道德勇气。但这些例子就够了，因为如果我们逐一叙述，所有的时间也不够。}\par

% structure: p0035 source=h2 target=subsection reason=relative-heading-rank; base=h2

\subsection{第十七章 编造这些故事的人应受责备}

\Quote{既然这些事可以清楚、有益而毫不损害虔诚地，以公开和直率的方式陈述，我奇怪你竟称那些人为明智和有智慧，他们却用弯曲的谜语隐藏了这些事，并用污秽的故事覆盖它们，从而，好像被邪灵驱使，欺骗了几乎所有人。因为要么这些事并非谜语，而是诸神真实的罪行，在这种情况下，它们不应被置于轻蔑之下，也不应将他们的这些需要完全作为模范摆在人面前；要么那些虚假归给诸神的事是用寓言陈述的，那么，\Person{阿皮翁}，你称为有智慧的人就错了，因为他们用不值的故事掩盖本身有价值的事，引导人犯罪，而且并非不羞辱他们所相信是神祇的那些。}\par

% structure: p0037 source=h2 target=subsection reason=relative-heading-rank; base=h2

\subsection{第十八章　同前}

\Quote{因此，你不要以为他们是智慧的人，而毋宁是邪灵，他们能用邪恶的故事掩盖高尚的行为，目的是让那些愿意效仿更好者的人，能够效仿这些所谓神祇的行为——就是昨天我在讲论中自由谈论的那些行为：弑父、杀子、各种乱伦、无耻的通奸和无数的污秽。其中最不虔敬的人，是那些希望这些故事被相信的人，以便他们自己做出同样的事时不会羞愧。如果他们有心行事虔敬，正如我刚才所说，即使神祇们真的做了那些歌咏他们的事，也应该用更体面的故事掩盖他们的不雅，而非像你所说，当神祇的行为是高尚的时，却把它们披上邪恶和不雅的形式，这些形式即使被解释，也只能通过大量辛劳才能理解；而当它们被某些人理解时，他们虽然因辛劳获免于受骗——其实不劳苦也能有这特权——而受骗的人却彻底沉沦了。（至于那些把寓言追溯到更尊贵源头的人，我不反对；例如，那些解释一个寓言时说智慧从\Person{宙斯}的头颅中生出的人。）总之，在我看来更可能的是：邪恶的人，剥夺了神祇的荣耀，竟敢传播这些侮辱性的故事。}\par

% structure: p0039 source=h2 target=subsection reason=relative-heading-rank; base=h2

\subsection{第十九章　这些寓言没有一个是连贯的}

\Quote{我们发现关于任何神祇的诗体寓言都与自身不一致。仅就宇宙的塑造而言，诗人们时而说自然是整个受造界的第一因，时而说是心智。因为，他们说，元素的最初运动和混合来自自然，但是心智的预见安排了它们的秩序。即使他们断言是自然塑造了宇宙，由于不能因作品中设计的痕迹而完全证明这一点，他们又织入心智的预见，从而能捕捉到最智慧的人。但我们对他们说：如果世界起源于自我运动的自然，它如何获得比例和形状？这只能来自监管的智慧，并且只能被知识理解，唯有知识能追踪这些事。另一方面，如果万物因智慧而存在并保持秩序，那些事物又如何由自我运动的偶然产生呢？}\par

% structure: p0041 source=h2 target=subsection reason=relative-heading-rank; base=h2

\subsection{第二十章　这些神祇其实是邪恶的术士}

\Quote{于是，那些选择对神圣事物作不名誉的寓言的人——例如，说\Person{墨提斯}被\Person{宙斯}吞下——陷入了两难，因为他们没有看到，在这些关于神祇的故事中，他们间接地教导了物理学，却否认了神祇本身的存在，把各种神祇变成宇宙各种实体的纯粹寓言式代表。因此更可能的是，这些人所颂扬的神祇是某种邪恶的术士，他们实际上是邪恶的人，却用魔术改变形状，犯下通奸，夺取生命，从而对不懂魔术的古人来说，因他们所行的事而被视为神祇；这些人的身体和坟墓在许多城镇仍可见到。}\par

% structure: p0043 source=h2 target=subsection reason=relative-heading-rank; base=h2

\subsection{第二十一章　他们的坟墓至今仍可见到}

\Quote{例如，如我先前已提到的，在\Place{高加索山脉}中，有一个叫\Person{克洛诺斯}的人的坟墓，他是个凶猛的君主，杀了自己的子女。这个人的儿子叫\Person{宙斯}，变得\Emph{比他的父亲更甚}；他凭魔术的力量被宣布为宇宙的统治者，犯下许多通奸，并惩罚了他的父亲和叔伯，然后死去；克里特人展示他的坟墓。在\Place{美索不达米亚}，一个叫\Person{赫利奥斯}的人葬于\Place{阿提尔}，一个叫\Person{塞勒涅}的人葬于\Place{卡尔赫}。一个叫\Person{赫尔墨斯}的人葬于\Place{埃及}；\Person{阿瑞斯}葬于\Place{色雷斯}；\Person{阿佛洛狄忒}葬于\Place{塞浦路斯}；\Person{阿斯克勒庇俄斯}葬于\Place{埃皮达鲁斯}；许多其他这类人的坟墓也可见到。}\par

% structure: p0045 source=h2 target=subsection reason=relative-heading-rank; base=h2

\subsection{第二十二章　因此，他们的同时代人并不视他们为神}

\Quote{因此，对正直思考的人来说，显然他们被承认为必死者。他们的同时代人，知道他们是必死的，在他们死后便不再留意；是漫长的时间给他们披上了神祇的光荣。你不必惊讶，生活在\Person{阿斯克勒庇俄斯}和\Person{赫拉克勒斯}时代的人受了骗，或\Person{狄俄尼索斯}或当时任何其他人的同时代人，甚至\Place{伊利昂}的\Person{赫克托耳}和\Place{琉刻}岛上的\Person{阿喀琉斯}，也被当地居民崇拜；\Place{奥浦斯}人崇拜\Person{帕特罗克洛斯}，\Place{罗得岛}人崇拜\Person{马其顿的亚历山大}。}\par

% structure: p0047 source=h2 target=subsection reason=relative-heading-rank; base=h2

\subsection{第二十三章　埃及人向一个人奉献神圣的敬礼}

\Quote{此外，在埃及人中，直至今日，有人在生前就被当作神受人敬拜。这确实是一种小不虔敬，即埃及人在人有生之年给予他神圣的尊荣；但最荒谬的是，他们敬拜鸟类、爬虫和各类走兽。因为大多数人既无理智，行事也无分寸。但我请求你们看看，最可耻的是：他们称为众神和众人之父的那一位，据他们说曾与勒达交合；他们中的许多人公开陈列这样一幅画，在上面写上宙斯的名字。为了惩罚这种侮辱，我希望他们画下他们现在的国王身处如此卑贱的拥抱中，就像他们对宙斯所大胆做的那样，并公之于众，这样，从一位暂时的君主，而且是一位凡人的愤怒中，他们或许能学会向应得尊荣者致敬。我说这些给你们，并非因为我自己已经认识真天主；但我很高兴地说，即使我不知道谁是天主，我想我至少清楚地知道天主是什么。}\par

% structure: p0049 source=h2 target=subsection reason=relative-heading-rank; base=h2

\subsection{第二十四章 什么不是天主}

\Quote{那么，首先，四大原始元素不可能是天主，因为它们有原因。混合、复合、生成、以及环绕可见宇宙的球体，都不是天主；在冥府中聚集的渣滓，漂浮其上的水，也不是天主；火质，以及从它延伸到我们大地的空气，也不是天主。因为四大元素，如果彼此分离，就不可能在没有一位伟大工匠的情况下混合在一起以产生生命。如果它们始终联合，即便如此，它们也是由艺术心智配合起来，以适应动物肢体和部分所需，使它们能够保持各自的比例，具有明确界定的形状，并且所有内部部分都能达到恰当的连贯性。同样，每个部分适合的位置也是由工匠心智非常优美地决定的。简而言之，在一个生物所必需的一切其他事物中，这个伟大的世界存在物毫无欠缺。}\par

% structure: p0051 source=h2 target=subsection reason=relative-heading-rank; base=h2

\subsection{第二十五章 宇宙是理智的产物}

\Quote{因此，我们不得不假设有一位无始的工匠，如果元素是分离的，他使之结合；如果它们原本在一起，他艺术地混合它们以产生生命，并从所有中完善出一件作品。因为一件完全智慧的作品，不可能在没有一个比它更大的心智的情况下被制造出来。也不能说爱是万物的工匠，或是欲望、权力或任何这类东西。所有这些都易变，而且在本质上短暂。那被另一者推动的，更不可能成为天主，尤其那被时间和本性所改变、可以被消灭的，更不是天主。}\par

% structure: p0053 source=h2 target=subsection reason=relative-heading-rank; base=h2

\subsection{第二十六章 伯多禄从凯撒勒雅抵达}

当我对阿皮翁说这些话时，伯多禄从凯撒勒雅走近了；在提洛，人们聚集起来，匆忙去迎接他，并一同表达对他来访的喜悦。阿皮翁退下了，只有阿努比翁和雅典诺多洛斯陪同；但我们其余的人都急忙去迎接伯多禄，我在城门第一个向他致意，并带他向客栈走去。我们到达后，遣散了众人；当他屈尊询问所发生的事时，我什么也没有隐瞒，而是告诉他西满的诽谤，他所变的怪异形状，以及在献祭宴后他送来的一切疾病，还有一些病人仍在提洛，而另一些人在我到达时就已随西满前往西顿，希望被他治愈，但我听说没有人被他治愈。我还告诉了伯多禄我与阿皮翁的辩论；他出于对我的爱，并想鼓励我，称赞并祝福了我。然后，用过晚餐后，他休息去了，这是他旅途劳顿所必需的。\par

\SourceRecord{Translated by Thomas Smith.；Ante-Nicene Fathers；Vol. 8.；Edited by Alexander Roberts, James Donaldson, and A. Cleveland Coxe.；Buffalo, NY: Christian Literature Publishing Co.,；1886.；<http://www.newadvent.org/fathers/080806.htm>.}

% structure: p0001 source=h1 target=section reason=page-title

\section{第七篇讲道}

% structure: p0002 source=h2 target=subsection reason=relative-heading-rank; base=h2

\subsection{第一章 \Person{伯多禄}向群众讲话}

\Emph{并且}我们停留在\Place{提洛}的第四天，黎明时分，\Person{伯多禄}外出，有不少周围的居民，连同许多\Place{提洛}本城的居民迎上前来，他们呼喊道，\Quote{天主藉着你怜悯我们，天主藉着你医治我们！}于是\Person{伯多禄}站在一块高石上，好让所有人都能看见他；他以一种虔敬的方式问候了他们，然后这样开始：——\par

% structure: p0004 source=h2 target=subsection reason=relative-heading-rank; base=h2

\subsection{第二章 \Person{西满}能力的缘由}

\Quote{天主，创造了诸天和整个宇宙，并不需要为那些愿意得救的人提供救恩的机会。因此，谁也不要因着看似邪恶的事情，轻率地指责祂对人缺乏慈爱。因为人不知道临到他们之事的结局，甚至猜疑结果将是邪恶的；但天主知道它们会转好。\Person{西满}的情况也是如此。他是\OriginalTerm{天主的左手}{left hand of God}的势力，有权柄加害那些不认识天主的人，因此他能使你们陷入疾病；但这些疾病，是借着\OriginalTerm{天主良善的眷顾}{good providence}允许临到你们的，你们因着寻求并找到了那能治愈者，就被迫借着身体的治愈而服从天主的旨意，并想到要相信，如此你们的灵魂也能像身体一样处于健康状态。}\par

% structure: p0006 source=h2 target=subsection reason=relative-heading-rank; base=h2

\subsection{第三章 疗愈}

\Quote{现在有人告诉我，他在献祭了一头公牛之后，在广场中间宴请了你们，你们因多饮了酒而沉醉，不仅与 \OriginalTerm{恶魔}{demons}，而且与它们的\OriginalTerm{它们的王子}{their prince}也成了朋友，就这样你们大多数人都被这些疾病抓住，不知不觉亲手给自己招来了毁灭之剑。因为恶魔决不会对你们有权势，除非你们先与它们的王子共餐。因为从起初，万有的创造者天主就为两个王子，\OriginalTerm{右手的}{right hand}和\OriginalTerm{左手的}{left hand}，立了一条法律，就是除非先与它们同桌吃饭，否则它们都无权对任何它们想要加害或加惠的人行使权力。因此，正如你们分享了\OriginalTerm{祭偶像的肉}{meat offered to idols}，就成了\OriginalTerm{邪恶王子}{the prince of evil}的仆人，同样，如果你们远离这些事，借着祂\OriginalTerm{右手的那位好王子}{the good Prince of His right hand}逃到天主那里，不用\OriginalTerm{祭献}{sacrifices}，只做祂所愿意的一切，来尊崇祂，你们要确知，不仅你们的身体要得医治，你们的灵魂也要变得健康。因为唯独祂能用左手毁灭，用右手复苏；唯独祂能打击，也能扶起跌倒的。}\par

% structure: p0008 source=h2 target=subsection reason=relative-heading-rank; base=h2

\subsection{第四章 黄金法则}

\Quote{所以，正如你们先前被先行者\Person{西满}所欺骗，灵魂上对天主成了死的，身体受到打击；同样，现在如果你们悔改，正如我所说的，并服从那些令天主喜悦的事，你们或许能使身体获得新力量，并恢复灵魂的健康。令天主喜悦的事是：向祂祈祷，向祂祈求，承认祂是万物的赐予者，并按着有区别的法律施予；远离\OriginalTerm{魔鬼的筵席}{the table of devils}，不尝\OriginalTerm{死肉}{dead flesh}，不摸血；洗净一切污秽；其余一言以蔽之——正如\OriginalTerm{敬畏天主的犹太人}{the God-fearing Jews}已经听到的，你们也要听到，并在众多身体中同心合意；愿每个人存心对\OriginalTerm{近人}{neighbour}行那些他希望别人为自己行的善事。你们都可以通过与自己进行这样的对话来发现什么是善：你不愿被人谋杀，就不要谋杀他人；你不愿你的妻子被人勾引，就不要犯奸淫；你不愿任何东西被人偷走，就不要偷别人的东西。因此，你们凭自己明白什么是合理的，并去做，你们就会成为天主所爱的，并获得医治；否则，在现世生活中你们的身体将受折磨，在来世你们的灵魂将受惩罚。}\par

% structure: p0010 source=h2 target=subsection reason=relative-heading-rank; base=h2

\subsection{第五章 \Person{伯多禄}起程往\Place{息冬}}

\Person{伯多禄}这样教导了他们并医治了他们几天之后，他们就\OriginalTerm{领受了圣洗}{were baptized}。此后，众人一同坐在广场上，穿着\OriginalTerm{麻布}{sackcloth}，撒\OriginalTerm{灰}{ashes}，因着他的其他奇事而忧伤，并忏悔他们从前的\OriginalTerm{罪恶}{sins}。\Place{息冬}的人听到后，也同样行了，并派人恳求\Person{伯多禄}，因为他们自己因疾病不能前来。\Person{伯多禄}在\Place{提洛}没有停留很多天；当他教导了所有居民，使他们摆脱各种疾病，并建立了一个\OriginalTerm{教会}{church}，立了与他同来的一位\OriginalTerm{长老}{elders}为管理该教会的\OriginalTerm{主教}{bishop}后，他就起程往\Place{息冬}去了。但\Person{西满}听说\Person{伯多禄}要来，就立刻与\Person{阿皮昂}和他的朋友们逃往\Place{贝鲁特}了。\par

% structure: p0012 source=h2 target=subsection reason=relative-heading-rank; base=h2

\subsection{第六章 \Person{伯多禄}在\Place{息冬}}

\Person{伯多禄}进入\Place{息冬}时，许多人用担架抬着病人，放在他面前。他对他们说：\Quote{请你们不要想，我求你们，我能做什么来医治你们，我是一个必死的人，自己也会遭受许多邪恶。但我不会拒绝向你们指出你们必须得救的道路。因为我从\OriginalTerm{真理的先知}{the Prophet of truth}那里学到了天主在创世以前预定的条件；也就是说，如果人行恶事，祂就预定他们必受\OriginalTerm{邪恶王子}{the prince of evil}的伤害；同样，如果人行善事，祂就颁令，凡信祂为\OriginalTerm{他们的医生}{their Physician}的人，他们的身体必得痊愈，灵魂必得安全。}\par

% structure: p0014 source=h2 target=subsection reason=relative-heading-rank; base=h2

\subsection{第七章 两条道路}

\Quote{那么，既然知道了这些善行与恶行，我向你们指明两条道路，并要指示你们哪些旅人迷失，哪些得救，因天主的引导。迷失者的道路是宽阔且极其平坦的——它毁了人却毫不困扰他们；但得救者的道路是狭窄、崎岖的，最终救赎那些历经艰辛行走其上的人。这两条道路由不信和信德掌管；那些选择了享乐、因此忘记了审判之日、行天主所不喜悦的事、不关心借圣言拯救自己的灵魂、也不急切寻求自身好处的人，走的是不信的道路。他们实在不知道，天主的计划不像人的计划；因为，首先，祂知道所有人的思念，所有人都必须交代自己的行为，甚至思念。那些努力正确理解却失足者的罪，远小于那些根本不追求好事之人的罪。因为天主喜欢：按人的标准估量，在认识善事上有错误的人，在略受惩罚后得以得救。但那些毫不关心认识更好道路的人，即使他们做了无数其他善事，如果没有遵守祂亲自设立的侍奉，就会被视为冷漠，受严厉惩罚，全然毁灭。}\par

% structure: p0016 source=h2 target=subsection reason=relative-heading-rank; base=h2

\subsection{第八章  天主设立的侍奉}

\Quote{祂设立的侍奉就是：只敬拜祂，只信赖真理的先知，领受圣洗以得罪赦，并借这纯净的洗礼由拯救之水重生归于天主；戒绝魔鬼的筵席，即祭偶像的食物、死尸、窒息而死的或被野兽捕获的动物、以及血；不再不洁地生活；交合后要沐浴；妇女要遵守洁净律；所有人都要清心寡欲，专务善行，远离恶事，期待全能天主所赐的永生，并借祈祷和恒切恳求来争取这永生。}这也是伯多禄对漆冬人所说的话。几天后，许多人悔改并相信了，得到了治愈。伯多禄建立了一个教会，并将与他同来的长老中的一位立为主教，然后离开了漆冬。\par

% structure: p0018 source=h2 target=subsection reason=relative-heading-rank; base=h2

\subsection{第九章  西满攻击伯多禄}

他刚到贝鲁特，就发生了一场地震；群众跑到伯多禄面前说：\Quote{救救我们，因为我们害怕会全部灭亡。}这时，西满与阿皮翁、阿努比翁、阿忒诺多罗斯及其同伴一起，大胆公开向群众喊叫，反对伯多禄：\Quote{朋友们，逃离这个人！他是个术士；相信我们，是他引起了这场地震；他给我们带来了这些疾病来恐吓我们，好像他自己就是天主一样。}西满和他的朋友们向伯多禄提出了许多此类虚假指控，仿佛他能行超乎人力之事。但当群众稍一安静，伯多禄便以惊人的勇气微微一笑，说：\Quote{朋友们，我承认，若天主愿意，我能做这些人所说的事；不仅如此，如果你们不信我所说的话，我甚至能把这城完全颠覆。}\par

% structure: p0020 source=h2 target=subsection reason=relative-heading-rank; base=h2

\subsection{第十章  西满被驱逐}

人们害怕了，答应遵行他所吩咐的一切。\Quote{那么，你们谁也不可，}伯多禄说，\Quote{与这些术士交往，或与他们有任何来往。}人们一听到这简洁的命令，就拿起棍棒追赶他们，直到把他们完全赶出城。患病和被鬼附的来到伯多禄脚前跪下。伯多禄看见这一切，急于使他们摆脱恐惧，便对他们说：——\par

% structure: p0022 source=h2 target=subsection reason=relative-heading-rank; base=h2

\subsection{第十一章  得救之道}

\Quote{即便我能引发地震、行我愿行的一切，我向你们保证，我不会毁灭西满和他的朋友们（因为我来不是为毁灭人），而是要使他成为我的朋友，好让他不再因诽谤我宣讲真理而阻碍许多人的得救。但如果你们相信我，他本人就是术士、诽谤者、对不认识真理者的邪恶使者。因此，他有能力把疾病带给罪人，因为有罪人本身助长了他对他们的权力。但我是造物主天主的仆人，是祂右边先知的弟子。所以，作为祂的宗徒，我宣讲真理：为服侍善人，我驱除疾病，因为我是祂的第二使者——首先疾病来临，之后是医治。借着那作恶的术士，你们因背叛天主而患病；借着我的宣讲，如果你们相信祂，必得痊愈。这样，你们既经验到祂的能力，便可以转向善行，使你们的灵魂得救。}\par

% structure: p0024 source=h2 target=subsection reason=relative-heading-rank; base=h2

\subsection{第十二章  伯多禄前往比布罗斯和的黎波里}

他说这些话时，众人都跪在他脚前。他举手向天，向天主祈祷，仅凭简单的祈祷就治好了他们所有人。他在贝鲁特没有停留多少天；在使许多人习惯了侍奉唯一的天主、为他们施洗、并从与他同来的长老中为他们立了一位主教之后，他去了比布罗斯。到达那里后，他得知西满没有等他们一天，就直接去了的黎波里，于是他只在那里停留了几天；在治愈了不少人、并用圣经训练他们之后，他循着西满的踪迹前往的黎波里，宁愿追赶他而不愿逃避他。\par

\SourceRecord{Translated by Peter Peterson.；Ante-Nicene Fathers；Vol. 8.；Edited by Alexander Roberts, James Donaldson, and A. Cleveland Coxe.；Buffalo, NY: Christian Literature Publishing Co.,；1886.；<http://www.newadvent.org/fathers/080807.htm>.}

% structure: p0001 source=h1 target=section reason=page-title

\section{讲道词第八}

% structure: p0002 source=h2 target=subsection reason=relative-heading-rank; base=h2

\subsection{第一章 伯多禄抵达特里波利斯}

\Emph{如今}，当伯多禄进入特里波利斯时，来自提洛、漆冬、贝里图斯和比布鲁斯的人们，他们渴望得到教导，还有许多附近的人，与他一同进入；此外，城里也有许多群众聚集，希望见到他。因此，在郊区，那些由他派出去的弟兄们与我们会合，他们去查明城中的详细情况，尤其是西满的所作所为，并来向他报告。他们迎接了他，并带他到马罗内斯的住所。\par

% structure: p0004 source=h2 target=subsection reason=relative-heading-rank; base=h2

\subsection{第二章 伯多禄的体贴}

但是，当他刚到住所门口时，他转过身来，向群众承诺第二天要和他们谈论宗教事宜。他进去后，先遣者为与他同来的人安排了住处。主人和接待者没有辜负那些寻求款待之人的期望。但伯多禄对此一无所知，我们请他进食，他说在与他同来的人安顿好之前，他自己不会进食。我们向他保证这件事已经办妥，因为所有人都因对他的爱戴而热切地接待了他们，以至于那些没有客人可招待的人感到极其难过。伯多禄听了这话，对他们的热忱慈善感到欣慰，祝福了他们，然后出去在海中沐浴，并与先遣者一同进食；随后，傍晚来临，他便睡了。\par

% structure: p0006 source=h2 target=subsection reason=relative-heading-rank; base=h2

\subsection{第三章 被打断的谈话}

大约在第二次鸡叫时，他醒来，发现我们已起床。我们总共十六人：伯多禄本人、我克莱孟、尼塞塔斯和阿奎拉，以及先于我们到达的十二人。他向我们致意后，说：\Quote{今天，不因外面的事忙碌，我们可以自由地彼此交谈。因此，我要告诉你们你们离开提洛后所发生的事；你们也要详细向我报告西满在这里的所作所为。}于是，当我们彼此以叙述回应时，一位朋友进来，向伯多禄报告说，西满得知他到来后，已动身前往叙利亚；而群众觉得这一夜如同一年之久，无法再等待他所定的日期，便站在门前，三五成群地谈论西满所提出的控告，以及西满如何引起他们的期望，并承诺当他来到时会指控伯多禄许多恶事，却在他得知伯多禄到来后连夜逃走了。\Quote{可是，}他说，\Quote{他们都急切地想听你讲话；我不知道他们从哪里听到某种传闻，说你今天要对他们讲话。因此，为了免得他们非常疲惫却被毫无理由地打发走，你自己知道该怎么办。}\par

% structure: p0008 source=h2 target=subsection reason=relative-heading-rank; base=h2

\subsection{第四章 蒙召的人多}

于是伯多禄对群众的热情感到惊讶，回答说：\Quote{弟兄们，你们看，我们主的话显然应验了。因为我记得他说：\InnerQuote{将有许多人从东方和西方、北方和南方来，在亚巴郎、依撒格和雅各伯怀中坐席。}但他又说：\InnerQuote{被召的人多，被选的人少。}因此，这些蒙召之人的到来已经应验。但是，既然这不是出于他们自己，而是出于那召叫他们并使他们前来的天主，单凭这一点他们就没有赏报，因为这并非出于他们自己，而是出于那在他们内运行的天主。然而，如果他们蒙召后行卓越之事，这乃是出于他们自己，因此他们将为此获得赏报。}\par

% structure: p0010 source=h2 target=subsection reason=relative-heading-rank; base=h2

\subsection{第五章 信德是天主的恩赐}

\Quote{因为连信从梅瑟的希伯来人，若不遵守他所说的话，也不能得救，除非他们遵守对他们所说的话。因为他们信从梅瑟并非出于自己的意愿，而是出于天主，他对梅瑟说：\InnerQuote{看，我要在云柱中来到你这里，为叫百姓听我与你说话，并且永远信服你。}因此，既然无论是希伯来人，还是从外邦人中蒙召的人，信从真理的导师都是出于天主，而卓越的行为却留给每个人凭自己的判断去行，赏报就公正地赐予那些行善的人。因为如果凭他们自己就能理解合理之事，就不需要梅瑟或耶稣的来临了。此外，信从导师并称他们为主，也无救恩可言。}\par

% structure: p0012 source=h2 target=subsection reason=relative-heading-rank; base=h2

\subsection{第六章 隐藏与启示}

\Quote{因此，耶稣对那些以梅瑟为导师的犹太人是隐藏的，而梅瑟对那些信从耶稣的人也是隐藏的。因为两者教导同一道理，天主接受信从两者中任何一位的人。但信从导师是为了遵行天主所说的话。我们主自己也说：\InnerQuote{父啊，天地的主宰，我感谢你，因为你将这些事向聪明和年老的人隐藏起来，却启示给婴孩。}因此，天主亲自向一些人隐藏了一位导师，因为他预知他们该做什么；而向另一些人启示了他，因为他们不知道该如何行。}\par

% structure: p0014 source=h2 target=subsection reason=relative-heading-rank; base=h2

\subsection{第七章 梅瑟与基督}

\Quote{因此，希伯来人因不认识耶稣而不会受谴责，这是由于那隐藏耶稣的一位，如果他们遵行梅瑟所\Emph{命令}的事，并不恨他们所不认识的那一位。同样，从外邦人中的人也不因不认识梅瑟而受谴责，这是由于那隐藏梅瑟的一位，只要他们也遵行耶稣所说的话，并不恨他们所不认识的那一位。另有一些人，若只称导师为主，却不做仆人的工作，便毫无益处。正因如此，我们的耶稣曾对一位多次称他为主却不遵行他所吩咐的人说：\InnerQuote{你们为什么称呼我主啊，主啊，却不照我所说的去做？}因为说话不能使人受益，惟有行动才可以。因此，无论如何都需要善工。此外，如果有人有幸认出两者教导同一道理，那人就在天主眼中算为富足，既明白新时事如旧时，也明白旧时事如新时。}\par

% structure: p0016 source=h2 target=subsection reason=relative-heading-rank; base=h2

\subsection{第八章 一大群会众}

当伯多禄这样说时，群众好像被某人召唤，涌入了伯多禄所在的地方。他看见一大群人，如同河流平稳的水流缓缓流向自己，便对马罗内斯说：\Quote{你这里有没有更好的地方，能容纳这人群？}然后马罗内斯带他到一个露天的园地，群众跟着。伯多禄站在一个不高的雕像基座上，以虔诚的方式问候群众后，知道站在旁边的许多人被魔鬼折磨，长期患有各种痛苦，\Emph{听见他们}悲戚地尖叫，\Emph{在他面前}恳求地倒下，就斥责他们，命令他们安静；并承诺在讲道后为他们治愈，便开始这样说道：\par

% structure: p0018 source=h2 target=subsection reason=relative-heading-rank; base=h2

\subsection{第九章 \Quote{向人类说明天主的道路}}

\Quote{当开始向那些完全无知、心智被我们的对手西满的指控所败坏的人讲解对天主的敬拜时，我认为有必要首先谈论那创造万物的天主的无瑕，从天主按其眷顾适时提供的契机开始，好让人们知道，为何许多人被许多魔鬼缠身、遭受奇怪的痛苦是有充分理由的，为的是在此显明天主的公义；而那些因无知而责备祂的人，现在可以通过善言善行学习他们应有的看法，并从先前的指控中回转，将无知作为他们邪恶僭妄的原因，以便得到宽恕。}\par

% structure: p0020 source=h2 target=subsection reason=relative-heading-rank; base=h2

\subsection{第十章 原始法律}

\Quote{但事情是这样。唯一良善的天主妥善地创造了万有，并将它们交给按祂肖像所造的人。那受造者因受造者的神性而呼吸，是一位真正的先知，知晓一切；为了那赐予他万有的父的荣耀，以及从他而生的子孙的得救，他作为真实的父亲，保持着对亲生儿女的爱，并盼望他们为了自己的益处爱天主并蒙天主所爱，就向他们指出了通往友谊的道路，教导他们人的哪些行为能使独一的天主、万有的主喜悦；向他们展示了祂所喜悦的事之后，为所有人制定了一条永久的法律，这法律既不能被敌人废除，也不能被任何不敬虔者败坏，也不隐藏在任何地方，反而能被所有人阅读。因此，他们因服从法律而一切丰裕——最美好的果实、长寿、无忧无虑、无疾病，无惧地赐予他们，空气清新。}\par

% structure: p0022 source=h2 target=subsection reason=relative-heading-rank; base=h2

\subsection{第十一章 人堕落的原因}

\Quote{但他们最初没有经历邪恶的经验，对善恩的恩赐毫无感知，因食物丰盛和奢侈而变得忘恩负义，甚至认为没有天主眷顾，因为他们没有通过先前的劳作得到善事作为义德的酬报，没有人遭受任何痛苦、疾病或其他需要；因此，正如通常因邪恶过犯而受苦的人，他们会寻找能医治他们的天主。但就在他们因无所畏惧和安逸奢侈而生出的轻慢之后，某种公义的惩罚临到他们，仿佛来自某种安排的和谐，将有益于他们的善事移除，反过来引入邪恶之事作为有益之事。}\par

% structure: p0024 source=h2 target=subsection reason=relative-heading-rank; base=h2

\subsection{第十二章 天使的变形}

\Quote{原来，居住在最低区域的天使，因人的忘恩负义而忧伤，请求进入人的生活，真正成为人，通过更多交往来定罪那些忘恩负义者，并给予每人适当的惩罚。当他们的请求得到准许后，他们变幻成各种本性；因为他们属于更类似神的实体，能轻易采取任何形式。于是他们变成宝石、美丽的珍珠、最美丽的紫色、精选的黄金以及一切最受珍视的物质。他们落入一些人的手中，另一些人的怀中，任由自己被偷走。他们也变成野兽、爬行动物、鱼、鸟以及任何他们喜欢的东西。你们当中的诗人，因无所畏惧，也歌唱这些事，好似它们真的发生，将许多不同的行为归于一人。}\par

% structure: p0026 source=h2 target=subsection reason=relative-heading-rank; base=h2

\subsection{第十三章 天使的堕落}

\Quote{但当他们采取这些形式，定罪那些偷窃他们的贪婪之人，并改变成人的本性，为要过圣洁的生活，显示这样生活的可能性，从而惩罚忘恩负义者时，因他们在各方面成了人，也参与了人的情欲，并被其制服，落入与妇女同居；与她们纠缠，沉溺于污秽，完全失去了最初的能力，无法返回本性的原始纯洁，他们的肢体偏离了火性的实体；因为火本身被情欲的重量熄灭，\Emph{变成}血肉，他们踏上了不敬虔的下行之路。因为他们自己披上了肉体的束缚，被强力捆绑，因此再也无法升到天上。}\par

% structure: p0028 source=h2 target=subsection reason=relative-heading-rank; base=h2

\subsection{第十四章 他们的发现}

\Quote{因为交合之后，她们被要求展示他们先前是什么，却因污秽后不能再做别的事而无法做到，但为了取悦她们的女主人，他们便展示大地的内部，即精选的金属：金、铜、银、铁等，以及所有最珍贵的宝石。与这些有魔力的宝石一起，他们传授了与每种事物相关的技艺，传授了魔法的发现，教导了天文学、草根的威力，以及一切人类心智无法发现的事物；还有金、银等的熔炼，以及衣物的各种染色。总之，凡是为装饰和愉悦妇女的一切，都是这些被束缚在肉身中的魔鬼的发明。}\par

% structure: p0030 source=h2 target=subsection reason=relative-heading-rank; base=h2

\subsection{第十五章 巨人}

\Quote{但从他们不圣洁的交合中，生出了非法的男子，身材远高于\Emph{常人}，后来被称为巨人；并非那些希腊亵渎神话所歌颂的龙足巨人与天主作战，而是性情野蛮，体型大于常人，因为他们出自天使；却又小于天使，因为他们生于女人。因此天主知道他们已野蛮化为残暴，世界不足以满足他们（因为世界是按人的比例和人类用途创造的），免得他们因缺乏食物而违反本性去吃动物，却又似乎无可指责，因为他们是不得已而为之，全能的天主便降下玛纳给他们，适合他们各种口味；他们尽情享用。但他们因私生子的本性，不满足于纯净的食物，只渴求血的味道。因此他们首先品尝了肉。}\par

% structure: p0032 source=h2 target=subsection reason=relative-heading-rank; base=h2

\subsection{第十六章 食人}

\Quote{而当时与他们在一起的人也渴望效仿。因此，虽然我们生来不善不恶，却变得\Emph{要么善要么恶}；一旦形成习惯，便难以改掉。但当无理性动物匮乏时，这些私生子也尝到了人肉。因为先尝了其他形式的肉，再吃与自己同类之肉，并不是很大的一步。}\par

% structure: p0034 source=h2 target=subsection reason=relative-heading-rank; base=h2

\subsection{第十七章 洪水}

\Quote{但流了许多血，纯洁的空气被不洁的蒸汽污染，使呼吸它的人生病，从而容易患病，从此人类过早死亡。而大地因此大大玷污，先就充满了带毒刺和致命的生物。因此，由于这些残暴的魔鬼，万事越来越糟，天主愿像抛弃邪恶的酵母一样抛弃它们，免得每个生于邪恶种子的世代，像以前一样不虔敬，使未来的世界没有得救的人。为此，祂警告了一位义人及其三个儿子，连同他们的妻子和孩子们，让他们在方舟里自救，便降下洪水的泛滥，使一切毁灭，好让被净化后的世界交给在方舟里得救的人，以开始第二次生命。事情就这样发生了。}\par

% structure: p0036 source=h2 target=subsection reason=relative-heading-rank; base=h2

\subsection{第十八章 对幸存者的法律}

\Quote{因此，已故巨人的灵魂比人类的灵魂更强大，正如它们也胜过身体一样，它们作为一个新种族，也被赋予了一个新名字。而对于世上幸存的人，天主通过一位天使制定了法律，告诉他们应如何生活。因为他们在种族上是私生子，出自天使的火和女人的血，因此容易渴望拥有自己的种族，他们被某种正义的法律预先阻止了。因为有一位天使被天主派到他们那里，向他们宣布祂的旨意，说道：——}\par

% structure: p0038 source=h2 target=subsection reason=relative-heading-rank; base=h2

\subsection{第十九章 对巨人或魔鬼的法律}

\Quote{这些事在无所不见的天主眼中看为好：你们不得统治任何人；不得扰乱任何人，除非有人自愿臣服于你们，敬拜你们，向你们献祭和奠酒，与你们共席，或做任何其他不当之事，或流人血，或尝死肉，或吃被野兽撕裂的、被切割的、被勒死的，或其他任何不洁之物。但那些投奔我法律的人，你们不仅不可触碰，还要尊敬他们，并从他们面前逃走。因为凡为义人的，就你们而论，他们所喜悦的事，你们必受约束而承受。但如果任何敬拜我的人走迷，或犯奸淫，或行邪术，或不洁地生活，或做任何其他我不喜悦的事，那么他们将根据我的命令，在你们或他人手中受罚。但当他们悔改时，我将判断他们的悔改是否配得赦免，然后宣判。因此，你们应当记住并遵行这些事，因为你们深知，连你们的念头也无法向祂隐藏。}\par

% structure: p0040 source=h2 target=subsection reason=relative-heading-rank; base=h2

\subsection{第二十章 自愿的俘虏}

\Quote{天使如此吩咐后便离开了。但你们仍不知道这条法律：凡敬拜魔鬼、向他们献祭、或与他们共席的人，都将受制于他们，并从他们那里受一切惩罚，如同在邪恶的主宰之下。而你们因不知晓这条\Emph{法律}，在他们的祭坛旁败坏了，饱食了祭献给他们的\Emph{食物}，就落在了他们的权下，却不知道你们的身体在各方面受到了伤害。但你们应当知道，魔鬼对任何人都没有权柄，除非那人先成为他们的同席者；因为连他们的首领也不能违反天主加于他们的法律，因此他对不敬拜他的人无权；但也没有人能从他那里得到任何想要的东西，也不会受到任何伤害，这从以下陈述可以得知。}\par

% structure: p0042 source=h2 target=subsection reason=relative-heading-rank; base=h2

\subsection{第二十一章 基督的诱惑}

\Quote{因为今世之王曾来就我们正义的君王，没有使用暴力，因为这事不在他的权能内，而是引诱和说服，因为被说服在于每个人的选择。因此，他作为今世之主，对来世之王说：\InnerQuote{今世万国都属我管辖；金、银和今世的一切奢华都在我权下。所以，你如果俯伏朝拜我，我就把这一切都赐给你。}他说这话，是知道一旦祂朝拜他，他就能控制祂，从而夺去祂将来的荣耀和国度。但祂知道一切，不仅不朝拜他，也不接受他所给的一切。因为祂为祂自己的人立下誓言：此后甚至不可触摸那些已交付给他的东西。因此祂回答说：\InnerQuote{你要敬畏上主你的天主，单要事奉祂。}}\par

% structure: p0044 source=h2 target=subsection reason=relative-heading-rank; base=h2

\subsection{婚宴}

\Quote{然而，不虔敬者的君王试图说服虔敬者的君王归顺他的计谋，却不能成功，就停止努力，决意在其余生中迫害祂。但你们由于无知于预定的法律，因恶行而处在他的权下。因此，你们身心灵被玷污，今生受痛苦和魔鬼的奴役，来世你们的灵魂必受惩罚。这不单是你们因无知而受的苦，也是我们民族中的一些人，因恶行被邪恶之君的权下所制，如同被邀请参加一位父亲为儿子娶亲所设的宴席，却不肯听从。但代替那些因耽延而不服从的人，为儿子娶亲的父亲已命令我们，藉着真理的先知，来到岔路口，即到你们这里，给你们穿上洁净的婚服，即\OriginalTerm{圣洗圣事}{baptism}，为赦免你们所犯的罪，并藉着悔改将好人带到天主的宴席，虽然起初他们被排除在宴席之外。}\par

% structure: p0046 source=h2 target=subsection reason=relative-heading-rank; base=h2

\subsection{散会}

\Quote{因此，如果你们愿意成为天主圣神的衣裳，就赶快先脱去你们卑劣的骄傲，那是不洁的灵和污秽的衣服。你们若不先在善工中\OriginalTerm{领受圣洗}{baptized}，就不能脱去它。这样，身心灵洁净，你们就能享受将来的永恒国度。所以不要信偶像，也不要与他们同赴不洁的宴席，不可杀人，不可奸淫，不可恨那不该恨的人，不可偷盗，不可行任何恶事；因为你们若在今生失去将来福乐的希望，就会受恶魔鬼和可怕痛苦的折磨，在来世将被永火惩罚。现今所说的已足够今日之用。其余的人中，你们有疾病的，留待治愈；其余的，你们愿意的，可以平安离去。}\par

% structure: p0048 source=h2 target=subsection reason=relative-heading-rank; base=h2

\subsection{病愈}

他这样说了以后，众人都留下来，有的为了得医治，有的为了看那些得痊愈的人。伯多禄只用按手在他们身上，祈祷，便治好了他们；以致立刻痊愈的人极其欢喜，旁观的人极其惊奇，赞美天主，怀着坚定的希望相信了，并与那些得治愈的人一同回到自己的家，受命次日清早聚集。他们离去后，伯多禄和他的同伴留在那里，用了饭，又睡眠休息。\par

\SourceRecord{Translated by Peter Peterson.；Ante-Nicene Fathers；Vol. 8.；Edited by Alexander Roberts, James Donaldson, and A. Cleveland Coxe.；Buffalo, NY: Christian Literature Publishing Co.,；1886.；<http://www.newadvent.org/fathers/080808.htm>.}

% structure: p0001 source=h1 target=section reason=page-title

\section{讲道集 第九篇}

% structure: p0002 source=h2 target=subsection reason=relative-heading-rank; base=h2

\subsection{第一章 伯多禄继续发言}

因此，次日伯多禄与同伴们一同外出，来到先前的地方，站定后，继续说道：\Quote{天主曾以洪水剪除了所有古代不敬虔的人，在众人中只找到一位虔诚者，便使他与三个儿子及儿媳在方舟中得救。由此可知，祂的本性是不关心众多恶人，也不忽视一位虔诚者的得救。因此，最大的不虔敬就是离弃唯一的万有之主，而去敬拜许多非神的假神，如同敬拜神一样。}\par

% structure: p0004 source=h2 target=subsection reason=relative-heading-rank; base=h2

\subsection{第二章 独一统治与多神统治}

\Quote{因此，当我阐述并向你们指明这是最大的罪，足以毁灭你们众人时，如果你们心中想：你们人数众多，并未被毁灭，你们就受骗了。因为你们有上古世界被淹没的榜样。然而他们的罪远比你们所犯的罪要小。因为他们只对同等人行恶，杀人或犯奸淫。但你们却对万有之神行恶，敬拜无生命的像以代替祂，或与祂一同敬拜，并将祂神圣的名归于各种无灵之物。首先，你们不幸的是不知道独一统治与多神统治的区别——独一统治产生和谐，而多神统治导致战争。因为一者不与自身争斗，而众多则彼此寻衅交战。}\par

% structure: p0006 source=h2 target=subsection reason=relative-heading-rank; base=h2

\subsection{第三章 诺厄家族}

\Quote{因此，洪水之后，诺厄立即与众多后裔和谐地继续生活了三百五十年，如同一位照着独一真神的形象为王的君王。但他死后，许多后裔觊觎王位，渴望统治，各人思虑如何实现；有的靠战争，有的靠诡计，有的靠说服，各有各的方式；其中一人出自含族，其子孙是梅斯伦，埃及人、巴比伦人和波斯人的部落就是从梅斯伦繁衍而来的。}\par

% structure: p0008 source=h2 target=subsection reason=relative-heading-rank; base=h2

\subsection{第四章 琐罗亚斯德}

\Quote{这一族中，在适当的时候生了一个人，他从事魔术，名叫内布罗德，他像巨人一样选择图谋反对天主的事。希腊人称他为琐罗亚斯德。洪水之后，他野心勃勃想统治，又是大魔术师，用魔术迫使统治世界的那个恶者的星辰将统治权作为礼物赐给他。但那位首领，拥有权柄，对他施压的人发怒，愤怒地倾泻了王国的火，为要使他服从，并惩罚那起初强迫他的人。}\par

% structure: p0010 source=h2 target=subsection reason=relative-heading-rank; base=h2

\subsection{第五章 英雄崇拜}

\Quote{因此，魔术师内布罗德被这从天上落到地上的闪电摧毁，为此他的名被改为琐罗亚斯德，因为那颗星的\OriginalTerm{活流}{\GreekText{ζῶσαν}} \OriginalTerm{星流}{\GreekText{ἀστέρος}}倾倒在他身上。但当时那些愚昧无知的人，以为因着天主的爱，他的灵魂被闪电召回，便埋葬了他身体的遗骸，并在波斯人那里，在那火降下的地方，用一座殿宇尊崇他的墓地，并敬拜他为神。受此例影响，其他人也将被闪电击毙的人当作天主所爱的人埋葬，用殿宇尊崇他们，并为死者立起他们形象的雕像。于是，各地统治者同样争相效仿，许多人尊崇他们所爱之人的坟墓，即或不死于闪电，也用殿宇和雕像尊崇，并点燃祭台，命令将他们当作神敬拜。久而久之，后代便以为他们真是神了。}\par

% structure: p0012 source=h2 target=subsection reason=relative-heading-rank; base=h2

\subsection{第六章 火崇拜}

\Quote{于是，就这样，最初的独一王国分裂成许多部分。波斯人首先从天上降下的闪电中取来火炭，用普通燃料保存，并尊崇这天火为神，火本身也以其最初敬拜者的身份，赐予他们第一个王国。之后，巴比伦人从那里的火中偷取火炭，安全地送往自己家中，并敬拜它，他们也依次统治了。埃及人照样行，用他们自己的方言称那火为\Emph{弗塔}，这词译作\Emph{赫淮斯托斯}或\Emph{奥西里斯}，他们中第一位统治者也以此名为名。那些在不同地方统治的人，也如此行，制造雕像，点燃祭台以尊敬火，他们中大多数被排除在王国外。}\par

% structure: p0014 source=h2 target=subsection reason=relative-heading-rank; base=h2

\subsection{第七章 祭宴狂欢}

\Quote{但他们并没有停止敬拜雕像，这是因为魔术师们的邪恶智慧，他们为雕像制造借口，这些借口竟有能力迫使他们进行愚蠢的崇拜。因为魔术师们通过仪式建立这些事，为他们指定了祭宴、奠酒、笛声和呐喊的节庆，藉此，愚昧无知的人受骗，即使他们的王国被夺走，也仍不放弃他们所接受的崇拜。他们竟如此因愉快而选择错误胜过真理。他们在祭宴饱足后嚎叫，他们的灵魂仿佛从深处，像通过梦境，预先警告他们将因这样的行为而受惩罚。}\par

% structure: p0016 source=h2 target=subsection reason=relative-heading-rank; base=h2

\subsection{第八章 最好的商品}

\Quote{世上许多崇拜形式已然过去，我们前来，如同好商人，将我们从祖先传承并保存下来的崇拜带给你们；像展示植物的种子一样，将之摆在你们的判断和权力之下。请选择你们看为好的。因此，如果你们选择我们的货物，你们不仅能逃脱魔鬼及其所加的痛苦，而且你们自己也能驱逐它们，并使它们降低到向你们哀求的地步，且永远享受未来的福乐。}\par

% structure: p0018 source=h2 target=subsection reason=relative-heading-rank; base=h2

\subsection{第九章 魔鬼如何获得对人的权柄}

\Quote{反之，既然你们正遭受由魔鬼加诸的怪异苦难，当你们离开身体时，你们的灵魂也将永远受罚；这并非由于天主的报复，而是因为恶行本身招致的审判。因为魔鬼藉着给予它们的食物获得能力，被你们亲手接纳进入你们的身体；长期潜伏在那里，与你们的灵魂相混合。由于那些不愿或甚至不想自助之人的疏忽，在他们的身体瓦解时，他们的灵魂既与魔鬼结合，就必然被它带到它所喜悦的任何地方。而最可怕的是，当万物终结时，魔鬼首先被投入炼净之火，与它混合的灵魂不得不受可怕的刑罚，而魔鬼却感到愉悦。因为灵魂由光构成，不能承受异质的火焰，因而受折磨；但魔鬼属于自身特性的实体，因而大为喜悦，成为它所吞没的灵魂的坚固锁链。}\par

% structure: p0020 source=h2 target=subsection reason=relative-heading-rank; base=h2

\subsection{第十章 如何驱逐它们}

\Quote{但魔鬼喜欢进入人身体的原因如下：它们是神灵，渴望肉食、饮料和性快感，但因是神灵，无法分享这些，又缺乏享受这些的器官，于是它们进入人的身体，以便获得器官为它们服务，得到它们所愿之物，无论是藉着人的牙齿获得肉食，还是藉着人的肢体获得性快感。因此，为了驱逐魔鬼，最有益的帮助是节制、斋戒和受苦受难。因为如果它们进入人身体是为了分享\Emph{快乐}，显然它们会被苦难驱逐。但有些更为恶毒的魔鬼，即使与身体一同受罚，仍停留在受罚的身体旁，因此需要通过祈祷和恳求转向天主，远离一切不洁的场合，好让天主的手触摸他，医治他，因他是纯洁而有信德的。}\par

% structure: p0022 source=h2 target=subsection reason=relative-heading-rank; base=h2

\subsection{第十一章 不信是魔鬼的堡垒}

\Quote{但在我们的祈祷中，必须承认我们已转向天主，并作证，不是证明魔鬼的冷漠，而是它的缓慢。因为一切为有信德者而成，无物为不信者而成。因此，魔鬼本身知道它们所占据之人的信德程度，便按比例衡量它们的停留。因此，它们永久地停留在不信者身上，在信德软弱者身上停留片刻；但在那些彻底相信并行善的人身上，它们连一刻也不能停留。因为灵魂因信德而变得如同水的性质，就能像泼水灭火星一样熄灭魔鬼。因此，每个人的努力就在于关心驱逐自身的魔鬼。因为当它们与人的灵魂混合时，就向每个人的心智暗示它们所喜悦的欲望，好让他忽略自己的救恩。}\par

% structure: p0024 source=h2 target=subsection reason=relative-heading-rank; base=h2

\subsection{第十二章 疾病理论}

\Quote{因此，许多人不知道他们如何受影响，就同意魔鬼所暗示的恶念，仿佛这些是他们自己灵魂的推理。因而他们变得不太积极去接近那些能拯救他们的人，也不知道他们自身被欺骗的魔鬼所俘虏。因此，潜伏在他们灵魂中的魔鬼诱导他们以为困扰他们的不是魔鬼，而是一种身体疾病，例如某些酸性物质、胆汁、痰、血过多、膜发炎或其他东西。但即使如此，情况也不会改变它属于魔鬼的事实。因为普遍而属土的灵魂，因各种食物而进入，由于过度饮食而过度摄取，它自身与那与它同类的灵（即人的灵魂）相结合；而食物的物质部分与身体结合，留下可怕的毒物。因此在各方面，节制都是卓越的。}\par

% structure: p0026 source=h2 target=subsection reason=relative-heading-rank; base=h2

\subsection{第十三章 魔鬼的欺骗}

\Quote{但有些恶毒的魔鬼以另一种方式欺骗。起初它们甚至不显露自己的存在，以免人们提防它们；但到了适当的时候，藉着愤怒、爱情或其他情感，它们突然用刀剑、绞索、悬崖或其他方式伤害身体，最后，正如我们所说，使那些与它们混合的被欺骗的灵魂受到惩罚，退入炼净之火。但另一些人，以另一种方式被欺骗，并不来接近我们，他们被恶毒魔鬼的煽动所引诱，仿佛这些事是神祇本身因他们忽略神祇而降下的，并且他们能够通过祭献与神祇和好，因此不需要来我们这里，反而要逃避和憎恨我们。同时，他们憎恨和逃避那些对他们有更大怜悯、跟随他们以便向他们行善的人。}\par

% structure: p0028 source=h2 target=subsection reason=relative-heading-rank; base=h2

\subsection{第十四章 更多诡计}

\Quote{因此，他们逃避和憎恨我们，是被欺骗了，不知道他们为何会策划出有害自身健康的事。因为我们不能强迫他们违背其意愿趋向健康，如今我们对他们没有这样的权能；他们自身也不能明白魔鬼的邪恶煽动，因为他们不知道这些邪恶的煽动从何而来。这些就是魔鬼所恐吓的人，魔鬼以它们愿现的形象显现。有时它们给患病者开出药方，从而从先前被欺骗的人那里接受神圣的尊荣。它们向许多人隐藏它们是魔鬼的事实，但不向我们隐藏，我们知晓它们的奥秘，知道它们为何做这些事——在梦境中改变自己来对付它们有权势的人；以及为何恐吓一些人，向另一些人发出神谕般的回应，要求他们献祭，并命令他们与它们一同吃喝，以便吞没他们的灵魂。}\par

% structure: p0030 source=h2 target=subsection reason=relative-heading-rank; base=h2

\subsection{第十五章 偶像的考验}

\Quote{因为如同可怕的毒蛇用气息吸引麻雀，这些魔鬼也通过食物和饮料，混杂于人的领悟中，在梦中根据偶像的形状改变自己，以增加错误。因为偶像既不是活的，也没有神圣的精神，而是显现的魔鬼滥用了形状。同样，有多少人曾在梦中被其他人看见，而他们醒着相遇并对比梦中所见，却不一致；所以梦不是启示，而是魔鬼或灵魂的产物，赋予当下的恐惧和欲望以形状。因为灵魂受惊时，会在梦中构想形状。但如果你认为偶像既然是活的，就能成就这些事，就把它们放在一架平衡的横梁上，在另一边放上等重的砝码，然后吩咐它们变重或变轻；如果这事成就，那么它们就是活的。但事实并非如此。即便如此，也不能证明它们是神，因为这可能是魔鬼的手指所为。连蛆虫都会动，但它们不被称为神。}\par

% structure: p0032 source=h2 target=subsection reason=relative-heading-rank; base=h2

\subsection{第十六章 魔鬼的权能}

\Quote{但是，每个人的灵魂按照自己的先入之见体现魔鬼的形状，并且被称为神的神并不显现，这从它们不向犹太人显现的事实可以明显看出。但有人会说，它们如何发布神谕，预言未来的事？这也是虚假的。但假设这是真的，这也不能证明它们是神；因为并非任何东西能预言就是神。例如，灵蛇能预言，却被我们当作魔鬼驱逐和赶走。但有人会说，它们为一些人治病。这是虚假的。但假设这是真的，这也不是神性的证据；因为医生也能治愈许多人，但他们不是神。但是，有人说，医生不能完全治愈他们负责的人，但这些魔鬼通过神谕治病。但魔鬼知道适合每种疾病的疗法。因此，作为熟练的医生，能治愈人类可治愈的疾病，并且也是预言者，知道每种疾病何时会自行痊愈，它们如此安排疗法，以便获得产生治愈的功劳。}\par

% structure: p0034 source=h2 target=subsection reason=relative-heading-rank; base=h2

\subsection{第十七章 他们的欺骗为何不被察觉}

\Quote{因为为什么它们通过神谕预言很久以后才治愈？如果它们是全能的，为什么不在不施用药的情况下就治愈？又为什么它们向一些祈求它们的人开处方，而对另一些（可能更合适的病例）却不给回应？因此，每当治愈即将自发发生时，它们就允诺，以便获得治愈的功劳；另一些人，生了病，祈祷后自然康复，就将治愈归功于他们所祈求的神，并向它们献祭。而那些祈祷后失败的人，则无法献祭。但如果死者的亲属或他们的孩子追问损失，你会发现失败多于成功。但没有人愿意因为羞耻或恐惧而控告它们；相反，他们掩盖他们认为它们应受的罪行。}\par

% structure: p0036 source=h2 target=subsection reason=relative-heading-rank; base=h2

\subsection{第十八章 体系的支柱}

\Quote{又有多少人伪造它们所给出的回应和所行的治愈，并用誓言加以确认！又有多少人为了报酬而把自己交给它们，假装遭受某些苦痛，从而宣扬自己的痛苦，并通过治疗手段得到康复，然后说它们曾通过神谕应许他们治愈，以便将原因归咎于这无知的敬拜！又有多少这样的事以前通过法术、解梦和占卜来进行！然而随着时间的推移，这些事已经消失了。现在又有多少人，希望得到这样的东西，使用符咒！然而，即使一件事是预言性或治疗性的，它也不是神性的。}\par

% structure: p0038 source=h2 target=subsection reason=relative-heading-rank; base=h2

\subsection{第十九章 受洗者的特恩}

\Quote{因为天主是全能的。因为祂是良善和公义的，现今容忍所有人，以便那些愿意的人，悔改他们所做的恶事，行善生活，在审判万物的日子得到配得的报酬。因此，现在开始由于正确的认识而服从天主，并反对你们邪恶的情欲和思想，以便你们能够恢复委托给人类的原始救恩性敬拜。因为这样，祝福将立刻向你们涌来，当你们接受了它们，你们从此将脱离邪恶的考验。但要感谢赐予者；永远作不可言喻之善事的君王，与和平的君王同在。但在现世生活中，在三倍有福的呼求下，在流动的河流、泉水或甚至海洋中洗涤，你们不仅能驱赶潜藏在你们内的邪灵；而且你们自己不再犯罪，无疑地相信天主，你们将驱逐他人身上的邪灵、可怕的魔鬼和可怕的疾病。有时，你们只要看它们一眼，它们就会逃跑。因为它们认识那些把自己交托给天主的人。因此，尊敬它们，它们恐惧地逃跑，正如你们昨天所看见的，当我讲道后为那些遭受这些疾病的人祈祷时，由于对敬拜的尊敬，它们大声喊叫，无法忍受短短的一小时。}\par

% structure: p0040 source=h2 target=subsection reason=relative-heading-rank; base=h2

\subsection{第二十章 \Quote{不是几乎，而是完全像我一样}}

\Quote{那么，不要以为我们因为与\Emph{你们}本性不同而因此不怕魔鬼。因为我们有相同的本性，但敬拜不同。因此，我们不仅比你们优越很多，而且完全优越，我们并不嫉妒你们成为\Emph{像我们一样}；相反，我们劝告你们，因为我们知道所有这些\Emph{魔鬼}都无比尊敬和惧怕那些与天主和好的人。}\par

% structure: p0042 source=h2 target=subsection reason=relative-heading-rank; base=h2

\subsection{第二十一章 魔鬼服从于信徒}

\Quote{因为，正如在凯撒的一位军官手下的士兵知道要尊敬那因授权者而获得权柄的人，所以指挥官对这个人说\InnerQuote{来}，他就来，对另一个人说\InnerQuote{去}，他就去；同样，那将自己交托给天主的人，若是忠信，当他只对魔鬼和疾病说话时，就蒙垂听；魔鬼就退避，尽管它们比命令它们的人强大得多。因为天主以不可言喻的能力，将祂所愿的人的心思置于权下。正如许多军官带着整营和整座城邑畏惧凯撒——虽然他只是一个人——每个人的心都急于为天主的旨意而尊崇所有人的肖像，万物因恐惧而被奴役，却不知道原因；同样，所有致病之灵，出于某种自然的方式，敬畏并逃避那投靠天主、并将正信如同祂的肖像带在心中的人。}\par

% structure: p0044 source=h2 target=subsection reason=relative-heading-rank; base=h2

\subsection{第二十二章 \Quote{反而要喜乐}}

\Quote{然而，尽管所有魔鬼连同各种疾病都在你面前逃遁，你不可只为此喜乐，而要因你们的名字藉着恩宠如同永生者被记录在天上而喜乐。同样，神圣的圣神也因人战胜死亡而喜乐；因为驱赶魔鬼有助于他人的得救。但我们说这话，并非否认我们应当帮助别人，而是说我们不可因此自高而忽略自己。有时，魔鬼也因那尊贵的名号在恶人面前逃遁，于是驱魔者和目睹者都受欺骗：驱魔者以为自己因义德而受尊荣，却不知魔鬼的诡计。因为魔鬼一面尊崇那名号，一面藉着逃跑使那恶人自以为义，从而诱骗他远离悔改。而旁观者见那驱魔者如同虔诚的人，便效法其生活方式，最终堕落。有时，魔鬼也佯装在不以天主之名发出的驱魔前逃遁，为要欺骗人，毁灭他们所欲毁灭的人。}\par

% structure: p0046 source=h2 target=subsection reason=relative-heading-rank; base=h2

\subsection{第二十三章 病者获愈}

\Quote{因此我们愿你们知道，除非有人自愿将自己交给魔鬼作奴隶，如我先前所说，魔鬼无权胜过他。因此，你们若选择敬拜唯一的天主，戒绝魔鬼的筵席，怀着\OriginalTerm{仁爱}{philanthropy}与义德持守贞洁，并为了赦罪而以\OriginalTerm{三重祝福的名号}{thrice-blessed invocation}受洗，尽己所能追求纯洁的完美，就能逃脱永罚，成为永福的继承者。}\par

说了这些话，他吩咐那些为疾病所苦的人前来。于是许多人前来，因昨日得愈者的经验而聚集。他按手在他们身上并祈祷，立刻治愈了他们，然后嘱咐他们和其他人更早前来，便沐浴进食，就寝了。\par

\SourceRecord{Translated by Peter Peterson.；Ante-Nicene Fathers；Vol. 8.；Edited by Alexander Roberts, James Donaldson, and A. Cleveland Coxe.；Buffalo, NY: Christian Literature Publishing Co.,；1886.；<http://www.newadvent.org/fathers/080809.htm>.}

% structure: p0001 source=h1 target=section reason=page-title

\section{讲道集 第十篇}

% structure: p0002 source=h2 target=subsection reason=relative-heading-rank; base=h2

\subsection{第一章 在\Place{的黎波里}的第三天}

\Emph{因此}，第三天在\Place{的黎波里}，\Person{伯多禄}清早起来，进入园中，那里有一个大水池，源源不断的水流入其中。他在那里沐浴后，祈祷完毕，便坐下；见我们围坐四周，热切注视他，似乎想听他讲些什么，便说：——\par

% structure: p0004 source=h2 target=subsection reason=relative-heading-rank; base=h2

\subsection{第二章 无知与错误}

\Quote{在我看来，无知者与犯错者之间有极大的区别。因为无知者在我看来，就像一个人不愿前往一座物产丰盛的城市，因为他不认识那里的美好事物；但犯错者则像一个人，确实知道了城中的美好事物，却在前往的途中放弃了大道，因此迷了路。所以，在我看来，拜偶像者与在敬拜天主方面有缺陷的人之间也有极大的区别。因为拜偶像者不认识永生，因此也不渴慕它；因为他们所不知道的，他们无法爱。但那些选择敬拜独一天主、并得知赐予善人的永生之人，若他们相信或做了任何不悦天主的事，就类似于那些从惩罚之城出来、渴望前往物产丰盛之城、却在路上偏离了正路的人。}\par

% structure: p0006 source=h2 target=subsection reason=relative-heading-rank; base=h2

\subsection{第三章 人是万物的主人}

当他这样向我们讲论时，我们中有一人进来，他受命作如下宣告，说道：\Quote{我的主 \Person{伯多禄}，门外有许多群众。}经他同意后，一大群人进来了。于是他起身，像前一天那样站在台基上，以虔诚的方式致候后，说：\Quote{天主创造了天地，并创造了其中的万物，正如真先知对我们所说的，人按天主的肖像和模样受造，被指定为万物的统治者和主人，我说，是在空气、地上和水中的万物；这从以下事实可知：人凭其才智降服空中之生物，从深渊中取出生物，猎取地上之生物，尽管它们比人强壮得多，我说的是大象、狮子以及诸如此类的动物。}\par

% structure: p0008 source=h2 target=subsection reason=relative-heading-rank; base=h2

\subsection{第四章 信德与本分}

\Quote{因此，当人正义时，他也超脱一切痛苦，因其不朽的身体不能体验任何疼痛；但当他犯罪时——正如我前天和昨天所显示的——他仿佛成了罪恶的奴隶，屈服于一切痛苦，因公义的审判而被剥夺了一切卓越的事物。因为既然离弃了施恩者，恩赐就不应留在忘恩者那里，这是合情合理的。因此，出于祂丰盛的慈悲，为了使我们不仅领受最初的，也领受将来的福乐，祂派遣了祂的先知。先知已托付我们告诉你们应当思想什么，应当做什么。所以，请选择吧；这事就在你们能力之内。因此，你们应当这样思想：敬拜创造万有的天主；你们若在心里接纳祂，就必从祂那里，连同最初的卓越事物，一同领受将来的永恒福乐。}\par

% structure: p0010 source=h2 target=subsection reason=relative-heading-rank; base=h2

\subsection{第五章 敬畏天主}

\Quote{因此，你们若像驱魔者那样，对潜伏在心中的可怕毒蛇说：\InnerQuote{你当敬畏上主你的天主，唯独事奉祂}，你们就能说服自己接受有益之事。从各方面说，敬畏祂独一者是有益的，不是作为不义的天主，而是作为公义的天主。因为人害怕不义者，怕被无理毁灭；但害怕公义者，是怕被抓住在罪中而受惩罚。因此，借着对祂的敬畏，你们能摆脱许多有害的恐惧。因为若你们不敬畏独一的主宰和创造主，你们就会成为一切邪恶的奴隶，伤害自己，我指的是魔鬼、疾病以及任何能伤害你们的事物。}\par

% structure: p0012 source=h2 target=subsection reason=relative-heading-rank; base=h2

\subsection{第六章 恢复天主肖像}

\Quote{因此，要满怀信心地接近天主，你们起初被造为万物的统治者和主人；你们身体上带有祂的肖像，你们心志中也同样带有祂审判的模样。既然你们像无理性的动物那样行事，从你们的灵魂中失去了人的灵魂，变得像猪一样，成了魔鬼的猎物。所以，若你们领受天主的法律，就成为了人。因为不能对无理性的动物说：\InnerQuote{不可杀人，不可奸淫，不可偷盗}，等等。因此，当被邀请时，不要拒绝返回你们最初的尊贵；因为若你们借着善工相似天主，这是可能的。你们因相似祂而被视为儿子，将恢复为万物的主人。}\par

% structure: p0014 source=h2 target=subsection reason=relative-heading-rank; base=h2

\subsection{第七章 偶像的无用}

\Quote{那么，开始脱去对虚妄偶像的有害恐惧，使你们能摆脱不义的捆绑。因为那些甚至作为仆人也不能有益于你们的，却成了你们的主人。我指的是无生命像的材料，它们在服事方面对你们毫无用处。因为它们不听、不看、不感觉，也不能移动。你们中有人愿意像它们那样看、像它们那样听、像它们那样感觉、像它们那样移动吗？愿天主阻止这样的事发生在任何带有天主肖像的人身上，即使他已失去了祂的模样。}\par

% structure: p0016 source=h2 target=subsection reason=relative-heading-rank; base=h2

\subsection{第八章 没有手造的神}

\Quote{因此，将你们那些金、银或任何其他材料的偶像还原为它们的本来性质；我是指还原为杯、盘和所有其他器皿，此类器具可能对你们有用；而最初赐予你们的那些美好之物将得以恢复。但也许你们会说：皇帝的法律不允许我们这样做。你们说得对，那只是法律，而虚妄的偶像本身毫无能力。既然如此，你们怎能把它们当作神呢？它们被人的法律所保护，被狗所看守，被众人所把持——如果它们是金、银或铜制的。至于木制或陶制的，则因其毫无价值而得以保存，因为没有人想要偷窃一个木质或陶制的偶像！因此，你们的偶像越贵重，它们就越处于危险之中。它们怎能是神呢？能被偷、被熔、被称重、被看守？}\par

% structure: p0018 source=h2 target=subsection reason=relative-heading-rank; base=h2

\subsection{第九章 \Quote{有眼，却看不见。}}

\Quote{可怜的人心啊，竟然畏惧比死人更死的东西！因为我甚至不能说它们曾经活过，除非它们是古人的坟墓。有时，一个人来到陌生之地，不知道他所见的庙宇是死人的纪念碑，还是属于所谓的神明；但若询问并得知它们属于神明，他便敬拜，却毫不羞愧：若他没有询问得知，便会因两者极其相似而将它们当作死人的纪念碑走过。然而，关于此类迷信，我不必提供过多证据。因为任何愿意的人都很容易明白：\Emph{偶像}本是虚无，除非有人视而不见。但现在至少听吧，它听不见，也明白它不明白。因为死人之手制造了它。如果制造者已死，他所造之物岂能不消亡？那你们为何敬拜一个全然无感觉的必朽者的作品？而有理性的人不敬拜动物，也不试图平息天主所造的元素——我是指天、日、月、闪电、海和其中所有——他们正当地判断不敬拜祂所造之物，而是敬畏造物主和维系者。因为连它们自己也乐于无人将属于造物主的尊荣归于它们。}\par

% structure: p0020 source=h2 target=subsection reason=relative-heading-rank; base=h2

\subsection{第十章 拜偶像乃是蛇的欺骗}

\Quote{因为唯独那一位的荣耀是卓越的：唯独祂是非受造的，而其他一切都是受造的。因此，正如不受造的特性是属天主的，凡受造的实在不是天主。所以，在一切之前，你们应当思想那欺骗之蛇在你们里面的恶毒提示，它以更好的理性为应许引诱你们，从你们的大脑蠕动到脊髓，并极其重视欺骗你们。}\par

% structure: p0022 source=h2 target=subsection reason=relative-heading-rank; base=h2

\subsection{第十一章 蛇为何引诱人犯罪}

\Quote{因为他知道那原始律法：若他能使你们相信所谓的神明，以致你们犯罪反对独一统御之善，你们的跌倒便成了他的获益。正因如此，他被判罚吃土，就有能力吞吃那因罪而化为尘土、成为泥土的人，你们的灵魂进入他火焰的肚腹。为使你们遭受这些，他暗示一切有害的念头。}\par

% structure: p0024 source=h2 target=subsection reason=relative-heading-rank; base=h2

\subsection{第十二章}

% structure: p0025 source=h2 target=subsection reason=relative-heading-rank; base=h2

\subsection{无知不恕}

\Quote{因为一切反对独一统御的欺骗性念头，都是他为了伤害你们而播撒在你们心中的。首先，为使你们不听虔敬的言论，从而驱散无知——这乃是邪恶的缘由——他以知识的假象引诱你们，首先给予并贯彻这种狂妄，即认为并错误地以为：若有人不听虔敬之言，就不受审判。因此，有些人受此欺骗，不愿听闻，以便保持无知，却不知无知本身就是一种足以致命的毒药。因为若有人无知地服下致命毒药，难道他不会死吗？同样，罪自然地毁灭罪人，即使他因无知而犯罪，不知何为正当。}\par

% structure: p0027 source=h2 target=subsection reason=relative-heading-rank; base=h2

\subsection{第十三章 对无知之人的谴责}

\Quote{但如果违背训诲要受审判，那么那些不肯接受敬拜祂的人，天主更要毁灭他们。因为那不愿学习、以免受审的人，其实已经受过审判：因为他知道他不愿听什么。因此，在洞悉人心的天主面前，这种想象不能作为托词。所以，要避免蛇所提示给你们的那些狡猾的念头。但若有人在真正的无知中结束此生，这指控将针对他：他活了这么久，却不知道谁是他食物的供应者；他像是一个无知、忘恩、极不配的仆人，被逐出天主的国。}\par

% structure: p0029 source=h2 target=subsection reason=relative-heading-rank; base=h2

\subsection{第十四章 多神论的比喻}

\Quote{再者，那可怕的蛇向你们提示这种假设，使你们思想并说出你们多数人所常说的话：我们知道有一位万有的主，但还有其他的神。因为正如有一位凯撒，但他下面有总督、巡抚、省长、千夫长、百夫长、十夫长；同样，有一位伟大的天主，如同有一位凯撒，同样，仿效较低级权柄的模式，有次等的诸神，虽然低于祂，却统治我们。因此，你们这些被这种观念如同可怕毒药所引诱的人——我是指此比喻的邪恶观念——要听，好知道什么是善、什么是恶。因为你们尚未看见，也不思想你们所说的话。}\par

% structure: p0031 source=h2 target=subsection reason=relative-heading-rank; base=h2

\subsection{第十五章 其不足为凭}

\Quote{因为如果你说，天主如同\Person{凯撒}一样拥有从属的权能——即那些被称为神的——你的比喻并没有达到目的。因为如果你遵循这个比喻，你必须知道，将\Person{凯撒}的名号给与他人——无论是执政官、总督、队长或任何其他人——是不允许的；而且，给出这名号的人不得存活，接受它的人也将被剪除。因此，按照你自己的比喻，天主的名字不得给与他人；而无论是诱惑人接受或给出这名字的人，都将被毁灭。既然如此，对人的这种侮辱尚且招致惩罚，那么那些称他人为神的人，岂不更当受永罚，因为他们侮辱了天主。这是理所当然的，因为你把你所能施加的一切侮辱加于那交托给你、为尊崇祂的独一统治而尊荣的名号。因为\Emph{天主}并非祂本来的名字；但你暂时领受了它，却侮辱了所赐之物，使这行为算作是针对真名而行的，正如你所使用的那样。但你却将它置于各种侮辱之下。}\par

% structure: p0033 source=h2 target=subsection reason=relative-heading-rank; base=h2

\subsection{第十六章 埃及人的神}

\Quote{因此，你们这些埃及人的魁首，自夸懂得气象学，并许诺判断星宿的本质，却因潜伏在他们心中的邪恶意见，尽可能地使那名号遭受各种不名誉。因为他们中有人教导崇拜一头称为\Person{阿匹斯}的牛，有人崇拜公山羊，有人崇拜猫，有人崇拜蛇；连鱼、洋葱、肚子里的咕噜声、公共下水道、牲畜的肢体，以及成千上万其他卑劣可憎之物，他们都称之为\Emph{神}。}\par

% structure: p0035 source=h2 target=subsection reason=relative-heading-rank; base=h2

\subsection{第十七章 埃及人为其体系的辩护}

\Person{伯多禄}说了这话，周围的人群大笑起来。于是\Person{伯多禄}对发笑的人群说：\Quote{你们嘲笑他们的行为，却不知道你们自己更成为他们的笑柄。但你们互相嘲笑对方的行为；因为受邪恶习俗引导而陷入欺骗，你们看不见自己的行为。不过我承认，你们有理由嘲笑埃及人的偶像，因为他们有理性，却崇拜无理性的动物，而且这些动物都是会死的。但请听他们嘲笑你们时所说的话。\InnerQuote{他们说，我们虽然崇拜会死的生物，但这些生物毕竟曾经有生命；而你们却尊崇从未活过的东西。}}\par

% structure: p0037 source=h2 target=subsection reason=relative-heading-rank; base=h2

\subsection{第十八章 回答埃及人}

\Quote{因此，这样回答他们：你们说谎，因为你们崇拜这些东西并不是为了尊崇真天主，否则你们所有人都会崇拜每一种形象，而不是像你们所做的这样。因为你们中那些认为洋葱是神明的，以及那些崇拜肚子咕噜声的，彼此争斗；同样，所有的人都偏爱某一种事物，而辱骂他人所偏爱的。而且，由于不同的判断，有人崇拜同一动物的这个肢体，有人崇拜那个肢体。此外，他们中那些仍然存有一丝正确理性的人，因明显的卑劣而感到羞愧，试图将这些推入寓意解释，希望以另一种妄想来确立他们致命的错误。但我们若在那里，就应当驳斥这些寓意解释，因为对寓意的愚蠢热情已经盛行到成为理解力的一大疾病。因为不必在身体的整个部分敷上膏药，而只应敷在患病的部分。既然你们嘲笑埃及人，表明你们没有染上他们的疾病，那么对于你们自己的疾病，我理当立刻为你们提供当前的治疗。}\par

% structure: p0039 source=h2 target=subsection reason=relative-heading-rank; base=h2

\subsection{第十九章 天主的独有属性}

\Quote{凡愿意崇拜天主的人，首先应当知道什么唯独属于天主的本性，不能归于其他任何事物，以便他注视天主的独有性，而不在其他事物中发现它，这样就不至于被诱骗而将神性归于其他事物。然而，这唯独属于天主：唯有祂既是万有的创造者，又是万有的至善者。那创造者，在能力上确实超越被造之物；那无限的，在量上超越有限的；在美善方面，祂是最美的；在幸福方面，祂是最有福的；在智慧方面，祂是最完美的。同样地，在其他方面，祂有着无可比拟的优越性。既然，如我所说，这本身——即作为万有的至善者——唯独属于天主，而包容万有的世界是由祂所造，那么祂所造的万物中没有一样能与祂同等相比。}\par

% structure: p0041 source=h2 target=subsection reason=relative-heading-rank; base=h2

\subsection{第二十章 世界或其任何部分都不能是天主}

\Quote{但世界，既然不是不可比拟、不可超越的，也不是在各方面完全无缺陷的，就不能是天主。如果整个世界因被造而不能是天主，那么它的各部分——我指你们称为神的那些，由金银铜石或其他任何材料制成，且由人手所造——岂不更应合理地不被称作天主？然而，让我们进一步看看，那可怕的毒蛇如何藉人的口毒害那些被其诱惑所引诱的人。}\par

% structure: p0043 source=h2 target=subsection reason=relative-heading-rank; base=h2

\subsection{第二十一章 偶像并非由神性之灵赋予生命}

\Quote{因为许多人说：我们敬拜的对象，我们并不敬拜其中的金、银、木、石。因为我们也知道这些不过是无生命的物质和必死之人的技艺。但住在它们里面的神灵，我们称之为天主。看那些如此说话之人的不道德！因为当那显现出来的东西被轻易证明为虚无时，他们就诉诸不可见之物，仿佛无法在非显明之事上被定罪。然而，他们部分同意我们，即他们偶像的一半不是天主，而是无感觉的物质。剩下的是，他们需要向我们证明，我们如何相信这些偶像有神圣之神。但他们无法向我们证明事实如此，因为它并非如此；当他们说他们看见了时，我们也不相信他们。我们将向他们提供证据，证明他们没有神圣之神，以便喜爱真理的人听了关于这些偶像是有灵的想法的驳斥后，可以远离那有害的迷惑。}\par

% structure: p0045 source=h2 target=subsection reason=relative-heading-rank; base=h2

\subsection{第二十二章 驳偶像崇拜}

\Quote{首先，如果你们敬拜它们当作有灵的，为什么你们也敬拜古代名人的坟墓，那些名人显然没有神圣之神？因此你们对此根本不说真话。但如果你们的敬拜对象真是有灵的，它们会自行移动，会有声音，会抖掉身上的蜘蛛，会赶走那些想要突袭和偷窃它们的人，会轻易抓获那些盗取供物的人。但现在它们不做任何这些事，反而像罪犯一样被看守，尤其是那些更昂贵的，如我们已说过的。还有，难道统治者为它们向你们征收赋税和税金，仿佛你们因它们大得利益？还有，难道它们不是常常被敌人掠夺，被打碎、分散？难道祭司们不比外面的敬拜者更多偷走许多供物，从而承认他们敬拜的无用？}\par

% structure: p0047 source=h2 target=subsection reason=relative-heading-rank; base=h2

\subsection{第二十三章 偶像崇拜的愚昧}

\Quote{不，有人会说；但它们被其预见所察觉。这是假的；因为有多少未被察觉？如果因为某些被捕获就声称它们有能力，那是错误。因为那些盗墓者，有些被发现，有些逃脱；但被捕者被察觉不是由于死者的能力。对于偷窃和盗窃神像者，我们也应如此结论。但有人会说，住在它们里面的神明不顾它们的偶像。那么，为什么你们照料它们，擦拭、清洗、打磨、加冠、献祭？所以请同意我，你们完全行事没有正当理性。因为你们像为死者哀悼一样，向你们的神献祭和奠酒。}\par

% structure: p0049 source=h2 target=subsection reason=relative-heading-rank; base=h2

\subsection{第二十四章 偶像的无能}

\Quote{然而，用恺撒和他手下权柄的比喻来称它们为管理者也并不协调；而你们却像我所说的，悉心照料它们，从各方面保养你们的偶像。因为它们毫无能力，什么也不做。所以告诉我们，它们管理什么？它们做了什么类似各地统治者所做的事？它们发挥了什么影响力，如同天主的星辰？它们能像太阳那样展示什么吗？还是你们在它们面前点灯？它们能像云彩降雨那样带来雨水吗？——它们连自己都不能移动，除非人抬着它们？它们能使土地为你们的劳作而丰饶吗？——你们却向它们供献祭品？因此它们什么也不能做。}\par

% structure: p0051 source=h2 target=subsection reason=relative-heading-rank; base=h2

\subsection{第二十五章 仆人成为主人}

\Quote{但如果它们能做什么，你们也不应当称它们为神：因为称元素为神是不对的，即便它们提供美好的事物；只有那安排它们为我们所用、命令它们服务于人的那一位——我们严格地称祂为天主，但你们没有察觉到祂的恩惠，却允许那些被分配给你们作仆人的元素统治你们。我何必谈论元素呢，你们不仅制造并敬拜无生命的偶像，而且屈尊在一切方面作为仆人服从它们？因此，由于你们错误的判断，你们成了魔鬼的奴仆。然而，通过承认天主本身，通过善行，你们可以再次成为主人，命令魔鬼如奴隶，并作为天主的儿子成为永恒国度的嗣子。}\par

% structure: p0053 source=h2 target=subsection reason=relative-heading-rank; base=h2

\subsection{第二十六章 病人得痊愈}

说了这话，他吩咐将附魔的和患病的人带到他面前；他们被带来时，他按手在他们身上，祈祷后，治好他们便遣散了，提醒他们和其余群众每天都要在他讲道的地方等候他。然后，当其他人退去后，伯多禄与那些愿意的人在那里的蓄水池中沐浴；随后，他吩咐在树荫浓密的树下地面上铺好桌子，为了荫凉，他按照我们的位份吩咐我们各自斜躺，于是我们进食。因此，按照希伯来人惯常的信德，他祝福并感谢天主，为这享用；我们面前还有很长的时间，他允许我们随意向他提问；虽然我们有二十人从四面八方向他提问，他使每个人都满意。现在夜晚降临，我们都随他进入住宿处最大的房间，我们都在那里睡了。\par

\SourceRecord{Translated by Peter Peterson.；Ante-Nicene Fathers；Vol. 8.；Edited by Alexander Roberts, James Donaldson, and A. Cleveland Coxe.；Buffalo, NY: Christian Literature Publishing Co.,；1886.；<http://www.newadvent.org/fathers/080810.htm>.}

% structure: p0001 source=h1 target=section reason=page-title

\section{讲道集十一}

% structure: p0002 source=h2 target=subsection reason=relative-heading-rank; base=h2

\subsection{第一章 清晨操练}

\Emph{因此}，在第四天，在\Place{的黎波里}，\Person{伯多禄}起身，发现我们醒着，向我们致意后，便出去到蓄水池那里，以便沐浴和祈祷；我们也跟随他这样做。当我们一起祈祷后，坐在他面前时，他就给我们讲论关于纯洁的必要性。此后，白天到来时，他让群众进来。当一大群人进来后，他照常向他们致意，并开始发言。\par

% structure: p0004 source=h2 target=subsection reason=relative-heading-rank; base=h2

\subsection{第二章 \Quote{尽力}}

\Quote{由于你们长久以来的疏忽，对你们自身造成损害，你们的心智在宗教方面生出了许多有害的观念，你们变得如同农夫疏忽所致的休耕地，所以你们需要长时间进行净化，好使你们的心智像接受好种子一样接受所传授的真道，不被邪恶的忧虑所窒息，以致无法结出能拯救你们的善工之果。因此，那些关心自己救恩的人应当更恒常地聆听，好使那长久积累的罪过，在余下的短暂时光里，得以用恒常的净化关怀来相称。既然没有人知道自己末时的界限，你们要赶紧拔除心中众多的荆棘；但不能一点点地拔，因为那样你们无法被净化，因为你们已经休耕很久了。}\par

% structure: p0006 source=h2 target=subsection reason=relative-heading-rank; base=h2

\subsection{第三章 \Quote{看，何等义愤}}

\Quote{但你们若不这样罚，就不会忍受为自己的净化付出极大的关怀，除非你们对自己发怒，并因那些事而惩罚自己——你们作为无用的仆人，曾被那些事所困，顺从了邪恶的情欲——好使你们能将义愤如火烧休耕地般注入你们的心智。因此，若你们没有义火，即对邪恶情欲的义愤，就请学习你们曾被什么美善所欺骗，被谁所蒙蔽，你们预备了怎样的惩罚；这样，你们的心智清醒，因那位派遣我们者的教导而点燃如同义火，才能烧尽情欲的邪恶。请相信我，如果你们愿意，你们能整顿一切。}\par

% structure: p0008 source=h2 target=subsection reason=relative-heading-rank; base=h2

\subsection{第四章 黄金律}

\Quote{你们是不可见的天主的肖像。因此，凡愿虔诚的人不可说偶像是天主的肖像，因而敬拜它们是正当的。因为天主的肖像就是人。凡愿对天主虔诚的人，应向人行善，因为人的身体带有天主的肖像。但并非所有人都已拥有祂的模样，而是善良灵魂的纯洁心智才有。然而，我们知道人是按天主的肖像和模样造的，我们告诉你们要对他虔诚，好使这恩惠算作对那肖像所属的天主所做的。因此，你们应当尊敬天主的肖像——即人——如此：给饥饿者食物，给口渴者饮料，给裸者衣服，照顾病人，给陌生人住所，探望狱中的人并尽己所能帮助他。简而言之，任何人希望自己得到什么美善，就应当同样供给有需要的他人，这样他就能被算作得到好报酬，如同对天主的肖像虔诚一样。同样，若他不愿做这些事，他将因忽视肖像而受惩罚。}\par

% structure: p0010 source=h2 target=subsection reason=relative-heading-rank; base=h2

\subsection{第五章 你们为我最小的兄弟所做的}

\Quote{难道可以说，为了对天主的虔诚，你们敬拜每一种形状，却事事伤害那真正是天主肖像的人，犯下谋杀、奸淫、偷窃，并在许多其他方面羞辱他吗？但你们连一件令人忧伤的恶事也不应做。然而如今你们做尽一切令人沮丧的事，因为不义也是痛苦。因此你们杀人、抢夺财物，以及任何你们自己所不愿承受的事。但你们被某种恶毒的爬虫诱惑走向恶意，藉着多神论教义的暗示，对真正的肖像——即人——不虔敬，却以为自己在对无感觉的事物虔诚。}\par

% structure: p0012 source=h2 target=subsection reason=relative-heading-rank; base=h2

\subsection{第六章 为何天主容许偶像崇拜的对象存留}

\Quote{但有人说，若天主不意欲这些事物存在，它们就不会存在，祂会除去它们。但我这样说：当所有人都表明对祂的偏爱时，这事就必发生，而现今的世界也将改变。然而，若你们愿祂如此行，以致没有一件受敬拜之物存留，请告诉我，现存之物中有什么是你们没有敬拜过的？难道你们中有人不是敬拜太阳，有人敬拜月亮，有人敬拜水，有人敬拜地，有人敬拜山，有人敬拜植物，有人敬拜种子，有人甚至敬拜人，如在埃及吗？因此，天主甚至必须不容忍任何人乃至你们存在，好使既无受敬拜者也无敬拜者。这正是潜伏在你们中的那可怕之蛇所愿望的，它毫不宽容你们。但事实并非如此。因为受敬拜之物并不犯罪；它只是被敬拜它的人所强迫。虽然众人作出不公的判决，但天主不会。因为受难者与安排者受同样惩罚并不公平，除非他甘愿接受那唯独属于至尊者的荣耀。}\par

% structure: p0014 source=h2 target=subsection reason=relative-heading-rank; base=h2

\subsection{第七章 \Quote{让两样一起长到收割}}

\Quote{但有人会说，敬拜者本人应由真天主除灭，以免别人再这样做。但你们并不比天主更聪明，不能作为比祂更谨慎者向祂献计。祂知道祂所做的事；因为祂对一切不虔敬者长久忍耐，如同仁慈而爱人的父亲，知道不虔敬的人会变为虔诚。而那些敬拜卑劣无感觉之物的敬拜者中，许多人变得清醒，停止敬拜这些事物并停止犯罪，许多希腊人得救而向真天主祈祷。}\par

% structure: p0016 source=h2 target=subsection reason=relative-heading-rank; base=h2

\subsection{第八章 自由与必然}

\Quote{但是，你会说，天主当初应该造得我们根本不会想到这些事。说这话的人不知道什么是自由意志，也不知道怎样才算是真正的善；因为一个人出于自己的选择而成为善的，才是真正的善；但若因被迫而由别人使之成为善的，便不是真正的善，因为他的善并非出于自己的选择。既然每个人的自由构成了真正的善，也显明了真正的恶，天主便藉各种机遇使各人心中产生友谊或敌意。然而，有人说，我们所想的一切都是祂使我们想的。且慢！你为何越说越亵渎？因为如果我们所想的一切都受祂影响，那你就是说祂是淫乱、情欲、贪婪和一切亵渎的根源。住口吧，你们这些本应赞美祂、将一切尊荣归于祂的人！不要说天主不要求任何尊荣；因为即使祂自己不要求什么，你也应当看什么是正当的，并以感恩的声音回应那位在一切事上施恩于你的天主。}\par

% structure: p0018 source=h2 target=subsection reason=relative-heading-rank; base=h2

\subsection{第九章　天主是忌邪的天主}

\Quote{但是，你会说，我们同时向祂和所有其他对象感恩，这样做更好。现在，当你说这话时，你并不知道那针对你所设的诡计。因为正如许多无能的医生应允治好一个病人，而那位真正能治好他的人却不施治，认为若他治好了他，别人会得功劳；同样，天主当被与许多无能为力者一同祈求时，祂就不施恩于你。怎么？有人会说，天主因此发怒吗？若祂治愈了，别人却得功劳？我回答：即使祂不愤慨，至少祂不会成为欺诈的同谋；因为当祂赐予恩惠时，那什么也不曾做的偶像却被归功。而且我也对你说，若那向无知偶像跪拜的人本性上未受损害，或许天主也会容忍这事。因此，你们要警醒，使你们能在救恩之事上获得合理的理解。因为天主毫无所需，祂自身既不缺什么，也不受损害；受益或受损的乃是我们自己。正如凯撒既不会因受辱骂而受损，也不会因受感谢而获益，而感谢者得安全，辱骂者遭毁灭；同样，赞美天主的人并未使祂获益，却救了自己；同样，亵渎祂的人并未损害祂，却自取灭亡。}\par

% structure: p0020 source=h2 target=subsection reason=relative-heading-rank; base=h2

\subsection{第十章　受造物为天主的正义复仇}

\Quote{但有人会说，天主与人之间的情况并不相同；我承认它们并不相同：因为对更大者行不敬的惩罚更大，对较小者犯罪的惩罚较小。既然天主是最大的，那么对祂不敬的人就要承受更大的惩罚，因为他是向更大的犯罪；这不是因为祂亲手自卫，而是因为整个受造界都向他发怒，并自然地报复他。因为对亵渎者，太阳不给他光，大地不给他果实，泉水不给他水，在阴府中那被设立为王的也不给灵魂安息；因为即使在现今世界秩序尚存之时，整个受造界也向他发怒。因此，云彩不降下足够的雨水，大地也不结果实，于是许多人丧亡；甚至空气本身也因愤怒而转为瘟疫之气。然而，我们所享有的一切美物，是祂出于仁慈迫使受造物为我们谋益。但是，对于你——那羞辱万物创造者的人——整个受造界都与你为敌。}\par

% structure: p0022 source=h2 target=subsection reason=relative-heading-rank; base=h2

\subsection{第十一章　灵魂的不朽}

\Quote{即使因身体的分解你逃脱了惩罚，你怎能藉朽坏逃脱你那不朽的灵魂呢？因为恶人的灵魂也是不朽的，对他们来说，没有这不朽倒更好。因为他们在不灭的火中受无尽折磨，永远不死，他们的痛苦永无终结。但或许你们中有人说：\InnerQuote{伯多禄，你恐吓我们。}那么请教我们，我们怎能对这事缄默，却又如实告诉你们？因为若不如此，我们无法告诉你们。但若我们说话，我们又被怀疑以虚假的理论恐吓你们。那么，我们如何能安抚那潜伏在你们灵魂中的恶蛇，它狡猾地以爱天主为伪装，播下敌视天主的怀疑呢？你们要与自己和好；因为为得救，你们必须以善行归向祂。你们心中不合理的私欲与天主为敌，因为它以智慧的幻觉增强了无知。}\par

% structure: p0024 source=h2 target=subsection reason=relative-heading-rank; base=h2

\subsection{第十二章　偶像的无益}

\Quote{另有些人说，天主不关心我们。这也是假的。因为若祂真不关心，就不会使祂的太阳升起，光照善人，也光照恶人；降雨给义人，也给不义的人。但另有些人说，我们比你们更虔诚，因为我们既敬拜祂，也敬拜像。我不认为，若有人对一位国王说：\InnerQuote{我把你与尸体和无用的粪土同等地尊崇}——我认爲他能因此获益。但有人会说：你把我们崇拜的对象称为粪土吗？我说是的，因为你们把它们用于崇拜而使它们对自己无益，而它们的材料或许本可用于其他用途，或用作肥料。但如今连这用途也不成，因为你们改变了它们的形状而崇拜它们。你们这些最邪恶的人，怎能说你们更虔诚呢？你们若持守这唯一的、无可比拟的罪，凭这罪就理应在那位真实者手中灵魂灭亡。因为正如一个儿子从父亲那里领受了许多恩惠，却将应归于父亲的光荣给与另一个不是他父亲的人，他当然被剥夺继承权；但若他按父亲的判断生活，并以此感谢父亲的恩惠，他就有正当理由成为继承人。}\par

% structure: p0026 source=h2 target=subsection reason=relative-heading-rank; base=h2

\subsection{第十三章　回答支持偶像崇拜的论点}

\Quote{但有人说，如果我们抛弃祖先传下来的崇拜对象，就是亵渎；因为那就像保护一笔存款。但按照这个原则，强盗或放荡者的儿子就不该清醒并选择更好的部分，以免行事亵渎，因与父母不同而犯罪！那么，那些说\InnerQuote{我们崇拜这些东西是为了不使他麻烦}的人是多么愚蠢啊！好像天主因祝福他的人而麻烦，因忘恩负义亵渎他的人而不麻烦？那么，为什么当雨水停止时，你只仰望天空，倾注祈祷和恳求；而当得到雨水时，你又很快忘记？因为当你收割庄稼或采集葡萄时，你将初果分给那些虚无的偶像，很快忘记了你的恩主天主；这样你就进入树林和庙宇，祭献并举行盛宴。因此你们中有人说，这些东西是为了欢乐和宴饮而精心设计的。}\par

% structure: p0028 source=h2 target=subsection reason=relative-heading-rank; base=h2

\subsection{第十四章  异教狂欢}

\Quote{哦，没有理解力的人！你们要正确判断所说的。因为如果有必要为了身体的休息而放松享乐，那么是在河流、树林和树丛中，那里有娱乐、欢宴和阴凉之处更好，还是在有鬼魔的疯狂、砍手、阉割、狂怒、疯癫、披头散发、呼喊、狂热和嚎叫，以及所有那些为了迷惑无思想的人而虚伪地做的事情更好，当你们向那些比死者更死的东西献上规定的祈祷和感谢时？}\par

% structure: p0030 source=h2 target=subsection reason=relative-heading-rank; base=h2

\subsection{第十五章  异教崇拜者在恶魔权下}

\Quote{你们为什么以这些事为乐？因为潜伏在你们里面的蛇，它在你们里面播下了无果的欲望，不会告诉你们，我要说出来并记录下来。情况就是这样。按照对天主的崇拜，宣告要清醒、要贞洁、要节制情欲、不偷窃他人的财物、正直、适度、无畏、温柔地生活；宁可节制自己的需要，也不通过不正当夺取他人财产来满足自己的需要。但在那些所谓的神祇那里，情况正好相反。你们否认一些事情，认为它们是\Emph{由你们所做的}，以便你们为自己的\Emph{你们的}义德而钦佩；然而，即使你们做了所有命令你们做的事，对天主的无知本身足以定你们的罪。但是，你们聚集在为他们奉献的地方，喜欢喝醉，你们点燃祭台，祭台散发的气味通过其影响力吸引盲目的和耳聋的精灵到那熏香的地方。这样，在场的人中，有些被灵感充满，有些被奇怪的魔鬼充满，有些沉溺于淫荡，有些偷窃和杀人。因为血的蒸发和酒的奠祭满足了这些潜伏在你们里面、使你们以那里所行的事为乐的污秽的精灵，并在梦中用虚假的幻象环绕你们，用无数的疾病惩罚你们。因为在所谓圣祭的幌子下，你们被可怕的魔鬼充满，它们狡猾地隐藏自己，毁灭你们，使你们不明白为你们设下的圈套。因为假借某种伤害、或爱、或愤怒、或悲伤、或用绳索勒死、或淹死、或从悬崖扔下、或自杀、或中风、或其他疾病，它们夺去你们的生命。}\par

% structure: p0032 source=h2 target=subsection reason=relative-heading-rank; base=h2

\subsection{第十六章  万事互相效力，叫爱天主的人得益处}

\Quote{但我们中没有一个人会遭受这样的事；相反，他们自己被我们惩罚，当他们进入任何人时，他们恳求我们让他们慢慢出去。但也许有人会说：即使一些崇拜天主的人也会遭受这样的痛苦。我说那是不可能的。因为我所说的崇拜天主的人，是真正虔诚的人，不是只在名义上是，而是真正实行所赐给他法律的行为。如果有人行事不虔，他就不虔诚；同样，如果另一个部落的人遵守法律，他就是犹太人；但不遵守法律的就是希腊人。因为犹太人信天主并遵守法律，凭借这种信德，他也能除去其他痛苦，尽管它们像山一样沉重。但不遵守法律的人显然是因为不信天主而叛逃；因此，作为非犹太人，而是罪人，他因自己的罪而服从于那些为惩罚罪人而规定的痛苦。因为按照起初所规定的天主旨意，惩罚公义地伴随着那些因过犯而崇拜他的人；这是为了，通过惩罚来清算他们的罪，就好像还债一样，他可以在公审判中把那些归向他的人纯洁地展现出来。因为正如恶人在这里享受奢华，丧失永福，同样，惩罚被降在犯罪的犹太人身上，作为清算，使他们在这里赎罪，以便在那里从永罚中得释放。}\par

% structure: p0034 source=h2 target=subsection reason=relative-heading-rank; base=h2

\subsection{第十七章  用爱心说诚实话}

\Quote{但你们不能这样说；因为你们不相信那时事情像我们所说的那样——我是指当有对所有人的报应时。因此，你们不知道什么是有益的，被暂时的快乐所引诱，未能抓住永恒的事物。因此，我们试图向你们展示有益的事，以便你们因对虔诚应许的确信，能通过善行与我们一同继承那没有悲伤的世界。在你们认识我们之前，不要向我们发怒，好像我们为你们所渴望的好处说了假话。因为那些被我们视为真实和良善的事，我们毫不迟疑地带给你们，相反，我们急于使你们成为我们看为良善之事的共同继承人。因为必须这样对不信的人说话。但我们在所说的中是否真的说真理，你们除非先以热爱真理的心去听，否则无法知道。}\par

% structure: p0036 source=h2 target=subsection reason=relative-heading-rank; base=h2

\subsection{第十八章  迷惑蛇}

\Quote{因此，就眼前之事而言，尽管潜伏在你内的蛇以千万种方式暗示邪恶的推理和阻碍，想要诱捕你，你便越发应当抵抗它，并勤勉地听从我们。因为你们曾受严重欺骗，如今正在思虑，理当知道如何驯服它。但别无他法。我所说的驯服，是指以理性对抗其邪恶的计谋，记住：它藉知识的许诺，起初将死亡带入了世界。}\par

% structure: p0038 source=h2 target=subsection reason=relative-heading-rank; base=h2

\subsection{第十九章 不是和平，乃是刀剑}

\Quote{因此，真理的先知知道世界大错特错，看见它站在邪恶一边，就不愿世界在错误中享有和平。所以直到末了，他反对所有与邪恶协和的人，将\Emph{真理}置于错误对面，仿佛向那些清醒的人投火，即对诱惑者的愤怒，这愤怒好比刀剑；他藉宣讲圣言，用知识摧毁无知，仿佛切割，将活人与死人分开。因此，当邪恶被合法的知识战胜时，战争便席卷了一切。因为为得救起见，顺服的儿子从不信的父亲中被分离出来，或父亲从儿子，或母亲从女儿，或女儿从母亲，亲戚从亲戚，朋友从同伴。}\par

% structure: p0040 source=h2 target=subsection reason=relative-heading-rank; base=h2

\subsection{第二十章 倘若它已经点燃呢？}

\Quote{且不要让任何人说：父母与子女分离、子女与父母分离，怎能是公义的？这是公义的，全然公义。因为倘若他们与他们同住，却无助于他们，反与他们一同灭亡，那么，愿意得救的人不愿被毁灭的人分开，岂不公义？而且，不是那些判断更正确的人想要分离，而是他们愿意与他们同住，并用更好的教导使他们获益；故此，不信的人不愿听从他们，便与他们作战，驱逐、迫害、憎恨他们。但遭受这些事的人，怜悯那些被无知捆缚的人，以智慧的教导为那些对他们图谋恶事的人祈祷，因为他们知道无知是他们犯罪的缘由。因为导师自己，被钉在\Emph{十字架上}时，曾向圣父祈祷，求赦免杀害他者的罪，说：\BibleQuote{父啊，赦免他们的罪，因为他们不知道他们所做的是什么。}他们因此也效法导师，在受苦时为那些图谋恶事的人祈祷，正如他们所学的。因此，他们并非因憎恨父母而被分离，因为他们不断祈祷，甚至为那些既非父母也非亲戚，而是仇敌的人祈祷，并努力爱他们，正如他们所受的命令。}\par

% structure: p0042 source=h2 target=subsection reason=relative-heading-rank; base=h2

\subsection{第二十一章——\Quote{我若是父亲，我的敬畏在哪里？}}

\Quote{但你告诉我，你如何爱你的父母？如果你总是以合乎正理的方式爱他们，我祝贺你；但如果你随意地爱他们，那就不是这样，因为你可能因小事成为他们的仇敌。但如果你明智地爱他们，请告诉我：父母是什么？你会说他们是我们的生命之源。那么，为何你们不爱万有的\Emph{源头}呢？如果你们真正明白而选择了这样做？但你们现在又会说：我们没有见过祂。那么，你们为何不寻求祂，反倒崇拜无感觉之物？但是什么？即使认识天主对你们是困难的，你们总不能不知道什么不是天主，以至于能推理出：天主不是木头、石头、铜，也不是任何由可朽坏之物制成的。}\par

% structure: p0044 source=h2 target=subsection reason=relative-heading-rank; base=h2

\subsection{第二十二章——\Quote{那些没有创造诸天的神祇。}}

\Quote{因为难道它们不是用铁雕刻的吗？难道雕刻的铁没有被火软化吗？难道火本身不是被水熄灭的吗？难道水的运动不是来自这神吗？难道这神的运行之始不是来自创造万有的天主吗？因为先知梅瑟曾如此说：\BibleQuote{起初天主创造了天地。大地是无形且荒芜的；黑暗笼罩深渊；天主的神在水面上运行。}这神，如同天主的手，遵从天主的命令造作万物，将光与黑暗分开，并在不可见的天之后铺展开可见的天，使高处由光明天使居住，低处由人和为他所用的一切受造物居住。}\par

% structure: p0046 source=h2 target=subsection reason=relative-heading-rank; base=h2

\subsection{第二十三章——\Quote{多给谁。}}

\Quote{因为为你，人啊，天主命令水退到地面，使大地能为你们结果实。祂开辟水道，为你们预备泉源，显露河床，使动物滋生；总之，祂为你们供应了一切。难道不是为你们，风才吹动，雨才降下，季节才变化以出产果实吗？再者，太阳、月亮和其他天体为你们完成出没；河流、池塘和一切泉源都为你们服务。因此，无知的人啊，既然给了你更大的尊荣，也为你这忘恩负义者预备了更大的火罚，因为你竟不肯认识那最当认识的一位。}\par

% structure: p0048 source=h2 target=subsection reason=relative-heading-rank; base=h2

\subsection{第二十四章——\Quote{由水而生。}}

\Quote{现在从低等事物中学习万有的起因：推理——水造成万物，水的运动由这神推动，这神之始来自万有的天主。你本该这样推理，以便藉理性达致天主，认识你的本源，并由首生之水重生，成为生你为不朽之父母的嗣子。}\par

% structure: p0050 source=h2 target=subsection reason=relative-heading-rank; base=h2

\subsection{第二十五章 善工须善行}

\Quote{所以，你要像儿子到父亲面前一样欣然前往，好使天主以无知为你罪过的借口。但若你被召唤后不肯来，或迟延，你将被天主公义的审判所毁灭，因你不愿，你便不被天主所愿。也不要以为，纵然你比所有曾有过的虔敬者更虔敬，但若未领洗，你就永远有希望。恰相反，正因如此，你要受更大的惩罚，因为你行了卓越的善工，却非以卓越的方式。行善只有在遵从天主的命令时才是卓越的。但你若不愿按祂的旨意领洗，你便是事奉自己的意愿，反对祂的旨意。}\par

% structure: p0052 source=h2 target=subsection reason=relative-heading-rank; base=h2

\subsection{第二十六章 圣洗圣事}

\Quote{但或许有人会说：用水受洗对虔敬有什么贡献？首先，因为你做了天主所喜悦的事；其次，你由水重生归于天主，因敬畏而改变你最初由情欲所生的本性，从而能获得救恩。否则，是不可能的。因为先知曾向我们起誓说：\InnerQuote{我实在告诉你们，你们若不由活水重生，因圣父、圣子及圣神之名，决不能进入天国。}所以，上前来吧。因为那里有从起初就慈悲的事物，浮在水上，借三圣呼求，拯救那些受洗者脱离将来的刑罚，并将受洗者在受洗后所行的善工作为礼品奉献给天主。因此，要逃向水，因为惟有水能熄灭烈火的威力。那现在不来的人，仍怀着争竞之灵，因此他不为使自己得救而接近活水。}\par

% structure: p0054 source=h2 target=subsection reason=relative-heading-rank; base=h2

\subsection{第二十七章 人人需要圣洗}

\Quote{所以，上前来吧，无论你是义人还是不义之人。因为你若是义人，只缺圣洗便可获得救恩。你若是不义之人，就来受洗，以赦免从前无知所犯的罪。对于不义之人，受洗后的善行必须与他\Emph{以前}的不虔敬程度相称。因此，无论你是义人还是不义之人，要赶快由天主重生，因为迟延带来危险——死亡预定之期并未显明；并要借善行显示你与那由水生你的父的相似。作为真理的爱好者，要尊崇真天主为父。祂的尊荣就在于你按祂——那义者——的意愿生活。那义者的旨意是叫你不行不义。不义之事就是杀害、仇恨、嫉妒之类，其形式很多。}\par

% structure: p0056 source=h2 target=subsection reason=relative-heading-rank; base=h2

\subsection{第二十八章 洁净}

\Quote{然而，还需对这些事加上一些与人无关、唯独与敬拜天主相关的事。我指的是洁净，即不与行经的妻子亲近，因为天主的律法如此命令。但还有其他事呢？若不加洁净于事奉天主的事上，你就会像粪蝇一样愉快地翻滚。因此，作为人，你拥有比无理性动物更多的东西，即理性，当用属天的理性洁净心灵脱离邪恶，并在沐浴中洗净身体。因为按照真理的洁净，不是身体的洁净先于心灵的洁净，而是洁净随善行而至。因为我们的导师也曾\Emph{论及}我们中间某些法利塞人和经师，他们自别于众人，作为经师比别人更通晓律法之事，但祂仍责备他们是伪善者，因为他们只洁净在人前显眼的事物，却忽略了心灵的洁净和那些唯独天主所见的事物。}\par

% structure: p0058 source=h2 target=subsection reason=relative-heading-rank; base=h2

\subsection{第二十九章 外表与内在的洁净}

\Quote{因此，祂用了这句值得纪念的话，说的是他们中的伪善者的实情，并非对所有人。因为祂曾对一些人说，应当服从他们，因为他们受托于梅瑟的座位。然而，对伪善者祂说：\BibleQuote{祸哉，你们经师和法利塞假善人！因为你们洗净杯盘的外面，里面却满是污秽。瞎眼的法利塞人！先洗净杯盘的里面，好使外面也洁净。}这确实如此：因为当理智被知识光照，门徒便能行善，随之而来的就是洁净；因为内在的明悟产生对外在身体良好的照料。正如对身体的疏忽不能产生对理智的照料，同样，洁净的人能洁净外在和内在。而那洁净外在事物，却是为得人的称赞而行，依靠旁观者的称赞，他便从天主一无所得。}\par

% structure: p0060 source=h2 target=subsection reason=relative-heading-rank; base=h2

\subsection{第三十章 \Quote{凡纯洁的事}}

\Quote{但有谁不明白，在女人经期不与妻子同房，而是保持洁净并沐浴，是更好的呢？并且在同房之后也应当沐浴。但若你不愿行这事，就当回想你曾事奉无知觉的偶像时，你追求的那些洁净的方面；并应感到羞愧：如今，当需要达到——我不说更多——而是达到那唯一而完全的洁净时，你反而更加怠惰。因此，要思想那造你的，你就会明白是谁将这种懈怠降在你身上。}\par

% structure: p0062 source=h2 target=subsection reason=relative-heading-rank; base=h2

\subsection{第三十一章——\Quote{你们做了什么比别人更多的事？}}

\Quote{但你们中有人会说：那么我们必须做我们仍为偶像崇拜者时所做的一切事吗？我告诉你们：不是一切，而是凡你们那时做得好的事，现在也必须做，并且做得更多；因为在谬误中行得正的事都依附于真理，正如在真理中行了不正的事，便是来自谬误。因此，要从各方领受属于你们自己的事物，而不要领受他人的；也不要说：若那些在谬误中的人做得好，我们就不必做。因为按此原则，若那拜偶像的人不杀人，我们就该杀人吗？}\par

% structure: p0064 source=h2 target=subsection reason=relative-heading-rank; base=h2

\subsection{第三十二章 \Quote{多给谁}}

\Quote{不，反而：如果那些在错误中的人不杀人，我们也不该发怒；如果那在错误中的人不犯奸淫，我们连最小程度上也不该贪恋；如果那在错误中的人爱那爱他的人，我们就该爱那恨我们的人；如果那在错误中的人借给那些有财的人，我们就该\Emph{给}那些没有的人。毫无疑问，我们——那些希望继承永生的人——应当比那些只认识今生的人所行的善事做得更好，知道如果他们的行为在审判之日与我们的行为相比较，发现同样良善，我们将会蒙羞，而他们将会丧亡，因为他们因错误而行事与自身相悖。我说我们会因此蒙羞，因为我们没有比他们做得更多，尽管我们比他们知道得更多。如果我们因表现出与他们相等的善行而没有更多而蒙羞，那么如果我们表现出比他们的善行更少，又该如何呢？}\par

% structure: p0066 source=h2 target=subsection reason=relative-heading-rank; base=h2

\subsection{第三十三章 南方女王和尼尼微人}

\Quote{但在审判之日，那些认识真理者的行为将与那些在错误中者的善行相比较，这不撒谎的那位亲自教导了我们，祂对那些疏忽前来听祂的人说：\InnerQuote{南方女王将同这一代人起来，并要定他们的罪；因为她从地极而来，要听撒罗满的智慧；看，这里有一位比撒罗满更大的，}你们却不信祂。又对那百姓中不愿因祂的宣讲而悔改的人说：\InnerQuote{尼尼微人将同这一代人起来，并要定他们的罪，因为他们听了约纳的宣讲而悔改了；看，这里有一位更大的，却无人相信。}因此，祂将那些在外邦人中\Emph{行善}的人与所有不虔敬的人相对照，为的是定那些拥有真宗教却不如那些在错误中的人行得好的罪，祂劝勉那些有理智的人不仅要在外邦人所行的各种卓越之事上与它们相等，而且要超过他们。这番话是因着尊重分离和交合后洗濯、以及不否认这种纯洁的必要性而引出的，尽管那些在错误中的人也做同样的事，因为在错误中行善却不获救的人，是为了定罪那些在敬拜天主中\Emph{行恶}的人；因为他们对纯洁的尊重是出于错误，而不是出于对真正万有的父和天主的敬拜。}\par

% structure: p0068 source=h2 target=subsection reason=relative-heading-rank; base=h2

\subsection{第三十四章 伯多禄的日常工作}

说了这些话后，他遣散了群众；并按照他的习惯，与最亲近的人一起吃了饭，然后去休息。他日复一日地如此行事和讲论，有力地巩固了天主的法律，以所谓的\OriginalTerm{创世记}{Genesis}挑战所谓的神祇，论辩说没有自动作用，而是世界按照天主眷顾被治理。\par

% structure: p0070 source=h2 target=subsection reason=relative-heading-rank; base=h2

\subsection{第三十五章 \Quote{你们要提防假先知。}}

然后满了三个月，他命令我禁食数日，然后带我到靠近海边的泉源，并在那如永流的水中为我授洗。因此当我们的弟兄因我天主所赐的重生而欢喜时，没过几天，他在全体教会面前转向长老们，嘱咐他们说：\Quote{我们的主和先知，那派遣我们的，向我们宣告说，那恶者与祂辩论了四十天，丝毫没有胜过祂，便应许要从他的臣仆中派遣宗徒来欺骗。因此，最重要的是，要记得躲避宗徒、教师或先知，凡不首先将他的宣讲与那被称为我主的兄弟的雅各伯的宣讲\Emph{加以}准确比较的——雅各伯被委托管理耶路撒冷希伯来人的教会——即使他带着见证人来到你们这里：免得那与主辩论了四十天却一无所胜的邪恶，后来像闪电从天上落到地上一样，派一个宣讲者来损害你们，正如他现在派了西满到我们这里来，在真理的伪装下，以主的名宣讲，撒播错误。因此那派遣我们的说：\InnerQuote{有许多人要到我这里来，外面披着羊皮，里面却是残暴的狼。凭他们的果实你们能认出他们来。}}\par

% structure: p0072 source=h2 target=subsection reason=relative-heading-rank; base=h2

\subsection{第三十六章 告别提洛波里}

说了这些话后，他打发先驱往叙利亚的安提约基雅，嘱咐他们立即在那里等他。当他们走了以后，伯多禄为许多被说服的群众驱除了疾病、痛苦和魔鬼，并在靠近海边的泉源中为他们授洗，举行了圣体圣事，又任命了接待他进入家中的玛若内斯——当时已成全——作他们的主教，并选立了十二位长老，指定了执事，安排了关于寡妇的事务，就教会的秩序讲了有益的公共道理，劝勉他们服从主教玛若内斯；三个月满了之后，他向腓尼基的提洛波里告别，启程往叙利亚的安提约基雅，所有的人都以应有的尊敬陪伴我们。\par

\SourceRecord{Translated by Peter Peterson.；Ante-Nicene Fathers；Vol. 8.；Edited by Alexander Roberts, James Donaldson, and A. Cleveland Coxe.；Buffalo, NY: Christian Literature Publishing Co.,；1886.；<http://www.newadvent.org/fathers/080811.htm>.}

% structure: p0001 source=h1 target=section reason=page-title

\section{第十二篇讲道}

% structure: p0002 source=h2 target=subsection reason=relative-heading-rank; base=h2

\subsection{第一章 两队}

\Emph{因此}，从腓尼基的的黎波里出发前往叙利亚的安提约基雅，同一天我们到达奥尔塔西亚，并在那里停留。由于它靠近我们刚离开的城市，几乎所有人都已听过宣讲，我们只在那里停留一天，启程前往安塔拉杜斯。由于有许多人与我们同行，伯多禄对尼塞特和阿基拉说：\Quote{因为与我们同行的大批群众，当我们进城时，给我们招来不小的嫉妒，我以为必须作出安排，既不让这些人因被阻止与我们同行而悲伤，也不让我们因过于显眼而落入恶人的嫉妒。因此，我希望你们，尼塞特和阿基拉，分成两队走在我的前面，秘密进入外邦人的城市。}\par

% structure: p0004 source=h2 target=subsection reason=relative-heading-rank; base=h2

\subsection{第二章 宣讲者与归化者之爱}

\Quote{我知道，你们被要求这样做时感到痛苦，因为要离开我两天的路程。因此，你们要知道，我们劝服者爱你们被劝服者，远超过你们爱我们这些劝服你们的人。因此，我们既然彼此相爱，就不应不合理地做所愿之事，而要竭尽所能，为安全着想。因为我宁愿，你们也知道，进入各省更著名的城市，停留数日并讲论。目前，你们先领路进入邻近的劳狄刻雅，两三天后，就我意愿而言，我会赶上你们。只由你们一人在城门接我，以免混乱，这样我们可以与你们一同进入，没有骚动。然后，在那里停留数日后，其他人将按次序替你们前往更远的地方，为我们准备住处。}\par

% structure: p0006 source=h2 target=subsection reason=relative-heading-rank; base=h2

\subsection{第三章 顺服}

伯多禄说了这话后，他们不得不顺从，说道：\Quote{主，做这事并不完全使我们忧伤，因为这是你的命令；首先，因为你被天主眷顾所拣选，配得在一切事上思考和劝善；其次，大部分时间我们只是因必须走在你前面而与你分离两天。那确实是很长时间见不到你，哦，伯多禄，但我们若不想到那些被派得更远的人将更加忧伤，他们被命令在每个城市更久地等你，因更长时间看不到你渴望的容颜而痛苦，我们就更觉难忍。我们虽不亚于他们痛苦，却不反对，因为你命令我们这样做是有益的。}就这样，他们说完后前行，受命在第一站对随行群众说话，要他们各城分开进入。\par

% structure: p0008 source=h2 target=subsection reason=relative-heading-rank; base=h2

\subsection{第四章 克肋孟的喜乐}

因此，当他们离去后，我克肋孟非常喜乐，因为他命令我留在他身边。于是我回答说：\Quote{我感谢天主，你没有像打发其他人那样打发我走，否则我会悲伤而死。}但他却说：\Quote{但什么呢？如果将来为了教导的缘故，有必要打发你走，你会因为与我分离一小会儿且为有益的目的而死吗？难道你不该将忍耐为你安排之事的责任印刻在心，并欣然顺从？你难道不知道，朋友虽身体分离，却在记忆中共在；而有些人虽身体在一起，却因缺乏记忆而灵魂远离其朋友？}\par

% structure: p0010 source=h2 target=subsection reason=relative-heading-rank; base=h2

\subsection{第五章 克肋孟的服务职分}

于是我回答说：\Quote{不要以为，主，我会愚蠢地忍受那种悲伤，而是有充足的理由。因为我将你，主，视为一切：父亲、母亲、兄弟、亲戚，你是天主使我获得救恩真理的媒介，我将你视为一切，就有了最大的安慰。此外，我害怕我年轻的本性欲火，担忧若被你留下（我虽年轻，但有如此决心，若不跟你走必引起天主的愤怒），我会被欲火战胜。但由于留在我身边对我来说更好更安全，因为我的理智合理地致力于尊敬你，因此我祈求永远与你同在。此外，我记得你在凯撒勒雅说过：\InnerQuote{如果有人愿意与我同行，让他虔诚地同行。}而用\Emph{虔诚地}你意指那些献身于敬拜天主的人不应在关于天主的事上使任何人忧伤，例如离开父母、亲爱的妻子或其他人。因此，我在各方面都是你合适的旅伴，若你赐予我最大的恩惠，就允准我履行仆人的职分。}\par

% structure: p0012 source=h2 target=subsection reason=relative-heading-rank; base=h2

\subsection{第六章 伯多禄的俭朴}

于是，伯多禄听了，微笑着说道：\Quote{那么，克莱孟，你认为如何？难道你不认为你因着极大的需要，已被置于我仆人的地位吗？因为除了你，谁还能照顾我那些众多华丽的束腰外衣，以及我所有的戒指和凉鞋的更换？谁又能预备那些可口又精致的佳肴，它们种类繁多，需要许多技艺高超的厨师，以及所有那些被急切寻求、为柔弱之人的胃口预备的东西，如同为某只巨大的野兽预备的一样？然而，这样的一种选择临到你，或许是因为你并不了解或不知道我的生活方式：我只用面包和橄榄，偶尔用些蔬菜；我所穿的就是这一件外衣和斗篷；我不需要任何那些东西，也不需要别的：因为即使这些，我也绰绰有余。因为我的心灵，看到那里所有永恒的美好事物，就不关注这里的一切。不过，我接受你的善意；我佩服并称赞你，因为你这样一位有教养的人，竟如此轻易地将你的生活方式顺应了你的需要。因为从童年起，我和我的兄弟安德肋，他也是在主内我的兄弟，不仅是在孤儿状态下被抚养长大，而且因着贫穷与不幸惯于劳作，所以很容易承受我们当前旅途的不便。因此，如果你听从我，你当允许我，一个劳动者，为你履行仆人的角色。}\par

% structure: p0014 source=h2 target=subsection reason=relative-heading-rank; base=h2

\subsection{第七章。\Quote{不是要受服事，而是要服事}}

但我听了这话，就颤抖哭泣，因为这样的一句话竟出自一个人，而这世代的所有人在知识和虔诚方面都不及他。但他见我哭泣，就问我流泪的原因。于是我说：\Quote{我犯了什么罪，以致你对我说了这样的话？}伯多禄回答说：\Quote{如果我说作你的仆人有错，那么你请求作我的仆人，是首先有过错的。}于是我说：\Quote{这两种情形并不相同；因为这样做确实很合我的身份；但你，天主的使者，拯救我们灵魂的人，对我这样做，却是可怕的。}伯多禄回答说：\Quote{我本应同意你的看法，但我们的主，祂为拯救整个世界而来，唯独祂极其尊贵，却取了仆人的身分，为说服我们不要羞于为我们的弟兄履行仆人的服务，无论我们的出身多么高贵。}于是我说：\Quote{如果我想在辩论中胜过你，那我是愚蠢的。不过，我感谢天主的眷顾，使我得以认你为父。}\par

% structure: p0016 source=h2 target=subsection reason=relative-heading-rank; base=h2

\subsection{第八章。家族历史}

于是伯多禄问道：\Quote{那么，你的家族中真的只有你一个人吗？}我回答说：\Quote{的确有许多伟人，他们是凯撒的亲属。因此凯撒亲自将他自己家族的一位妻子赐给我的父亲，我父亲是他的义兄弟；她生了我们三个儿子，两个在我之前，他们是双胞胎，彼此十分相像，正如我父亲告诉我的。但我几乎不认识他们，也不认识我们的母亲，只带着他们模糊的印象，如同在梦中一般。我母亲名叫玛提狄亚，父亲名叫法乌斯督；两个兄弟一个叫法乌斯提诺，另一个叫法乌斯提尼雅诺。在我，他们的第三个儿子出生后，我母亲看见一个异象——这是我父亲告诉我的——\Emph{这异象告诉她}，除非她立刻带走她的双生儿子，离开罗马城流亡十二年，否则她和他们必遭毁灭性的命运而死。}\par

% structure: p0018 source=h2 target=subsection reason=relative-heading-rank; base=h2

\subsection{第九章。失踪的人}

\Quote{因此，我父亲因疼爱孩子，便为他们适当地配备了男女仆人，将他们送上船，连她一起送到雅典去读书，只留下我这个独生子在身边安慰他；为此我非常感恩，因为那异象没有命令我也随母亲离开罗马城。过了一年，我父亲寄钱到雅典给他们，同时打听他们的近况。但那些去办这事的人没有回来。第三年，我父亲忧心忡忡，又派人带着供给品前去，他们在第四年回来了，带来消息说，他们没有见到我母亲和兄弟，他们从未到达雅典，也没有找到任何同行者的踪迹。}\par

% structure: p0020 source=h2 target=subsection reason=relative-heading-rank; base=h2

\subsection{第十章。寻找者的失落}

\Quote{我父亲听了这事，因过度悲伤而茫然不知所措，不知该去哪里寻找他们，便常常带着我下到港口，向许多人打听，可有谁在四年前见过或听说过海难。有人说在这里，有人说在那里。他又问他们是否见过一具女尸带着\Emph{两}个孩子被冲上岸。当他们告诉他在许多地方见过许多尸体时，我父亲听了呻吟起来。但他情急之下，问了不合情理的问题，试图寻找如此广阔的海域。不过，这是情有可原的，因为出于对所寻找之人的爱，他以虚幻的希望为食。最后，他将我托付给监护人，在我十二岁时将我留在罗马，他自己含泪下到港口，登船出发去寻找。从那天直到今日，我没有收到他的任何信，也不知道他是死是活。但我更猜想他已死在某处，或因悲伤过度，或因船难而亡。证据就是，至今已二十年，我没有听到任何关于他的确切消息。}\par

% structure: p0022 source=h2 target=subsection reason=relative-heading-rank; base=h2

\subsection{第十一章。义人的苦难}

但伯多禄听了这话，因同情而流泪，立即对在场的诸位说：\Quote{若是一位恭敬天主的人遭受了这样的苦难，如同这人的父亲所遭受的，他必会立即将原因归于他恭敬天主，而归咎于那恶者。外邦人遭受苦难也是如此；我们这些恭敬天主的人却不知道。但我很有理由称他们为可怜，因为他们在今生被缠住，却得不到你们所怀的盼望。因为那些在恭敬天主中遭受苦难的人，是为赎免自己的过犯而受苦。}\par

% structure: p0024 source=h2 target=subsection reason=relative-heading-rank; base=h2

\subsection{第十二章 一次愉快的旅行}

伯多禄这样说了之后，我们中有一人冒昧地以大家的名义邀请他第二天一早乘船到对面的\Place{阿剌多斯岛}（相距恐怕不到三十斯塔狄亚），去看那里两根非常粗大的葡萄木柱子。于是仁慈的伯多禄同意了，说道：\Quote{你们下船后，不要许多人一起去看你们想看的东西；因为我不希望居民注意你们。}于是我们出发航行，不久便到达了那岛。下船后，我们走到有葡萄木柱子的地方，并且连同它们一起观看了\Person{菲狄亚斯}{Phidias}的几件作品。\par

% structure: p0026 source=h2 target=subsection reason=relative-heading-rank; base=h2

\subsection{第十三章 一个忧伤心灵的妇人}

但伯多禄独自一人认为不值得观看那里的景物；他注意到一个妇人坐在门外，不停地乞讨以维持生计，便对她说：\Quote{妇人，你肢体是否有残缺，以致甘心受此羞辱——我指乞讨——，而不愿用天主赐给你的双手劳动，以获取每日的食粮？}她叹息着回答：\Quote{但愿我有手能劳动！但现在它们只剩下手的形状，因我啃咬而变得麻木无用。}伯多禄接着问：\Quote{你为何遭受如此可怕的痛苦？}她回答：\Quote{软弱的心灵，此外无他。若我有男子的心志，早有悬崖或水池，我可跳下，从折磨我的不幸中得安息。}\par

% structure: p0028 source=h2 target=subsection reason=relative-heading-rank; base=h2

\subsection{第十四章 基肋阿得有乳香}

伯多禄说道：\Quote{那么如何？妇人，你以为自尽者能免于惩罚吗？那些如此死去之人的灵魂，在\OriginalTerm{阴府}{Hades}中不是因自杀而受更重的刑罚吗？}她说：\Quote{但愿我能相信灵魂在阴府中确实活着；那么我会喜爱死亡，轻视刑罚，只为能见我所渴想的儿子们一小时！}伯多禄说：\Quote{是什么事使你忧伤？妇人，我愿知道。若你告诉我，以此恩惠为报，我必使你确信灵魂在阴府中活着；并且，代替悬崖或水池，我将给你一种药，使你活着或死去都无痛苦。}\par

% structure: p0030 source=h2 target=subsection reason=relative-heading-rank; base=h2

\subsection{第十五章 妇人的故事}

于是那妇人，不理解那模棱两可的话，因这承诺而喜悦，开始这样说道：——\Quote{若我谈及我的家族和祖国，我想无人会相信。但你得知这些，除了我在痛苦中因啃咬而毁了我的双手的原因之外，又有什么要紧呢？但我还是要向你述说我的事，尽你所能听的。我出身极为高贵，因某权贵之安排，嫁给了他的一个亲戚。我先产下一对孪生子，后又生一子。但我丈夫的兄弟完全疯了，贪恋我这极为贞洁的不幸女子。我既不愿顺从我的情人，也不愿向我丈夫揭露他兄弟对我的爱，我如此思量：我不可因犯奸淫而玷污自己，不可羞辱我丈夫的床榻，不可使兄弟相争，更不可使整个家族（一个大家族）蒙受众人的羞辱。我认为最好暂时带着我的孪生孩子离开那座城，直到那不洁的爱情（他无耻地谄媚我）止息。至于另一个儿子，我留在他父亲身边，作为对他的一点安慰。}\par

% structure: p0032 source=h2 target=subsection reason=relative-heading-rank; base=h2

\subsection{第十六章 海难}

\Quote{然而，为使事情如此安排，我决定编造一个梦：似乎有人在夜间站在我身旁，如此说：\InnerQuote{妇人，立即带着你的孪生孩子离开这城一段时候，直到我吩咐你返回此地；否则你和你的丈夫及所有儿女将立即惨死。}我于是这样做了。因我一将这假梦告诉丈夫，他惊慌起来，便用船将我和我的两个儿子以及奴仆、婢女、大量金钱送往雅典，要我教育孩子，直到，他说，那位赐神谕者喜悦我回到他那里。但可怜的我，在航行中带着我的孩子，被狂风吹到此地，船在夜间碎裂，我遇了难。其余的人都死了，只有我这个不幸的人被大浪抛到岩石上。我坐在岩石上时，因希望找到孩子还活着，而阻止了自己当时跳入深水（我本可轻易做到），因为我的灵魂已被海浪灌醉。}\par

% structure: p0034 source=h2 target=subsection reason=relative-heading-rank; base=h2

\subsection{第十七章 徒劳的寻找}

\Quote{但天破晓时，我大声呼喊，悲惨地哀号，环顾四周，寻找我不幸孩子们的尸体。因此居民们怜悯我，见我赤身露体，先给我穿上衣服，然后潜入深处，寻找我的孩子们。当他们什么也没找到时，几位好客的妇女来安慰我，每个人都讲述自己的不幸，希望我从类似的不幸中得到安慰。但这反而让我更悲痛，因为我说我还没有坏到能从别人的不幸中得到安慰。于是，当许多人邀请我接受她们的款待时，一位贫穷的妇女极力强求我进入她的小屋，对我说：\InnerQuote{妇人，鼓起勇气，因为我的丈夫，他是一名水手，也死在海上，当时他正值青春年华；自那以后，虽然许多人向我求婚，但我宁愿守寡，为失去丈夫而悲伤。我们将共同拥有我们双手所能赚取的一切。}}\par

% structure: p0036 source=h2 target=subsection reason=relative-heading-rank; base=h2

\subsection{第十八章. 祸不单行}

\Quote{为了不冗述不必要的细节，我因为她的爱夫之情而与她同住。不久之后，我的双手因啃咬而变得无力；收留我的那位妇女被某种疾病缠身，困在屋里。自此，妇女们先前的同情减退了，我和那屋里的妇女都无助。我在这里坐了很久，如你所见，乞讨；我所得到的一切都带给我的同伴，供我们维持生计。关于我的事就到此为止。至于其余的，有什么妨碍你履行承诺，把那毒药给我，好让我也给她——那个想死的人；这样，如你所说，我也能逃离生命？}\par

% structure: p0038 source=h2 target=subsection reason=relative-heading-rank; base=h2

\subsection{第十九章. 推诿}

当那妇人这样说时，\Person{伯多禄}似乎因许多思虑而犹豫不决。但我走上前说：\Quote{我找了你好久。现在怎么样了？}但\Person{伯多禄}命令我带路，并到船上等他；因他命令时不可违抗，我便照做了。但\Person{伯多禄}后来将整件事告诉我，他心中稍起疑心，便询问那妇人说：\Quote{请告诉我，妇人，你的家族、你的城邑、以及你孩子们的名字，然后我就给你那毒药。}但她受到逼迫，不愿说话，却又急切想得到毒药，便狡猾地以假乱真。她说自己是\Place{厄弗所}人，丈夫是\Place{西西里}人；同样地，她改变了三个孩子的名字。于是\Person{伯多禄}以为她说的是真话，便说：\Quote{唉！妇人，我原以为今天会给你带来极大的喜乐，我猜你正是我想起的那个人，我听说过她的事并准确知道。}但她恳求他说：\Quote{请告诉我，我求你，好让我知道妇女中是否有比我更悲惨的人。}\par

% structure: p0040 source=h2 target=subsection reason=relative-heading-rank; base=h2

\subsection{第二十章. \Person{伯多禄}的叙述}

于是\Person{伯多禄}，不知道她说了假话，出于怜悯开始告诉她真相：\Quote{有一个跟随我的青年，渴慕宗教言论，是罗马公民，他告诉我，他有一个父亲和两个孪生兄弟，他都失去了他们的踪迹。因为，}他说，\Quote{我的母亲，正如我父亲告诉我的，因见到异象，带着双胞胎子女暂时离开罗马城，以免遭遇厄运；她与他们一同离去后便杳无音讯；而她的丈夫，那青年的父亲，去寻她，也失去了踪迹。}\par

% structure: p0042 source=h2 target=subsection reason=relative-heading-rank; base=h2

\subsection{第二十一章. 真相大白}

当\Person{伯多禄}这样说时，那妇人一直留心听着，突然昏厥如痴呆。但\Person{伯多禄}走近她，扶住她，劝她克制自己，说服她承认自己的状况。但她浑身无力，如醉酒一般，转过头来，\Emph{头}能够承受那所期望的喜乐，搓着脸说：\Quote{那少年在哪里？}他此时已看透整个事情，说：\Quote{你先告诉我；否则你不能见他。}她急切地说：\Quote{我就是那少年的母亲。}于是\Person{伯多禄}说：\Quote{他叫什么名字？}她说：\Quote{克肋孟。}于是\Person{伯多禄}说：\Quote{正是他，他刚才与我说话，我命他在船上等我。}她俯伏在\Person{伯多禄}脚前，恳求他赶快到船上去。\Person{伯多禄}说：\Quote{你若与我约定，我就照办。}她说：\Quote{我什么都愿意做；只求你让我见到我唯一的孩子。因为我在他身上似乎看到了我死在这里的两个孩子。}\Person{伯多禄}说：\Quote{你们见到他时，要保持安静，直到我们离开这岛。}她说：\Quote{我会的。}\par

% structure: p0044 source=h2 target=subsection reason=relative-heading-rank; base=h2

\subsection{第二十二章. 失而复得}

因此，\Person{伯多禄}拉着她的手，把她带到船上。但我看见他拉着那妇人的手，笑了，走近，自愿代替他扶她，以示尊敬。但我一触到她的手，她便发出母性的呼喊，紧紧地拥抱我，急切地吻我，像吻自己的儿子。但我对整件事一无所知，把她当作疯女人推开。然而，出于对\Person{伯多禄}的尊重，我克制了自己。\par

% structure: p0046 source=h2 target=subsection reason=relative-heading-rank; base=h2

\subsection{第二十三章. 好客的报偿}

但伯多禄说：\Quote{哎呀！我的儿子克莱孟，你在做什么？竟要推开你的生身母亲？} 我一听这话，便哭了起来，倒在已昏倒的母亲旁边，亲吻了她。因为一经告诉我，我隐约记起了她的容貌。这时，大批群众跑来观看这女乞丐，互相传告说她的儿子认出了她，而且他是个有地位的人。正当我们想立刻带母亲离开那岛时，她对我说：\Quote{我深深渴望的儿子，我应当向那款待我的妇人告别，她穷困且身体衰弱，正躺在屋里。} 伯多禄听到这话，在场的群众也都敬佩那妇人的好心。伯多禄立刻吩咐几个人去把她连睡榻一起抬来。榻一抬到放下，伯多禄当着全体群众说：\Quote{如果我是真理的使者，为使旁观的众人相信，叫他们知道有一位创造世界的天主，就让这妇人立刻痊愈站起来吧。} 伯多禄还在说话时，那妇人已痊愈站起，跪在伯多禄面前，亲吻她亲爱的同伴，并询问这一切是怎么回事。于是她简短地向她讲述了相认的整个经过，使听者惊奇。随后，我的母亲看见她的女主人被治好，也请求自己能获得痊愈。伯多禄按手在她身上，也治好了她。\par

% structure: p0048 source=h2 target=subsection reason=relative-heading-rank; base=h2

\subsection{第二十四章　一切安排妥当}

接着，伯多禄论及天主和应事奉他的道理，最后这样总结：\Quote{若有人愿意准确学习这些事，就让他到安提约基雅来，我已决定在那里住一段时间，学习有关他得救的事。因为你们若惯于为经商或作战而离乡背井，来到远地，就不该为永生的缘故，不愿走三天的路程。} 伯多禄讲完话后，我当着全体群众的面，送给那被治好的妇人一千银币作为生活费用，并把她托付给城中一位品行良好的首领，他欣然主动承担起照料的责任。此外，我又分钱给许多别的妇人，并感谢那些曾安慰过我母亲的人，然后带着母亲、伯多禄和其余的同伴，乘船前往安塔拉杜斯；我们就这样继续前往住处。\par

% structure: p0050 source=h2 target=subsection reason=relative-heading-rank; base=h2

\subsection{第二十五章　博爱与友谊}

到达后，我们吃了饭，并按惯例谢了恩，时间还有余，我便对伯多禄说：\Quote{我的主伯多禄，我母亲做了件博爱的事，纪念了她的女主人。} 伯多禄回答说：\Quote{克莱孟，你果真认为你母亲对待那位在她遭海难后收留她的妇人，是行了博爱之事，还是说这话是为了大大赞扬你的母亲？如果你说的是真话，而不是出于恭维，那么在我看来，你似乎还不知道博爱的伟大何在。博爱就是：无论对谁，只因他是人，便对他怀有仁爱，而不受肉体说服的影响。但我甚至不敢称那位在你母亲遇难后收留她的女主人为博爱者；因为她是被怜悯所推动，被说服而成为一位遭海难、为子女忧伤、陌生、赤身露体、穷困潦倒、极其哀叹厄运的女子的恩人。因此，她处于这种情形，谁见了，即使是不敬神的人，岂能不怜悯她？所以，在我看来，连那位接待陌生人的妇人也没有做博爱的事，而是因怜悯她无数的苦难才出手相助。更何况你的母亲，当她境遇好转、回报女主人时，她所做的不是博爱的事，而是友谊的事！因为友谊与博爱大有分别，友谊出于回报。而博爱，除却肉体说服，爱每一个人，只因他是人，并施恩于他。因此，如果她怜悯了女主人，同时也怜悯并善待那些曾经错待她的仇人，那她才是博爱的；但如果她因某一原因而友好或敌对，因另一原因而敌对或友好，那么这样的人是某种品质的朋友或敌人，而不是人本身的朋友或敌人。}\par

% structure: p0052 source=h2 target=subsection reason=relative-heading-rank; base=h2

\subsection{第二十六章　何谓博爱}

于是我回答说：\Quote{那么，你难道不认为那位接待陌生人的妇人是博爱的吗？她善待了一个她不认识的陌生人。} 伯多禄说：\Quote{我固然可以称她为怜悯的，但不敢称她为博爱的，正如我不能称一位母亲为慈爱的（\OriginalTerm{慈爱}{philoteknic}），因为她对子女的爱是出于生产的痛苦和养育之情。同样，情人也因与情妇的交往和愉悦而满足，朋友因友谊的回报，怜悯的人也因不幸而感动。然而，怜悯的人接近博爱的人，因为他并不追求收受任何东西，而施予善意。但他还不是博爱的人。} 我说：\Quote{那么，怎样做才算是博爱的人呢？} 伯多禄回答说：\Quote{我看你急切想听什么是博爱的工作，我不妨告诉你。博爱的人就是连仇人也善待的人。你且听：博爱是阴阳兼有的；其阴性部分称为\Emph{怜悯}，阳性部分称为\Emph{爱近人}。但每个人都彼此为近人，并不只是这个人或那个人；因为善与恶、友与敌都是人。因此，实践博爱的人，应当效法天主，施恩于义人和不义之人，如同天主自己在这现世赐阳光和天空给众人一样。但你若只善待善人，而恶待恶人，甚至惩罚他们，那你就是从事审判官的工作，而不是努力持守博爱。}\par

% structure: p0054 source=h2 target=subsection reason=relative-heading-rank; base=h2

\subsection{第二十七章 谁能审判}

于是我说：\Quote{那么，连天主，如你所教导的，将要审判的那一位，也不是慈爱的。}伯多禄说：\Quote{你提出一个矛盾；因为祂将审判，正是因此祂是慈爱的。因为那爱那些受委屈的人并同情他们的人，会为受委屈者报复那些委屈他们的人。}于是我说：\Quote{那么，如果我也善待善人，并惩罚那些伤害他人的作恶者，我难道不是慈爱的吗？}伯多禄答道：\Quote{如果你同时拥有知识和审判的权柄，你会做得正确，因为你已领受了权柄去审判天主所造的人，并且你的知识无误地将一些人定为义人，将一些人定为不义者。}于是我说：\Quote{你说得正确又真实；因为任何人若没有知识，就不可能正确审判。因为有时有些人看似良善，却暗中行恶；而有些善人因仇敌的控告而被视为恶人。但即使一个人审判，拥有拷问和审查的权力，也未必能完全公正地审判。因为有些凶手经受住了酷刑，脱身而被视为无罪；而另一些无罪的人未能承受酷刑，却对自己作假供，并作为有罪者受到惩罚。}\par

% structure: p0056 source=h2 target=subsection reason=relative-heading-rank; base=h2

\subsection{第二十八章 审判的难度}

于是伯多禄说：\Quote{这些是平常的事；现在请听更重大的。有些人的罪或善行一部分是他们自己的，一部分是别人的；但每个人都应该为自己的罪受罚，为自己的功德得赏。然而，除了那唯一拥有全知的先知，没有人能知道一个人所做的哪些是出于他自己，哪些不是；因为一切都被视为是他做的。}于是我说：\Quote{我想知道，人的恶行或善行，哪些是自己的，哪些是别人的。}\par

% structure: p0058 source=h2 target=subsection reason=relative-heading-rank; base=h2

\subsection{第二十九章 善人的苦难}

于是伯多禄答道：\Quote{真理的先知曾说过：\InnerQuote{善事必然要来，但他说，藉着谁而来，那人是有福的；同样，恶事也必然要来，但藉着谁而来，那人是有祸的。}但如果恶事藉着恶人而来，善事由善人带来，那么每个人必须是出于自己而成为善或恶，并且出于他所预定的，以便后来善或恶的到来；这些出于他自己的选择，并非是出于天主的眷顾而安排从他而来。既然如此，这就是天主的审判：那如同通过战斗，经历一切灾祸而找到无可指责的人，被堪当承受永生；因为那些凭自己意志持守善道的人，被那些凭自己意志持守恶道的人所试探，受迫害、被仇恨、被毁谤、被图谋、被击打、被欺骗、被控告、被折磨、被羞辱——忍受所有这些事情，由此看来，他们理当被激怒并趋向报复。}\par

% structure: p0060 source=h2 target=subsection reason=relative-heading-rank; base=h2

\subsection{第三十章 绊倒事必来}

\Quote{但主知道那些不义地做这些事的人因他们先前的罪而有罪，并且邪恶的灵藉着有罪的人做这些事，祂已劝告我们要怜悯人，因为他们是人，并且是藉着罪成为邪恶的工具；祂将这\Emph{劝告}赐给祂的门徒，作为要求仁爱的表现，并且尽我们所能，免除作恶者的定罪，以便节制的人如同帮助醉酒的人一样，藉着祈祷、斋戒和祝福来帮助他们，不抵抗，不报复，免得迫使他们犯更多的罪。因为当一个人被判定要承受苦难时，他没有理由对带来苦难的人发怒；因为他应当想：即使他没有虐待我，但因为我要受苦，必定也要藉着另一个人受苦。那么，既然我被判定无论如何都要受苦，我为何要向施行者发怒呢？然而，更进一步：如果我们这些善人为报复而对恶人做同样的事，那么我们做的就是恶人所做的事，只不过他们先做，我们后做；而且，如我所说，我们不应发怒，因为知道在天主的眷顾中，恶人惩罚善人。因此，那些对他们的惩罚者怀恨的人，是犯罪，因为他们藐视天主的使者；而那些尊敬他们、并抵挡那些想伤害他们的人，是对这样安排的天主敬虔。}\par

% structure: p0062 source=h2 target=subsection reason=relative-heading-rank; base=h2

\subsection{第三十一章 \Quote{然而，他们并非此意。}}

我回答说：\Quote{那么，那些作恶的人是无罪的，因为他们藉着天主的审判来委屈义人。}于是伯多禄说：\Quote{他们确实犯了大罪，因为他们已将自己交付给罪。因此，知道这事，\Emph{天主}从他们中拣选了一些人来惩罚那些正义地悔改了先前之罪的人，以便义人先前所行的恶事可以藉着这惩罚得以赦免。但对于那些惩罚人、又想虐待他们、且不肯悔改的恶人，允许他们虐待义人，以填满他们自己的惩罚。因为没有天主的旨意，连一只麻雀也不能掉进网里。同样，义人的头发也被天主数过了。}\par

% structure: p0064 source=h2 target=subsection reason=relative-heading-rank; base=h2

\subsection{第三十二章 金科玉律}

\Quote{义人为了合理之事与本性相争。例如，爱那些爱自己的人，是所有人的本性。但义人还努力爱仇敌，祝福诽谤他的人，甚至为仇敌祈祷，同情那些对他行不义的人。因此，他也避免行不义，祝福咒骂他的人，原谅打击他的人，服从迫害他的人，向不向他致敬的人致敬，与那些没有的人分享他所拥有的，劝导对他发怒的人，与他的仇敌和解，劝诫悖逆的人，教导不信的人，安慰哀悼的人；困苦时，他忍耐；被忘恩负义地对待时，他不发怒。但他全心爱邻人如己，不怕贫穷，反而因与一无所有者分享自己的财物而成为贫穷。他也不惩罚罪人。因为爱邻人如己的人，知道自己犯罪时不愿受罚，因此也不惩罚犯罪的人。他愿意被赞美、祝福、尊敬，并所有罪过得到赦免，他也如此对待邻人，爱邻人如己。总之，他为自己所愿望的，也为邻人所愿望。因为这是天主和先知的法律，这是真理的教义。而对每个人的这种完美的爱，就是\OriginalTerm{爱人之心}{philanthropy}的男性部分；其女性部分是同情，即喂养饥饿者、给口渴者水喝、给赤身者衣服穿、探访病人、接待旅客、探望并尽力帮助被囚者，总之，同情不幸者。}\par

% structure: p0066 source=h2 target=subsection reason=relative-heading-rank; base=h2

\subsection{第三十三章 畏惧与爱}

但我听了这话，说：\Quote{这些事确实不可能做到；但对仇敌行善，忍受他们的一切无礼，我认为不可能是人的本性。}然后\Person{伯多禄}回答说：\Quote{你说得对；因为\OriginalTerm{爱人之心}{philanthropy}，作为不朽的原因，是付出巨大代价才赐下的。}然后我说，\Quote{那么，如何可能在人心中得到它呢？}\Person{伯多禄}回答说：\Quote{亲爱的\Person{克莱孟}，得到它的方法是这样的：若有人被说服，认为仇敌暂时虐待他们所恨的人，却成为那些人从永罚中得救的原因，他便会立即热切地爱他们如恩人。但得到它的方法，亲爱的\Person{克莱孟}，只有一个，就是敬畏天主。因为敬畏天主的人起初不能爱邻人如己；因为这样的秩序不会进入灵魂。但因着敬畏天主，他能够做那些有爱之人所做的事；因此，当他行爱德的事时，新娘\Quote{爱}就被带到新郎\Quote{畏惧}面前。于是这位新娘生出爱人之心的思想，使拥有者成为不朽，如同天主的准确肖像，其本质不会腐朽。}他这样给我们讲解爱人之心的教义时，夜晚降临，我们就去休息了。\par

\SourceRecord{Translated by Peter Peterson.；Ante-Nicene Fathers；Vol. 8.；Edited by Alexander Roberts, James Donaldson, and A. Cleveland Coxe.；Buffalo, NY: Christian Literature Publishing Co.,；1886.；<http://www.newadvent.org/fathers/080812.htm>.}

% structure: p0001 source=h1 target=section reason=page-title

\section{讲道词 第十三篇}

% structure: p0002 source=h2 target=subsection reason=relative-heading-rank; base=h2

\subsection{第一章  往劳狄刻雅之行}

天刚破晓，\Person{伯多禄}进来，说道：\Quote{\Person{克莱孟}、他的母亲\Person{玛蒂狄雅}和我的妻子，必须立刻坐上马车。}他们便立刻照办了。当我们匆忙赶往\Place{巴兰纳艾}的路上，母亲问起父亲的景况；我说：\Quote{父亲外出寻你和我的双胞胎兄弟\Person{福斯提诺}和\Person{福斯提尼亚诺}，如今下落不明。但我猜想他必已去世多年，或因船难丧生，或因迷路失踪，或因悲伤憔悴。}她听了这话，泪如雨下，因悲伤而叹息；但因找到我而生的喜悦，多少减轻了她回忆的痛苦。于是我们一同下行至\Place{巴兰纳艾}。次日，我们到了\Place{帕尔托斯}，又从那里到\Place{加巴拉}；第三日，我们抵达\Place{劳狄刻雅}。看！城门前\Person{尼凯塔}和\Person{阿桂拉}迎接我们，拥抱我们，带我们到住所。\Person{伯多禄}见这城市美丽而宏伟，便说：\Quote{值得在此停留数日；因为一般而言，人口众多之处最能为我们提供所求的人。}\Person{尼凯塔}和\Person{阿桂拉}问我那位陌生的妇人是谁；我说：\Quote{我的母亲，是天主借着我主\Person{伯多禄}赐我认出的。}\par

% structure: p0004 source=h2 target=subsection reason=relative-heading-rank; base=h2

\subsection{第二章  \Person{伯多禄}向\Person{尼凯塔}和\Person{阿桂拉}述说\Person{克莱孟}及其家史}

我说了这话，\Person{伯多禄}便向他们简述了全部经过——他们先行出发后，我\Person{克莱孟}向他讲述了我的家世，母亲以梦的假托带着双胞胎孩子出行；此外，父亲外出寻她；然后，\Person{伯多禄}本人听了这事后，前往岛上，遇见那妇人，见她乞讨，便问她为何如此；接着确认了她的身份、她的生活方式、那伪造的梦，以及她孩子们的名字——即我（留在我父亲身边）的名字，以及与她同行的双胞胎孩子的名字，她以为他们已葬身深海。\par

% structure: p0006 source=h2 target=subsection reason=relative-heading-rank; base=h2

\subsection{第三章  \Person{尼凯塔}和\Person{阿桂拉}被认出}

\Person{伯多禄}作完这概要叙述后，\Person{尼凯塔}和\Person{阿桂拉}惊讶地说：\Quote{这果真是真的吗，宇宙的统治者和主？还是做梦？}\Person{伯多禄}说：\Quote{除非我们睡着了，这当然是真的。}他们便沉思片刻，然后说道：\Quote{我们是\Person{福斯提诺}和\Person{福斯提尼亚诺}。从你谈话开始，我们便彼此注视，心中诸多猜测，不知所说是否与我们有关；因为我们想到生活中多有巧合。所以我们沉默不语，心中怦怦直跳。但当你讲完叙述时，我们清楚看出你的话是指我们，于是我们便承认了自己的身份。}说着，他们泪流满面，冲进去见母亲；虽见她正睡着，却仍想拥抱她。但\Person{伯多禄}阻止他们说：\Quote{让我先带你们，把你们引见给你们的母亲，免得她因过度突然的喜悦而失去理智，因她正在睡梦中，心神被睡眠束缚。}\par

% structure: p0008 source=h2 target=subsection reason=relative-heading-rank; base=h2

\subsection{第四章  母亲不可与儿子一同进食；其理由}

母亲睡足醒来后，\Person{伯多禄}立刻开始先与她谈论\Emph{真正的}虔诚，说道：\Quote{妇人，我希望你明白我们宗教的生活方式。我们敬拜独一的天主，祂创造了你所见的这个世界；我们遵守祂的法律，其主要诫命是单单敬拜祂，尊崇祂的圣名，孝敬父母，保持贞洁，虔诚生活。此外，我们不与所有人无分别地共处；也不与外邦人同席吃饭，因为我们不能与他们一同进食，因为他们生活不洁。但当我们说服他们拥有真思想、遵循正确行为，并以三次可称颂的呼号为他们施洗后，我们才与他们同住。因为即使是我们的父亲、母亲、妻子、孩子、兄弟或任何其他有自然亲情的人，我们也不敢与他一同吃饭；因为我们的宗教迫使我们做出区分。因此，不要因为你的儿子不与你一同吃饭而视为侮辱，直到你也拥有与他相同的见解、遵循与他相同的行为方式为止。}\par

% structure: p0010 source=h2 target=subsection reason=relative-heading-rank; base=h2

\subsection{第五章  \Person{玛蒂狄雅}愿受洗}

她听了这话，便说：\Quote{那么，有什么阻止我今天受洗呢？因为在我见到你之前，我已因这想法而远离所谓的神祇：尽管我几乎每日向他们献上许多祭品，他们却未在我困苦时帮助我。至于奸淫，何须多言？因我富有时，并未因奢侈而陷入此罪；后来的贫穷也未能迫使我犯罪，因为我坚守贞洁，视之为最大的美丽，为此我陷入如此大的苦难。但我绝不以为我主\Person{伯多禄}你不知，最大的诱惑正是在一切顺遂时来临。因此，若我在顺境中保持贞洁，也不在沮丧时放纵享乐。是的，你切莫以为我的灵魂如今已摆脱痛苦，尽管因认出\Person{克莱孟}而得到些许安慰。因我失去两个孩子的悲伤仍涌上心头，多少遮蔽了我的喜乐；我悲伤，并非因他们葬身大海，而是因他们灵魂与肉身一起毁灭，未拥有对天主的真正虔诚。此外，我从\Person{克莱孟}得知，他们的父亲、我的丈夫外出寻我和儿子们，这么多年音讯全无；他必已去世。因那可怜人，以贞洁爱着我，又疼爱子女；因此这老人，失去我们所有至亲，心碎而死。}\par

% structure: p0012 source=h2 target=subsection reason=relative-heading-rank; base=h2

\subsection{第六章　儿子们向母亲表明自己的身份}

母亲的这番话，儿子们听了，再也无法克制自己，便依\Person{伯多禄}的劝勉，起身拥抱她，泪如雨下，连连亲吻她。母亲却说：\Quote{这是什么意思？}\Person{伯多禄}回答说：\Quote{妇人，鼓起勇气，振奋精神，好享受你的孩子们；因为这就是你的儿子\Person{法乌斯蒂诺}与\Person{法乌斯蒂尼亚诺}，你以为他们已葬身海底了。至于他们如何在你认为已死于那个最悲惨的夜晚之后仍然活着，以及为什么现在一个名叫\Person{尼塞塔}，另一个名叫\Person{阿奎拉}，他们自己会告诉你；因为我们和你一样，还不知道这事。}\Person{伯多禄}这样说时，我母亲因过度欢喜而昏厥，几乎死去。等我们使她苏醒过来，她坐起身，定了定神，说：\Quote{我亲爱的孩子们，请行个好，告诉我们那悲惨的夜晚之后发生了什么事。}\par

% structure: p0014 source=h2 target=subsection reason=relative-heading-rank; base=h2

\subsection{第七章　尼塞塔讲述自己的遭遇}

于是\Person{尼塞塔}——今后要改称为\Person{法乌斯蒂诺}——开始说道：\Quote{就在那个你知道船破碎的夜里，我们被一些人捞起，他们不惧怕在海上干那强盗的营生。他们将我们放在一条小船上，沿着海岸行进，时而划船，时而派人去取食物，最终将我们带到\Place{凯撒勒雅·斯特拉托尼}。在那里，他们用饥饿、恐惧和殴打折磨我们，免得我们轻率地泄露任何他们不愿让人知道的事；此外，他们还换了我们的名字，并成功地将我们卖掉。买我们的是一个犹太人的归依者，一个非常可敬的女人，名叫\Person{尤斯塔}。她收养我们如同己出，热心地用希腊人的一切学问教养我们。但我们随着年龄增长，变得谨慎，对她的信仰深怀依恋，并认真修习文化，以便与其他民族辩论时，能说服他们认清错误。我们还精确研究了哲学家的学说，尤其是最无神论的——我指的是\Person{伊壁鸠鲁}和\Person{皮浪}的学说——以便更好地驳斥他们。}\par

% structure: p0016 source=h2 target=subsection reason=relative-heading-rank; base=h2

\subsection{第八章　尼塞塔险些被西满术士欺骗}

\Quote{我们与一个术士\Person{西满}一同被抚养长大；因与他友好交往，我们险些被引入歧途。现在有关于某人的传闻，说他显现时，众多虔诚的人要在他的王国里生活，免受死亡与痛苦。但这事，母亲，将在适当的时候更详尽地解释。然而，当我们即将被\Person{西满}引入歧途时，我们主人\Person{伯多禄}的一位朋友，名叫\Person{匝凯}，来警告我们不要被那术士欺骗；\Person{伯多禄}到来后，他便将我们带到\Person{伯多禄}面前，让他给我们充分的信息，并使我们确信那些与虔诚有关的事。因此，我们恳求您，母亲，也分享那已赐予我们的福分，好让我们能围坐同一张餐桌！母亲，这就是您以为我们死了的原因。在那个灾难性的夜晚，我们被海盗从海上捞起，但您以为我们已丧命。}\par

% structure: p0018 source=h2 target=subsection reason=relative-heading-rank; base=h2

\subsection{第九章　母亲为自己和女主人请求圣洗}

\Person{法乌斯蒂诺}说完后，我们的母亲俯伏在\Person{伯多禄}脚下，恳求他差人将她自己和她的女主人接来，立即为她们付洗，好使，她说，\Quote{在我找回孩子们之后，一天也不要错过与他们一同进食。}我们与母亲一同提出同样的请求时，\Person{伯多禄}说：\Quote{你们怎么想？难道我一人如此无情，竟不愿让你们今天与母亲一同进餐，为她付洗吗？但她在受洗前必须先禁食一天。而且只是一天，因为她单纯地说了些为自己辩护的话，在我看来这足以表明她的信德；否则，她的净化必须持续许多天。}\par

% structure: p0020 source=h2 target=subsection reason=relative-heading-rank; base=h2

\subsection{第十章　玛蒂迪亚正确看重圣洗}

于是我说：\Quote{请告诉我们，她说了什么使她的信德显明出来。}\Person{伯多禄}说：\Quote{她请求她的女主人和恩人一同受洗。因为她若没有首先亲身感到圣洗是一份伟大的恩赐，就不会为她所爱的人祈求这一恩赐。正因如此，我谴责许多人，他们受洗后，声称自己有信德，却做不出与信德相称的事；也不督促他们所爱的人——我指的是他们的妻子、儿女或朋友——去领受圣洗。因为如果他们相信天主在接受圣洗后以善工赐予永生，他们必会立即督促他们所爱的人去受洗。但你们中有人会说：\InnerQuote{他们的确爱他们，关心他们。}那是胡扯。难道他们不是最肯定地，当他们看见所爱的人患病、或被引向死亡的道路、或遭受任何其他考验时，就为他们悲伤怜悯吗？所以，如果他们相信永恒的火等待着那不敬拜天主的人，他们就不会停止劝诫他们，或为他们的不信而深深忧伤，若看见他们不顺从，深知惩罚正等待着他们。但现在我要派人去请那女主人，问她是否愿意接受我们传布的律法；然后根据她的心态，我们将做当做的事。}\par

% structure: p0022 source=h2 target=subsection reason=relative-heading-rank; base=h2

\subsection{第十一章　玛蒂迪亚无意中已禁食一天}

\Quote{但是，既然你的母亲对圣洗的功效有真实信心，让她至少在领洗前一天斋戒。}但她发誓说：\Quote{在过去两天里，当我向那位妇人讲述所有与相认有关的事件时，由于过度喜乐，我无法进食；只在昨天喝了一点水。}\Person{伯多禄}的妻子也发誓为她作证，说：\Quote{她确实没有尝过任何东西。}\Person{阿桂拉}——他以后应被称为\Person{福斯蒂尼亚努斯}——说：\Quote{因此，没有任何事阻止她受洗。}\Person{伯多禄}微笑着回答：\Quote{那不是因圣洗本身而守的斋戒。}\Person{福斯蒂努斯}回答说：\Quote{或许天主不愿在相认之后将我们的母亲与我们分开一整天，已经预先安排了这次的斋戒。因为她在无知之时保持贞洁，做了真宗教所教导的事，所以如今或许天主也安排她为了真圣洗而提前斋戒一日，好使她自与我们相认的第一天起，就能与我们一同用餐。}\par

% structure: p0024 source=h2 target=subsection reason=relative-heading-rank; base=h2

\subsection{第十二章 难题的解决}

\Person{伯多禄}说：\Quote{不要让邪恶借着上智安排和你们对母亲的孝爱而找到借口来支配我们；相反，你们今天继续守斋，我也与你们一同守，明天她将受洗。此外，今天的时辰也不适合施行圣洗。}于是我们都同意如此。\par

% structure: p0026 source=h2 target=subsection reason=relative-heading-rank; base=h2

\subsection{第十三章 伯多禄论贞洁}

当天晚上，我们都受益于\Person{伯多禄}的训导。他以我们母亲所遭遇的事为引子，向我们指明贞洁的结果如何美好，而奸淫的结果如何灾难性，并且自然会为整个种族带来毁灭，即使不是迅速的，至少也是缓慢的。\Quote{甚至到了这样的程度，}他说道，\Quote{贞洁的行为中悦天主，以至于祂在此生为这事赐予一些小的恩惠，甚至赐给那些在谬误中的人；因为来世的救恩只赐给那些因信赖祂而领洗、并贞洁正义行事的人。你们自己已经在你们母亲的事上看到，贞洁的结果终究是好的。因为她若犯了奸淫，或许已经被除掉了；但天主怜悯她，因她行为贞洁，将她从威胁她的死亡中解救出来，并将她失去的子女归还给她。}\par

% structure: p0028 source=h2 target=subsection reason=relative-heading-rank; base=h2

\subsection{第十四章 伯多禄的讲论（续）}

\Quote{但有人会说：\InnerQuote{有多少人因贞洁而丧亡！}是的，但这是因为他们没有察觉危险。因为那察觉到自己在爱某人或被人所爱的女子，应立即回避与他的一切交往，如同回避烈火或疯狗一样。你们的母亲正是如此行的，因为她真心视贞洁为祝福；因此她被保全了，并与你们一同获得了对永恒国度的圆满认识。愿意贞洁的女子，应当知道她被邪恶嫉妒，并且因为爱的缘故，许多人窥伺她。那么，如果她藉着对贞洁的坚定持守而保持圣洁，她将战胜所有诱惑，并得救；然而，即使她做了一切对的事，却一旦犯下奸淫之罪，她也必须受罚，正如先知所说的。}\par

% structure: p0030 source=h2 target=subsection reason=relative-heading-rank; base=h2

\subsection{第十五章 伯多禄的讲论（续）}

\Quote{承行天主旨意的贞洁妻子，是祂起初创造的佳美回忆；因为天主是唯一，祂为一男创造了一女。她若不忘记自己的被造，将未来的惩罚放在眼前，不忽视永恒福乐的失落，她就更为贞洁。贞洁女子乐于那些愿意得救的人，对虔诚者来说是虔诚的榜样，因为她乃是善生的模范。愿意贞洁的女子，断绝一切毁谤的机会；但即使她被敌人毁谤，虽然不给他任何借口，她仍是有福的，并由天主为她复仇。贞洁女子渴慕天主、爱慕天主、中悦天主、光荣天主；她对人也不提供毁谤的机会。贞洁女子以她的美名使教会芬芳，以她的虔诚光荣教会。此外，她也是其师长的赞美，并在他们的贞洁上成为助手。}\par

% structure: p0032 source=h2 target=subsection reason=relative-heading-rank; base=h2

\subsection{第十六章 伯多禄的讲论（续）}

\Quote{贞洁女子以天主之子为净配而得装饰。她以圣光为衣。她的美在于灵魂的良序；她散发馨香，即美名。她穿着美丽的衣服，即端庄。她佩戴珍贵的珍珠，即贞洁的话语。她光芒四射，因为她的心智已被灿烂照亮。她注视美丽的镜子，因为她注视天主。她使用美丽的化妆品，即对天主的敬畏，她用这敬畏劝诫自己的灵魂。女子的美丽不因她有金链，而因她已从短暂的私欲中得释放。贞洁女子深为大君王所渴慕；她已被祂追求、守护和爱恋。贞洁女子不提供被人渴慕的机会，除自己的丈夫以外。贞洁女子若被他人渴慕，就感到忧伤。贞洁女子从心底爱自己的丈夫，拥抱、安慰、取悦他，向他如奴仆般服从，在一切事上顺从他，除非会违背天主。因为那顺从天主的女子，无需守卫便可灵魂贞洁、身体纯洁。}\par

% structure: p0034 source=h2 target=subsection reason=relative-heading-rank; base=h2

\subsection{第十七章 伯多禄的讲论（续）}

\Quote{因此，凡使妻子远离对天主的敬畏的丈夫，是愚昧的；因为那不敬畏天主的妻子，并不畏惧自己的丈夫。如果她不敬畏那看见不可见之物的天主，她怎能在她不可见的选择中保持贞洁呢？而且，那不来集会聆听使人贞洁的言语的女子，又怎能贞洁呢？她又怎能得到劝诫呢？如果她未被告知天主即将来临的审判，未曾完全确信永恒的惩罚是那轻微快感的代价，她无人看守，又怎能贞洁呢？因此，另一方面，即使她不愿意，也要迫使她常来聆听那使人贞洁的言语，是的，要劝诱她这样做。}\par

% structure: p0036 source=h2 target=subsection reason=relative-heading-rank; base=h2

\subsection{第十八章　伯多禄的言论（续）}

\Quote{最好的是，如果你愿意牵着她的手前来，好使你自己也成为贞洁的；因为你将会渴望成为贞洁的，好你能体验神圣婚姻的圆满成果；如果你渴望的话，你将不会犹豫成为父亲，爱自己的子女，并被自己的子女所爱。凡希望有贞洁妻子的人，他自己也应当是贞洁的，尽他对妻子应尽的责任，与她一同用餐，与她交往，与她一同去听那使人贞洁的话语，不使她悲伤，不轻易与她争吵，不使自己被她厌恶，尽可能为她提供一切美好的事物；若他没有这些，就用抚爱来弥补不足。贞洁的妻子不期待被抚爱，她承认丈夫是自己的主，在他贫穷时忍受他的贫穷，在他饥饿时与他一同饥饿，在他旅行时与他一同旅行，在他悲伤时安慰他；即使她拥有丰厚的嫁妆，她也如同什么也没有一样顺从他。但如果丈夫娶了贫穷的妻子，就让她把贞洁视为丰厚的嫁妆。贞洁的妻子在饮食上节制，以免如此娇养的身体的疲倦将灵魂拖向非法的欲望。此外，她绝不同青年男子单独相处，且怀疑老年人；她远离无序的欢笑，唯独将自己交托于天主；她不受迷惑；她喜爱聆听圣善的话语，却远离那些不是为了产生贞洁而说的话。}\par

% structure: p0038 source=h2 target=subsection reason=relative-heading-rank; base=h2

\subsection{第十九章　伯多禄的言论结束}

\Quote{天主是我的见证：一次通奸如同许多谋杀一样恶劣；其可怕之处在于，其所带来的恐怖和不敬并不被看见。因为，当鲜血流淌时，尸体躺卧在那里，所有人都被这事件的恐怖所震惊。但通奸所造成的灵魂的谋杀，虽然更为可怕，却因不被人看见，并不会使胆大妄为者的冲动有丝毫减弱。人啊，要知道使你保持生命的是谁的嘘气，你就会不愿它被玷污。唯有通奸玷污了天主的嘘气。因此，它把玷污它的人拖入火中；因为它急忙将侮辱者交付永罚。}\par

% structure: p0040 source=h2 target=subsection reason=relative-heading-rank; base=h2

\subsection{第二十章　伯多禄对玛提狄雅说话}

当\Person{伯多禄}说这话时，他看见善良贞洁的\Person{玛提狄雅}因喜乐而流泪；但以为她因过去所受的许多苦而悲伤，便说：\Quote{鼓起勇气，妇人！因为许多人因通奸受了许多苦，你却因贞洁受苦，因此你没有死。但如果你死了，你的灵魂本会得救。你因贞洁离开了你的故乡\Place{罗马}城，但你藉此找到了真理，即永恒国度的冠冕。你在深渊中经历了危险，却没有死；即使你死了，那深渊本身对因贞洁而死的你，也会成为你灵魂得救的圣洗。你与子女分离了片刻；然而这些，你丈夫的真正后裔，已在更好的境况中被找到。当你饥饿时，你乞求食物，却没有以奸淫玷污你的身体。你使身体遭受折磨，却拯救了你的灵魂；你逃离了奸夫，以免玷污你丈夫的床榻；但因你的贞洁，知道你的逃离的天主，将代替你丈夫的位置。你忧伤孤独，暂时失去了丈夫和子女，但这些迟早都会因死亡——人预定的命运——而失去。然而，你甘愿因贞洁而失去他们，远胜于日后仅仅因罪恶而不情愿地灭亡。}\par

% structure: p0042 source=h2 target=subsection reason=relative-heading-rank; base=h2

\subsection{第二十一章　同题续}

\Quote{那么，你的最初的境况是忧苦的，这远更好。因为当情况如此时，它们不会使你过度悲伤，因为你希望它们会过去，并且它们因对更好境况的期待而带来喜乐。但最重要的是，我希望你知道贞洁多么中悦天主。贞洁的妇人是天主的拣选，天主的喜悦，天主的荣耀，天主的子女。贞洁是如此伟大的恩赐，以至于若没有那条法律说连义人也不可未受洗而进入天主的国，或许连迷途的外邦人也能仅仅因贞洁而得救。因此，我极为那些因躲避圣洗而守贞洁——从而选择没有美好希望的贞洁——的迷途者感到惋惜。因此他们不得救；因为天主的法令清楚写明：未受洗的人不能进入祂的国。}当他说了这些话和更多的话之后，我们就去睡了。\par

\SourceRecord{Translated by James Donaldson.；Ante-Nicene Fathers；Vol. 8.；Edited by Alexander Roberts, James Donaldson, and A. Cleveland Coxe.；Buffalo, NY: Christian Literature Publishing Co.,；1886.；<http://www.newadvent.org/fathers/080813.htm>.}

% structure: p0001 source=h1 target=section reason=page-title

\section{第十四篇讲道}

% structure: p0002 source=h2 target=subsection reason=relative-heading-rank; base=h2

\subsection{第一章 玛提狄雅在海中受洗}

\Emph{很}早以前，\Person{伯多禄}醒了，来到我们这里，唤醒我们，说：\Quote{让\Person{福斯提诺}和\Person{福斯提尼亚诺}，连同\Person{克肋孟}和全家，陪我一起去，使我们能到海边某个隐蔽处，在那里给她付洗而不引人注意。}于是，我们来到海边，他在一些岩石间给她付了洗，那里既无风也无尘。但我们兄弟们，连同我们的兄弟和另外几个人，因妇女们而回避，沐浴后，又回到妇女们那里，带她们同行，这样我们到了一个隐秘处祈祷。然后\Person{伯多禄}因人多，先打发妇女们前行，吩咐她们从另一条路回寓所，只允许我们男子陪伴我们的母亲和其余妇女。我们就去了寓所，等待\Person{伯多禄}到来时，彼此谈论。几小时后\Person{伯多禄}来了，为圣体圣事擘饼，并在上面加盐，先给了我们的母亲，然后给了我们她的儿子们。我们就与她一同进食，赞美天主。\par

% structure: p0004 source=h2 target=subsection reason=relative-heading-rank; base=h2

\subsection{第二章 伯多禄迟到的原因}

于是，终于，\Person{伯多禄}看见人群已经进入，便坐下，叫我们坐在他旁边，首先讲述了为什么他在洗礼后打发我们先行，以及他自己为何迟归。他说原因如下：\Quote{当你们上来的时候，}他说，\Quote{一个老人，一个工人，同你们一起进来，出于好奇而隐藏自己。他之前曾窥探我们，正如他后来自己承认的，为了看我们在进入那隐蔽处时做什么，然后他悄悄出来，跟随着我们。在一个合适的地方走近我，对我说：\InnerQuote{我跟您已经很久，想和您谈谈，但我怕您会生我的气，好像我被好奇心驱使；但现在，如果您愿意，我会告诉您我认为的真理。}我回答说：\InnerQuote{请告诉我们您认为好的，我们会认可您的行为，即使您说的并不真的好，因为您出于善意，急于陈述您认为好的事物。}}\par

% structure: p0006 source=h2 target=subsection reason=relative-heading-rank; base=h2

\subsection{第三章 老人不信有天主或眷顾}

\Quote{老人开始如下说话：\InnerQuote{当我看见你们在海中沐浴后退到隐秘处时，我上前偷偷窥看你们进入隐秘处的目的，看见你们祈祷，我就退下了；但可怜你们，我等着，等你们出来时好和你们说话，并说服你们不要被引入歧途。因为没有天主，也没有眷顾，但万事都服从于命星。对此我深信不疑，因为我自己所经历的事，并且我长久以来仔细研究这门学问。所以，不要受骗，我的孩子。因为无论你祈祷与否，你都必须承受命星所分派给你的。因为如果祈祷能有什么作用或好处，我自己现在就会处于更好的境况了。现在，除非我这破旧的衣服误导你，你不会拒绝相信我的话。我过去曾经富有，向众神献了许多祭，慷慨施舍给穷人。然而，尽管我祈祷并虔诚行事，我仍无法逃脱我的命运。}我说道：\InnerQuote{你遭受了什么灾难？}他回答说：\InnerQuote{我现在不必告诉你；或许到最后你会知道我是谁，我的父母是谁，我落入了何等窘迫的境地。但此刻我希望你完全确信，万事都服从于命星。}}\par

% structure: p0008 source=h2 target=subsection reason=relative-heading-rank; base=h2

\subsection{第四章 伯多禄反驳命星的论证}

\Quote{\InnerQuote{如果万事都服从于命星，而您深信这是事实，那么您的想法和劝告就与您自己的意见相悖。因为如果连思想都无法违反命星，您为何徒劳劝我做不能做的事？不仅如此，即使命星存在，也不要急于说服我不敬拜那位也是星辰之主，因祂的意愿使某事不发生，那事就成为不可能。因为服从者必须始终服从统治者。然而，至于敬拜普通神祇的事，如果命星掌权，那是多余的。因为没有任何事情发生违反命运之意，他们自己也不能做什么，因为他们服从于他们自己的普遍命星。如果命星存在，对此就有这样一个反对意见：即那并非首先的却掌权；或者\Emph{换句话说}，非受造者不能服从，因为非受造者，作为非受造的，没有比它自身更古老的事物。}}\par

% structure: p0010 source=h2 target=subsection reason=relative-heading-rank; base=h2

\subsection{第五章 对命星的实际反驳}

\Quote{当我们这样谈话时，一大群人聚集在我们周围。于是我望着人群说：\Quote{我和我的家族从祖先那里继承了崇拜天主的传统，并且我们有一条诫命，不可关注\InnerQuote{起源}（我是指占星术的科学）；因此我不曾留意它。正因如此，我毫无占星术的技艺，但我要陈述我所擅长的事。既然我无法通过与\InnerQuote{起源}相关的科学来反驳\InnerQuote{起源}，我愿用另一种方式证明：这个世界的事务是由天主的眷顾所管理的，并且每个人都将按其行为得到赏罚。至于这赏罚是现在还是将来发生，对我来说无关紧要；我所断言的是：每个人无疑都会收获其行为的果实。证明\InnerQuote{起源}不存在的方法如下：如果在你们当中有人失明、手残、脚瘸，或身体某部分有缺陷，并且完全无法医治，远远超出医学的范围——连占星家也声称无法治愈的病例，因为久远的岁月中从未发生过这样的治愈——然而，我向天主祈祷，就能治愈它，尽管\InnerQuote{起源}从未能将其矫正。既然这样，那些亵渎创造万物的天主的人，岂不犯了罪吗？}\newline{}\newline{}老人回答说：\InnerQuote{那么，说万物都受\InnerQuote{起源}支配，这是亵渎吗？}\newline{}\newline{}我回答说：\InnerQuote{毫无疑问是的。因为如果人类的所有罪恶、一切不虔和放荡的行为都源自星辰，而星辰又是天主为完成这工作而设立的，成为所有邪恶的成因，那么一切的罪恶就都追溯到了那位将\InnerQuote{起源}置于星辰中的天主身上。}}\par

% structure: p0012 source=h2 target=subsection reason=relative-heading-rank; base=h2

\subsection{第六章　老人以个人经验反对\Person{伯多禄}的论证}

\Quote{老人回答说：\InnerQuote{你说得对，然而，尽管你的论证无与伦比，我仍因个人的知识而无法赞同。因为我曾是占星家，先住在\Place{罗马}；后来与一位\Person{凯撒}家族的人结为朋友，我精确地查考了他本人和他妻子的\InnerQuote{起源}。追踪他们的历史，我发现所有事迹都确实按照他们的\InnerQuote{起源}实现了，因此我不能屈服于你的论证。因为她的\InnerQuote{起源}的排列是那种使妇女犯奸淫、爱上自己的奴隶、并死于异乡水中的类型。而这事确实发生了：她爱上了自己的奴隶，无法忍受羞辱，便与他私奔，匆忙逃往异邦，与他同床，并死于海中。}}\par

% structure: p0014 source=h2 target=subsection reason=relative-heading-rank; base=h2

\subsection{第七章　老人讲述自己的故事}

\Quote{我回答说：\InnerQuote{那么你怎么知道那个逃跑并定居异乡的女人嫁给了奴隶，并且婚后死了呢？}\newline{}\newline{}老人说：\InnerQuote{我十分确信这是真的，并非因为她嫁给了他——我甚至不知道她爱上了他；但在他离开后，她丈夫的兄弟告诉了我整个关于她激情的故事，以及他如何作为一位高尚的人行事，不愿因兄弟之谊玷污其床榻，以及那个可怜的女人（她并不可指责，因为她受\InnerQuote{起源}的驱使不得不做和忍受这一切）如何渴望他，却又对他和他的斥责感到畏惧，以及她如何策划了一个梦——真假我无法分辨；因为他说她曾说：\InnerQuote{有一个人站在我面前异象中，命令我立即带着我的孩子们离开\Place{罗马}城。} 然而丈夫希望她能同他的儿子们一起得救，便立即将他们送往\Place{雅典}受教育，由母亲和奴隶们陪伴，而他自己则保留了第三个也是最年轻的儿子；因为梦中发出警告的那一位允许这个儿子与父亲同住。过了很久，他没有收到她的信，便亲自多次派人去\Place{雅典}，最终带着我——作为他最忠实的朋友——一同去寻找她。旅途中我竭尽全力与他一同奔波，因为他记得在他昔日兴旺之时，我曾分享他的一切，他爱我胜过其他朋友。最后我们从\Place{罗马}本身出发，航行了这些\Place{叙利亚}地区，在\Place{塞琉西亚}登陆；登陆后没过几天，他因心碎而死。但我来到这里，从那天至今一直靠双手劳动谋生。}}\par

% structure: p0016 source=h2 target=subsection reason=relative-heading-rank; base=h2

\subsection{第八章　老人提供关于\Person{克莱孟}的父亲\Person{法乌斯托}的信息}

\Quote{当老人说完这些，我从他所说的话中得知，他所说的那位死去的老人正是你们的父亲。然而，我不想在与你商议之前就将你们的情况告诉他。但我查到了他的住处，也把我的住处指给他看；为了确认\Emph{我的猜测是对的}，我只问了他一个问题：\InnerQuote{那位老人叫什么名字？}他说：\InnerQuote{\Person{法乌斯托}}。\InnerQuote{他的双胞胎儿子叫什么名字？}他回答：\InnerQuote{\Person{法乌斯蒂诺}和\Person{法乌斯蒂尼安诺}}。\InnerQuote{第三个儿子叫什么名字？}他说：\InnerQuote{\Person{克莱孟}}。\InnerQuote{他们的母亲叫什么名字？}他说：\InnerQuote{\Person{玛提迪雅}}。因此，出于同情，我与他一同流泪，然后遣散了众人，来到你这里，为的是在我们一同吃过饭后与你商议。但我不想在吃饭前把这件事告诉你，免得你因悲伤而陷于忧愁，在就连天使也欢庆的圣洗之日仍然忧郁。}\newline{}\newline{}听了\Person{伯多禄}的这些话，我们都与母亲一同哭泣。但他看见我们流泪，便说：\Quote{现在，你们每个人因敬畏天主，要勇敢地承受所说的话：因为你们的父亲并非今天才死，而是很久以前就死了，正如你们猜测的那样。}\par

% structure: p0018 source=h2 target=subsection reason=relative-heading-rank; base=h2

\subsection{第九章　\Person{法乌斯托}本人出现}

当\Person{伯多禄}说这话时，我们的母亲再也无法忍受，便喊道：\Quote{哀哉！我的丈夫！你爱我们，却自决而死，而我们还活着，得见光明，刚刚吃了食物。}这一声喊叫还未停歇，看！那老人进来了，同时想询问喊叫的原因，他看着那妇人说：\Quote{这是什么意思？我见到谁了？}他走到她面前，注视着她，也更仔细地注视她，拥抱了她。但他们几乎因突如其来的喜悦而昏厥，想要彼此交谈，却因那未满足的喜悦而说不出话来，因为他们被言语不能所困。但不久之后，我们的母亲对他说：\Quote{我现在拥有你，福斯托，你是我最亲爱的人。那么你怎么还活着，我们刚才听说你死了？但这是我们的儿子：法斯提诺、法斯蒂尼亚诺和克莱孟。}她这样说时，我们三个人都扑到他身上，亲吻他，以一种相当模糊的方式唤起了我们对他的形象的记忆。\par

% structure: p0020 source=h2 target=subsection reason=relative-heading-rank; base=h2

\subsection{第十章 福斯托向伯多禄解释他的叙述}

\Person{伯多禄}看见这事，说：\Quote{你是福斯托，这妇人的丈夫，她孩子的父亲吗？}他说：\Quote{我是。}\Person{伯多禄}说：\Quote{那么，你怎么向我讲述你自己的经历，仿佛是在讲别人的事，告诉我你的劳苦、忧愁和埋葬？}我们的父亲回答说：\Quote{我出身于\Person{凯撒}的家族，不想被认出，因此我假托别人的名字编造了叙述，以免被人察觉我是谁。因为我知道，如果我被认出，当地的地方官们会知道这事，就会把我召回以满足\Person{凯撒}，并将那种我早已决心告别的前繁荣强加给我。因为我不能沉溺于奢华的生活，而我已对自己宣告了最严厉的谴责，因为我深信我是我所爱的人死亡的起因。}\par

% structure: p0022 source=h2 target=subsection reason=relative-heading-rank; base=h2

\subsection{第十一章 论星命学}

\Person{伯多禄}说：\Quote{但关于星命学，你肯定它的时候是仅仅在演戏，还是真心坚持它存在？}我们的父亲说：\Quote{我不会对您说假话。我真心坚持星命学存在。因为我对这门学问并非外行；相反，我与一位最优秀的星象家交往，一个名叫\Person{安努比翁}的埃及人，他在我旅行开始时成为我的朋友，并向我透露了我妻子和孩子的死亡。}\Person{伯多禄}说：\Quote{你现在难道不是被事实说服，星命学的教义没有坚实基础吗？}我的父亲回答说：\Quote{我必须向你摆出我心中出现的一切想法，好让你听了之后能理解你的反驳。我确实知道，星象家既会犯很多错误，也常常说真话。因此我猜想，他们之所以说真话，是因为他们精确了解这门学问，而他们的错误则是无知的结果；所以我推测这门学问有坚实基础，但星象家们自己之所以说假话，纯粹是由于无知，因为他们无法绝对准确地知道所有事情。}\Person{伯多禄}回答说：\Quote{考虑一下，他们所说真话是否并非偶然，他们是否在并不精确了解事情的情况下就做出宣告。因为当许多预言被说出来时，有些当然会成真。}那老人说：\Quote{那么，怎样才能完全确信这一点呢，即星命学是否有可靠的基础？}\par

% structure: p0024 source=h2 target=subsection reason=relative-heading-rank; base=h2

\subsection{第十二章 克莱孟承担讨论}

当两人都沉默时，我说：\Quote{既然我精确知道这门学问，而我们的主和我们的父亲却不在这种状态，我希望能有\Person{安努比翁}本人在这里，在我的父亲面前与他讨论。因为这样，当一位实际了解此学科的人与另一位同样知情的人进行讨论时，事情才能公开。}我们的父亲回答说：\Quote{那么，在哪里可能遇见\Person{安努比翁}呢？}\Person{伯多禄}说：\Quote{在\Place{安提约基雅}，因为我得知\Person{西满·玛古斯}在那里，\Person{安努比翁}是他形影不离的同伴。那么，当我们去那里时，如果遇到他们，讨论就可以进行。}于是，当我们讨论了许多主题，为相认而欢喜，并感谢天主时，夜幕降临了，我们便去睡了。\par

\SourceRecord{Translated by James Donaldson.；Ante-Nicene Fathers；Vol. 8.；Edited by Alexander Roberts, James Donaldson, and A. Cleveland Coxe.；Buffalo, NY: Christian Literature Publishing Co.,；1886.；<http://www.newadvent.org/fathers/080814.htm>.}

% structure: p0001 source=h1 target=section reason=page-title

\section{\WorkTitle{讲道集第十五篇}}

% structure: p0002 source=h2 target=subsection reason=relative-heading-rank; base=h2

\subsection{第一章：\Person{伯多禄}希望劝化\Person{福斯托斯}}

天破晓时，我们的父亲同我们的母亲和他的三个儿子进入\Person{伯多禄}所在的处所，向他致意后坐下。然后我们也应他的请求同样坐下；\Person{伯多禄}看着我们的父亲，说：\Quote{我希望你与你的妻子和孩子同心合意，这样你在此世就能与他们一同生活，而在灵魂与身体分离后的来世，你也将与他们在一起，没有忧愁。因为你们不能彼此交往，难道不是令你极为悲伤吗？}我们的父亲说：\Quote{确实如此。}\Person{伯多禄}说：\Quote{那么，如果现在彼此分离使你们痛苦，而且毫无疑问等待你的惩罚是死后你们不能与彼此同在，那么你，一个智慧之人，竟因自己的见解与家人分离，你的悲伤将有多大呢？他们也必定因意识到永恒的惩罚在等待你，因为你持有与他们不同的见解，并否认已经确立的真理，而感到更加痛苦。}\par

% structure: p0004 source=h2 target=subsection reason=relative-heading-rank; base=h2

\subsection{第二章：聆听\Person{伯多禄}论据的理由}

我们的父亲说：\Quote{但并非如此，我亲爱的朋友，灵魂在阴间不会受罚，因为灵魂一离开身体就消散为空气了。}\Person{伯多禄}说：\Quote{在我们说服你相信这一点之前，请回答我：你难道不认为自己没有悲伤，是因为你不相信\Emph{未来的}惩罚吗？而他们完全相信惩罚存在，必定为你忧虑。}我们的父亲说：\Quote{你说得有理。}\Person{伯多禄}说：\Quote{那么，你为什么不通过同意他们的宗教来使他们从对你的最大忧虑中解脱出来呢？我不是说出于恐惧，而是出于善意，聆听并判断我所说的，是否属实。如果真相如我们所说，那么在此世你将与最亲爱的人一同生活，在另一世界你将与他们一同安息；但如果你在检验论据时证明我们所说的是虚构的故事，那么你将做一件好事，因为你会把你的朋友争取到你这边，并终止他们依赖虚假希望的倾向，使他们摆脱虚假的恐惧。}\par

% structure: p0006 source=h2 target=subsection reason=relative-heading-rank; base=h2

\subsection{第三章：信仰的障碍}

我们的父亲说：\Quote{你的话显然很有道理。}\Person{伯多禄}说：\Quote{那么，是什么阻止你来到我们的信仰呢？请告诉我，以便我们从此开始讨论。因为阻碍有很多。信者被商业、公务、耕种、忧虑等事务阻碍；不信者（你也是其中之一）被以下观念阻碍：认为根本不存在的诸神确实存在，或者认为万物受命运支配，或受偶然支配，或者认为灵魂是会死的，或者认为我们的教义是虚假的，因为没有天主的眷顾。}\par

% structure: p0008 source=h2 target=subsection reason=relative-heading-rank; base=h2

\subsection{第四章：\Person{福斯托斯}及其家庭生活中所见的天主眷顾}

\Quote{但我坚持认为，从你所遭遇的事来看，万事皆由天主的眷顾安排，而你与家人分离这么多年也是天主的安排；因为如果他们与你在一起，或许就不会聆听真宗教的教义，所以安排你的孩子与母亲一同旅行，遭遇海难，被认为已死，并被贩卖；此外，他们受教于希腊学识，尤其是无神论学说，以便作为熟悉它们的人能够更好地驳斥它们；此外，他们还皈依了真宗教，并与我联合，以便帮助我宣讲；再进一步，他们的兄弟\Person{克莱孟}应在同一地点遇见他们，从而他母亲被认出，并通过她的治愈而完全信服了正确的敬拜天主；不久之后，双胞胎彼此认出，那天与你会面，你重新获得你的亲人。因此，我认为，如此迅速汇聚的各种情况，好像来自四面八方，以完成一个计划，不可能没有上智安排的指引。}\par

% structure: p0010 source=h2 target=subsection reason=relative-heading-rank; base=h2

\subsection{第五章：真宗教与哲学的区别}

我们的父亲开始说：\Quote{不要以为，我最亲爱的\Person{伯多禄}，我不思考你所宣讲的教义。我确实在想。但在昨夜，当\Person{克莱孟}热切地敦促我同意你所宣讲的真理时，我最后回答说：\InnerQuote{为什么？因为有什么新的诫命能比古人嘱咐我们遵守的更多呢？}而他温和地微笑道：\InnerQuote{父亲，真宗教的教义与哲学的教义之间有很大区别；因为真宗教从预言得到证明，而哲学给我们提供优美的语句，似乎从推测提出证明。}说完这句话，他举了一个例子，将你向他解释的仁爱之道放在我们面前，那在我看来似乎非常不公正，我将告诉你原因。他认为，把另一边脸也伸给打你耳光的人，把外衣也交给拿你外套的人，和强迫你走一里路却走两里，以及诸如此类，是正确的。}\par

% structure: p0012 source=h2 target=subsection reason=relative-heading-rank; base=h2

\subsection{第六章：仁爱}

伯多禄回答说：\Quote{你曾认为最公正之事为不公正。若你愿意，是否听我说？} 我的父亲说：\Quote{全心全意。} 伯多禄说：\Quote{你以为如何？假设有两个彼此为敌的国王，他们的国家彼此隔绝；假设其中一国的臣民在另一国境内被捕，因而被判死罪：如今若他被打一顿而不处死，便得以免罚，那么赦免他的人岂不是仁爱之人吗？} 我们的父亲说：\Quote{确实如此。} 伯多禄说：\Quote{再假设同一个人偷了某人属于他或别人的东西；若被抓获，他赔偿双倍，而非承受应得的惩罚，即赔偿四倍并被处死，因为是在敌方领土被抓获；那么你认为，接受双倍赔偿并赦免他死刑的人，岂不是仁爱之人吗？} 我们的父亲说：\Quote{他显然如此。} 伯多禄说：\Quote{那么如何？居于他人王国，且是敌对的邪恶君主的国中之人，难道不应该为了活命而取悦众人，当受强力时更加退让，向不招呼的人打招呼，与仇敌和好，不与发怒的人争吵，将自己的财产慷慨给予一切求取者，以及诸如此类之事吗？} 我们的父亲说：\Quote{他若宁要生命而不要它们，便有理由忍受一切。}\par

% structure: p0014 source=h2 target=subsection reason=relative-heading-rank; base=h2

\subsection{\Emph{第七章 比喻的解释；今世与来世}}

伯多禄说：\Quote{那么，那些你认为受到不公正待遇的人，难道不正是违法者吗？既然他们身处另一个王国，他们所有的一切不正是通过欺诈得来的吗？而那些被认为行为不义的人，却在善待那敌对王国的每个臣民，允许他们拥有财产。因为这些财物属于那些选择了现世的人。他们如此仁慈，竟允许他人存活。这就是比喻；现在请听实际真理。显现\Emph{于世}的真理先知教导我们，万有的创造者和天主将两个王国赐予了两个——善与恶；授予恶者现世的主权连同法律，使他（即\Emph{它}）有权惩罚行为不义的人；而给予善者未来的永恒。但祂使人自由，有能力将自己交付给他所偏好的，无论是现世的恶还是未来的善。那些选择现世的人有权富有，纵情奢华，享乐，以及做一切能做的事。因为他们不会拥有任何未来的好处。但那些决心接受未来国度祝福的人，无权将此地之物视为己有，因为它们属于异国的君王，唯独水和面包，以及为维持生命而流汗得来的物品（因为他们不得自杀），还有一件外衣，因为不得赤身露体，以避全视之天。}\par

% structure: p0016 source=h2 target=subsection reason=relative-heading-rank; base=h2

\subsection{\Emph{第八章 现世与来世}}

\Quote{那么，若你希望获得准确的说法，请听。你方才说他们受到不公正待遇的那群人，其实自己行不义；因为选择了未来祝福的人，在现世与恶人一同生活，享受许多与恶人相同的享乐——如生命本身、光、面包、水、衣服，以及其他类似之物。但那些你认为是行不义的人，却不会在来世与善人一同生活。} 我们的父亲回答：\Quote{现在你让我相信，行不义的人自己受到不公正待遇，而遭受不公正的人却远为有利，整个事情在我看来仍然是最不公正的交易；因为那些似乎行不义的人允许了许多事物给那些选择了未来祝福的人，而那些似乎受到不公正待遇的人却在行不义，因为他们没有在另一个世界中，给予那些在现世给了他们祝福的人同样的好处，正如那些人所给予他们的。} 伯多禄说：\Quote{这完全不公正，因为每个人有权选择现世或未来的好处，无论大小。他凭自己的判断和意愿选择，没有受不公正待遇——我的意思是，即使他的选择落到了小事上也不受不公，因为大事也在他的选择范围之内，正如小事一样。} 我们的父亲说：\Quote{你说得对；因为一位希腊智者曾说过：\InnerQuote{过失在于选择者——天主无可指责。}}\par

% structure: p0018 source=h2 target=subsection reason=relative-heading-rank; base=h2

\subsection{\Emph{第九章 财产即是过犯}}

\Quote{你能否也解释一下这件事？我记得克肋孟对我说，我们遭受伤害和磨难是为了赦免我们的罪。} 伯多禄说：\Quote{这完全正确。因为我们这些选择了未来事物的人，就我们比别人拥有更多财物，无论是衣服、食物、饮料，或任何其他东西而言，我们拥有罪，因为我们本不应拥有任何东西，正如我刚才向你解释的。对我们所有人来说，财产即是罪。剥夺这些财产，无论以何种方式发生，就是除去罪。} 我们的父亲说：\Quote{这似乎合理，正如你解释那两位君王的两个分界线，并且每个人都可以在其权下选择任何他所愿的。但为何磨难被降下，或我们公正地忍受它们？} 伯多禄说：\Quote{极其公正；因为得救者的分界线，如我所说，是无人应拥有任何东西，但许多人拥有许多财产，或者说罪，因此天主的无比仁爱向那些心地不纯洁的人降下磨难，为了使他们对天主有一点爱，他们可以借着暂时的惩罚，从永远的惩罚中得救。}\par

% structure: p0020 source=h2 target=subsection reason=relative-heading-rank; base=h2

\subsection{\Emph{第十章 贫穷未必是义}}

我们的父亲说：\Quote{这是怎么回事呢？我们不是看见许多不虔敬的人贫穷吗？那么，这些人就因此属于得救的人吗？} 伯多禄说：\Quote{绝非如此。因为那种渴求不应得之物的贫穷，是不蒙悦纳的。所以有些人在他们自己的选择上是富有的，虽然在实际财富上是贫穷的，他们因渴望拥有更多而受惩罚。但一个人并不因为他恰好贫穷就无可置疑地是义人。因为就实际财富而言，他可能是一个乞丐，但他可能渴望甚至去做他最不应该做的事。因此，他可能崇拜偶像，或是亵渎者或奸淫者，或他可能不分青红皂白地生活，或发假誓，或说谎，或过不信者的生活。但我们的老师宣告有信德的穷人是蒙福的；他这样做，不是因为他们给了什么，因为他们什么也没有，而是因为他们没有犯罪，仅仅因为他们没有给施舍，因为他们没有什么可以给，就不应被定罪。} 我们的父亲说：\Quote{确实，就讨论的主题而言，一切似乎都很正确；因此我决定按顺序听完你的全部论点。}\par

% structure: p0022 source=h2 target=subsection reason=relative-heading-rank; base=h2

\subsection{第十一章 应许的真宗教之阐述}

伯多禄说：\Quote{既然你们从此热切地渴望学习与我们宗教相关的事，我应当按顺序加以解释，从天主本身开始，并表明我们应唯独称祂为天主，既不应称其他者为神，也不应视他们为神，而违背此事的人将受到永恒的惩罚，因为他显出了对万有之主的最大不虔敬。} 说了这话，他就按手在那些受苦难折磨、患病和被魔鬼附身的人身上；并祈祷着，治好了他们，解散了群众。然后他这样进去，吃了惯常的食物，就去睡了。\par

\SourceRecord{Translated by James Donaldson.；Ante-Nicene Fathers；Vol. 8.；Edited by Alexander Roberts, James Donaldson, and A. Cleveland Coxe.；Buffalo, NY: Christian Literature Publishing Co.,；1886.；<http://www.newadvent.org/fathers/080815.htm>.}

% structure: p0001 source=h1 target=section reason=page-title

\section{讲道集第十六篇}

% structure: p0002 source=h2 target=subsection reason=relative-heading-rank; base=h2

\subsection{第一章　西满愿与伯多禄讨论天主的唯一性}

在拂晓时分，伯多禄外出，来到他惯常讲道的地方，看见聚集了一大群人。正当他要开始讲论时，他的一位执事进来说：\Quote{西满已从\Place{安提约基雅}来了，他在傍晚时就动身，因听说你承诺要讲论天主的唯一性；他准备与伊壁鸠鲁派的\Person{雅典诺多录}一同来听你的讲论，为要公开反驳你为天主唯一性所提出的一切论据。}执事刚说完这话，看！西满本人进来了，由\Person{雅典诺多录}和其他几位朋友陪同。伯多禄尚未开口，西满便抢先说道：——\par

% structure: p0004 source=h2 target=subsection reason=relative-heading-rank; base=h2

\subsection{第二章　同一主题续篇}

\Quote{我听说你昨天向\Person{法乌斯托}承诺，今日要按部就班地提出你的论据，从宇宙之主开始，证明我们应当说唯有祂是天主，我们既不该说也不该想有别的神，因为违背这一点的将受永罚。但我最惊异的是，你竟妄想使一位智者、一位年高德劭之人转变为你的思想。但你的图谋不会得逞；更何况我在此，能彻底驳斥你的谬论。因为倘若我不在场，这位明智的老人或许会被引入歧途，因他对犹太人中间公认的典籍并无批判性的了解。目前我将略去许多，以便更快地驳斥你承诺要证明的东西。所以，请在我们这些了解圣经的人面前，开始讲说你承诺要说的话吧。但倘若你惧怕我们的反驳，不愿在我们面前履行承诺，这本身就足以证明你是错误的，因为你竟敢在了解圣经的人面前发言。不过，我何必等待你告诉我呢？在场的这位老人便是你承诺的有力见证。}说完这话，他转向我的父亲说：\Quote{告诉我，万民中最卓越的人，这个人不是承诺今日要向你证明天主是唯一的，我们不该说或想有任何别的神，并且违背这一点的要受永罚，犯了最严重的罪吗？那么，你拒绝回答我吗？}\par

% structure: p0006 source=h2 target=subsection reason=relative-heading-rank; base=h2

\subsection{第三章　讨论的方式}

我们的父亲说：\Quote{西满，如果\Person{伯多禄}先否认他曾做过承诺，你或许有理由向我索要证词。但此刻，我却不羞于说出我应说的话。我认为你是在愤怒中想要开始讨论。这种心态下，你说话和听你说话都不适宜；因为我们不再帮助走向真理，而是观看一场争斗。如今，我从希腊文化中学到寻求\Emph{真理}者应当如何行事，我将提醒你。你们每人陈述自己的观点，轮流发言。因为如果只有伯多禄愿意阐述他的思想，而你对你的思想保持缄默，那么你提出的某个论据可能同时摧毁你和他二人的意见；而你们二人虽然被这一论据击败，却不会显得被击败，只有阐述自己意见的那一位显得被击败；而未曾阐述自己意见的那一位，虽然同样被击败，却不会显得被击败，甚至被认为得胜了。}西满回答说：\Quote{我按你说的做；但我担心你并非爱好真理的裁判者，因为你已被他的论据所左右。}\par

% structure: p0008 source=h2 target=subsection reason=relative-heading-rank; base=h2

\subsection{第四章　法乌斯托的成见更偏向西满而非伯多禄}

我们的父亲回答说：\Quote{不要强迫我在不经判断的情况下同意你，以便我显得是个爱好真理的裁判者；但若你希望我告诉你实情，我的成见是偏向你的意见的。}西满说：\Quote{你不知我的意见为何，怎能如此？}我们的父亲说：\Quote{这很容易知道，我告诉你。你承诺要证明\Person{伯多禄}在维护天主的唯一性上是错误的；但如果有人承诺要证明维护天主唯一性的人是错误的，那么很明显，作为正确的一方，他并不持有同样的意见。因为如果他持有与完全错误之人同样的意见，那么他自己也是错误的；但如果他在提出证明时持有相反的意见，他就是正确的。所以，你说维护天主唯一性的人是错误的，这并不正确，除非你相信有许多神。而我相信有许多神。因此，在讨论之前我持有与你相同的意见，我自然是偏向你的。为此，你不必为我担忧，伯多禄却应该，因为我的意见仍与他相反。所以希望在你们的讨论之后，作为一个爱好真理、抛开成见的裁判者，我将赞同那得胜的教义。}我父亲这样说时，群众中不由自主地爆发出一阵欢呼的喃喃声。\par

% structure: p0010 source=h2 target=subsection reason=relative-heading-rank; base=h2

\subsection{第五章　伯多禄开始讨论}

伯多禄于是说：\Quote{我愿按我们讨论的裁判官所说的去做；立刻毫不迟延地陈述我对天主的意见。我断言有一位天主创造了天和地，以及其中的一切。不可以说或想有别的任何神。}西满说：\Quote{但我坚持认为犹太人中间相信的圣经说有许多神，而天主对此并不发怒，因为祂在祂的圣经中亲自提到了许多神。}\par

% structure: p0012 source=h2 target=subsection reason=relative-heading-rank; base=h2

\subsection{第六章　西满援引旧约证明有许多神}

\Quote{例如，在法律的起初话语中，祂显然说到他们与自己相似。因为经上记载，当第一个人从天主领受命令，可以吃园中所有树上的果子，但不可吃知善恶树上的果子时，蛇藉着女人说服了他们，藉着他们将成为诸神的应许，使他们举目仰望；然后，当他们这样仰望时，天主说：\BibleQuote{看，亚当成了我们中的一个。}当蛇说：\BibleQuote{你们将如同诸神一样}时，它清楚地说话是基于相信诸神存在；更何况天主也加上了祂的见证，说：\BibleQuote{看，亚当成了我们中的一个。}因此，说有许多诸神的蛇并没有说谎。又如圣经说：\BibleQuote{不可咒骂诸神，也不可诅咒你百姓的首领}，这指出了许多诸神，甚至连诅咒他们也不可以。但在别处又记载：\BibleQuote{难道有另一个神敢进去，从另一民族中为自己夺取一个民族，如同我上主天主所做的那样？}当祂说：\BibleQuote{难道有另一个神敢？}祂是在假设有其他神存在的前提下说话。还有别处：\BibleQuote{那没有创造天地的诸神，应当灭亡}，仿佛那些创造天地的诸神不会灭亡。在另一处，当它说：\BibleQuote{你要谨慎，免得你去事奉你祖先不认识的其他诸神}，它说话如同有其他诸神存在，不可跟随他们。又如：\BibleQuote{其他诸神的名号不可挂在你的口上}，这里提到了许多诸神，它不希望他们的名号被说出来。又如经上记载：\BibleQuote{你的天主是上主，祂是万神之神。}又如：\BibleQuote{上主啊，在众神中谁能像你？}又如：\BibleQuote{天主是万神之主。}又如：\BibleQuote{天主站在众神的会中，在众神中施行审判。}因此，我惊奇，既然经上有这么多段落证明有许多诸神，你却断言我们既不应说也不应认为有许多诸神。最后，如果你对如此明确所说的话有任何反驳，请在众人面前说出来。}\par

% structure: p0014 source=h2 target=subsection reason=relative-heading-rank; base=h2

\subsection{第七章 伯多禄引用旧约证明天主的唯一性}

\Quote{我要简短答复你所说的话。那经常提到诸神的法律，它自己向犹太群众说：\BibleQuote{看，上主你的天主拥有诸天和天上的一切，}这暗示即使有诸神，他们也在祂之下，即在犹太人的天主之下。又说：\BibleQuote{上主你的天主，祂是天上地下的天主，除祂以外没有别的神。}别处圣经又向犹太群众说：\BibleQuote{上主你的天主是万神之神，}所以即使有诸神，他们也在犹太人的天主之下。别处圣经论到祂说：\BibleQuote{天主，伟大而真实，不看情面，也不受贿赂，祂为孤儿寡妇伸冤。}圣经称犹太人的天主为伟大而真实，并施行审判，就表明其他的神是渺小而不真实的。但另一处圣经又说：\BibleQuote{我指着我的永生起誓，上主说，除我之外没有别的天主。我是首先的，我是以后的；除我以外没有天主。}又说：\BibleQuote{你要敬畏上主你的天主，只事奉祂。}又说：\BibleQuote{以色列，你要听：上主你的天主是唯一的主。}此外还有许多经文以誓言证实天主是唯一的，除祂以外没有天主。因此，我惊奇，既然有这么多经文证明天主是唯一的，你却说是许多。}\par

% structure: p0016 source=h2 target=subsection reason=relative-heading-rank; base=h2

\subsection{第八章 西蒙与伯多禄继续讨论}

西蒙说：\Quote{我与你最初的约定是，我应从圣经证明你是错误的，因为你坚持我们不应说有许多诸神。因此我引用了许多经文章节，表明神圣的圣经本身就在说有许多诸神。}伯多禄说：\Quote{那说到有许多诸神的圣经也劝告我们说：\BibleQuote{其他诸神的名号不可挂在你的口上。}所以，西蒙，我并没有违反圣经。}西蒙说：\Quote{伯多禄，请你听我说。在我看来，你似乎有罪，因为你说话反对他们，而圣经说：\BibleQuote{不可咒骂\Emph{诸神}，也不可诅咒你百姓的首领。}}伯多禄说：\Quote{我并没有犯罪，西蒙，因为我按照圣经指出他们的灭亡；因为经上记载：\BibleQuote{那没有创造天地的诸神，应当灭亡。}祂这样说，并不是像你所解释的那样，好像\Emph{有些}创造了天地而不会灭亡。因为在圣经的开头清楚宣告：\BibleQuote{起初天主创造了天地。}它并没有说\InnerQuote{诸神}。别处又说：\BibleQuote{穹苍显示祂手的化工。}另一处也记载：\BibleQuote{诸天将要灭没，你却永远长存。}}\par

% structure: p0018 source=h2 target=subsection reason=relative-heading-rank; base=h2

\subsection{第九章 西蒙试图证明圣经自相矛盾}

\Quote{我从圣经引用了明确的段落证明有许多诸神；而你，作为答复，从同一圣经提出了同样多或更多段落，表明天主是唯一的，祂是犹太人的天主。当我说我们不应咒骂诸神时，你接着指出那创造者是唯一的，因为那些没有创造的诸神将要灭亡。作为对我主张应该认为有诸神的答复（因为圣经也这样说），你表明我们不应说出他们的名字，因为同一圣经告诉我们不可提起其他诸神的名号。既然这些圣经有时说有许多诸神，有时又说只有一位；有时说它们不应被咒骂，有时又说应当咒骂；结果我们该得出什么结论呢？岂不是圣经本身误导我们吗？}\par

% structure: p0020 source=h2 target=subsection reason=relative-heading-rank; base=h2

\subsection{第十章 伯多禄对圣经明显矛盾的解释}

伯多禄说：\Quote{他们不会引人误入歧途，却会揭露并光照那像蛇一样潜伏在每人心中对抗天主的邪恶性情。因为圣经就像多种多样的模型摆在每人面前。于是，每人都有自己的性情如同蜡块，研读圣经并在其中找到一切，他就按照自己的意愿塑造对天主的观念，如同我所说的，把自己的性情像蜡一样盖在上面。因此，既然每人在圣经中都找到自己希望关于天主的意见，那么，他——\Emph{\Person{西蒙}}——就从圣经中塑造出许多神的形状，而我们则塑造了那真实存在者的形状，从我们自身的模样认识了真实的典型。因为我们的灵魂确实穿戴着祂的肖像而得永生。如果我抛弃这灵魂的父亲，它也会抛弃我去接受公义的审判，通过这胆大妄为的行为显明不义；既然它来自一位公义者，它会公义地抛弃我；因此，就灵魂而言，我在受罚之后将被毁灭，因为抛弃了来自它的帮助。但若另有别的神，首先让他穿上另一种形式、另一种形状，好让我通过身体的新形状认识那新神。但如果他改变形状，难道他就因此改变了灵魂的实体吗？但如果他也改变了灵魂，那么我就不再是我自己，在形状和实体上都变成了另一个。所以，若有另一位，就让他创造别的吧。但并没有。因为若曾有，他早就创造了。既然他没有创造，那么就让这位不存在者离开那真实存在者吧。因为他什么也不是，只存在于西蒙的意见中。除了创造我的那一位，我不接受任何别的神。}\par

% structure: p0022 source=h2 target=subsection reason=relative-heading-rank; base=h2

\subsection{第十一章 \BibleRef{创 1:26}被\Person{西蒙}引用}

\Person{西蒙}说：\Quote{既然我看到你经常提到创造你的天主，那就听我讲讲你如何甚至对祂不敬。因为显然有两位创造者，正如圣经所说：\BibleQuote{天主说：让我们照我们的肖像，按我们的模样造人。} 现在\InnerQuote{让我们造}意味着两个或更多，肯定不是一位。}\par

% structure: p0024 source=h2 target=subsection reason=relative-heading-rank; base=h2

\subsection{第十二章 伯多禄对经文的解释}

\Person{伯多禄}回答说：\Quote{祂是一位，曾对祂的智慧说：\BibleQuote{让我们造一个人。}但祂的智慧是祂自己常常与之欢愉如同自己的神的那一位。它如同灵魂与天主结合，却被祂伸展如同手，塑造宇宙。为此，一个人被造，从他而出了女性。它在种类上是统一，却是双重，因为通过扩展和收缩，统一被认为是双重。所以我做得对，将一切荣耀归于一位天主，如同归于父母。} \Person{西蒙}说：\Quote{那么怎样？即使圣经说有别的神，你也不接受这个意见吗？}\par

% structure: p0026 source=h2 target=subsection reason=relative-heading-rank; base=h2

\subsection{第十三章 圣经的矛盾意在考验读经者}

\Person{伯多禄}回答说：\Quote{如果圣经或先知说到神祇，他们这样做是为了考验听者。因为经上记载：\BibleQuote{如果你们中间兴起一位先知，显神迹和奇事，那神迹和奇事也应验了，他对你们说：\InnerQuote{我们去随从并事奉你们祖先所不认识的其他神祇。} 你们不可听从那位先知的话；你们的手要首先向他投石，将他处死。因为他已试图引你们离开上主你们的天主。但如果你们心里说：\InnerQuote{他怎能行那神迹或奇事？} 你们应当知道：那考验你们的，考验你们是要看看你们是否敬畏上主你们的天主。} \InnerQuote{那考验你们的，考验你们}这句话指的是早期时代；但在遷徙到巴比伦之后似乎有所不同。因为天主知道一切，可以用许多论点证明，祂不会为了自己知道而考验，因为祂预知一切。但如果你愿意，让我们讨论这一点，我将证明天主预知。但已经证明祂不知道的观点是错误的，且这是为了考验我们而写的。因此，我们，西蒙，既不会被圣经也不会被任何人引入歧途；我们也不会被欺骗而承认许多神，也不会同意任何反对天主的话。}\par

% structure: p0028 source=h2 target=subsection reason=relative-heading-rank; base=h2

\subsection{第十四章 其他被称为神的受造物}

\Quote{因为我们自己也晓得，天使被圣经称为神——例如，那位在荆棘中说话（\BibleRef{出 3:2}）、并与雅各伯摔跤的（\BibleRef{创 32:24-30}）——这名称也同样应用于那位生于厄玛奴耳、被称为大能的天主（\BibleRef{依 7:14}）。是的，就连梅瑟对法郎也成了神，尽管他实际是人（\BibleRef{出 7:1}）。外邦人的偶像也是如此。但我们只有一位天主，就是那创造万物、安排宇宙的，祂的儿子是基督。顺从基督，我们从圣经学会辨别什么是假的。此外，我们的祖先以虔诚的心确立了祂是真正天主的坚定信仰，并将这信仰传给了我们，使我们知道，若有人说什么反对天主的话，那就是谎言。我还要附加一句必要的话：如果事情不像我所说的，那么，愿我和所有喜爱真理的人，在赞美造我们的天主之事上遭遇危险。}\par

% structure: p0030 source=h2 target=subsection reason=relative-heading-rank; base=h2

\subsection{第十五章 基督不是天主，而是天主子}

\Person{西蒙}听了这话说：\Quote{既然你说我们连那显神迹和奇事的先知也不该信，若他说有别的神，而且你知道他甚至招致死刑，因此你们的老师也因显神迹和奇事而被合理剪除。} \Person{伯多禄}回答说：\Quote{我们的主既没有断言除了万有的创造者之外还有神，也没有自称是天主，而是有理由称那称祂为那位安排宇宙的天主之子的人为有福。} \Person{西蒙}回答说：\Quote{那么，你不认为那来自天主的本身就是天主吗？} \Person{伯多禄}说：\Quote{请告诉我们这怎么可能；因为我们不能肯定这事，因为我们没有从祂本人那里听到。}\par

% structure: p0032 source=h2 target=subsection reason=relative-heading-rank; base=h2

\subsection{第十六章 非受生者与受生者必然彼此不同}

\Quote{除此以外，圣父的特性是不受生，而圣子的特性是受生；但受生者不能与\OriginalTerm{无始}{unbegotten}或自生者相比。}西满说：\Quote{难道不是因它的起源而相同吗？}伯多禄说：\Quote{若某人在一切方面不与另一个人相同，就不能对那人应用所有相同的称号。}西满说：\Quote{这是在断言，并非证明。}伯多禄说：\Quote{怎么，你没有看到，如果一个人碰巧是自生或无始的，他们就不能被称作相同；也不能断言那受生者与他生者\OriginalTerm{同实体}{of the same substance}？也要学习这个：人的身体有灵魂是不朽的，它们被天主的嘘气所覆蔽；它们出自天主，是同实体的，但它们不是神。但如果它们是神，那么所有已死、尚存、及将要出生之人的灵魂都是神。但若你出于争辩的精神坚持认为这些也是神，那么基督被称为神又有什么了不起呢？因为祂只拥有所有都拥有的。}\par

% structure: p0034 source=h2 target=subsection reason=relative-heading-rank; base=h2

\subsection{第十七章 天主的性体}

\Quote{我们称那一位为天主，祂的特性不能属于任何别的性体；因为，正如祂被称为\OriginalTerm{无界限者}{Unbounded}，因为祂在每一方面都是无限的，那么必然，\InnerQuote{无界限者}这一称呼不是其他任何人的特性，因为另一位不能同样无限。但若有人说可能，他错了；因为两个在每一方面无限的事物不能共存，因为一方被另一方界定。因此，在事物的本性中，\OriginalTerm{无始}{unbegotten}者是唯一的。但若祂拥有\OriginalTerm{形象}{figure}，即使这样，形象也是唯一且无比的。因此祂被称为\OriginalTerm{至高者}{Most High}，因为祂高于万有，使宇宙服从于祂。}\par

% structure: p0036 source=h2 target=subsection reason=relative-heading-rank; base=h2

\subsection{第十八章 天主的名字}

西满说：\Quote{这个\InnerQuote{天主}一词是祂\OriginalTerm{不可言喻的名}{ineffable name}吗，所有人都使用，因为你如此坚持一个名字不可给予别人？}伯多禄说：\Quote{我知道这不是祂不可言喻的名，而是人类约定中给予的一个；但若你将之给予另一位，你也会将不被使用的名字给予另一位；而且是故意。被使用的名是不被使用的名的前驱。这样，冒犯甚至被归给尚未说出的，正如对已知的尊敬被传给未晓的。}\par

% structure: p0038 source=h2 target=subsection reason=relative-heading-rank; base=h2

\subsection{第十九章 人身上的\OriginalTerm{天主的形象}{shape of God}}

西满说：\Quote{我想知道，伯多禄，你是否真的相信人的形象是按天主的形象塑造的？}伯多禄说：\Quote{我真的非常确信，西满，情况正是如此。}西满说：\Quote{死亡如何能解散身体，当它已被如此伟大的\OriginalTerm{印记}{seal}印上？}伯多禄说：\Quote{这是\OriginalTerm{公义的天主}{just God}的形象。那么，当身体开始不义地行动时，它里面的形象就逃离，于是身体被解散，因为形象消失，以便不义的身体不再拥有公义的天主的形象。然而，解散并不发生在印记本身，而是在被印的身体上。但被印的没有那印者就不能解散。因此，没有审判就不允许死亡。}西满说：\Quote{有什么必要给从地土中升起的人赋予这样一位存在的形象呢？}伯多禄说：\Quote{这样做是由于造人的天主的爱。因为，就\OriginalTerm{实体}{substance}而言，一切都优于人的肉体——我指以太、太阳、月亮、星辰、空气、水、火——总之，所有为服务人而造的一切——然而，尽管在实体上更优越，它们甘心情愿地服事实体次等的，因那优越者的形象。因为正如他们尊崇一个国王的泥像，已是尊崇那国王本人——泥像恰具有其形象——同样，整个受造界喜乐地服事从土而出的人，因看到这尊崇归于天主。}\par

% structure: p0040 source=h2 target=subsection reason=relative-heading-rank; base=h2

\subsection{第二十章 天主的特性}

\Quote{看，西满，你所希望说服我们忘恩负义的那位天主的特性，而大地仍承载着你，或许是想看看谁敢抱持与你相似的见解。因为你是第一个胆敢做他人不敢做的：你是第一个说出我们初次听到的话。我们首先且唯独看到了天主的无限忍耐，祂容忍你这般大不敬，而这位天主不是别的，正是世界的创造者，你胆敢对祂行不敬。然而，大地没有裂开，火也没有从天上降下并出去烧灭人，雨也没有倾注，成群的野兽没有从丛林中被派来，而毁灭性的义怒也没有开始向我们显现，因那犯了罪的，仿佛成了\OriginalTerm{属灵的奸淫}{spiritual adultery}，这比肉体的更坏。因为，假如从前惩罚罪恶的不是那创造天地的天主，那么现在，当祂被极度亵渎时，祂理应施加最严厉的惩罚。但相反，祂忍耐，召叫悔改，将那些以毁灭不敬者告终的箭镞贮存在祂的库房里，当祂坐下审判祂的人时，祂要像放出活物一样释放它们。因此，让我们畏惧公义的天主，人的身体带着祂的形象为得尊荣。}\par

% structure: p0042 source=h2 target=subsection reason=relative-heading-rank; base=h2

\subsection{第二十一章 西满许诺求教于基督的训诲；伯多禄遣散群众}

当伯多禄说了这话，西满回答说：\Quote{既然我看出你巧妙地暗示，那些书中所记载的反对\OriginalTerm{世界的创造者}{framer of the world}的话并非真实，明天我将从你老师的言论中证明，他曾断言\OriginalTerm{世界的创造者}{framer of the world}并非至高的天主。}西满说完便出去了。但伯多禄对聚集的群众说：\Quote{即使西满不能在有关天主的事上对我们造成其他伤害，他至少阻止了你们聆听那能净化灵魂的话语。}伯多禄这样说时，人群中起了许多窃窃私语，说：\Quote{有什么必要让他进来，在这里宣讲亵渎天主的话呢？}伯多禄听见了，便说：\Quote{但愿那些为试探人而反对天主的教义仅限于西满！因为正如主所说的，将有假宗徒、假先知、异端、贪图首位者，我猜想，他们始于亵渎天主的西满，将与他一同持守相同的反对天主的意见。}他含着泪说了这话，用手召唤群众到他跟前；他们来了，他就给他们覆手祈祷，然后遣散他们，嘱咐他们第二天更早来。说完，他叹息着进去，没有进食就睡了。\par

\SourceRecord{Translated by James Donaldson.；Ante-Nicene Fathers；Vol. 8.；Edited by Alexander Roberts, James Donaldson, and A. Cleveland Coxe.；Buffalo, NY: Christian Literature Publishing Co.,；1886.；<http://www.newadvent.org/fathers/080816.htm>.}

% structure: p0001 source=h1 target=section reason=page-title

\section{讲道集第十七篇}

% structure: p0002 source=h2 target=subsection reason=relative-heading-rank; base=h2

\subsection{第一章 西满来到伯多禄面前}

第二天，因此，正当\Person{伯多禄}要与\Person{西满}进行讨论时，他比平日更早起身祈祷。祈祷结束后，\Person{匝凯}进来，说：\Quote{\Person{西满}坐在外面，正与大约三十名他自己的特别门徒交谈。}\Person{伯多禄}说：\Quote{让他说下去，直到群众聚集，然后我们这样开始讨论。我们将听他说的一切，然后针对此准备答复，再出去发言。}果然如此发生。\Person{匝凯}因此出去，不久后又进来，将\Person{西满}针对\Person{伯多禄}的演说传达给他。\par

% structure: p0004 source=h2 target=subsection reason=relative-heading-rank; base=h2

\subsection{第二章 西满反对伯多禄的演说}

现在他说：\Quote{他指控你，\Person{伯多禄}，是邪恶的仆人，拥有强大的法术能力，以比偶像崇拜更坏的方式迷惑人的灵魂。为证明你是术士，他似乎向我提出以下证据，说道：\InnerQuote{我察觉到这一点，就是当我来与他讨论时，我记不住我自己默想的任何一个词。因为当他发言时，我的心思忙于回忆我想要在与他会议时说的话，我完全听不到他说什么。既然我只在他面前而不是任何其他人面前经历此事，岂不显然是我受了他的法术影响？}至于他的教义比偶像崇拜更坏，我能向任何有理解力的人清楚说明。因为没有其他益处，只有使灵魂摆脱各类形像。因为当灵魂在眼前呈现一个形像时，它被恐惧束缚，因忧虑而消瘦，唯恐遭受某种灾祸；并且发生变化，落在那魔的影响之下；在那魔的影响下，它在众人眼中似乎聪明。}\par

% structure: p0006 source=h2 target=subsection reason=relative-heading-rank; base=h2

\subsection{第三章 西满对伯多禄的控告}

\Quote{\Person{伯多禄}在承诺使你们有智慧的同时，对你们这样做。因为，在以宣布唯一天主为借口下，他似乎使你们摆脱了许多无生命的形像，这些形像丝毫无损于崇拜它们的人，因为它们本身就被眼睛看见是由石头、铜、金或其他无生命材料制成的。因此，灵魂因为知道所见之物毫无意义，无法因可见之物而被恐惧同等程度地束缚。反而通过骗人的教导仰望一位可怖的天主，它所有的自然根基都被颠覆了。我说这话，并非劝你们崇拜形像，而是因为\Person{伯多禄}似乎使你们的灵魂摆脱可怕的形像，却通过更可怕的形像使你们每个人的心智疯狂，引进一种带有形状的天主，而且是一个极其公义的天主——一个伴随着令沉思的灵魂感到可怕可畏之物的形像，能完全摧毁健全心智的精力。因为心智在如此风暴中，就像被狂风搅动的深渊，翻腾而黑暗。因此，如果他要造福你们，就不要在看似化解来自无生命形状的温和恐惧时，反而引进天主的可怕形状。但天主有形状吗？如果有，祂就有形象。如果有形象，如何不受限制？如果受限制，祂就在空间内。如果在空间内，祂就小于包含祂的空间。如果小于任何事物，祂如何大于一切、超越一切、至高一切？情况就是这样。}\par

% structure: p0008 source=h2 target=subsection reason=relative-heading-rank; base=h2

\subsection{第四章 声称基督的教导与伯多禄的教导不同}

\Quote{而且他实际上不相信他的老师所宣讲的教义，这是明显的，因为他宣扬与老师相反的教义。因为他曾对某个人说，据我所知，\InnerQuote{不要称我为善，因为只有一位是善的。}如今谈到那位善者，他不再谈论那位公义者，就是圣经所宣告的那位，祂杀死又使活着，——杀死那些犯罪的人，使那些按照祂的旨意生活的人活着。但他并没有真正称那位世界的创造者为善，这对任何能思考的人来说都是明显的。因为世界的创造者为祂所造的\Person{亚当}所知，也为取悦祂的\Person{哈诺客}所知，为祂眼中看为正义的\Person{诺厄}所知；同样为\Person{亚巴郎}、\Person{依撒格}、\Person{雅各伯}所知；也为\Person{梅瑟}、百姓和全世界所知。但是\Person{耶稣}，\Person{伯多禄}本人的老师，来了并说：\InnerQuote{除了子，没有人认识父，正如除了父，也没有人认识子，以及子愿意启示的那些人。}如果，那么，是子本人亲自临在，那么从祂出现之时起，祂开始向那些祂愿意的人启示那位不为所有人所知的祂。因此，父对所有在祂之前生活的人是未知的，故此不能是那位为所有人所知的祂。}\par

% structure: p0010 source=h2 target=subsection reason=relative-heading-rank; base=h2

\subsection{第五章 耶稣的教导自相矛盾}

\Quote{耶稣这样说，甚至与自己也不一致。因为有时祂引用圣经中的其他言论，把天主描述为可敬畏而公义的，说：\InnerQuote{不要害怕那杀害肉身，而不能杀害灵魂的；但要害怕那能把灵魂和肉身都投到火狱里的。是的，我告诉你们，要害怕祂。}但祂断言祂确实应被敬畏，因为祂是公义的天主，那些受不义之苦的人向祂呼号，这在祂解释的一个比喻中显明了，说：\InnerQuote{如果不义的法官因不断受恳求而这样做了，何况天父岂不更要为那些昼夜向祂呼号的人伸冤吗？或者你们以为，因为祂容忍他们，祂就不会做吗？是的，我告诉你们，祂必会做，而且要快快地做。}那称天主为施行报应和赏报的天主的人，把祂描述为本质上公义，而非良善。此外，祂感谢天地的主。但如果祂是天地的主，就被承认为世界的造物主；如果是造物主，那么祂就是公义的。因此，当祂时而称祂为良善，时而称祂为公义时，祂在这一点上与自己不一致。但祂智慧的弟子昨日提出了第三点，即真实的看见比异象更令人满意，却不知道真实的看见可以是人的，而异象则明确来自神性。}\par

% structure: p0012 source=h2 target=subsection reason=relative-heading-rank; base=h2

\subsection{第六章 伯多禄出去答复西满}

\Quote{伯多禄，西满站在外面向群众说了这些话以及类似的话，在他看来他似乎搅乱了多数人的心思。所以你立刻出去，用真理的力量驳倒他的虚假言论。}匝凯说完这话，伯多禄按惯例祈祷后便出去，站在前一天讲话的地方，按他宗教所规定的习俗向群众致意后，开始说道：\Quote{我们的主耶稣基督，祂是真先知（我将在适当时候确凿证明），关于那些与真理有关的事，祂作了简明的宣告，这有两个原因：第一，因为祂习惯于对虔诚的人讲话，他们有足够的知识，能够相信祂以宣告方式所陈述的意见；因为祂的话并不陌生于他们惯常的思维方式；第二，因为祂传道的时间有限，祂没有采用论证的方法，以免把所有有限的时间都花在辩论上，因为这样祂可能会忙于解决少数几个需要脑力劳动才能理解的问题，而不会大量地赐给我们那些与真理有关的言论。因此，祂随自己的意愿陈述意见，像是对能理解祂的人讲话，我们也属于这些人。每当我们对祂所说的任何话有不理解之处（这是很少发生的事），我们就私下问祂，这样祂所说的任何话都不至于让我们不懂。}\par

% structure: p0014 source=h2 target=subsection reason=relative-heading-rank; base=h2

\subsection{第七章 具有天主形状的人}

\Quote{因此，知道我们明白祂所说的一切，并能提供证明，祂就派我们到无知的 \OriginalTerm{外邦人}{Gentiles}那里，为他们施洗以赦免罪过，并命令我们先教导他们。祂的诫命中，第一条也是最大的一条，就是敬畏主天主，单单事奉祂。但祂要我们敬畏那位天主，祂的天使是我们中间最卑微的信徒的天使，他们常在天上观看圣父的面。因为祂有形状，祂有每个肢体，主要是为了美丽，而不是为了使用。因为祂没有眼睛用于观看；祂从各方面都能看见，因为祂的身体无与伦比地比我们内在的视觉之灵更明亮，祂比一切更辉煌，以至太阳的光在祂面前算为黑暗。祂也没有耳朵用于听；因为祂从各方面听、感知、移动、运作、行动。但祂为了人而有最美丽的形状，使心地纯洁的人能看见祂，好让他们因所受的苦而喜乐。因为祂按自己的形状造了人，如同用最大的印模，使人成为万物的统治者和主人，使万物都服从他。因此，祂断定自己就是宇宙，而人是祂的肖像（因为祂本身是不可见的，但祂的肖像——人是可见的），那愿意崇拜祂的人，便尊敬祂可见的肖像，即人。所以，任何人向人做什么，无论是好是坏，都视为向祂做的。因此，从那发出来的审判，将按各人的行为给予报应。因为祂为自己的肖像伸冤。}\par

% structure: p0016 source=h2 target=subsection reason=relative-heading-rank; base=h2

\subsection{第八章 天主的形状：由此驳斥西满的异议}

\Quote{但是有人会说：\InnerQuote{若祂有形状，那么祂也有形象，并且处于空间之中；但若祂处于空间之中，并由于较小而被空间包围，祂怎能超越一切？若祂有形象，祂怎能无所不在？}我首先要向提出这些异议的人说的是：圣经劝我们持有这种情感并相信关于祂的这些陈述；我们知道这些宣言是真实的，因为我们的主耶稣基督为它们作证，我们奉祂的命令必须向你们提供证据证明事实如此。但首先我要谈论空间。天主之空间是\OriginalTerm{非存在者}{non-existent}，而天主是\OriginalTerm{存在者}{existent}。但非存在者不能与存在者相比。因为空间怎能是存在者？除非它是另一个空间，如天、地、水、空气，以及任何其他填充虚空的物体——虚空之所以称为虚空，是因为它是无。因为\InnerQuote{无}是它更恰当的名称。因为所谓虚空除了像一个空无一物的容器，即容器本身之外，还能是什么？但既是虚空，它本身不是空间；空间是虚空的所在，若它确实是容器。因为存在者必定存在于非存在者之中。但我所说的非存在者，正是有些人称为空间的东西，即无。但既是无，它怎能与存在者相比，除非通过表达相反，说它是非存在者，而称非存在者为空间？但即使它是某种东西，我手边有许多例子，但我只满足于一个例子，以表明包围者并不毫无疑问地优于被包围者。太阳是一个圆形形象，完全被空气包围，但它照亮空气，温暖空气，分割空气；若太阳离开空气，空气便陷入黑暗；无论太阳从哪一部分离开，空气都变得寒冷如死；但太阳升起时再次照亮它，温暖它，以更大的美丽装饰它。这样做是通过给予自己的一份，尽管它的实体是有限的。那么，有什么能阻止天主，作为这一切及其他万物的创造者和主宰，拥有形象、形状和美丽，并且使这些品质的沟通从祂自身无限地延伸？}\par

% structure: p0018 source=h2 target=subsection reason=relative-heading-rank; base=h2

\subsection{第九章 天主是宇宙的中心或心脏}

\Quote{因此，真正存在的一位天主，以卓越的形象主持一切，是上方和下方之物的心脏，双重地，从祂如同从中心发出赋予生命和无形体的能力；整个宇宙连同星辰和天界区域、空气、火，以及任何其他存在，被证明是一个在高度上无限、深度上无底、宽度上不可测量的实体，将赋予生命和智慧的本性从祂延伸到\OriginalTerm{三种无限}{three infinites}。因此，从祂各方面发出的这个无限必须存在，以那超越一切并如此拥有形象的一位为中心；因为无论祂在哪里，祂就如同在无限的中心，是宇宙的界限。从祂开始的延伸具有\OriginalTerm{六种无限}{six infinites}的本性；其中一种随着祂向上穿透高处，另一种向下穿透深处，另一种向右，另一种向左，另一种向前，另一种向后；祂自己看着如同各方面都相等的数目，在六个时间间隔中完成了世界，祂自己是安息，并拥有无限的未来时代作为祂的肖象，祂是元始和终结。因为六种无限在祂内结束，并从祂接受向着无限的延伸。}\par

% structure: p0020 source=h2 target=subsection reason=relative-heading-rank; base=h2

\subsection{第十章 天主的本质与形象}

\Quote{这是\OriginalTerm{七的奥迹}{hebdomad}。因为祂自己是整体的安息，将自己作为安息赐给那些在微小尺度内效法祂伟大的人。因为祂是唯一的，有时可理解，有时不可理解，\Emph{有时可限制}，\Emph{有时不可限制}，从祂发出延伸到无限的。因为祂是这样可理解和不可理解，近和远，在这里和那里，作为唯一的存在者，并分享那各方面无限的心智，结果灵魂呼吸并拥有生命；如果它们与身体分离并发现对祂的渴望，它们就被带入祂的怀中，如同冬日山上的雾气，被太阳的光芒吸引，不朽地奔向太阳。因此，若我们用我们的心凝视祂美丽的形象，我们内应当产生怎样的情感！但另一方面，\Emph{谈论美丽}是荒谬的。因为美丽不能脱离形象而存在；人也不能被吸引去爱天主，甚至不能认为自己能看见祂，如果天主没有形态。}\par

% structure: p0022 source=h2 target=subsection reason=relative-heading-rank; base=h2

\subsection{第十一章 对天主的敬畏}

\Quote{但有些与真理无关，并将精力用于服事邪恶的人，以光荣天主为借口，说祂没有形体，为的是，祂既无形状又无形体，便无人能看见祂，以致无人渴慕祂。因为理智看不见天主的形体，就缺少祂。然而，人若没有可以投靠、倚靠的对象，又怎能祈祷呢？因为如果遇不到任何阻力，他便会坠入虚空。\InnerQuote{是的，他说，我们不应敬畏天主，而应爱慕祂。}我同意；但每件善行中行善的意识就能成就这件事。而行善源于敬畏。但他说，敬畏将死亡注入灵魂。不，我断言，它并非注入死亡，而是唤醒灵魂，并使之皈依。也许，禁止敬畏天主的命令可能是正确的，如果我们人类不敬畏许多其他事物的话；例如，同类对我们的阴谋、野兽、蛇、疾病、痛苦、魔鬼，以及成千上万的其他祸患。那么，谁要求我们不敬畏天主，就让他救我们脱离这些，使我们不再敬畏它们；但如果他不能，他又何必嫉妒我们，使我们能以对公义者的敬畏这一种恐惧，从千百种恐惧中得救，并藉着对祂的一点信德，就能除去自己和他人身上的千百种痛苦，反而得到福乐的交换，并且由于敬畏那洞察万有的天主而不再作恶，甚至在此生也能安享和平呢？}\par

% structure: p0024 source=h2 target=subsection reason=relative-heading-rank; base=h2

\subsection{第十二章 敬畏天主与爱慕天主}

\Quote{如此，对那真正是主的事奉，感恩的事奉，使我们脱离所有其他主人的事奉。因此，如果有人能在不敬畏天主的情况下脱离罪恶，就让他不要敬畏；因为，在爱慕祂的影响下，人不能做令祂不悦的事。因为，一方面，经上记载我们要敬畏祂，我们也奉命令爱慕祂，为使我们每个人都能使用适合其本性的处方。所以，当敬畏祂，因为祂是公义的；但无论你是敬畏祂还是爱慕祂，都不要犯罪。但愿任何敬畏祂的人能战胜不法的欲望，不贪图他人之物，施行仁慈，清醒，行义！因为我看见有些人在敬畏祂方面不完全，却犯下许多罪。因此，让我们敬畏天主，不仅因为祂是公义的；因为祂正是出于对受冤者的怜悯，才对行不义者施以惩罚。正如水熄灭火焰，敬畏也扑灭行恶的欲望。教导无畏的人自己并不敬畏；但不敬畏的人，不相信将有审判，反而增强自己的情欲，行邪术，并指控他人行他自己所行的事。}\par

% structure: p0026 source=h2 target=subsection reason=relative-heading-rank; base=h2

\subsection{第十三章 感官证据与超自然视觉之对比}

\Person{西满}听了这话，打断他，说：\Quote{我知道你这些话是针对谁说的；但为了不花时间讨论我不想讨论的问题，重复同样的话来驳斥你，请回答我们简洁陈述的问题。你声称你亲眼看见了你的老师的教义和作为，亲耳听见了，因此你充分理解了它们，并且不可能有其他任何人通过异象或显现获得类似的东西。但我要证明这是假的。一个人亲耳听见某人说话，并不能完全确信所说之事的真实性；因为他的心智必须考虑自己是否错误，因为那人表面上也是人。但显现不仅呈现对象在眼前，而且激发看见者的信心，因为它来自天主。现在，请先回答这个。}\par

% structure: p0028 source=h2 target=subsection reason=relative-heading-rank; base=h2

\subsection{第十四章 感官的证据比超自然视觉更可信}

\Person{伯多禄}说：\Quote{你提出要针对一个论点发言，却回答了另一个。因为你的主张是，一个人通过异象听见时，能更充分地认识\Emph{并获得信心}，比亲耳听见时更完全；但当你着手处理时，你却试图说服我们，通过异象听见的人比亲耳听见的人更确定。最后，你断言，正因如此，你比我更满意地知道了\Person{耶稣}的教义，因为你是通过异象听见祂的话的。但我要回答你最初提出的主张。先知，正因为他是先知，首先对客观所说的事提供确定的信息，就能被满怀信心地相信；并且，他预先被知道是真先知，按门徒所愿接受审问和询问，他就回答。但信任异象或显现和梦境的人是不可靠的。因为他不知道自己在信任谁。因为那可能是一个恶鬼，或一个欺骗之灵，在言语中假装成不是自己的样子。但如果有人想要询问显现者是谁，他可以说任何自己喜欢的话。就这样，像恶者一样闪现，停留任意长的时间，最终消失，并不留在询问者身边，如他所愿地供他咨询。因为靠梦境看见的人不能随意询问任何他想知道的事。因为思考并非睡着者的特权。因此，我们想在清醒时获得信息，却在梦中询问别的事；或者不经询问，听见与我们无关的事，并从睡梦中醒来，垂头丧气，因为我们既没有听见也没有询问我们渴望知道的事。}\par

% structure: p0030 source=h2 target=subsection reason=relative-heading-rank; base=h2

\subsection{第十五章 论来自梦境的证据}

西满说：\Quote{你若主张显现并非总是揭示真理，然而无论如何，来自天主的异象和梦境，在它们所要传达的事上，并不说假话。}伯多禄说：\Quote{你说得好：来自天主的，不说假话。但看见的人是否看见了来自天主的梦，这是不确定的。}西满说：\Quote{如果看见异象的人是义人，他就看见了真实的异象。}伯多禄说：\Quote{你说得对。但谁是义人呢？如果他需要异象才能学到该学的、做该做的？}西满说：\Quote{请允许我这一点：只有义人能看见真实的异象；然后我再回答另一个问题。因为我断定，不虔敬的人不会看见真实的梦。}伯多禄说：\Quote{这是假的；我可以不用圣经、也用圣经证明它；但我不想说服你。因为一个倾向爱上坏女人的男人，不会改变心意去关心与另一个各方面都好的女人合法结合；但有时他们因先入之见而爱较差的女子，尽管知道有另一个更优越的女子。由于某种类似的心态，你也是无知的。}西满说：\Quote{放下这个话题，讨论你答应要讲的事。因为在我看来，不虔敬的人绝不可能以任何方式从天主领受梦境。}\par

% structure: p0032 source=h2 target=subsection reason=relative-heading-rank; base=h2

\subsection{第十六章 只有恶鬼向不虔敬者显现}

伯多禄说：\Quote{我记得我答应要证明这一点，并从圣经和圣经之外提出证据。现在请听我说。我们知道有许多人——若你容我如此说；你若不容，我可以请在场的人作裁判——他们崇拜偶像、犯奸淫、各样犯罪，却仍看见真实的异象和梦境，其中一些还有魔鬼的显现。因为我主张，凡人的眼睛不能看见圣父或圣子的无形之形，因为它被极大的光照亮。因此，不是因为天主嫉妒，而是因为祂怜悯，才不能被已变成血肉的人看见。因为看见天主的人不能存活。因为光之过度会消融看见者的血肉；除非借着天主隐秘的能力，血肉变成光的本性，才能看见光；或光的本体变成血肉，才能被血肉看见。因为看见圣父而不经历任何改变的能力，唯独属于圣子。但义人也将同样看见天主；因为在死人复活时，当他们身体方面被改变成光，并变得像天使一样，他们将能看见祂。最后，如果有一位天使被派来，要被人看见，他就变成血肉，以便能被血肉看见。因为没有人能看见无形的能力，不仅圣子的不能，连天使的也不能。但若有人看见显现，他就该知道这是恶鬼的显现。}\par

% structure: p0034 source=h2 target=subsection reason=relative-heading-rank; base=h2

\subsection{第十七章 不虔敬者看见真实的梦与异象}

\Quote{但显然，不虔敬的人看见真实的异象和梦境，我能用圣经证明这一点。最后，律法上记着：亚比米勒如何不虔敬，想要与义人亚巴郎的妻子同房玷污她，以及他如何在睡中从天主那里听到命令，如圣经所说，不可碰她，因为她与丈夫同住。法郎也是不虔敬的人，看见了关于饱满和干瘦麦穗的梦，若瑟解释这梦来自天主。拿步高崇拜偶像，下令将崇拜天主的人投入火中，却看见一个涵盖整个世界时代的梦。谁也不能说：没有不虔敬的人在清醒时看见异象。那是假的。拿步高本人下令将三人投入火中，当他往火窑里看时，看见了第四位，并说：\InnerQuote{我看见第四位好像天主子。}然而，尽管他们看见显现、异象和梦境，他们仍是不虔敬的。因此，我们不能绝对地推断，看见异象、梦境和显现的人无疑就是虔诚的。因为对虔诚的人来说，真理自然而纯地涌现在他的理智中，不是通过梦境形成，而是通过理智赐予善人的。}\par

% structure: p0036 source=h2 target=subsection reason=relative-heading-rank; base=h2

\subsection{第十八章 启示的本质}

\Quote{同样，圣子也由圣父启示给我。因此我知道启示的含义，因我亲自体验过。正当主说：\InnerQuote{他们说我是谁？}我听到有人说这个，有人说那个，我心里便说出（我不知道怎么说的）：\InnerQuote{你是生活的天主之子。}祂称我有福，指出是父启示给我；从那时起，我知道启示是无须教导、无须显现和梦境的认知。事实正是如此。因为在天主安置在我们内的灵魂中，有一切真理；但真理被天主的手遮盖和揭示，天主的工作是照每人通过其知识所应得的。而借着外在的显现和梦境来宣告某事，这证明它不是来自启示，而是来自愤怒。最后，律法上记着：天主发怒，对亚郎和米黎盎说：\InnerQuote{如果你们中有一位先知，我将在异象和梦境中显示给他，但对我仆人梅瑟却不是这样；我要在外表上与他说话，而不是通过梦境，就像一个人与自己的朋友说话一样。}你看，愤怒的宣告是通过异象和梦境，而对朋友的宣告却是面对面、在外表上，而不是通过谜语、异象和梦境，如同对敌人一般。}\par

% structure: p0038 source=h2 target=subsection reason=relative-heading-rank; base=h2

\subsection{第十九章 反对伯多禄是不合理的}

\Quote{如果，那么，我们的\Person{耶稣}在异象中显现给你，使你认识祂，并向你说话，那是以愤怒对待仇敌的方式；正因如此，祂才通过异象、梦境或外在的启示向你说话。但人能通过显现而适于受教吗？如果你会说，\InnerQuote{这是可能的}，那么我问，\InnerQuote{为什么我们的导师与那些清醒的人同住并讲论整整一年？} 而且，当你告诉我们祂向你显现时，我们怎能相信你的话？当你持守与祂教导相悖的意见时，祂又如何向你显现？但倘若你曾见过祂，受祂教导，并在一小时之内成为祂的宗徒，那么就宣扬祂的话语，解释祂的言论，爱祂的宗徒，不要与我这曾与祂同行的人争辩。因为你现在站在与我，这坚如磐石、教会的基础直接对立的位置。如果你不反对我，你就不会控告我，并诽谤我所宣扬的真理，为的是不让人相信我从主亲耳所听到的话，仿佛我是显然被定罪和名声败坏的人。但如果你说我是被定罪的，你就是在控告那向我启示基督的天主，并且咒骂那因启示而称我有福的那一位。但是，如果你真的愿意为真理工作，首先要从我们这里学习我们从祂所学到的，并成为真理的门徒，与我们同工。}\par

% structure: p0040 source=h2 target=subsection reason=relative-heading-rank; base=h2

\subsection{第二十章 另一个讨论题目}

当\Person{西满}听了这话，就说：\Quote{我绝不会成为他或你的门徒。因为我不是不知道我所应当知道的；但我是作为学习者所提出的询问，是为了看你们能否证明实际的看见比显现更清晰。但你只是凭你自己的喜好说话；你没有证明。而且，明天我将探讨你关于天主的意见，你认为祂是世界的创造者；在我与你的讨论中，我将表明祂不是至高的，也不是善的，而且你们的导师也说过我现在所说的同样的话；我将证明你没有理解祂。} 说了这话他就走了，不愿意听他对他所提出的主张可能有的回应。\par

\SourceRecord{Translated by James Donaldson.；Ante-Nicene Fathers；Vol. 8.；Edited by Alexander Roberts, James Donaldson, and A. Cleveland Coxe.；Buffalo, NY: Christian Literature Publishing Co.,；1886.；<http://www.newadvent.org/fathers/080817.htm>.}

% structure: p0001 source=h1 target=section reason=page-title

\section{第十八篇讲道}

% structure: p0002 source=h2 target=subsection reason=relative-heading-rank; base=h2

\subsection{第一章 \Person{西蒙}主张世界的创造者不是至高的天主}

天刚亮时，\Person{伯多禄}出来讲论，\Person{西蒙}抢在他前头说：\Quote{昨天我离开时，答应你们今天回来，并在讨论中证明世界的创造者不是至高的天主，而是另一位至高的、独一的善者，祂至今仍不为人所知。那么，请立刻告诉我：你认为世界的创造者与立法者是同一位吗？如果祂是立法者，祂就是公义的；但如果祂是公义的，祂就不是善的。但如果祂不是善的，那么\Person{耶稣}所宣扬的就是另一位，当祂说：\InnerQuote{不要称我为善；因为只有一位是善的，就是那在天上的父。}然而，一位立法者不能既公义又善良，因为这两种品质并不协调。}然后\Person{伯多禄}说：\Quote{请先告诉我们，依你之见，哪些行为使人成为善人，哪些使人成为义人，这样我们才能对准同一个目标发言。}然后\Person{西蒙}说：\Quote{你先说说你认为什么是善，什么是义。}\par

% structure: p0004 source=h2 target=subsection reason=relative-heading-rank; base=h2

\subsection{第二章 善与义的定义}

\Person{伯多禄}说：\Quote{为了不浪费时间在争辩中，我公平地要求你回答我的命题，我却愿意如你所愿，回答自己提出的问题。那么我认为，慷慨施予\Emph{善物}的人是善的，正如我看见世界的创造者所做的那样，祂将太阳赐给善人，也将雨水赐给义人和不义的人。}然后\Person{西蒙}说：\Quote{这非常不公，祂将同样的东西赐给义人和不义之人。}\Person{伯多禄}说：\Quote{那么，请你反过来告诉我们，怎样的行为才使祂成为善的。}\Person{西蒙}说：\Quote{应当由你来陈述。}\Person{伯多禄}说：\Quote{我会的。按你的看法，将同样的东西赐给善人和义人，也赐给恶人和不义之人，祂甚至不是公义的；但如果你说祂将善赐给善人，将恶赐给恶人，那么你就有理由称祂为公义。那么，如果祂不采取将暂时之物赐给恶人（如果他们或许会悔改），并将永恒之物赐给善人（如果他们至少保持\Emph{善}），祂将采取怎样的行为呢？这样，通过赐给众人，但奖赏更优秀者，祂的公义是善的；并且更显忍耐的是，祂自愿赦免悔改者的罪，将永生赐予行善者。但在末后审判，按各人所应得的报应各人，祂是公义的。因此，如果这是正确的，就承认吧；但如果你认为不正确，就驳斥它。}\par

% structure: p0006 source=h2 target=subsection reason=relative-heading-rank; base=h2

\subsection{第三章 天主既是善的又是公义的}

\Person{西蒙}说：\Quote{我一次说过：\InnerQuote{每一位立法者，着眼于正义，就是公义的。}}\Person{伯多禄}说：\Quote{如果善者的本分不是设立法律，而是公义者的本分是设立法律，那么世界的创造者就既是善的又是公义的。祂是善的，因为很明显从\Person{亚当}到\Person{梅瑟}的时代，祂并没有设立成文的法律；但从\Person{梅瑟}到现在，祂有成文的法律，因此祂也是公义的。}\Person{西蒙}说：\Quote{请从你老师的言论中向我证明同一个人可以是善的又是公义的；因为在我看来，一位既是善的又是公义的立法者似乎是不可能的。}\Person{伯多禄}说：\Quote{我要向你解释善本身如何是公义的。我们的老师亲自先对那位问祂\InnerQuote{我该做什么才能承受永生？}的法利赛人说：\InnerQuote{不要称我为善；因为只有一位是善的，就是那天上的父。}然后立刻又说了这些话：\InnerQuote{但如果你愿意进入生命，就遵守诫命。}当法利赛人问\InnerQuote{什么诫命？}时，祂就指给他律法的诫命。如果祂是在指向另一位善者，祂就不会把他引向那公义者的诫命了。诚然，我承认公义和善是不同的，但你不明白同一位可以是善的又是公义的。因为祂是善的，在于如今祂对悔改者长久忍耐，并接纳他们；但祂是公义的，当祂作为审判者时，将按各人的行为报应各人。}\par

% structure: p0008 source=h2 target=subsection reason=relative-heading-rank; base=h2

\subsection{第四章 未启示的天主}

\Person{西蒙}说：\Quote{那么，如果世界的创造者（祂也塑造了\Person{亚当}）是已知的，并且也被那些按律法为义的人所知，甚至也被义人和不义之人以及全世界所知，你的老师却在所有这些之后说：\InnerQuote{除了圣子，没有人认识父，除了父，也没有人认识圣子，以及圣子愿意启示的人。}但祂不会说这句话，除非祂宣告了一位仍不为人知的父，律法称祂为至高者，祂从未发出任何或善或恶的言语（正如\Person{耶肋米亚}在\BibleBook{}{哀歌}中所见证的）；祂将列国限定为七十种语言，按照进入埃及的以色列子孙的数目，并按照这些民族的边界，将希伯来人赐给祂的圣子（祂也称为主，并安排天地），作为祂的产业，并将祂定义为万神之神，即那些接受其他民族为产业的诸神之神。因此，从所有所谓的神祇产生了给各自管辖区域的律法，这些区域包括其他民族。同样，也从万主之主的圣子那里产生了在希伯来人中建立的律法。这种状态被确定下来：如果任何人投靠某一位的律法，他就属于他所服从的律法的那一位的管辖。没有人知道那最高的、未启示的父，正如他们不知道圣子是祂的圣子一样。因此，此刻你自己将未启示的至高者的属性归给圣子，却不知道他是圣子，是\Person{耶稣}的父，而\Person{耶稣}在你们中间被称为基督。}\par

% structure: p0010 source=h2 target=subsection reason=relative-heading-rank; base=h2

\subsection{第五章 \Person{伯多禄}怀疑\Person{西蒙}的诚实}

当西满说了这些话后，伯多禄对他说：\Quote{你能作证这是你的信仰吗——祂自己，我不是指你现在所说不被启示的那一位，而是你相信却不肯承认的那一位？因为当你用一个东西定义另一个东西时，你是在胡说。因此，如果你呼求祂作证你相信你所说的话，我就会回答你。但如果你继续与我讨论你不相信的东西，你就迫使我在打空气。}西满说：\Quote{我是从你的一些门徒那里听说\Emph{这是真理}。}伯多禄说：\Quote{不要作假见证。}西满说：\Quote{不要责备我，最无礼的人。}伯多禄说：\Quote{只要你不说出是谁说的，\Emph{我就肯定}你在说谎。}西满说：\Quote{假设我自己编造了这些道理，或者我从别人那里听来，你给我回答。因为如果它们不能被推翻，那么我就知道这是真理。}伯多禄说：\Quote{如果是人的发明，我不回答；但如果你被这假设所束缚，认为它是真理，请你向我承认这一点，然后我自己可以就此事说点什么。}西满说：\Quote{那么，一言为定，这些道理在我看来是真实的。如果你有什么反驳的话，请给我你的回答。}\par

% structure: p0012 source=h2 target=subsection reason=relative-heading-rank; base=h2

\subsection{第六章 启示的本性}

伯多禄说：\Quote{如果是这样，你行事最不虔敬。因为如果安排天地属于子，由祂按自己的意愿将未被启示的父启示给任何人，那么正如我所说的，你行事最不虔敬，竟将父启示给祂没有启示的人。}西满说：\Quote{但祂自己希望我启示祂。}伯多禄说：\Quote{你不明白我的意思，西满。但请听并理解。当说子将启示给祂所愿之人时，意思是那人不是通过教导，而是唯独通过启示来认识祂。因为启示就是那秘密隐藏在人所有心中的东西，由祂\Emph{天主的}自己的意愿、不用任何言语而\Emph{公开}揭示。于是知识临到一个人，不是因为受教导，而是因为明白。但理解的人无法将其演示给别人，因为他自己不是通过教导领受的；也不能启示它，因为他自己不是子，除非他宣称自己是子。但你并不是常存的子。因为如果你是子，你必定知道那些配得如此启示的人。但你不知道他们。因为如果你知道他们，你行事就会像知道的人那样。}\par

% structure: p0014 source=h2 target=subsection reason=relative-heading-rank; base=h2

\subsection{第七章 西满承认自己的无知}

西满说：\Quote{我承认我没有理解你所说的\InnerQuote{你行事就会像知道的人那样}是什么意思。}伯多禄说：\Quote{如果你没有理解，那么你就不能知道每个人的心思；如果你对此无知，那么你就不知道那些配得启示的人。你不是子，因为子知道。因此祂将\Emph{祂}启示给所愿之人，因为他们配得。}西满说：\Quote{不要受骗。我知道那些配得的人，我也不是子。但我还没有理解你赋予\InnerQuote{祂将\Emph{祂}启示给所愿之人}这句话的含义。但我说我不理解，不是因为我不知道，而是因为我知道在场的人不理解，为的是让你更清楚地说出来，好让他们明白我们进行这场讨论的理由。}伯多禄说：\Quote{我不能更清楚地说出来：请解释你赋予这句话的含义。}西满说：\Quote{我没有必要陈述你的意见。}伯多禄说：\Quote{显然，西满，你不理解它，但你又不愿承认，以免被发现在无知中，从而被证明不是常存的子。因为你暗示了这一点，尽管你不愿直说；而且，我虽不是先知，却是真先知的弟子，从你的暗示中清楚知道你的愿望。因为尽管你连明确说出的话都不理解，却仍想自称子来反对我们。}西满说：\Quote{我要除去你的一切借口。我承认我不明白\InnerQuote{子将\Emph{祂}启示给所愿之人}这句话可能是什么意思。因此请更清楚地说出它的意思。}\par

% structure: p0016 source=h2 target=subsection reason=relative-heading-rank; base=h2

\subsection{第八章 启示的工作唯独属于子}

伯多禄说：\Quote{既然至少在表面上，你已承认你不理解它，请回答我所提的问题，你就会明白这句话的意思。告诉我，你是否主张子，无论祂是谁，是公义的，还是祂不公义？}西满说：\Quote{我主张祂至为公义。}伯多禄说：\Quote{既然祂是公义的，为什么祂不向所有人启示，而只向祂所愿的人？}西满说：\Quote{因为祂是公义的，只愿向配得的人启示。}伯多禄说：\Quote{那么，祂岂不知道每个人的心思，以便向配得的人启示？}西满说：\Quote{当然祂必须知道。}伯多禄说：\Quote{因此，施行启示的工作被唯独限于祂，是合理的，因为唯有祂知道每个人的心思；而它没有给你，因为你甚至不能理解我们对你说的话。}\par

% structure: p0018 source=h2 target=subsection reason=relative-heading-rank; base=h2

\subsection{第九章 西满如何承受他的揭露}

当\Person{伯多禄}说这话时，群众都称赞。但\Person{西满}这样被揭露，羞得面红耳赤，揉着额头说：\Quote{那么，他们宣称我这个行邪术的，甚至我这个能推理的人，竟被\Person{伯多禄}打败了吗？并非如此。但如果有人推理，即使被带走打败，他仍保留内心的真理。因为辩护者的软弱与被征服者所拥有的真理并不相同。但我向你们保证，我已判定所有在场的人都配得认识那未显明的圣父。因此，因为我公开地向他们启示他，你自己出于嫉妒，对我这个愿意造福他们的人发怒。}\par

% structure: p0020 source=h2 target=subsection reason=relative-heading-rank; base=h2

\subsection{第十章 \Person{伯多禄}答\Person{西满}}

\Person{伯多禄}说：\Quote{既然你这样说以取悦在场的群众，那么我就对他们说话，不是要取悦他们，而是要告诉他们真理。请告诉我，你怎么知道所有在场的人都配得？他们当中没有一个人赞同你对论题的阐述；因为你们向我鼓掌喝彩，这是反对你的，不是赞同你的人，而是赞同我的人，他们向我鼓掌是因为我说了真理。但是，既然公正的天主判断每个人的心思——这种教义你声称是真的——他就不会愿意通过左手把这赐给右手的人，正如有人从强盗那里接受任何东西，他自己就有罪。所以，因此，他不愿意他们接受你所带来的东西；但他们是藉着那为此工作而被分别出来的圣子接受启示。因为圣父将启示给谁才是合理的呢？不正是给祂的独生子吗？因为祂知道祂配得这样的启示。所以，这是一件不是人能教导或学习的事，而是必须由那不可言说的手启示给配得知道的人。}\par

% structure: p0022 source=h2 target=subsection reason=relative-heading-rank; base=h2

\subsection{第十一章 \Person{西满}声称说出真实想法}

\Person{西满}说：\Quote{如果作战的人使用自己的武器，对胜利很有帮助；因为一个人所爱的，他才能真诚地捍卫，而用真诚热情捍卫的事物，其中含有非凡的力量。因此今后我将向你们陈述我的真实意见。我坚持认为有一种未显明的能力，不为所有人所知，甚至造物主自己也不知道，正如\Person{耶稣}自己也曾宣称的那样，尽管他不知道自己在说什么。因为当人说很多话时，有时会碰巧说出真理，却不知自己所说的是什么。我指的是他所说的这句话：\BibleQuote{没有人认识父。}}\Person{伯多禄}说：\Quote{你不可再声称你知道他的教义。}\Person{西满}说：\Quote{我不自称相信他的教义；但我讨论的是他偶然说对的地方。}\Person{伯多禄}说：\Quote{为了不给你任何逃跑的借口，我将按照你所希望的方式与你讨论。同时，我请众人作证，你还不相信你刚才所说的话。因为我知道你的意见。为了不让你以为我没有说实话，我将阐述你的意见，好让你知道你在与一个熟悉这些意见的人讨论。}\par

% structure: p0024 source=h2 target=subsection reason=relative-heading-rank; base=h2

\subsection{第十二章 \Person{伯多禄}阐述\Person{西满}的意见}

\Quote{我们，\Person{西满}，并不主张从那伟大的能力（也称为统治的能力）中派出了两位天使，一位创造世界，另一位颁布法律；也不主张他们各自到来时，因自己所做的事而宣称自己是唯一的创造者；也不主张有一位站立者、将要站立者并受反对者。要知道你连在这个问题上也不相信。如果你说有一种未显明的能力，那能力就是充满无知的。因为它没有预知它所派出的天使的忘恩负义。}\Person{西满}因\Person{伯多禄}说这话而极其恼怒，打断他的谈话，说：\Quote{你说的是什么胡话，你这个胆大包天、最无耻的人，竟然在众人面前公然揭露秘密教义，让人轻易学会？}\Person{伯多禄}说：\Quote{你为什么吝啬不让在场的听众获益呢？}\Person{西满}说：\Quote{那么你认为这样的知识是有益的吗？}\Person{伯多禄}说：\Quote{我认为是有益的：因为认识错误的教义是有益的，这样你就不会因无知而堕入其中。}\Person{西满}说：\Quote{你显然不能回答我先前提出的命题。我坚持认为，连你的老师也肯定有一位未显明的圣父。}\par

% structure: p0026 source=h2 target=subsection reason=relative-heading-rank; base=h2

\subsection{第十三章 \Person{伯多禄}对经文的解释}

\Person{伯多禄}说：\Quote{我要答复你希望我讲的那段经文，即：\BibleQuote{除了子，没有人认识父；除了父，也没有人认识子；只有子愿意启示给谁，才认识祂。}首先，我感到惊奇的是，虽然这句话有无数种解释，你却选择了那个非常危险的立场，认为这句话是针对造物主及其属下一切人的无知而说的。因为，首先，这句话可以适用于所有犹太人，他们以为\Person{达味}是基督的父，而基督自己是\Person{达味}之子，却不知道祂是天主子。因此，说\BibleQuote{没有人认识父}是合宜的，因为他们不承认祂是天主，反而声称\Person{达味}是祂的父；而附加的评论说也没有人认识圣子，也是完全正确的，因为他们不知道祂是圣子。至于\BibleQuote{子愿意启示给谁}，也正确；因为祂从起初就是圣子，唯独被派定将启示赐给祂愿意赐给的人。这样，第一个人\Person{亚当}必定曾听说过祂；蒙\Emph{天主}喜悦的\Person{哈诺客}必定认识祂；义人\Person{诺厄}必定熟悉祂；祂的朋友\Person{亚巴郎}必定理解祂；\Person{依撒格}必定感受到祂；与祂角力的\Person{雅各伯}必定相信祂；而所有在百姓中配得的人都领受了启示。}\par

% structure: p0028 source=h2 target=subsection reason=relative-heading-rank; base=h2

\subsection{第十四章 \Person{西满}被驳倒}

\Quote{但是，如果你说，由于祂如今通过耶稣向所有人启示，因而能够认识祂，那么你岂不是在说一件最不义的事吗？因为你说，这些不认识祂的人，就是世上的七根支柱，并且能够取悦最公义的天主，而如今这么多来自万邦的不虔敬之人，却完全认识祂？难道那些超越一切的人，不配认识祂吗？那不善的怎么会是公义的呢？除非你愿意把\Quote{善}的名号，不给那善待行义者的人，而是给那爱不义者的人，即使他们不信，也向他们启示祂不愿启示给义人的奥秘。但这样的行为，既不配一个善者，也不配一个义者，而是配一个已经仇恨虔敬者的人。西门，你这站立者，岂不正是你胆敢说出这些前所未有的话吗？}\par

% structure: p0030 source=h2 target=subsection reason=relative-heading-rank; base=h2

\subsection{第十五章 论玛窦福音11:25}

西门对此感到恼怒，说：\Quote{责备你的老师吧，祂曾说：\InnerQuote{天地的主，我感谢祢，因为祢将这些事瞒住了智慧人，而启示给了吃奶的婴孩。}}伯多禄说：\Quote{这话不是那样说的；但我姑且按你所喜欢的方式来谈它。我们的主，即使祂说过\InnerQuote{智慧人被瞒住的事，父却启示给了婴孩}，也绝不能因此被认为指出了另一位不同于创造世界者的天主和父。因为祂所说的那些隐藏的事，可能是造物主（\OriginalTerm{德穆革}{Demiurge}）自己的事；因为依撒意亚说：\BibleQuote{我要开口说比喻，我要说出从创世以来隐藏的事。}那么，你承认先知并不无知于那些隐藏的事吗？而耶稣却说这些事对智慧人是隐藏的，却启示给了婴孩。如果祂的先知依撒意亚并不无知于它们，造物主（\OriginalTerm{德穆革}{Demiurge}）又怎能无知呢？但是，我们的耶稣实际上并没有说\Quote{隐藏的事}，而是说了一句听起来更严厉的话；因为祂说：\Quote{祢将这些事瞒住了智慧人，却启示给了吃奶的婴孩。}现在，\Quote{祢瞒住了}这个词暗示他们曾经知道这些事；因为天国的钥匙，即奥秘的知识，曾在他们那里。}\par

% structure: p0032 source=h2 target=subsection reason=relative-heading-rank; base=h2

\subsection{第十六章 这些事被公义地隐藏于智慧人}

\Quote{不要说祂对智慧人行了不虔敬的事，把这些事隐藏起来。这种想法远非我们所能接受。因为祂没有行不虔敬的事；而是因为他们隐藏了天国的知识，自己既不进去，也不让那些愿意进去的人进去，因此，公义地，正如他们将道路向那些愿意的人隐藏一样，奥秘也同样向他们隐藏，好使他们自己体验到他们对别人所做的事，并且他们用什么尺度量给人，也用什么尺度量给他们。因为对于配得认识的人，他所不知道的应当给予他；但对于不配的人，即使他似乎在别的事上有任何东西，也要被夺去，即使他在其他事上是智慧的；并且要给予那配得的人，即使他们在作门徒的时间上是婴孩。}\par

% structure: p0034 source=h2 target=subsection reason=relative-heading-rank; base=h2

\subsection{第十七章 通往天国的道路并未向以色列子民隐藏}

\Quote{但若有人说没有什么是向以色列子民隐藏的，因为经上记载：\BibleQuote{以色列啊，没有一件事逃过你的眼目（因为不要说，雅各伯啊，道路向我隐藏了），}他应当明白，属于天国的事曾向他们隐藏，但通往天国的道路，即生活方式，却没有向他们隐藏。为此祂说：\Quote{因为你们不要说，道路向我隐藏了。}但所谓\Quote{道路}指的是生活方式；因为梅瑟说：\BibleQuote{看，我把生命之路和死亡之路摆在你们面前。}教师也与之一致地说：\BibleQuote{你们要从窄而狭的门进去，通过它你们将进入生命。}另一次，有人问祂：\BibleQuote{我该做什么才能承受永生？}祂向他指出了法律的诫命。}\par

% structure: p0036 source=h2 target=subsection reason=relative-heading-rank; base=h2

\subsection{第十八章 论依撒意亚先知书1:3}

\Quote{从依撒意亚以天主的位格所说的话：\BibleQuote{但是以色列不认识我，我的百姓不明白我，}不能推断出依撒意亚指出了另一位不同于那被认识的天主；而是他意指那被认识的天主在另一种意义上不被认识，因为百姓犯罪，对那被认识的天主的公义属性无知，并且幻想他们不会受到善天主的惩罚。因此，在他说了\BibleQuote{但是以色列不认识我，我的百姓不明白我}之后，他接着说：\BibleQuote{祸哉！犯罪的民族，罪恶深重的百姓。}因为他们因无知于祂的公义而不惧怕，正如我所说，他们便罪恶深重，以为祂仅仅是善的，因此不会因他们的罪而惩罚他们。}\par

% structure: p0038 source=h2 target=subsection reason=relative-heading-rank; base=h2

\subsection{第十九章 旧约中对天主的误解}

\Quote{有些人如此犯罪，是因为因祂的善而想象不会有审判。但另一些人则走了相反的路。因为他们以为圣经中那些反对天主、不义而虚假的话是真实的，所以他们不认识祂真正的神性和权能。因此，他们相信祂无知、喜悦谋杀、因祭献的礼物而释放恶人；不仅如此，他们还相信祂欺骗、说谎、行一切不义的事，他们自己就做类似于他们的天主所做的事，并以此犯罪，却声称自己行事虔诚。所以，他们不可能改变为更好，并且受警告也不留意。因为他们不惧怕，因为通过这样的行为，他们变得像他们的天主。}\par

% structure: p0040 source=h2 target=subsection reason=relative-heading-rank; base=h2

\subsection{第二十章 旧约中的某些部分是为考验我们而写的}

\Quote{但有人可能有充分理由主张，这句话是针对那些认为祂是这样的人说的：\BibleQuote{除了圣子，没有人认识圣父，正如除了圣父，没有人认识圣子一样。}这是合理的。因为如果他们认识，就不会因信赖那些针对天主而写的、其实是为了考验我们而写的书籍而犯罪。但是，为了更清楚地向他们显明他们错误的缘由，祂也在某处说：\BibleQuote{因此你们错了，因为你们不知道圣经的真义，所以你们也不认识天主的德能。}因此，凡是愿意得救的人，都必须像导师所说的那样，成为考验我们的书籍的审判官。因为祂这样说：\Quote{成为经验丰富的钱庄主。}现在需要钱庄主，是因为赝品与真品混在一起。}\par

% structure: p0042 source=h2 target=subsection reason=relative-heading-rank; base=h2

\subsection{第二十一章 西满对伯多禄对待圣经的方式感到惊讶}

\Quote{但愿这事远离我，以及爱我的人，就是听你的讲论。事实上，只要我不知道你对圣经持有这些见解，我就容忍你，与你讨论；但现在我退出了。真的，我一开始就该退出，因为我听你说：\InnerQuote{我本人不相信任何人说任何反对创造世界的天主的话，不管是天使、先知、圣经、司铎、导师，还是任何其他人，即使他行出标记和奇迹，即使他在空中发出灿烂光芒，或透过神视或梦境做出启示。}那么，无论是好是坏，谁能改变你的心意，使你持有与你的决定不同的意见，因为你如此固执地坚守自己的决定？}\par

% structure: p0044 source=h2 target=subsection reason=relative-heading-rank; base=h2

\subsection{第二十二章 伯多禄敬拜独一的天主}

当西满这样说并准备离开时，伯多禄说：\Quote{再听这一句话，然后你随便去哪里。}于是西满转身留下，伯多禄说：\Quote{我知道你当时听见我说\InnerQuote{无论谁说什么反对创造世界的天主，我都不相信他。}时何等惊讶。但现在请听一些更进一步的、更大的事。如果创造世界的天主实际上正如圣经所描述的性格，并且如果不知为何祂是无比邪恶的，比圣经所能描绘的或任何其他人所能想象的更邪恶，我仍然不会放弃单独敬拜祂并遵行祂的旨意。因为我希望你明白并确信，那对自己的创造者没有情感的人，永远不能对另一个人有情感。如果他对另一个人有情感，那是违反本性的，并且他不知道这种对不义者的热情来自那恶者。他甚至不能持久地保持它。而且，如果有一位在创造者（\OriginalTerm{德缪哥}{Demiurge}）之上，祂既然是善的，会更加欢迎我，因为我爱自己的圣父；祂不会欢迎你，因为祂知道你抛弃了自己的天然创造者：因为我称祂为圣父，并非受更大希望的影响，也不关心合理的事。因此，即使你找到一个比祂更高超者，祂也知道你终有一天会抛弃祂；尤其因为祂不是你的圣父，既然你抛弃了那位真正是你的圣父的那位。}\par

% structure: p0046 source=h2 target=subsection reason=relative-heading-rank; base=h2

\subsection{第二十三章 西满离开}

\Quote{但你会说：\InnerQuote{祂知道没有别的天主在祂之上，因此祂不能被抛弃。}多亏没有别的天主；但祂知道你的心态倾向于忘恩负义。但如果祂明知你忘恩负义仍接纳你，而明知我感恩却不接纳我，那么按照你自己的说法，祂是不体谅的，行为不合理。因此，西满，你不晓得你是邪恶的仆役。}西满回答说：\Quote{那么，邪恶从何而来？请告诉我们。}伯多禄说：\Quote{既然今天你是第一个离开，并且声明你今后不会听我这个亵渎者的话，那么明天如果你确实愿意学习，我将向你解释这件事，并允许你随意问我任何问题，不作任何争辩。}西满说：\Quote{我将按我看为好的去做。}说了这话，他就走了。然而，那些与他一同进来的人没有一人与他一同出去；反而俯伏在伯多禄脚前，恳求他们因被西满迷惑而得到宽恕，并后悔后得到接纳。但伯多禄接纳了那些悔改的人以及其余的群众，按手在他们身上，祈祷，并医治他们中的病人；然后遣散他们，劝他们黎明时早早回来。说了这话，他与密友进去，为立即休息做了通常的准备，因为已是傍晚。\par

\SourceRecord{Translated by James Donaldson.；Ante-Nicene Fathers；Vol. 8.；Edited by Alexander Roberts, James Donaldson, and A. Cleveland Coxe.；Buffalo, NY: Christian Literature Publishing Co.,；1886.；<http://www.newadvent.org/fathers/080818.htm>.}

% structure: p0001 source=h1 target=section reason=page-title

\section{讲道集第十九篇}

% structure: p0002 source=h2 target=subsection reason=relative-heading-rank; base=h2

\subsection{第一章 \Person{西满}着手证明世界的创造者并非无可指责}

第二天，伯多禄比平时出来得更早。他看见西满和许多其他人在等他，就向群众致意，开始讲道。但他刚开始讲，西满就打断了他，说：\Quote{请跳过您这些冗长的开场白，直接回答我向您提出的问题。因为我从开头听到的话中看出，您除了用尽一切办法表明创造者自己是惟一无可指责的天主之外，没有其他目的。——既然，正如我所说，我看出您如此坚决地要维护这一点，以至于您胆敢宣称圣经中那些明显反对他的部分为虚假，因此我今天决定证明，他既是万物的创造者，就不可能无可指责。但这一证明我现在就可以开始，只要您回答我提出的问题。}\par

% structure: p0004 source=h2 target=subsection reason=relative-heading-rank; base=h2

\subsection{第二章 魔鬼的存在得到肯定}

\Quote{您认为有魔鬼吗，还是没有？因为如果您说没有，我可以从许多陈述中，特别是从您导师的陈述中，向您证明有；但如果您诚实地承认魔鬼存在，那么我将按照这一信念说话。} 伯多禄说：\Quote{我无法否认我导师的断言。因此我承认魔鬼存在，因为我的导师，他在一切事上都说了真理，曾多次肯定他存在。例如，他承认曾与他交谈，并被他试探了四十天。而且我知道他在别处说过：\BibleQuote{若撒殚驱逐若撒殚，他是自相纷争，他的国怎能站得住呢？} 他还指出，他看见魔鬼像闪电一样从天上坠落。另外，他在别处说：\BibleQuote{那撒坏种子的仇人是魔鬼。} 又说：\BibleQuote{不可给魔鬼留下余地。} 此外，在劝告时，他说：\BibleQuote{你们的话，是就说是，不是就说不是；多余的话，便是出于魔鬼。} 还有，在他教导我们的祈祷中，我们读到：\BibleQuote{救我们免于凶恶。} 在另一处，他许诺要对那些不虔敬的人说：\BibleQuote{你们到外面黑暗里去，那原是圣父为魔鬼和他的使者预备的。} 我不再赘述，我知道我的导师常说有魔鬼。因此我也同意认为他存在。所以，如果您今后要按这一信念说些什么，就请说吧，如您所承诺的。}\par

% structure: p0006 source=h2 target=subsection reason=relative-heading-rank; base=h2

\subsection{第三章 伯多禄拒绝讨论关于魔鬼的某些问题}

西满说：\Quote{既然您凭圣经的见证诚实地承认魔鬼存在，那么请告诉我们，他是如何产生的——如果他是产生的——被谁产生，以及为什么。} 伯多禄说：\Quote{请您原谅，西满，如果我不敢肯定没有记载的事。但如果您说这已被记载，请证明它。但既然没有记载，您无法证明，我们为什么要冒险就未记载的事发表自己的见解呢？因为如果我们过于大胆地谈论天主，要么是我们不相信自己将受审判，要么是我们相信只凭行为受审判，而不凭信仰和言语受审判。} 但西满明白伯多禄指的是他自己的狂妄，便说：\Quote{请允许我冒这个险；但您不要以您所断定为亵渎的话为借口而退缩。因为我察觉您想退避，以免在大众面前被驳倒，有时假装害怕听亵渎的话，有时又坚持说，既然没有记载魔鬼如何、被谁、为何产生，我们就不应胆敢多于圣经断言。因此，您作为一个虔诚的人，只肯定他存在。但通过这些计谋，您欺骗了自己，不知道如果仔细探究魔鬼是亵渎，那么罪责在于我这个指控者，而不在于您这位天主的辩护者。如果所探究的事不在圣经中，因此您不愿探究，那么有些令人满意的方法也能像圣经一样有效地向您证明所寻求的事。例如，难道不是必然的：您所肯定的魔鬼要么是有起源的，要么是无起源的吗？}\par

% structure: p0008 source=h2 target=subsection reason=relative-heading-rank; base=h2

\subsection{第四章 关于魔鬼起源的假设}

伯多禄说：\Quote{必然如此。} 西满说：\Quote{因此，如果他有起源，那么他正是由那位创造万物的天主所造，要么作为动物而生，要么是实质性地被发出，并且是由外部元素的混合而产生。因为造他的物质，无论是有生命的还是无生命的，要么在他之外，要么他是通过天主本身、或通过自身、或从无中产生、或仅是相对之物、或永恒存在。我想我已经清楚地指出了所有可能的途径，沿着其中一条，我们必定能找到他。因此我们必须沿着每一条途径寻找他的起源；当我们找到他的作者时，我们必须认识到他应受责备。或者您怎么看？}\par

% structure: p0010 source=h2 target=subsection reason=relative-heading-rank; base=h2

\subsection{第五章 天主允许魔鬼存在并不应受责备}

伯多禄说：\Quote{依我之见，即便是明显他由天主所造，那造他的造物主也不应被指责；因为或许会发现他所做的事是完全必要的。但另一方面，若证明他并非受造，而是永远存在，那么在此事上造物主也不应被指责，因为祂比所有\Emph{其他一切}更好，即使祂没有能力终结一位无始的存在——因其本性不容如此；或者，若祂有能力却未除灭他，认为终结一位没有开端的存在是不义的，并宽恕那本性邪恶者，因为他即使愿意也无法变成别的样子。但如果愿意行善却无能为力，在此情形下祂仍是善的，因祂有善意，虽无能力；而无能力时，祂依然是最有能力的，因为能力并未留给另一位。但如果另有其人有能力却未完成，则必须承认，他既有能力却不完成，他就是邪恶的，没有除灭他，仿佛以他所行的事为乐。但若连他也无能为力，那么那虽不能却愿意尽力施恩于我们的，便是更好的。}\par

% structure: p0012 source=h2 target=subsection reason=relative-heading-rank; base=h2

\subsection{第六章 \Person{伯多禄}指责\Person{西满}比魔鬼更坏}

西满说：\Quote{等你讨论完我提出的一切主题，我就要给你指出邪恶的原因。然后我也要回答你刚才所说的话，并证明你所断言无可指责的那位天主是应该受指责的。} 伯多禄说：\Quote{既然我从你一开始所说的话中看出，你无非是想要把天主当作邪恶的制造者而加以指责，我就决心随你走你喜欢的任何道路，并证明天主完全无可指责。} 西满说：\Quote{你这样说是因为你爱你所认为认识的天主；但你并不正确。} 伯多禄说：\Quote{而你，作为邪恶的人，憎恨你所不认识的天主，说出亵渎的话。} 西满说：\Quote{你要记住，你曾把我比作邪恶的制造者。} 伯多禄说：\Quote{我承认，我把你比作那邪恶者是我错了；因我不得不这样做，因为我找不到一个与你相等或比你更坏的人。因此我把你比作那邪恶者；因为你碰巧比那邪恶的制造者还要坏得多。因为没有人能证明那邪恶者曾说过反对天主的话；但在场的我们众人都看见你胆敢反对祂。} 西满说：\Quote{寻求真理的人不应当在任何方面迎合任何人而有违真实。因为他为何要查问呢？我问，为何？因为我也不能放弃对事物的精确探究，把全部时间用来赞美我所不认识的那位天主。}\par

% structure: p0014 source=h2 target=subsection reason=relative-heading-rank; base=h2

\subsection{第七章 \Person{伯多禄}怀疑\Person{西满}甚至不相信有天主}

伯多禄说：\Quote{你并非有福到能赞美祂，你甚至不能行这样的善事；因为那样你就会充满祂。因为我们的师傅——祂总是说真话——这样说过：\BibleQuote{心里充满什么，口里就说什么。}因此，你充满邪恶的意念，因无知而反对那唯一美善的天主。你尚未为你敢于说出的这些话而遭受应得的苦难，就或者以为不会有审判，或者竟以为连天主也没有。因此，你不明白祂的如此忍耐，竟越来越疯狂。} 西满说：\Quote{不要以为你能吓倒我，使我不敢探究你那些例子的真实性。因为我如此渴慕真理，为了真理我不怕冒险。如果你对我一开始所提的命题有什么话要说，现在就说吧。}\par

% structure: p0016 source=h2 target=subsection reason=relative-heading-rank; base=h2

\subsection{第八章 \Person{伯多禄}着手讨论魔鬼的起源}

伯多禄说：\Quote{既然你强迫我们在精确探究天主的安排之后，斗胆陈述它们，而且是向那些不能完全理解自己同胞安排的人陈述，那么至少为了在场的人，我不保持沉默——这本来是最虔诚的做法——而要讨论你希望我谈论的主题。我同意你的看法，相信有一位邪恶之君，关于他的起源，圣经既未敢说真话，也未敢说假话。但让我们以多种方式追问他如何存在——如果事实上他确实存在；并且在呈现的观点中，让我们挑选那最恭敬的一种，因为在可能成立的意见中，那\Emph{基于一个原则}——即我们应当将更恭敬之事归于天主——的意见会被自信地采纳；而且尤其如此，如果当所有其他假设被排除后，仍有一个充分且风险较小的假设。但我向你保证，在我继续探究之前，探究中的每一种方法都能表明唯有天主无可指责。}\par

% structure: p0018 source=h2 target=subsection reason=relative-heading-rank; base=h2

\subsection{第九章 关于魔鬼起源的各种理论}

\Quote{但是，正如你所说，如果恶者是被造的，那么他要么是被生为动物，要么是祂实质地发出的，要么是外面合成的，要么他的意志是通过组合产生的；要么他是从不存在的东西中、没有组合和天主的意志而存在的；要么他是天主从完全不存在的地方造的；要么质料（无生命的或有生命的）在他之外，他从中产生；要么他塑造了自己，要么他是天主造的，要么他是一个相对的东西，要么他永远存在：因为我们不能说他不存在，因为我们已同意认为他存在。} 并且\Person{西满}说：\Quote{你很好地概括了所有解释其存在的方法。现在我的任务是考察这些不同的观点，并表明造物主是可指责的。但你的职责是证明，如你所承诺的，祂完全无可指责。但我怀疑你是否能够做到。因为，首先，如果魔鬼是从天主生出的，如同动物，那么他的恶与发出他的那一位相同。} 并且\Person{伯多禄}说：\Quote{并非如此。因为我们看见许多良善的人是邪恶孩子的父亲，而另一些邪恶的人是良善孩子的父亲，还有另一些邪恶的人既产生良善的孩子，也产生\Emph{邪恶的}孩子，还有一些良善的人既有邪恶也有良善的孩子。例如，首先被造的人产生了不义的加音和义人亚伯尔。} 对此\Person{西满}说：\Quote{你在谈论天主时使用人类的例子，这是愚蠢的。} 并且\Person{伯多禄}说：\Quote{那么，请你向我们谈论天主而不使用人类的例子，但同时又让你所说的能被理解；但你做不到。}\par

% structure: p0020 source=h2 target=subsection reason=relative-heading-rank; base=h2

\subsection{第十章 绝对天主完全无法为人理解}

\Quote{例如，那么，你在开始时说了什么？如果恶者是从天主所生，与祂同实体，那么天主就是邪恶的。但是当我从你亲自引用的例子中向你表明，邪恶来自良善，良善来自邪恶时，你并不接受这个论证，因为你说这个例子是人类的。因此，我现在不接受\Quote{受生}一词可用于天主；因为生育是人的特性，而不是天主的特性。不仅如此；天主也不能是善或恶，正义或不义。祂甚至不能拥有理智、生命或任何其他存在于人中的属性；因为所有这些都属于人。如果我们对天主的研究中，不能将属于人的善属性归给祂，那么我们就不能对天主有任何思想或任何陈述；我们所能做的就是研究一点——即，祂的意志是什么，祂亲自允许我们领悟的意志，以便我们在审判时，对于那些我们明知却未遵守的法律，无可推诿。}\par

% structure: p0022 source=h2 target=subsection reason=relative-heading-rank; base=h2

\subsection{第十一章 人的属性应用于天主}

\Person{西满}听到这话，说：\Quote{你不会通过羞愧迫使我保持沉默关于祂的实体，而仅仅探究祂的意志。因为有可能思考并谈论祂的实体。我是说，从属于人的善属性。例如，生命和死亡是属于人的；但死亡不是天主的属性，而是生命，永生。此外，人可以是邪恶的，也可以是良善的；但天主只能是无可比拟的善。而且，不必过分延长这个话题，人的较好属性是天主的永恒属性。} 并且\Person{伯多禄}说：\Quote{告诉我，\Person{西满}，生邪恶和良善，以及行邪恶和良善，是人的属性吗？} \Person{西满}说：\Quote{是的。} \Person{伯多禄}说：\Quote{既然你做出这个断言，我们必须将人的较好属性归给天主；因此，当人生育邪恶和良善时，天主只能生育良善；当人行邪恶和良善时，天主只喜乐于行善。这样，\Emph{对于天主}，我们要么不主张任何人的属性并保持沉默，要么合理地应将最好的善属性归给祂。因此，祂是一切善的唯一原因。}\par

% structure: p0024 source=h2 target=subsection reason=relative-heading-rank; base=h2

\subsection{第十二章 天主产生了恶者，但并非邪恶}

\Person{西满}说：\Quote{如果，那么，天主只是善的原因，我们还能想什么，除了某些其他本原产生了恶者；或者邪恶是非受生的？} 并且\Person{伯多禄}说：\Quote{没有其他权能产生了恶者，邪恶也不是非受生的，正如我将在结论中表明；因为现在我的目标是证明，如我在开头所承诺的，天主在各方面都是无可指责的。我们已经承认，天主以无可比拟的方式拥有属于人的较好属性。因此，祂也有可能成为四种物质——我指的是热、冷、湿、干——的生产者。这些物质起初是单纯未混合的，自然在欲望上是中性的；但它们由天主产生，并在外部混合，自然会成为一个活的存在，拥有自由抉择去毁灭那些邪恶的人。因此，既然万有都是从祂受生的，恶者并非来自其他根源。他也没有从创造万有的天主（在他内不可能存在邪恶）那里获得他的邪恶，因为这些物质是由祂以中性状态产生的，并仔细地彼此分离；当它们通过他的技艺在外部混合时，通过意志产生了毁灭那些邪恶的人的欲望。但良善不能被产生的邪恶所毁灭，即使它想要这样做：因为它只对犯罪的人行使它的能力。因此，由于对每位的品性无知，他试图攻击他，并定罪他，然后惩罚他。} \Person{西满}说：\Quote{天主能够混合元素，并按照祂所希望的任何倾向进行混合，为什么祂不使每个组合都具有偏好良善的倾向？}\par

% structure: p0026 source=h2 target=subsection reason=relative-heading-rank; base=h2

\subsection{第十三章 天主是魔鬼的制造者}

伯多禄说：\Quote{如今我们的目的确实是要表明恶者是如何、由谁产生的，既然他确实产生了；但当我们完成眼前这个论题之后，我们将表明他是否无辜地被产生。然后我将说明他是如何、因何产生的，并且我将充分说服你，他的创造者是无辜的。我们说过，四种基质是由天主产生的。这样，通过混合它们的那位的意愿，按照祂的愿望，恶的选择就产生了。因为如果它是违背祂的决定、或者从另一种基质或原因产生的，那么天主的意志就不坚固了：因为也许即使祂不愿意，恶的首领也会不断出现，与祂的意愿作战。但这是不可能的。因为没有生物，尤其是能提供引导的生物，能出于偶然产生：因为一切被产生之物必须由某位产生。}\par

% structure: p0028 source=h2 target=subsection reason=relative-heading-rank; base=h2

\subsection{第十四章 物质是永恒的吗？}

西满说：\Quote{但如果物质与祂同样永恒、拥有同等能力，作为祂的敌人产生阻碍祂意愿的首领呢？}伯多禄说：\Quote{如果物质是永恒的，那它就不是任何人的敌人：因为永恒存在之物是不受苦难的，不受苦难的是有福的，而有福的不能容纳仇恨，因为它因着永恒的被造，不惧怕失去任何东西。但物质为何不是更爱创造者，当它显然发出它的果实来滋养所有由祂所造之物时？它为何不畏惧祂为更优越者，如同通过地震颤抖而承认，又如当师傅乘船其上并命令平静时，它立刻服从并变得安静呢？什么！难道邪魔不是因对祂的恐惧和尊敬而出去，其中一些还渴望进入猪群吗？但它们先恳求祂，然后才出去，显然因为没有祂的许可，它们甚至没有能力进入猪群？}\par

% structure: p0030 source=h2 target=subsection reason=relative-heading-rank; base=h2

\subsection{第十五章 罪是恶的原因}

西满说：\Quote{但如果它没有生命，却具有产生善恶的本性呢？}伯多禄说：\Quote{根据这个说法，它既非善也非恶，因为它不是凭自由选择行动，没有生命且无知觉。因此在这件事上可以清楚地察觉到，它如何没有生命却如同有生命一样产生，无知觉却显然在动物和植物中塑造出艺术的形状。}西满说：\Quote{什么！如果天主亲自赋予它生命，那么祂不就成了它所产生之恶的原因吗？}伯多禄说：\Quote{如果天主按祂自己的意愿赋予它生命，那么是祂的圣神在产生它，它就不再是与天主敌对或同等力量的；或者说，凡由祂所造之物不可能不符合祂的愿望。但你会说，祂自己是邪恶的原因，因为祂通过它亲自产生恶。那么，你说的恶是哪种呢？毒蛇、致命植物、邪魔，还是任何其他能扰乱人的东西？——这些东西若不是人犯了罪，就不会有害，为此死亡进入了。因为如果人无罪，蛇的毒就不会起作用，有害植物的活性也不会，邪魔的扰乱也不会，人也不会自然有其它任何苦难；而是因着罪失去不死性，如我所说，人变得能承受一切苦难。但如果你说，为何人的本性起初被造为可死的？我告诉你，因为自由意志；因为如果我们不能死，作为不死之物，我们就不能因自愿的罪受惩罚。这样，因着我们不受苦难，义德会更加削弱，如果我们是出于选择而邪恶；因为那些怀有恶意目的的人，因着不能受苦而无法受惩罚。}\par

% structure: p0032 source=h2 target=subsection reason=relative-heading-rank; base=h2

\subsection{第十六章 为何恶者被赋予权能}

西满对此说：\Quote{关于恶者，我还有一点要说。无疑，既然天主从无中造了他，他在这方面是恶的，尤其因为天主本可造他为善，在造他时赋予他绝不能选择邪恶的本性。}伯多禄说：\Quote{说天主从无中造了他、赋予他选择的能力，这就像我们上面所说的：制造了这样一种以恶为乐的性情，祂自己似乎就成了所发生之事的原因。但既然两种说法有一个解释，我们以后会说明为何祂使他以恶人的毁灭为乐。}西满说：\Quote{如果祂也使天使成为自愿行动者，而恶者离开了义的状态，为何他被赋予统领的职位？难道不是显然，那这样荣耀他的人以恶人为乐，因为他这样荣耀了他？}伯多禄说：\Quote{如果天主在他反叛时按律法设立他统治那些像他的人，命令他对犯罪者施加惩罚，祂并非不义。但如果情况是，即使在他反叛之后祂仍荣耀了他，那么荣耀他的那一位预见了他的用处；因为这荣耀是暂时的，并且恶者被恶者统治、罪人被他惩罚，是正当的。}\par

% structure: p0034 source=h2 target=subsection reason=relative-heading-rank; base=h2

\subsection{第十七章 魔鬼没有与天主同等的权能}

\Person{西满}说：\Quote{如果祂永远存在，那唯一统治权\Emph{天主}的事实岂不是被摧毁了，因为另有一种能力，即与物质相关的能力，与祂一同统治？} \Person{伯多禄}说：\Quote{如果它们在本质上不同，它们的能力也不同，上级统治下级。但如果它们本质相同，那么它们能力相等，同样地善或恶。但显然它们能力不等；因为造物主将物质塑造成祂愿意塑造成的世界的形态。那么，是否可能主张它作为实体永远存在？物质不就像是天主的仓库吗？因为不可能主张曾有一个时期天主一无所有，而祂一直是它的唯一统治者。因此，祂是永恒的独一统治者；正因如此，可以公正地说，属于那存在、统治且是\Emph{永恒的}一位。} \Person{西满}说：\Quote{那么怎样？那邪恶者造了自己吗？而天主是如此良善，明知他会成为恶的原因，却在他起源时没有毁灭他，那时他尚未完全形成，本可以被毁灭？因为如果他突然而完全地出现，那么他因此便与造物主开战，作为突然出现且拥有同等能力者。}\par

% structure: p0036 source=h2 target=subsection reason=relative-heading-rank; base=h2

\subsection{第十八章　魔鬼是一种关系吗？}

\Person{伯多禄}说：\Quote{你所说的不可能；因为如果他逐渐形成，祂本可以凭自己的自由选择将他像敌人一样剪除。而且事先知道他会形成，祂就不会允许他作为善存在，除非祂知道因他而产生有益之物。他也不可能突然、完整地凭借自己的能力出现。因为不存在者不能塑造自己；他既不能从无中变成完整，也没有人能公正地说他有实体，以致如果他被生出来，就永远同等能力。} \Person{西满}说：\Quote{那么他仅仅是一种关系，并因此邪恶？——具有伤害性，如水对火有害，但对适时干旱的土地有益；如铁对耕地有益，但对谋杀有害；情欲在婚姻方面不是恶，但在奸淫方面是恶；如谋杀是恶，但对谋杀者就其目的而言却是善；欺骗是恶，但对欺骗者却是愉悦；其他类似事物同样善恶。这样，既非恶，也非善；因为一方产生另一方。因为看似有害的事，难道不是让行善者喜悦，却惩罚受者吗？虽然一个人出于自爱，用尽一切手段满足自己似乎不义，但另一方面，一个人因自爱而受到公正法官的严厉惩罚，难道对谁来说不是不公正的吗？}\par

% structure: p0038 source=h2 target=subsection reason=relative-heading-rank; base=h2

\subsection{第十九章　有些行为真正邪恶}

\Person{伯多禄}说：\Quote{一个人应当通过自我克制来惩罚自己，当他的情欲想要匆忙去伤害他人时，要知道\Emph{那}邪恶者能毁灭邪恶者，因为他从一开始就获得了对他们的权能。这还不是对那些作恶者的恶；但他们的灵魂在毁灭之后仍受惩罚，你认为这实在残酷，是对的，尽管那预定为恶的人会说是对的。所以，正如我所说，我们应当避免为了一时之快而伤害他人，以免为了一点快乐而使我们陷入永罚。} \Person{西满}说：\Quote{那么，难道事物没有天生好坏之分，区别只是通过法律和习俗产生的吗？\Emph{难道不是吗？}波斯人的习惯是娶自己的母亲、姐妹和女儿，而与别的女人结婚则被视为最野蛮而禁止？所以，如果没有确定什么是恶，就不可能所有人都期待天主的审判。} \Person{伯多禄}说：\Quote{这不能成立；因为所有人都清楚与母亲同居是可憎的，即使波斯人，不过是整体的一小部分，在坏习俗的影响下未能认识到他们可憎行为的不义。同样，不列颠人公开在众人面前同居，却不以为耻；有些人吃别人的肉，并不感到厌恶；另一些人吃狗肉；还有人做其他不可言说的事。因此，我们不应以因习惯而偏离了自然作用的感知来做出判断。因为被杀是恶，即使所有人都否认；因为没有人愿意自己遭受它，在偷窃的事上，没有人因自己的惩罚而欢喜。那么，如果根本没有人承认这些是罪，那也仍然有必要期待对罪的审判。} \Person{伯多禄}说了这些话，\Person{西满}回答说：\Quote{那么，你认为这就是关于那邪恶者的真理吗？告诉我。}\par

% structure: p0040 source=h2 target=subsection reason=relative-heading-rank; base=h2

\subsection{第二十章　痛苦与死亡是罪的结果}

伯多禄说：\Quote{我们记得我们的主和导师曾命令我们说：\InnerQuote{为我和我家的儿子保存这些奥秘。}因此，祂也私下里给门徒解释天国的奥秘。但你们与我们争辩，并且只审查我们的言论，看它们是否真实，对于你们，陈述隐藏的\Emph{真理}是不敬的。但为了使旁观者不以为我是在找借口，因为我无法回答你们提出的断言，我将先反问一个问题：如果存在一种无痛苦的状态，那么\InnerQuote{恶者曾存在}这句话是什么意思？}西满说：\Quote{这话没有意义。}伯多禄说：\Quote{那么，恶与痛苦和死亡是一回事吗？}西满说：\Quote{似乎是。}伯多禄说：\Quote{那么，恶并不是永远存在的，甚至根本不能实体性地存在；因为痛苦和死亡属于偶性的类别，两者都不能与持久的力量共存。因为痛苦若不是和谐的断裂又是什么？死亡若不是灵魂与身体的分离又是什么？因此，当和谐存在时，就没有痛苦。因为死亡根本不属于实体存在的事物：因为死亡，如我所说，不过是灵魂与身体的分离；当这发生时，身体——本性上不能感觉——就分解了；但灵魂——能够感觉——仍然活着，并实体性地存在。因此，当和谐存在时，就没有痛苦，没有死亡，甚至没有致命的植物或毒蛇，也没有任何结局是死亡的事物。因此，在永生统治的地方，一切事物都会显得是被理性地创造的。当为了义，人通过基督和平统治的盛行而成为不朽时，情况就会如此：那时他的构造将被如此良好地安排，以至于不会\Emph{产生}剧烈的冲动；而且他的知识也将无误，以至于他不会\Emph{误认}恶为善；他将不受痛苦，因此他不会死亡。}\par

% structure: p0042 source=h2 target=subsection reason=relative-heading-rank; base=h2

\subsection{第二十一章 情欲、愤怒、忧伤的用途}

西满说：\Quote{你说得对；但在现今世界，人岂不似乎能感受到各种情感——例如，情欲、愤怒、忧伤等等？}伯多禄说：\Quote{是的，这些属于偶性的事物，而不是永远存在的事物，并且会发现它们现在对灵魂有益。因为情欲，由于创造万物的那位的美善旨意，被造就为存在于生物中，以便人被它引导去交合，从而增加人类，从其中选出一群适合永生的高级存在。如果没有情欲，没有人会麻烦自己去与妻子交合；但现在，为了快乐，并且似乎满足自己，人实现了祂的旨意。如果人将情欲用于合法的婚姻，他并非不敬；但如果他奔向奸淫，他就不敬，并且受惩罚，因为他滥用了一个好的法令。同样，愤怒被天主自然地点燃在我们内，以便我们被它引导去抵御伤害。然而，如果有人无节制地放纵它，他行事不义；但如果他适度使用它，他就做得对。此外，我们能够忧伤，以便我们在亲属、妻子、儿女、兄弟、父母、朋友或其他人的死亡时感到同情；因为如果我们不能同情，就会不人道。同样，所有其他情感都会发现是为我们而设的，只要考虑它们存在的原因。}\par

% structure: p0044 source=h2 target=subsection reason=relative-heading-rank; base=h2

\subsection{第二十二章 无知的罪}

西满说：\Quote{那么，为什么有些人夭折，周期性疾病发生？而且还有恶魔的攻击、疯狂以及其他各种能大大惩罚的痛苦？}伯多禄说：\Quote{因为人凡事随从自己的快乐，不在适当的时候交合；因此，精子的播种不合时宜，自然会产生许多祸害。因为他们应当反思，正如适合种植和播种的季节已被确定，同样，交合的适当日子也被指定，这些日子必须谨慎遵守。因此，一个精通梅瑟教导的人，责备百姓的罪，称他们为新月和安息日的儿子。然而在创世之初，人活得很长，没有疾病。但当他们因疏忽而忽略了适当时间的观察时，后续的儿女由于无知在不该交合的时候交合，使他们的孩子处于无数的痛苦中。因此，我们的导师，当我们问祂关于那个生来瞎眼、后来恢复视力的人，这人或他的父母是否犯了罪，使他生来瞎眼时，祂回答说：\InnerQuote{不是他犯了罪，也不是他的父母，而是为叫天主的德能藉着他彰显出来，以医治无知的罪。}（\BibleRef{若 9:3}）事实上，这样的痛苦是由于无知而来；例如，不知道何时应当与妻子同房，比如她是否洁净了经期。你之前提到的痛苦是无知的结果，而不是由于任何已犯的恶行。此外，给我一个不犯罪的人，我将给你一个不受苦的人；你会发现他不仅自己不受苦，还能治好别人。例如，梅瑟因他的虔诚，一生不受痛苦，并且籍着他的祈祷，他治好了因自己的罪而受苦的埃及人。}\par

% structure: p0046 source=h2 target=subsection reason=relative-heading-rank; base=h2

\subsection{第二十三章 人类生活中命运的不平等}

西满说：\Quote{让我承认这是实情：世人的命运不平等，在您看来岂非极不公平？因为一人贫困，另一人富裕；一人患病，另一人健康；人生中还有无数类似的差异。}伯多禄说：\Quote{西满，您没有察觉吗？您又一次观察过头了。我们正在讨论恶，您却岔开话题，引入了世上出现的反常现象。但我也会就此发言。世界是一件巧妙设计的工具，为使那将永远存在的男性，女性可以生出永远正义的儿子。但若没有需要帮助的穷人，他们在这里就不能成为完全虔诚的人。同样地，有病人，使他们有照顾的对象。其他的苦难也可作类似的解释。}西满说：\Quote{那些处境卑微的人不就不幸吗？因为他们遭受痛苦，好使别人成义。}伯多禄说：\Quote{如果他们的卑微是永恒的，那么他们的不幸就非常大。但人的卑微和升高是按着命运；对自己的命运不满的人可以申诉，并依法审理他的案件，他就能换另一种生活方式。}西满说：\Quote{您说的这命运和申诉是什么意思？}伯多禄说：\Quote{您现在要求另一个论题的解释；但若您允许，我们可以向您说明，人如何重生，改变根源，按法律生活，并获得永久的救恩。}\par

% structure: p0048 source=h2 target=subsection reason=relative-heading-rank; base=h2

\subsection{第二十四章 西满受法乌斯图斯责备}

西满听到这话，说：\Quote{不要以为我一边问您，一边在每个论题上同意您，然后转向下一个，仿佛我确信前一个为真；我装作向您的无知让步，好让您继续下一个论题，以便了解您全部的无知之后，不是凭猜测，而是根据充分的知识定您的罪。现在请允许我退下三天，然后我会回来，证明您一无所知。}西满说了这话，正要出去时，我父亲说：\Quote{西满，请听我说一会儿，然后您随便去哪里。我记得，在讨论开始之前，您曾指责我有偏见，尽管您对我还不了解。但现在，我依次听了你们的讨论，判断伯多禄占优势，并把说真理的功劳归给他，我看我判断得是否正确，是否有见识；还是并非如此？因为如果您说我判断正确，却不同意，那么您显然有偏见，因为您承认失败后却不愿同意。但如果我坚持伯多禄在讨论中占优势是不正确的，那么请说服我们，我们如何判断错误；否则您将停止在众人面前与他讨论，因为您总是被打败并同意，结果您自己的灵魂将受苦，被定罪，并因自己的良心而蒙羞，即使您不在全体听众面前感到羞愧，那将是最大的折磨；因为我们确实看到您被征服了，也亲耳听到您自己的口承认了。最后，所以我认为您不会按承诺回到讨论中来；但为了显得不是被打败，您在离开时承诺会回来。}\par

% structure: p0050 source=h2 target=subsection reason=relative-heading-rank; base=h2

\subsection{第二十五章 西满退场；索福尼亚斯请伯多禄陈述他对恶的真正看法}

西满听到这话，气得咬牙切齿，默默地走了。但伯多禄（因为白天还有相当长的时间）按手在众多人身上治愈他们，然后打发他们离开，与更亲密的朋友一同进屋坐下。他的随从中有一位名叫索福尼亚斯的，说：\Quote{天主受赞颂，啊，伯多禄，祂拣选了您，教导您为安慰善人。事实上，您与西满讨论得有尊严且非常耐心。但我们恳求您向我们谈论恶；我们期待您会向我们陈述您自己对恶的真实信念——但不是现在，而是明天，如果您觉得合适的话：因为我们顾惜您，因您讨论而感到疲劳。}伯多禄说：\Quote{我希望你们知道，凡乐意做事的人，即使在劳苦中也得到安息；但那个不做自己愿意之事的人，即使休息也感到极度疲惫。所以，当你们让我谈论我所喜欢的题目时，你们真是给了我极大的安息。}于是，我们满意于他的态度，并因他的疲劳而顾惜他，请求他将讨论推迟到夜里，那时他习惯对他的真诚朋友们讲话。我们一同进食盐，然后去睡觉。\par

\SourceRecord{Translated by James Donaldson.；Ante-Nicene Fathers；Vol. 8.；Edited by Alexander Roberts, James Donaldson, and A. Cleveland Coxe.；Buffalo, NY: Christian Literature Publishing Co.,；1886.；<http://www.newadvent.org/fathers/080819.htm>.}

% structure: p0001 source=h1 target=section reason=page-title

\section{第二十篇讲道}

% structure: p0002 source=h2 target=subsection reason=relative-heading-rank; base=h2

\subsection{第一章 伯多禄愿意满足索福尼亚}

在夜间，伯多禄起身叫醒了我们，然后像往常一样坐下，说：\Quote{你们可以问我任何喜欢的问题。} 索福尼亚首先开始对他说话：\Quote{您愿意为我们这些渴望学习的人解释关于恶的真正真理吗？} 伯多禄说：\Quote{我在与西满的讨论中已经解释过了；但因为我在与其它话题结合时陈述了关于它的真理，所以没有完全清楚地说出来；因为许多看似与真理同等重要的话题却为大众提供了某种关于真理的知识。因此，如果我现在陈述从前与西满一起曾陈述过的话题，不要以为你们\Emph{没有}受到与他同等的尊重。} 索福尼亚说：\Quote{您说得对；因为如果您现在为我们从当时讨论的许多话题中分离出来，您将使真理更加显明。}\par

% structure: p0004 source=h2 target=subsection reason=relative-heading-rank; base=h2

\subsection{第二章 两个时代}

伯多禄说：\Quote{因此，请听关于恶者的和谐真理。天主设定了两个王国，建立了两个时代，决定将现今的世界交给恶者，因为它微小且迅速消逝；但祂应许为善者保存未来的时代，因为它将是伟大而永恒的。因此，祂以自由意志创造了人，并具有倾向自己所愿行为的能力。人的身体由三部分组成，源自女性；因为它有淫欲、愤怒和悲伤，以及由此而来的一切。但精神并非单一，而是由三部分组成，源自男性；它能够推理、知识和恐惧，以及由此而来的一切。每个三元组都有一个根源，因此人是两种混合物的组合：女性与男性。因此，在他面前摆了两条道路——顺从法律与不顺从法律；并且建立了两个王国——一个称为天国，另一个称为现今在地上作王的人的国。同时也设立了两位君王，其中一位被拣选按照法律统治现今暂时的世界，他的构成使他对恶人的毁灭感到喜乐。但另一位善者，是未来时代的君王，喜爱人的整个天性；但由于不能在现今世界放胆行事，他劝告有益的事，如同一个试图隐藏自己真实身份的人。}\par

% structure: p0006 source=h2 target=subsection reason=relative-heading-rank; base=h2

\subsection{第三章 善者与恶者的作为}

伯多禄说：\Quote{但在这二者中，\Emph{一位}因天主的命令对另一位采取暴力。此外，每个人有权服从他所喜欢的任何一个，以行善或作恶。但如果有人选择行善，他就成为未来善王的产业；但如果有人作恶，他就成为现在恶者的奴仆，恶者因他的罪经过公正审判获得权柄管辖他，并希望在未来的时代之前\Emph{使用}它，因在今世惩罚他而喜乐，因而仿佛满足自己的私欲，他完成了天主的旨意。而另一位，被造成以权柄管辖义人为乐，当他找到义人时，就极其喜悦，用永生拯救他；他也仿佛满足自己，将他由此感受到的满足归于天主。现在，每个不义的人都有能力悔改并获得拯救；而每个义人可能因自己生涯末了所犯的罪而受惩罚。此外，这两位领袖是天主敏捷的双手，渴望抢先完成祂的旨意。但事实如此，法律也以天主的口气说过：\BibleQuote{我要杀死，也要使人活；我要击伤，也要医治。} 因为确实，祂杀死也使人活。祂通过左手杀死，即通过恶者，他被造成以折磨不虔敬者为乐。祂通过右手拯救并施恩，即通过善者，他被造以义人的善行和得救为乐。现在，他们的本质并非出于天主之外：因为没有其他最初源头。的确，他们也不是作为动物从天主发出的，因为他们与祂同心；也不是偶然自发地违抗祂的旨意而出现，因为那样祂权能的最大运用就会被破坏。而是从天主发出了四个原初元素——热与冷，湿与干。因此，祂是每个本质的父亲，但不是由元素组合可能产生的性情；因为当这些元素从外部结合时，性情作为孩子在其中产生。那么，恶者效忠天主直到现今世界的终结，可以改变其构成而成为善，因为他确实不是单一趋向罪恶的均质本质。因为他并不更多作恶，尽管他是恶的，因为他获得了合法折磨的权柄。}\par

% structure: p0008 source=h2 target=subsection reason=relative-heading-rank; base=h2

\subsection{第四章 人因无知而犯罪}

当伯多禄这样说时，米加——他自己也是门徒之一——问道：\Quote{那么，人犯罪的原因是什么呢？} 伯多禄说：\Quote{这是因为他们无知，不知道当审判时他们必为自己的恶行受罚。因此，他们拥有淫欲——如我在别处所说——为了生命的延续，以任何偶然的方式满足它，也许是通过败坏男童，或通过其他谄媚的罪。因为由于他们的无知，如前所述，他们被无所畏惧驱使，以非法的方式满足淫欲。因此，天主不是恶的，祂正当地将淫欲置于人内，使生命得以延续；但最不虔敬的是那些恶劣地使用淫欲之善的人。同样的道理也适用于愤怒：如果有人正当地使用它，就像在他能力范围内，他就是虔诚的；但如果超过限度，自己施行审判，他就是不虔敬的。}\par

% structure: p0010 source=h2 target=subsection reason=relative-heading-rank; base=h2

\subsection{第五章 索福尼亚主张天主不能产生与自己不相似的事物}

索福尼亚斯又说：\Quote{我主伯多禄，您的极大忍耐给予我们胆量，为求准确而向您提出许多问题。因此我们满怀信心地四处询问。我记得，西满昨天在同您的讨论中曾说，那恶者若是出于天主，便因而具有与差遣他的那一位相同的本体，他本该是善的，而不是恶的。但您回答说情况并非总是如此，因为许多恶子生于善父母，正如从亚当生了两个不同的儿子，其中一个为恶，另一个为善。西满指责您使用了人的例子，您便回答说，既然如此，我们根本不该承认天主产生；因为这也是一个人的例子。我，索福尼亚斯，承认天主产生，但我不同意祂产生恶的东西，即使人中善者生出恶子。您不要以为我毫无理由地将人的某些属性归于天主，而拒绝将其他属性归于祂；我承认祂产生，但不同意祂产生与祂不相似的东西。因为人，正如你所料，生出于性情方面不像自己的儿子，原因如下：人由四元素构成，根据年岁的变化而多样地改变身体；因此，人的身体发生适当的增减变化，每个季节破坏和谐的混合。当混合并非总是精确地保持同一状态时，种子有时带一种混合，有时带另一种而被释放；接着，根据季节的混合，随之而来的是或善或恶的性情。但在天主身上，我们不能设想任何这类事情；因为祂是不变的、永远存在的，无论何时祂愿意发出，所发出的东西在各方面都必然与产生者处于同一状态——我指的是在本体和性情方面。但若有人执意主张祂是可变的，我不知道他怎能主张祂是不死的。}\par

% structure: p0012 source=h2 target=subsection reason=relative-heading-rank; base=h2

\subsection{第六章 天主改变自身的能力}

伯多禄听了这话，思考片刻，说道：\Quote{我认为任何人谈论恶事都不可能不承行那恶者的旨意。因此，我知道这点，却不知该怎么办——是该沉默还是说话。因为我若沉默，就会招致众人的嘲笑，因为我自称宣扬真理，却不知恶的解释。我若陈述己见，又恐怕寻求恶事完全不中悦天主，因为只有寻求善事才中悦祂。然而，在回答索福尼亚斯的陈述时，我要使自己的思想更加清楚。我同意他这样的看法：我们不该将人的一切属性归于天主。例如，人没有可转换的身体，故不能转换；但他们具有一种可通过年岁季节的流逝而变化的性质。但天主并非如此；因为祂藉着内在的圣神，以一种无法形容的能力，变成祂所喜欢的任何形体。这一点更容易相信，因为空气从祂接受了这样的性质，被无形精神渗透而化为露水，凝聚而成水，水凝结而成石头和泥土，石头碰撞而燃起火。按照这样的变化和转化，空气先变成水，最后通过转化而成为火，潮湿变成其天然的对立面。为什么？天主不是把梅瑟的棍杖变成动物，使它成为一条蛇，然后又把它变回棍杖吗？并且藉着这根被转化的棍杖，他把尼罗河的水变成血，后来又变回水。是的，甚至人——本是尘土——祂藉着嘘气变成血肉，又使他变回尘土。梅瑟本人本是血肉，岂非转化为极灿烂的光明，以致以色列子民不能直视他的脸？更何况，天主完全能够将祂自己转化为祂所愿意的任何事物。}\par

% structure: p0014 source=h2 target=subsection reason=relative-heading-rank; base=h2

\subsection{第七章 对人不能改变自身这一异议的答复}

\Quote{但或许你们中有人认为，某物可受一种影响而成为某种东西，受另一种影响而成为另一种东西，但没有人能将自己改变成他所希望的任何东西；并且，改变是属于那按本性必老必死者的特性，我们对不朽的存在不应怀有这种想法。难道不是那些不老、具有火一样体质的天使变成了血肉吗？例如，那些款待亚巴郎的天使，人们洗他们的脚，仿佛他们是同样体质的人。还有，与雅各伯——一个人——角力的，是一位天使，他转化为血肉，以便能与他近身搏斗。同样，在他自愿角力之后，又变回自己本来的形态；当他化为火焰时，并没有烧毁雅各伯的宽腱，却使他发热并使他瘸了。现在，那不能变成其他东西（无论它想变成什么）的，是有死的，因为它受其本性的支配；但那能随时变成他愿意的任何东西的，是不死的，因为他能恢复新的状态，因为他控制着自己的本性。因此，天主的能力更是如此：无论何时，祂都能将身体的本体改变为祂所愿意的任何东西；并且藉着所发生的改变，祂发出那在本体上相似、但在能力上不平等的东西。因此，发出者能将祂所发出的东西变成不同的本体，并能再次变回自己的本体；而被发出者，既是由于从发出者发出的改变而产生的，作为祂的孩子，没有发出者的意愿就不能变成任何别的东西，除非他自己愿意。}\par

% structure: p0016 source=h2 target=subsection reason=relative-heading-rank; base=h2

\subsection{第八章 善者与恶者的起源不同}

当伯多禄这样说时，米该亚——他也是陪伴他的同伴之一——说：\Quote{我也愿从你学习，善者是否以与恶者相同的方式产生？但若它们以相似方式产生，则依我看来，它们便是兄弟。}伯多禄说：\Quote{它们并非以相似方式产生：因为你无疑记得我起初说过，邪恶者身体的本体，其起源为四重，是由天主精心拣选并派遣的；但当它按照派遣者的旨意在外在结合时，因这结合便产生了那以恶为乐的性情；由此你可见，那由祂派遣、并且始终存在的四重起源的本体，是天主的孩子；但那偶然产生、以恶为乐的性情，是在本体被他外在结合时才出现的。因此，这性情既非由天主所生，也非由任何其他者所生，也非由祂所派遣，也非自发出现，也非如结合前的本体那样始终存在，而是根据天主的旨意，因外在结合而偶然出现的。我们常说，这是必然的。但善者是从天主最美的变化中受生，并非因外在结合而偶然出现，祂真是圣子。然而，既然这些教义未经书写，只凭推论向我们确证，我们绝不可视之为绝对确切的情况。因为若我们反其道而行，我们的心智便会停止探索真理，以为已完全领悟。因此，要记住这些事；因为我不可向所有人陈述这些，只可向那些经过考验最为可信的人陈述。我们也不应鲁莽地彼此坚持这些主张，也不可胆敢说话，仿佛你们精确知晓隐秘真理的发现，而只应在沉默中反思它们；因为，若有人偶然说某事是如此，说话者便会犯错，他将因敢于对那应受沉默尊崇之事说话——甚至只是对自己——而受罚。}\par

% structure: p0018 source=h2 target=subsection reason=relative-heading-rank; base=h2

\subsection{\Emph{第九章　为何公义的天主将邪恶者派定于恶人之上}}

当伯多禄这样说时，拉匝禄——他也是他的追随者之一——说：\Quote{请向我们解释这和谐，如何能合理：邪恶者由公义的天主派定作为不虔敬者的惩罚者，而他自己后来却同他的天使及罪人一起被投入下层黑暗中：因为我记得导师本人曾这样说。}伯多禄说：\Quote{我确实承认，邪恶者不作恶，因他成就了所赐予他的律法。虽然他具有邪恶的性情，但出于对天主的敬畏，他行无不义的事；但他控告真理的导师以诱捕不谨慎的人，他本人被称为控告者（魔鬼）。但我们无误导师的声明，即他和他的天使，与受骗的罪人一同进入下层黑暗中，可作如下解释。邪恶者按其构成得到以黑暗为乐的份，便喜悦与他的同伴天使们一同下到塔尔塔洛斯的黑暗；因为黑暗是火所喜爱的。但人的灵魂是纯净光明的点滴，被不同类的火本体吸收；它们没有可死的本性，便按功过受罚。但若那引人行恶的首领不因不以黑暗为乐而被投入黑暗，那么他那以恶为乐的性情就不能因另一种结合而改变为向善的倾向。这样，他就会被判定与善者同在，更是因为他虽具有以恶为乐的构成，但因敬畏天主，他没有作任何违反天主法律命令的事。圣经不是通过一个神秘的暗示，藉着大司祭亚郎的棍杖变成蛇，又再变回棍杖，指出邪恶者的构成以后会改变吗？}\par

% structure: p0020 source=h2 target=subsection reason=relative-heading-rank; base=h2

\subsection{\Emph{第十章　为何有些人相信，有些人不信}}

之后，若瑟——他也是他的追随者之一——说：\Quote{你所说的一切都是对的。也请教我这个，因为我渴望知道：为什么当你向众人讲同样的言论时，有些人相信，有些人不信？}伯多禄说：\Quote{这是因为我的言论并非咒语，使每个听见的人都必然毫不犹豫地相信。有些人相信，有些人不信，这事实向明智者指出了自由意志。}他说这话时，我们都赞美了他。\par

% structure: p0022 source=h2 target=subsection reason=relative-heading-rank; base=h2

\subsection{\Emph{第十一章　阿皮翁与安努比翁到达}}

当我们正要去吃饭时，有人跑进来说：\Quote{阿皮翁·普雷斯托尼塞斯刚与安努比翁从安提约基雅来到，他与西满同住。}我父亲听见这事，很高兴，对伯多禄说：\Quote{你若允许，我要去问候阿皮翁和安努比翁，他们从小就是我的朋友。因为或许我能说服安努比翁与克肋孟讨论星命学。}伯多禄说：\Quote{我允许你，并称赞你尽朋友的本分。但现在你看，在天主的眷顾下，从各方汇聚了有助于你完全确证的事，使和谐圆满。但我这样说，是因为安努比翁的到来对你有利。}我父亲说：\Quote{实在，我看正是如此。}说完这话，他便往西满那里去了。\par

% structure: p0024 source=h2 target=subsection reason=relative-heading-rank; base=h2

\subsection{\Emph{第十二章　福斯托以西满的面貌出现在他的朋友们面前}}

现在，我们所有与\Person{伯多禄}在一起的人整夜互相询问，因为从所说的话中获得的喜悦和快乐，我们一直醒着。但当黎明终于开始破晓时，\Person{伯多禄}看着我及我的兄弟们，说：\Quote{我很困惑，不知道你父亲在干什么。}他正说着这话，我们的父亲进来了，听见\Person{伯多禄}在谈论他；看见他不高兴，便上前搭话，为在外过夜而道歉。但我们看他时都惊异：因为我们看见\Person{西满}的形貌，却听见我们父亲\Person{法乌斯托}的声音。当我们逃避他、厌恶他时，我们的父亲因受到我们如此严厉和敌意的对待而惊讶。但只有\Person{伯多禄}看见他本来的形状，对我们说：\Quote{你们为什么恐惧地躲避自己的父亲？}但我们和我们的母亲说：\Quote{我们看到面前的是\Person{西满}，却有着我们父亲的声音。}\Person{伯多禄}说：\Quote{你们只认出他的声音，那不受魔法影响；但我的眼睛也不受魔法影响，我能看见他真正的形状，他不是\Person{西满}，而是你们的父亲\Person{法乌斯托}。}然后，他看着我的父亲，说：\Quote{他们看见的不是你真正的形状，而是我们最致命的敌人、最不虔敬的人\Person{西满}的形状。}\par

% structure: p0026 source=h2 target=subsection reason=relative-heading-rank; base=h2

\subsection{第十三章 西满的逃亡}

正当\Person{伯多禄}这样说时，有一个先前去了\Place{安提约基雅}的人进来了，他从\Place{安提约基雅}回来，对\Person{伯多禄}说：\Quote{我希望您知道，我的主，\Person{西满}在\Place{安提约基雅}公开行了许多奇迹，并称您为术士、变戏法的\Emph{和杀人犯}，激起了他们对您的仇恨，以至于每个人都渴望，如果您到那里，他们就要吃你的肉。因此，我们先前去的那些人，连同那些表面上被您派到\Person{西满}那里的我们的弟兄，看见那座城对您狂怒，便秘密会面，商量我们该做什么。确实，当我们非常困惑时，百夫长\Person{科尔乃略}到了，他被皇帝派到该省的总督那里。他就是我们的主在\Place{凯撒勒雅}治愈的那人，当时他被魔鬼附身。我们秘密派人去请他，告诉他我们沮丧的原因，恳求他的帮助。他非常爽快地答应，如果我们协助他的努力，他就会惊吓\Person{西满}，迫使他逃跑。当我们都答应会欣然做一切事时，他说：\InnerQuote{我要通过许多朋友散布消息，说我秘密来捉拿他；我要假装在搜索他，因为皇帝处死了许多术士，并得到了关于他的情报，派我来搜寻他，好惩罚他，如同惩罚他之前的术士；而你们那边与他在一起的人必须向他报告，仿佛他们从秘密渠道听说，我被派来捉拿他。也许当他从他们那里听到时，会惊慌而逃跑。}当我们曾打算做别的事，但事情却以如下方式发生。因为他从许多陌生人那里听到消息，那些人通过秘密告知他而大大取悦了他，也从我们那些假装依附他的弟兄那里听到，并把它当作他自己追随者的意见，他便决定退去。于是匆匆离开\Place{安提约基雅}，带着\Person{阿特诺多鲁斯}来到这里，正如我们所听说的。因此我们建议您暂时不要进那座城，直到我们弄清他们能否在他不在时忘记他对您提出的指控。}\par

% structure: p0028 source=h2 target=subsection reason=relative-heading-rank; base=h2

\subsection{第十四章 西满导致法乌斯托形貌的改变}

当先前的报告者做了报告后，\Person{伯多禄}看着我的父亲，说：\Quote{你听到了，\Person{法乌斯托}；你形貌的改变是由术士\Person{西满}造成的，如今已清楚了。因为他以为皇帝的\Emph{一个仆人}要捉拿他惩罚他，他便害怕而逃跑了，把你变成他自己的形貌，好使你被处死时，你的孩子们会悲伤。}我的父亲听到这话，便哭泣哀叹，说：\Quote{你猜对了，\Person{伯多禄}。因为\Person{阿努比翁}，我的挚友，向我暗示了他的计谋；但我没有相信他，我真是\Emph{不幸的人}，因为我应受这苦。}\par

% structure: p0030 source=h2 target=subsection reason=relative-heading-rank; base=h2

\subsection{第十五章 法乌斯托的悔改}

当我的父亲这样说后，不久\Emph{\Person{阿努比翁}来了}，向我们宣布\Person{西满}的逃亡，以及他如何在那天夜里赶往\Place{犹太}。他发现我们的父亲在哀哭，并悲叹说：\Quote{哎呀，哎呀！不幸的人！当有人告诉我他是术士时，我没有相信。我真是可悲的人！我被妻子儿女认出只有一天，就很快回到我先前无知时的悲惨境况。}我的母亲哀哭着，拔自己的头发；我们因父亲的变形而痛苦呻吟，无法理解这究竟是怎么回事。但\Person{阿努比翁}站着哑口无言，看到并听到这些事；而\Person{伯多禄}当着众人的面对我们，他的孩子们说：\Quote{相信我，这是你们的父亲\Person{法乌斯托}。所以我劝你们要把他当作父亲来照顾。因为天主将恩赐某个时机，让他脱去\Person{西满}的形状，再清楚地显出你们父亲的形状。}说着这话，他看着我的父亲，说：\Quote{我允许你问候\Person{阿皮翁}和\Person{阿努比翁}，因为你声称他们是你的童年朋友，但我不允许你与术士\Person{西满}交往。}\par

% structure: p0032 source=h2 target=subsection reason=relative-heading-rank; base=h2

\subsection{第十六章 西满为何将自己的形貌给予法乌斯托}

我父亲说：\Quote{我犯了罪，我承认。} 安努比昂说：\Quote{我同他一起恳求你饶恕这位高尚良善、受欺骗的老人：因为这不幸的人成了那个臭名昭著的家伙的玩物。但我要告诉你事情是怎样发生的。这位善良的老人来向我们致意。但就在那时，我们这些在场的人恰好在听西蒙说话，他想在那天晚上逃走，因为他听说有人奉皇帝之命到\Place{劳迪塞亚}来捉拿他。但法乌斯塔斯进来时，他把自己的怒气\Emph{转向}他，这样对我们说：\InnerQuote{等他来了，叫他同你们一起吃饭；我要准备一种油膏，当他吃过晚饭，他可以用一点涂在脸上，然后他在众人眼中就会变成我的模样。但我用某种植物的汁液涂你们，这样你们就不会被他的新形象欺骗；但在其他人眼中，法乌斯塔斯看上去就像西蒙。}}\par

% structure: p0034 source=h2 target=subsection reason=relative-heading-rank; base=h2

\subsection{第十七章·安努比昂为法乌斯塔斯服务}

\Quote{当他事先这样说时，我说：\InnerQuote{那么，你现在期望从这样的计谋中得到什么好处？}西蒙说：\InnerQuote{首先，那些寻找我的人，当他们捉住他，就会放弃对我的搜索。但如果他被皇帝处死，巨大的悲痛将降临到他的孩子身上，他们离开了我，逃到\Emph{伯多禄}那里，现在帮助他工作。}现在，伯多禄，我向你承认真相：我由于对\Emph{西蒙}的恐惧，未能将此事告知\Emph{法乌斯塔斯}。但西蒙没有给我们私下交谈的机会，\Emph{以免}我们中有人可能向他揭露西蒙的邪恶计谋。西蒙于是在半夜起身，由阿皮翁和阿特诺多鲁斯护送，逃往\Place{犹太}。然后我假装生病，以便在他们走后留下来，我可以让法乌斯塔斯立即回到他自己的人那里，如果他可能有机会，通过藏在你们中间而避开注意，以免被那些搜寻西蒙的人当作西蒙抓住，死于皇帝的愤怒。因此，在深夜，我把他送到你那里；出于对他的挂虑，我夜间来看他，打算在护送西蒙的人返回之前回去。} 他望着我们，说：\Quote{我，安努比昂，看到你父亲的真实形象；因为正如我先前告诉你们的，我被西蒙亲自涂了油，使我眼中能看到法乌斯塔斯的真实形象。因此，我极为惊异，惊叹于西蒙的魔力，你们站着却不认识自己的父亲。} 我们的父亲、母亲和我们自己都为这共同的祸患哭泣，安努比昂也出于同情与我们一同哭泣。\par

% structure: p0036 source=h2 target=subsection reason=relative-heading-rank; base=h2

\subsection{第十八章·伯多禄应许恢复法乌斯塔斯的本来面目}

然后伯多禄向我们承诺恢复我们父亲的形象，并对他说：\Quote{法乌斯塔斯，你听到了我们目前的情况。因此，当包裹你的欺骗性形状对我们有用，并且你协助我做了我要你做的事之后，我就会恢复你的真实形象，但你必须先执行我的命令。} 我父亲说：\Quote{我将非常乐意做一切我能做的事；只是请恢复我对自己家人的本来面目。} 伯多禄回答说：\Quote{你亲耳听到，那些在我之前的人从安提约基雅回来，说西蒙曾在那里，并强烈煽动群众反对我，称我为魔术师、杀人犯、骗子、变戏法的，以至于那里所有的人都渴望吃我的肉。那么，你要照我告诉你的去做。你将把克肋孟留在我这里，你将与你的妻子、你的儿子法乌斯蒂努斯和法乌斯蒂亚努斯先我们一步前往安提约基雅。还有一些我会认为有助于推进我的计划的人将陪同你。}\par

% structure: p0038 source=h2 target=subsection reason=relative-heading-rank; base=h2

\subsection{第十九章·伯多禄对法乌斯塔斯的指示}

\Quote{\Emph{当你}与这些人一起在\Place{安提约基雅}时，你看起来像西蒙，要公开宣告你的\Emph{悔改}，说：\InnerQuote{我西蒙向你们宣告：我承认我以前关于伯多禄所说的一切都是\Emph{彻头彻尾的谎言}；\Emph{他不是}骗子，也不是杀人犯，也不是变戏法的；我之前出于愤怒对他所说的一切恶事都不是真的。因此，我本人恳求你们，我曾是你们对他仇恨的原因，请停止恨他；因为他是那位为拯救世界而被天主派遣的真先知的真宗徒。因此我也劝你们相信他所宣讲的；因为如果你们不这样做，你们的整个城市将被彻底摧毁。现在我希望你们知道我为何向你们作此忏悔。今夜天主的使者狠狠地鞭打了我这个不敬神的人，因为我是真理宣告者的敌人。因此我恳求你们，即使我以后再来，试图说任何反对伯多禄的话，也不要听我。因为我向你们承认，我是魔术师，我是骗子，我是变戏法的。但或许通过悔改，我可以擦去我从前所犯的罪。}}\par

% structure: p0040 source=h2 target=subsection reason=relative-heading-rank; base=h2

\subsection{第二十章·法乌斯塔斯、他的妻子和儿子们准备前往安提约基雅}

当\Person{伯多禄}提议此事后，我父亲说：\Quote{我知道你想要什么，因此不必费心。因为当我到达那地方时，我必定会好好照管，就你的事作我该作的陈述。} \Person{伯多禄}又建议道：\Quote{那么，当你看到这城从对我的仇恨中转变，并渴望见我时，请将此事通知我，我就会立刻到你那里去。当我到达那里时，当天我就除去现在笼罩你的怪异外貌（\OriginalTerm{怪异外貌}{strange shape}），并使你原本的样貌清楚可见，为你自己的人和所有其他人看见。} 说完这话，他让他的儿子们，我的兄弟们，以及我们的母亲\Person{玛提迪亚}与他同行；他也命令一些更亲密的熟人与他同去。但我母亲不愿与他同去，说道：\Quote{我若与\Person{西满}的外貌交往，似乎成了淫妇；但若我必须与他同行，我不可能与他在同一榻上躺卧！但我不知道我是否会被说服与他同去。} 正当她极不情愿去时，\Person{安努比翁}敦促她，说：\Quote{请相信我和\Person{伯多禄}，以及这声音本身，这确实是你的丈夫\Person{法斯托斯}，我爱他不亚于你。而我自己\Emph{将与他同行}。} \Person{安努比翁}说这话后，我们的母亲答应与他同行。\par

% structure: p0042 source=h2 target=subsection reason=relative-heading-rank; base=h2

\subsection{第二十一章 \Person{阿皮翁}与\Person{雅典诺多罗斯}为寻\Person{法斯托斯}而返}

但\Person{伯多禄}说：\Quote{天主以最令人满意的方式安排我们的事务；因为我们有占星家\Person{安努比翁}与我们同在。因为当我们到达\Place{安提约基雅}时，他日后将就出生命盘（\OriginalTerm{出生命盘}{genesis}）讲论，作为朋友向我们提供他真正的见解。} 如今，在午夜过后，我父亲与\Person{伯多禄}吩咐同行的人以及\Person{安努比翁}急忙赶往邻近的\Place{安提约基雅}。第二天清早，在\Person{伯多禄}出去讲论之前，护送\Person{西满}的\Person{阿皮翁}和\Person{雅典诺多罗斯}返回\Place{劳狄刻雅}寻找我父亲。但\Person{伯多禄}查明此事，敦促他们进来。他们进来坐下后，说：\Quote{法斯托斯在哪里？} \Person{伯多禄}答道：\Quote{我们不知道；因为从昨晚他去你们那里以后，他的亲属再未见过他。但昨天早晨\Person{西满}来找他；当我们没有回答他时，他似乎有些恍惚，因为他自称是法斯托斯；但未被相信，他就哭泣哀叹，威胁要自杀，然后冲了出去，朝向海的方向。}\par

% structure: p0044 source=h2 target=subsection reason=relative-heading-rank; base=h2

\subsection{第二十二章 \Person{阿皮翁}与\Person{雅典诺多罗斯}返回\Person{西满}处}

当\Person{阿皮翁}和同来的人听闻此事，他们嚎啕哀叹，说：\Quote{你们为何没有接待他？} 同时\Person{雅典诺多罗斯}想对我说：\Quote{那是你的父亲法斯托斯；} \Person{阿皮翁}抢先说道：\Quote{我们听一人说，\Person{西满}找到了他，并敦促他\Emph{与自己同行}，而法斯托斯本人也恳求他，因为他不想在儿子们成为犹太人后见到他们。我们听到这事，为了他的缘故，前来寻找他。但他不在这里，显然那告诉我们信息的人说了实话，我们从他听到并已转告你们。} 而我\Person{克莱孟}察觉\Person{伯多禄}的计谋，他想在他们心中引起怀疑，以为他打算在他们中间寻找那老人，好使他们害怕并逃走，于是配合他的计划，对\Person{阿皮翁}说：\Quote{听我说，我最亲爱的\Person{阿皮翁}。我们曾渴望将我们自己看为好的东西给他，当作我们的父亲。但如果他自己不愿接受，反而恐惧地逃离我们，我要说一句相当冷酷的话：\InnerQuote{我们也不在乎他。}} 我说这话后，他们离去，似乎被我冷酷激怒；据我们次日所知，他们跟随\Person{西满}去了\Place{犹太}。\par

% structure: p0046 source=h2 target=subsection reason=relative-heading-rank; base=h2

\subsection{第二十三章 \Person{伯多禄}前往\Place{安提约基雅}}

如今，十天过去后，\Emph{我们中的一人}从我们父亲那里来，向我们宣告我们的父亲如何公开\Emph{站在}以\Emph{西满的}外貌（\OriginalTerm{外貌}{shape}）\Emph{出现}，控诉他；又如何通过赞扬\Person{伯多禄}使整个\Place{安提约基雅}城渴望见他；结果，所有人都说他们急于见他，并且有些人因对\Person{伯多禄}的极高爱戴，因以为他是西满而愤怒，想要动手抓\Person{法斯托斯}，以为他是西满。因此他害怕被处死，已派人请求\Person{伯多禄}立刻前来，若想见到他活着，并在全城渴望最强烈时在适当的时候出现在城中。\Person{伯多禄}听到这事，召集众人商议，并委派一位随从担任主教；在\Place{劳狄刻雅}停留三天施洗并治愈后，他赶往邻近的\Place{安提约基雅}城。阿们。\par

\SourceRecord{Translated by James Donaldson.；Ante-Nicene Fathers；Vol. 8.；Edited by Alexander Roberts, James Donaldson, and A. Cleveland Coxe.；Buffalo, NY: Christian Literature Publishing Co.,；1886.；<http://www.newadvent.org/fathers/080820.htm>.}
